\documentclass[11pt]{article}
\usepackage[T1]{fontenc}
\usepackage{lmodern}
\usepackage[margin=1.08in]{geometry}
\usepackage{amsmath,amssymb,amsthm,mathtools}
\usepackage{enumitem,microtype}
\usepackage{tikz}
\usetikzlibrary{arrows.meta,calc,positioning}
\usepackage[numbers,sort&compress]{natbib}
\usepackage[hidelinks,unicode]{hyperref}
\usepackage{xurl}
\usepackage{bookmark}
\numberwithin{equation}{section}
\newtheorem{theorem}{Theorem}[section]
\newtheorem{proposition}[theorem]{Proposition}
\newtheorem{lemma}[theorem]{Lemma}
\newtheorem{corollary}[theorem]{Corollary}
\theoremstyle{definition}
\newtheorem{definition}[theorem]{Definition}
\theoremstyle{remark}
\newtheorem{remark}[theorem]{Remark}
\newcommand{\R}{\mathbb R}
\newcommand{\dd}{\,\mathrm d}
\DeclareMathOperator{\tr}{tr}
\DeclareMathOperator{\divg}{div}
\DeclareMathOperator{\Div}{div}
\DeclareMathOperator{\sym}{sym}
\DeclareMathOperator{\Sym}{Sym}
\DeclareMathOperator{\Vol}{Vol}
\DeclareMathOperator{\scal}{Scal}
\DeclareMathOperator{\Area}{Area}
\DeclareMathOperator{\dist}{dist}
\DeclareMathOperator{\Ric}{Ric}
\DeclareMathOperator{\Hess}{Hess}
\DeclareMathOperator{\supp}{supp}
\DeclareMathOperator{\Scal}{Scal}
\DeclareMathOperator{\grad}{grad}
\DeclareMathOperator{\Id}{Id}
\newcommand{\Achi}{A_\chi}
\newcommand{\Ad}{A_d}
\newcommand{\cT}{\mathcal T}
\newcommand{\gbar}{\bar g}
\newcommand{\ghat}{\widehat g}
\newcommand{\cB}{\mathcal B}
\newcommand{\cP}{\mathcal P}
\newcommand{\cD}{\mathcal D}
\newcommand{\cW}{\mathcal W}
\newcommand{\eps}{\epsilon}
\newcommand{\omegaThree}{\omega_3}
\newcommand{\inner}[2]{\langle#1,#2\rangle}
\setlength{\emergencystretch}{1.5em}
\allowdisplaybreaks[1]

\newcommand{\N}{\mathbb N}
\newcommand{\Z}{\mathbb Z}
\newcommand{\cE}{\mathcal E}
\newcommand{\Mext}{M_{\mathrm{ext}}}
\newcommand{\abs}[1]{\lvert#1\rvert}
\newcommand{\norm}[1]{\lVert#1\rVert}
\newcommand{\AH}{\mathrm{AH}}
\newcommand{\II}{\mathrm{II}}
\newcommand{\tf}{\circ}
\usepackage{mathrsfs,bm,booktabs,longtable,array,needspace}
\makeatletter
\def\input@path{{four/}{neutral/}{maximal/}}
\makeatother
\setcounter{tocdepth}{1}
\pdfinfoomitdate=1
\pdftrailerid{}
\pdfsuppressptexinfo=15
\hypersetup{pdftitle={Boundary graph deformations for the spacetime Penrose inequality},pdfauthor={OpenAI}}
\title{Boundary graph deformations for the spacetime Penrose inequality}
\author{OpenAI}
\date{September 27, 2026}
\begin{document}
\maketitle
\begin{abstract}
We construct graph and conformal deformations of trapped initial-data
exteriors in three and four spatial dimensions. The resulting metrics
have nonnegative scalar curvature, strictly negative inner mean curvature,
a lower bound protecting every enclosing cut, and an arbitrarily small
upper error in ADM energy. We also give a direct four-dimensional maximal
vacuum construction that retains a noncompact decaying second fundamental
form. These constructions give boundary routes to the corresponding
numerical Penrose inequalities.
\end{abstract}
\tableofcontents
\section{Introduction}\label{b:introduction}

The spacetime Penrose inequality compares the invariant ADM mass
$m_{\mathrm{ADM}}=\sqrt{E^2-|P|^2}$ with the area needed to enclose a trapped
region, where $E$ and $P$ are the energy and linear momentum of
asymptotically flat initial data.  The deformations below control energy
in a chosen end; for general weak data, the invariant-mass conclusion
uses a separate passage to rest ends, where $P=0$.
Reducing the problem to the Riemannian Penrose inequality already
exposes the main difficulty.  A deformation must create nonnegative scalar
curvature while controlling both quantities in the comparison: energy
at infinity and the area of every enclosing cut.  If the deformation is
performed on an exterior with a genuine inner boundary, its boundary
flux is a third quantity that must be controlled.

This paper constructs such boundary deformations in three and four
spatial dimensions.  It also develops a direct four-dimensional maximal
construction that retains a noncompact decaying second fundamental form.
The common output is a smooth complete Riemannian exterior with
nonnegative scalar curvature, strictly negative inner mean curvature,
a metric lower bound protecting all enclosing areas, and an arbitrarily
small upper error in energy.  A minimal enclosure then supplies the
boundary to which the Riemannian Penrose inequality applies.

To explain the area comparison, let $(M,g,K)$ be an initial data exterior
with compact boundary and one asymptotically Euclidean end.  The tensor
$K$ is the second fundamental form.  Orient an inner boundary toward
the end, let $H$ be its mean curvature, and write
$\theta_+=H+\tr_{\mathrm{tan}}K$ for its future expansion.
A full enclosing cut is the entire compact boundary of a connected
closed outer domain that contains the distant end and has interior in
the open exterior.  Every boundary component counts, including any
part coincident with the original boundary.  Denote the infimum of their
$g$-areas by $A_*(g)$.  The infimum need not be attained.
The precise asymptotic and constraint hypotheses for each construction
are stated at its beginning.

Both the area and the horizon predicate matter. In spatial dimension
three, Ben-Dov's example
has an outermost future apparent horizon $S$ with
$|S|>16\pi m_{\mathrm{ADM}}^2$~\cite{BenDov2004}.
Carrasco--Mars construct outer area-minimizing generalized apparent
horizons, defined by $H=|\tr_{\mathrm{tan}}K|$, with the same strict area
inequality~\cite{CarrascoMars2010}.
These concern distinct formulations.  We use full enclosing area and
specify the expansion condition in each theorem.

The area comparison uses an elementary point that remains valid for all
these cuts.  In spatial dimension $d$, if $h\ge c g$ as quadratic forms
for some $c>0$, then the restriction to each tangent $(d-1)$-plane gives
\[
 |\Gamma|_h\ge c^{(d-1)/2}|\Gamma|_g
 \quad\hbox{and hence}\quad
 A_*(h)\ge c^{(d-1)/2}A_*(g).
\]
No selection of a minimizing component and no convergence of minimizers
is needed for this inequality.  We use it after constructing a metric
of the form
\[
 \widehat g=e^{4t/(d-2)}\bigl(g+l(t)^2\dd f^2\bigr),
 \qquad t\ge-\varepsilon.
\]
Here $f$ is the graph height, $t$ is the conformal unknown, and $l>0$
is the warping factor.
The graph term is nonnegative on every tangent plane.  Thus the lower
bound for $t$ protects the full enclosing area before the graph is
known to have bounded slope.

\subsection{The constructions and their relationship}

The first two constructions start with strict dominant-energy data,
a strictly future-trapped boundary, and a prepared end with compactly
supported $K$. The third keeps a maximal tensor tail. The following
map separates these geometric outputs from their original-data
applications; the cited statements give the full hypotheses.
Here $\omega_3=|S^3|=2\pi^2$.

\begin{center}
\small
\renewcommand{\arraystretch}{1.2}
\begin{tabular}{@{}p{.25\linewidth}p{.30\linewidth}p{.37\linewidth}@{}}
\toprule
Construction & Geometric output & Original-data consequence \\
\midrule
Prepared dimension three; compact $K$
& Theorem~\ref{n3:prepared:output}
& $m_{\mathrm{ADM}}\ge\sqrt{A_*/(16\pi)}$ for weak decay $q>1/2$;
Theorem~\ref{n3:intro:exterior-inequality} \\[3pt]
Prepared dimension four; compact $K$
& Theorem~\ref{four:thm:boundary-deformation}
& $m_{\mathrm{ADM}}\ge\frac12(A_*/\omega_3)^{2/3}$ for weak decay $q>1$;
Corollary~\ref{b:four-weak} \\[3pt]
Direct maximal dimension four; decaying $K$
& Proposition~\ref{four:maximal:prop:geometric-output}
& $E\ge\frac12(|S|_g/\omega_3)^{2/3}$ in the maximal-vacuum application;
Theorem~\ref{four:maximal:thm:main} \\
\bottomrule
\end{tabular}
\end{center}

In the first two applications, weak decay means
$g-\delta=O_2(r^{-q})$ and $K=O_1(r^{-1-q})$.
The exteriors satisfy the dominant energy condition, integrable
constraint densities, finite timelike ADM charges $E>|P|$, and weak
future trapping. Their separate rest preparations control all enclosing
cuts. In dimension three, the prepared energies tend to
$m_{\mathrm{ADM}}$ and the full enclosing-area infima converge. In
dimension four, the distant-end replacement supplies a one-sided lower
area bound; strictification is removed at each fixed repaired end before
the replacement radius tends to infinity. These inputs are
\cite[Proposition~7.13 and Proposition~8.8/Lemma~8.9, respectively]{NumericalCompanion}.

The direct maximal theorem has its own approximation and retains the
tensor tail. Its stated application has $q>1$, ball-exterior topology,
a connected spherical marginally outer trapped boundary ($\theta_+=0$),
a torus action, $P=0$, $E>0$,
outermostness, and outer area minimization against all enclosures.
Here $P=0$ identifies energy with invariant mass, and outer area
minimization identifies $A_*$ with $|S|_g$. The strict geometric output
in the middle column has its own hypotheses; the theorem records this
particular original-data application. The broader numerical comparison
in Corollary~\ref{b:four-weak} uses the distinct end-replacement route.

The difference at infinity is concrete.  With compactly supported $K$,
the distant trace equation is a homogeneous graph equation with a
positive height term.  When $\tr_gK=0$, the retained tensor contributes a
term quadratic in $\nabla f$, which becomes a drift along the solution.
The equation remains homogeneous and still forces exponential graph
decay at fixed deformation parameter.  The scalar equation behaves
differently: its zero-graph limit retains $|K|^2/2$.
Section~\ref{b:transition} derives this term and the tail weights needed
to control it.  That calculation identifies both the shared argument
and the additional preparation, threshold geometry and estimates required
for the maximal route.

\subsection{Why the inner boundary matters}

The scalar-curvature identity separates the original constraint densities,
a nonnegative quadratic expression, and the divergence of a vector field
$V$.  Prescribing this divergence produces the required curvature sign.
At infinity, the integral of $V$ measures the ADM energy change.  A
boundary condition for $V$ also prescribes the mean curvature of the
deformed inner boundary.  The divergence theorem therefore links the
two ends of the construction.

The useful boundaries are threshold trapped-region frontiers.  For
general $K$, two signs of the tensor produce black and white faces,
whose names distinguish their expansion equations and the corresponding
Dirichlet heights of the graph.  The collars outside these faces supply
geometric barriers for the graph before uniform ellipticity is available.
They first confine its scaled height to the required interval and then
force a signed normal slope.  That slope makes the possible negative
inner flux exponentially small relative to the positive bulk terms.
The all-black maximal construction has its own threshold argument,
using maximality in place of the two-sided trace budget.

This order is essential.  Uniform ellipticity is an outcome of the
geometric height, floor and slope estimates; it is not assumed to prove
them.  Once the coefficients lie in a bounded elliptic range, local
scalar estimates give the boundary regularity needed for a compact
continuation map.  The four-dimensional proof retains the critical
interpolation step and the interface measures produced by reflection.
Each application verifies its own boundary conditions and completes
its own degree and bootstrap argument.

There are two kinds of constants throughout.  Polynomial bounds in the
large deformation parameter $N$ are used in the final flux comparison.
Classical regularity constants may depend arbitrarily on fixed $N$.
One first exhausts the outer radius with $N$ fixed, then lets $N$ tend
to infinity using only the polynomial estimates.  This distinction
prevents a local compactness argument from being mistaken for a
uniform estimate in the final mass limit.

\subsection{Relation to earlier methods}

Penrose's proposed inequality arose from the relation between gravitational
collapse, horizons and total energy~\cite{Penrose1973}.  In the
time-symmetric setting the second fundamental form vanishes and the
dominant-energy condition becomes nonnegative scalar curvature.
Geroch's smooth-flow idea and the conditional argument of Jang--Wald
preceded the weak-flow theory~\cite{Geroch1973,JangWald1977}; see also
the historical discussion in~\cite[p.~359]{HuiskenIlmanen2001}.
Huisken--Ilmanen's weak inverse-mean-curvature flow proved the
three-dimensional Riemannian inequality for each connected component
of the horizon~\cite[Main Theorem, p.~356]{HuiskenIlmanen2001};
Bray's conformal flow gives the bound for the total area of a disconnected
horizon~\cite[Theorems~1 and~19]{Bray2001}.  Bray--Lee extended the numerical
conformal-flow theorem to dimensions below eight
\cite[Theorem~1.4]{BrayLee2009}.  The final comparison
here uses the numerical clauses of these theorems after constructing
their complete Riemannian inputs.

The graph strategy comes from Jang's equation and the Schoen--Yau
positive-energy argument~\cite{Jang1978,SchoenYau1981}. The earlier
Schoen--Yau proof establishes the positive mass theorem in the Riemannian
and maximal-data cases by minimal surfaces~\cite{SchoenYau1979}.  A classical Jang graph can develop
a cylindrical end at an apparent horizon; Metzger analyses this blow-up
and its exponential cylindrical convergence near a strictly stable
component~\cite{Metzger2010}.
Conformal normalization of such a graph can lose the horizon-area
comparison.  Bray and Khuri developed the warped-graph approach to
address this obstruction~\cite{BrayKhuri2010,BrayKhuri2011}; their 2010
work also proves a sharp spherical case under its stated hypotheses.
The general warped scalar identity already retains a trace discrepancy
before the generalized Jang equation is imposed
\cite[Identity~9, p.~580]{BrayKhuri2011}.

Han--Khuri prove existence and quantitative boundary blow-up for
prescribed warping factors~\cite{HanKhuri2013}.  Their later coupling
to conformal flow gives a reduction contingent on solving the coupled
system with the required boundary behavior~\cite{HanKhuri2018}.
The direct conclusion of that argument concerns ADM energy, as its
Remark~2 explains.  Here the coupled equations prescribe a nonzero
divergence source, a trace depending increasingly on graph height, and
a penalty near the conformal lower bound.  The boundary analysis proves
global existence for these equations and supplies the area and energy
comparisons needed for the subsequent limits.

This distinction concerns the equations, not only their interpretation.
Jaracz proves radial nonexistence for a specified generalized-Jang and
zero-divergence system~\cite{Jaracz2023}; that result does not apply to
the nonzero-source system studied here. His spacetime Poisson method
also provides a predecessor for conformal strictification
\cite{Jaracz2025StrictDEC}. The maximal construction below solves its
own problem with nonzero inner Neumann data and proves the energy and
full enclosing-area comparisons needed to remove the strict
approximation.

Trapped-region boundaries and their stability enter through the work of
Andersson--Metzger and Andersson--Eichmair--Metzger
\cite{AnderssonMetzger2009,AnderssonEichmairMetzger2011}.
Eichmair's Perron and almost-minimizing constructions provide the
original higher-dimensional method~\cite{Eichmair2010}; the precise
total-region theorem used below is the one recorded by
Andersson--Eichmair--Metzger.  For the stability operator of marginally
outer trapped surfaces (MOTS) and its principal-eigenfunction
characterization of stability we use
Andersson--Mars--Simon~\cite{AnderssonMarsSimon2008}.  The local arguments
below explain the additional threshold continuity, support transport
and collar choices used by the deformation.  The analytic part combines
scalar elliptic estimates with the Leray--Schauder degree principle
\cite{LeraySchauder1934}; the particular boundary homotopies and their
admissible open sets are constructed explicitly.

Recent spacetime results also distinguish the relevant scopes.
Khuri--Kunduri prove the sharp inequality for cohomogeneity-one
$\mathrm{SU}(n+1)$-invariant data in spatial dimensions $2(n+1)\ge4$,
with their outermost future-or-past horizon condition
\cite[Theorem~1.1]{KhuriKunduri2025}.  The direct maximal theorem here
has different symmetry and outer-area-minimization hypotheses.
Allen--Bryden--Kazaras--Khuri prove a three-dimensional invariant-mass
bound by a universal suboptimal constant times the square root of
minimum enclosing area, with additional trace decay and a specified
horizon selection~\cite[Theorem~1.1 and Remark~1.2]{AllenBrydenKazarasKhuri2025}.

Two companion manuscripts supply specified inputs: scalar algebra and
weak-end preparations from~\cite{NumericalCompanion}, and full-perimeter
enclosure geometry from~\cite{RiemannianCompanion}. Their exact hypotheses
and normalization are checked at each use. The numerical comparisons
at the end of the boundary constructions use Bray and Bray--Lee directly.

\subsection{Reading the proof}

The three routes retain their own geometry, equations, global bounds and
end comparison. Shared arguments are proved at their first use in the
three-dimensional route: threshold continuity, stable collars and exact
support transport in Section~\ref{n3:domains:section}; the passage from
lower tests to sublevel supports in Lemma~\ref{n3:apriori:sublevel-support};
and scalar measure-source regularity and degree on varying open sets in
Lemmas~\ref{b:lem:natural-growth} and~\ref{b:lem:moving-degree}.
Each later use verifies its tensor, coefficients and boundary conditions.

For dimension three, Sections~\ref{n3:deformation:section}--\ref{n3:limit:section}
proceed from the coupled equations and conformal floor through geometric
bounds, boundary existence, and the complete-end comparison. Dimension
four follows the same order in
Sections~\ref{four:sec:geometry}--\ref{four:sec:existence}, with its own
compact total-region construction and the critical interpolation in
Theorem~\ref{four:thm:axial-regularity}. Corollary~\ref{b:four-weak}
then gives its separate weak-decay transfer.

Section~\ref{b:transition} explains the maximal tensor tail before
Sections~\ref{four:maximal:sec:preparation}--\ref{four:maximal:sec:comparison}
construct the direct route. Its pointwise scalar and boundary identities
come from Section~\ref{four:sec:system}, its axial estimate from
Theorem~\ref{four:thm:axial-regularity}, and its enclosure input from
Proposition~\ref{four:input:enclosure}. The maximal part supplies its own
strict approximation, noncompact drift bounds, scalar solve, reflected
boundary checks and final flux estimate.

\section{A boundary deformation in three dimensions}\label{n3:intro:section}

The aim of this part is to replace an initial-data exterior by a
Riemannian exterior while keeping a genuine inner boundary. Three
quantities have to be controlled together: scalar curvature, the area
of every enclosing cut, and the ADM energy. The inner boundary creates
a flux term in the energy balance. We choose its geometry and the
boundary conditions so that this term is absorbed by a positive bulk
integral.

\subsection{Initial data, normals, and enclosing area}

We work in three spatial dimensions, with zero cosmological constant and
\(G=c=1\). For a Riemannian metric \(g\) and a symmetric covariant
two-tensor \(K\), define the constraint densities by
\begin{equation}\label{n3:intro:constraints}
 16\pi\mu=R_g+(\tr_gK)^2-|K|_g^2,\qquad
 8\pi J=\Div_g\bigl(K-(\tr_gK)g\bigr).
\end{equation}
Here \(J\) is a covector field. The dominant energy condition is
\begin{equation}\label{n3:intro:dec}
                         \mu\ge |J|_g.
\end{equation}
All initial data considered here are smooth, including at a compact
boundary when one is present.

On an asymptotically flat end with Euclidean coordinates \(x\) and
\(r=|x|\), our basic assumptions are
\begin{equation}\label{n3:intro:decay}
 g_{ij}-\delta_{ij}=O_2(r^{-q}),\qquad
 K_{ij}=O_1(r^{-1-q}),\qquad q>\tfrac12.
\end{equation}
The notation \(O_j(r^{-a})\) includes the corresponding coordinate
derivative bounds through order \(j\). We assume that \(\mu\) and
\(|J|_g\) are integrable and that the following ADM limits exist and are
finite:
\begin{align}
 E&=\frac1{16\pi}\lim_{r\to\infty}
   \int_{S_r}(\partial_jg_{ij}-\partial_i g_{jj})n^i
                    \dd A_\delta,\label{n3:intro:energy}\\
 P_i&=\frac1{8\pi}\lim_{r\to\infty}
   \int_{S_r}\bigl(K_{ij}-(\tr_gK)g_{ij}\bigr)n^j
                    \dd A_\delta.\label{n3:intro:momentum}
\end{align}
The normal and area form in these definitions are Euclidean. Whenever
\(E>|P|_\delta\), write
\begin{equation}\label{n3:intro:rest-mass}
                         m=\sqrt{E^2-|P|_\delta^2}.
\end{equation}

For a two-sided surface \(\Sigma\) with specified unit normal \(\nu\),
we use
\[
 H_\Sigma=\Div_\Sigma\nu,\qquad
 \theta_+(\Sigma)=H_\Sigma+\tr_\Sigma K.
\]
Thus Euclidean spheres have positive mean curvature for the normal
toward infinity. A future marginally outer trapped surface, abbreviated
MOTS, satisfies \(\theta_+=0\). A weakly future outer trapped surface
satisfies \(\theta_+\le0\). Our spacetime sign convention is
\begin{equation}\label{n3:intro:k-sign}
 K(Y,Z)=\mathbf g(\mathbf\nabla_Y n,Z),
\end{equation}
where \(n\) is the future unit timelike normal. Equivalently, in Gaussian
normal coordinates the spatial metric has normal derivative \(2K\).

\begin{definition}[Exterior and enclosing area]\label{n3:intro:exterior-class}
An exterior is a connected orientable smooth three-manifold \(\Omega\)
with nonempty compact smooth boundary, complete as a metric space with
its boundary included, and with one asymptotically flat coordinate end
whose complement is compact.
An enclosing cut is the entire compact smooth embedded two-sided
intrinsic boundary of a connected smooth end-containing domain $D$
closed in $\overline\Omega$, whose intrinsic interior lies in
$\operatorname{int}\Omega$. It separates the entire inner boundary
from the distant end. The inner side includes the original obstacle
and every bounded complementary pocket. All components of the cut
and every portion coincident with the inner boundary are counted.
Its normal points into $D$, toward the end. We put
\[
 a_g(\partial\Omega)=
  \inf\{|\Gamma|_g:\Gamma\text{ is an enclosing cut}\}.
\]
When using perimeter compactness, we take the closure of this class
among filled inner sets and include the area of any frontier coinciding
with the obstacle. Smooth outward approximation gives the same
infimum. Enclosing cuts and minimizing sets are taken with bounded
complementary pockets filled. Auxiliary perimeter arguments may use
competitors containing the inner obstacle before filling: filling a bounded
pocket deletes its frontier, cannot increase perimeter, and leaves the enclosing
infimum unchanged.
\end{definition}

The definition permits coincidence because first variations of the
minimum need not be realized by a cut lying strictly outside the
obstacle. It permits disconnected cuts because neither the minimizing
hull nor an intermediate trapped boundary need be connected. The smooth-cut and full-perimeter formulations have the same infimum;
we use the independent full-enclosure geometry of~\cite[Lemma~2.1]{RiemannianCompanion}.

\subsection{Prepared data and the geometric output}

The construction itself starts from the following explicit geometric
conditions. Its hypotheses do not require that the data have been
obtained by a particular end-replacement procedure.

\begin{definition}[Prepared three-dimensional data]\label{n3:prepared:data}
Let $(\Omega_0,g,K)$ be an exterior in
Definition~\ref{n3:intro:exterior-class}, with inner boundary $S_0$.
Suppose $K$ is compactly supported and, on its single coordinate end,
$g-\delta=O_j(r^{-1})$ for every fixed $j\ge0$. Assume the physical
constraint densities in Equation~\eqref{n3:intro:constraints} are
integrable, the ADM energy is $E>0$, and the ADM momentum is zero.
There are $0<\delta<1$, $c>0$, and a smooth radius extension
$\varrho\ge1$, equal to $r$ sufficiently far out, such that
\[
 R_g=O(r^{-3-\delta}),\qquad
 8\pi(\mu-|J|_g)\ge c\varrho^{-3-\delta}.
\]
The boundary is strictly future trapped with its normal into the
exterior: $\theta_+(S_0)<0$ on every component. Its full enclosing-area
infimum is denoted by $A_*=a_g(S_0)>0$.
\end{definition}

\begin{theorem}[Three-dimensional boundary deformation]
\label{n3:prepared:output}
For every prepared exterior $(\Omega_0,g,K)$ there is a connected
one-ended subexterior $\Omega\subset\Omega_0$ whose smooth compact
boundary $B$ separates $S_0$ from infinity, with the following property.
For each sufficiently small fixed $\epsilon>0$ and every sufficiently
large integer $N$, there is a smooth metric $\widehat g_{\epsilon,N}$
on $\overline\Omega$, complete with its boundary included, for which
\begin{align}
 R_{\widehat g_{\epsilon,N}}&\ge0,
 &H_{\partial\Omega}(\widehat g_{\epsilon,N})&<0,
 \label{n3:prepared:curvature}\\
 \widehat g_{\epsilon,N}&\ge e^{-4\epsilon}g,
 &E_{\widehat g_{\epsilon,N}}&\le E+\eta_{\epsilon,N},
 \label{n3:prepared:comparison}
\end{align}
where the boundary normal points into $\Omega$ and, at fixed
$\epsilon$, $\eta_{\epsilon,N}\to0$. More precisely, for constants
$C,a$ depending on the fixed data and $\epsilon$,
\[
 0\le\eta_{\epsilon,N}\le
 C(1+N)^a\bigl(e^{-\epsilon N}+e^{-\epsilon N/2}\bigr).
\]
The new end has $\widehat g_{\epsilon,N}-\delta=O_2(r^{-1})$ and
$R_{\widehat g_{\epsilon,N}}=O(r^{-3-\delta'})$ for some
$\delta'>0$. Every full enclosing cut $T$ in $\Omega$ satisfies
\[
 |T|_{\widehat g_{\epsilon,N}}\ge e^{-4\epsilon}A_*.
\]
\end{theorem}

The construction begins by selecting maximal black and white trapped
regions. Their outer collars prescribe opposite signs for the graph
slope, but the same negative mean-curvature condition for the final
metric. Section~\ref{n3:domains:section} proves the required geometry,
including the passage from weak supports to smooth enclosures.
Section~\ref{n3:deformation:section} specifies the trace and divergence
equations and excludes a lower conformal barrier. The estimates of
Section~\ref{n3:apriori:section} then control the solutions before
uniform ellipticity is available. Section~\ref{n3:existence:section}
supplies the boundary solve and degree argument. Finally,
Section~\ref{n3:limit:section} computes the end charge and proves
Theorem~\ref{n3:prepared:output}.

The algebraic inputs from~\cite[Propositions~4.1--4.2 and Lemma~4.3]{NumericalCompanion} are its scalar
identity, weighted coercivity, and the pointwise algebraic clause of
its floor lemma. Their
proofs are independent of its numerical conclusion. The boundary
construction here retains the signed-collar estimates, boundary flux,
homotopy, and degree that a filled construction does not require.

\subsection{Consequence for weakly decaying original data}

\begin{theorem}[Minimum-enclosing-area inequality]\label{n3:intro:exterior-inequality}
Let \((\Omega,g,K)\) be an exterior as in
Definition~\ref{n3:intro:exterior-class}. Assume
Equations~\eqref{n3:intro:constraints}--\eqref{n3:intro:decay}, integrability
of \(\mu\) and \(|J|_g\), and the finite ADM limits
\eqref{n3:intro:energy}--\eqref{n3:intro:momentum}. Orient \(\partial\Omega\)
by the normal into \(\Omega\), and suppose
\(\theta_+(\partial\Omega)\le0\). If \(E>|P|_\delta\), then
\begin{equation}\label{n3:intro:penrose}
          \sqrt{E^2-|P|_\delta^2}
               \ge \sqrt{\frac{a_g(\partial\Omega)}{16\pi}}.
\end{equation}
No extension of the data across \(\partial\Omega\) is required.
\end{theorem}


The last assertion will follow from the boundary construction and the
weak three-dimensional preparation theorem of~\cite[Proposition~7.13]{NumericalCompanion}.
That preparation supplies compact tensor support, strict trapping and
the weighted margin in Definition~\ref{n3:prepared:data}; its energies
converge to $\sqrt{E^2-|P|^2}$ and its full enclosing areas converge
to the original infimum. We will use only this preparation theorem,
not the companion's numerical inequality. In particular, the ADM
normalizations in Equations~\eqref{n3:intro:energy} and
\eqref{n3:intro:momentum} remain unchanged.

\section{Trapped domains, collars, and weak supports}
\label{n3:domains:section}

We need a boundary on which the graph height can be constant and its
slope has a definite sign. A future trapped face supplies one sign; a
past trapped face supplies the other. We first choose maximal regions
of these two types, then construct outward collars and compatible
height offsets. The original obstacle stays on their excluded side,
so every resulting enclosing cut still controls its full area.
Throughout, the data satisfy
Definition~\ref{n3:prepared:data}.  Write $S_0$ for their inner
boundary.  Thus $S_0$ is strictly future outer trapped, $K$ has compact
support, and the single end has the all-order $O(r^{-1})$ bounds stipulated
in that definition.  In particular, sufficiently large coordinate spheres
have positive outward mean curvature.  The function
\begin{equation}\label{n3:domains:margin}
 \mathfrak m(x)=8\pi\bigl(\mu(x)-|J(x)|_g\bigr)
\end{equation}
is positive, has a positive lower bound on every compact set, and has
the positive $r^{-3-\delta}$ lower bound on the end, with $0<\delta<1$.  Let $A_*$ denote the prepared data's
filled enclosing perimeter infimum.

\subsection{The bounding-region convention}

For a symmetric tensor $Q$ and a two-sided surface $\Sigma$ with a
specified unit normal $\nu$, put
\begin{equation}\label{n3:domains:expansion}
 \Theta_Q(\Sigma,\nu)
   =H_\Sigma(\nu)+\tr_\Sigma Q.
\end{equation}
For the boundary of a domain the specified normal always points
\emph{out of that domain}.  Notice both elementary identities
\begin{equation}\label{n3:domains:shift-reversal}
 \Theta_{Q-(c/2)g}=\Theta_Q-c,
 \qquad
 \Theta_{-Q}(\Sigma,-\nu)=-\Theta_Q(\Sigma,\nu).
\end{equation}
The factor $1/2$ in the first formula is the reciprocal of the surface
dimension.

We use the following geometric theorem.  Its boundary parts may be
disconnected.

\begin{theorem}[MOTS barriers and the total trapped region]
\label{n3:domains:region-theorem}
Let $(N,g,Q)$ be a smooth compact three-dimensional initial data set
with boundary $\Gamma_-\sqcup\Gamma_+$.  Assume $\Gamma_+$ is nonempty
and has positive expansion with its normal out of $N$.  The inner part
$\Gamma_-$ may be empty; when it is present, assume its expansion is
negative with its normal into $N$.

Consider all smooth domains whose inner boundary is $\Gamma_-$ and
whose remaining boundary $\Sigma$ lies in the interior of $N$ and
satisfies $\Theta_Q(\Sigma)\leq0$.  They avoid $\Gamma_+$, and their
normal on $\Sigma$ points out of the domain.  If their union is
nonempty, its closure has a smooth compact embedded stable MOTS
frontier in the interior, separating it from $\Gamma_+$.  This closure
is itself the closure of such a domain and contains every domain in
the defining class.  If $\Gamma_-$ is nonempty, the two strict barriers
also give a nonempty smooth embedded enclosing MOTS.
\end{theorem}

The total-region assertion is
\cite[Definitions 7.1--7.2 and Theorem 7.3]{AnderssonMetzger2009}; the
barrier assertion is Theorem 3.1 there, with stability furnished by
Theorem 4.1.  The preliminary conventions in Section 2 of that paper
allow an empty inner boundary.  That paper uses open regions;
taking their closures preserves the smooth frontier and its expansion,
giving the closed-region formulation above.  These results require no energy
condition on $Q$.  We therefore apply them to the shifted tensors in
Equation~\eqref{n3:domains:shift-reversal}, even though these tensors
need not satisfy the dominant energy condition.  A component of the
complement of a smooth trapped domain which meets no designated outer
boundary may be filled: its entire boundary is then deleted and the
remaining boundary has the same expansion.  The maximal region is
therefore filled in this relative sense.

We use compact inner fillings only to express enclosing competitors.
One may attach a smooth collar behind $S_0$, extend the one-sided
metric smoothly, and complete the inner side by a fixed compact
filling.  All competitors contain that filling and only their
boundaries in the original exterior are counted.  This introduces no
geometric condition behind $S_0$.  For perimeter minimization with
this smooth obstacle, the relevant regularity input is
\cite[Regularity Theorem 1.3(iii)]{HuiskenIlmanen2001}: the minimizing
boundary is $C^{1,1}$ and is smooth away from contact.  We retain the filled enclosing class and its smooth approximation
from the independent enclosure geometry of~\cite[Lemmas~2.1--2.2]{RiemannianCompanion}.  The trapped regions below are not
assumed to minimize perimeter.

\subsection{Choosing the two maximal regions}

Fix a large sphere $S_{R_0}$ outside $\supp K$ such that all coordinate
spheres with radius at least $R_0$ have positive outward mean
curvature.  A compact smooth domain with $\Theta_{\pm K}\leq c\leq0$
cannot reach this part of the end.  Indeed, at a point where its
boundary maximizes the coordinate radius, it lies inside the tangent
coordinate sphere and has the same outward normal.  The graph
second-derivative comparison gives
$H_{\partial D}\geq H_{S_r}>0$ there; also $K=0$ there.  This
contradicts its assumed expansion.  The same argument applies to an
enclosing domain with prescribed inner obstacle, since the maximum
radius is attained on its free boundary.  Consequently all the
negative-threshold regions used below lie in one fixed compact radial
range.  They may be constructed using any sufficiently distant
sphere as the outer barrier, without dependence on that choice.

For
\begin{equation}\label{n3:domains:black-interval}
 \max_{S_0}\Theta_K(S_0)<c<0,
\end{equation}
let $\mathcal B_c$ be the closed total trapped region for the
condition $\Theta_K\leq c$, with $S_0$ as the required inner
boundary and with its inner filling understood.  It is nonempty:
a sufficiently short outward collar of $S_0$ supplies a strict seed.
Theorem~\ref{n3:domains:region-theorem}, applied to $K-(c/2)g$, gives
its smooth stable frontier with expansion exactly $c$.  The regions
$\mathcal B_c$ increase with $c$.  Their complement toward the end
has no bounded component, since such a component could be filled to
enlarge the total trapped region.

We will choose $c_b<0$ close to zero and set
$\mathcal B=\mathcal B_{c_b}$.  In the complement of $\mathcal B$,
use the full black frontier and a distant coordinate sphere as
\emph{outer} barrier faces.  With the normal out of this complement,
the black frontier has expansion for $-K$ equal to $-c_b>0$.
For $c<0$ close to zero, let $\mathcal W_c$ be the closed total
trapped region of compact smooth domains in that complement with
\begin{equation}\label{n3:domains:white-condition}
 \Theta_{-K}=H-\tr_\Sigma K\leq c.
\end{equation}
There is no required inner boundary in this application of
Theorem~\ref{n3:domains:region-theorem}.  The shifted expansion of a
reversed black face is $-c_b-c>0$, and that of a distant sphere is
also positive.  Thus the theorem applies whenever the defining union
is nonempty.  When the union is empty we simply put
$\mathcal W_c=\varnothing$.  When nonempty, it has smooth stable
frontier of white expansion $c$ and lies a positive distance from all
black faces.  The family $\mathcal W_c$ is increasing, with the black
region held fixed.

To control limits of the deformation, we choose thresholds at which these
regions are right-continuous.  The following lemma includes empty regions.

\begin{lemma}[Right-continuous thresholds]\label{n3:domains:right-continuity}
Let $E_c$, $c$ in an interval, be an increasing family of closed
subsets of a fixed compact metric space.  Outside a countable set of
parameters,
\begin{equation}\label{n3:domains:hausdorff-right}
 E_{c_j}\longrightarrow E_c\quad\hbox{in Hausdorff distance whenever}
 \quad c_j\downarrow c.
\end{equation}
Here continuity at $E_c=\varnothing$ means that $E_{c'}$ is empty
for every sufficiently close $c'>c$.
\end{lemma}

\begin{proof}
Choose a countable dense set $\{x_i\}$ and a number $L$ larger than
the diameter of the compact space.  Set
$d_c(x)=\min\{L,\operatorname{dist}(x,E_c)\}$, with the value $L$
when $E_c$ is empty.  For each $i$, $c\mapsto d_c(x_i)$ is a
bounded nonincreasing real function, so its right discontinuities
form a countable set.  Avoid their countable union.  Then
$d_{c_j}(x_i)\to d_c(x_i)$ for all $i$ whenever $c_j\downarrow c$.
Every $d_c$ is $1$-Lipschitz.  A finite net chosen from the dense set
therefore upgrades this convergence to uniform convergence on the
compact space.

For nonempty $E_c$, if points $y_j\in E_{c_j}$ stayed a fixed
distance from $E_c$, uniform convergence evaluated at $y_j$ would
contradict $d_{c_j}(y_j)=0$.  The opposite Hausdorff inclusion follows
from $E_c\subset E_{c_j}$.  For empty $E_c$, uniform convergence to
$L$ is incompatible with any nonempty $E_{c_j}$, since its distance
function has a zero.  This proves the stated convention as well.
\end{proof}

Choose $c_b$ at such a right-continuity point of the black family,
then choose $c_w<0$ at such a point of the white family for this
fixed black region.  Both choices can be made arbitrarily close to
zero.  Write $\mathcal W=\mathcal W_{c_w}$, and let $\Omega$ be the
closure of the component of the complement of
$\mathcal B\cup\mathcal W$ which reaches the selected end.  Its
boundary $B$ is a finite disjoint union of complete components of
the two smooth frontiers.  Write
\begin{equation}\label{n3:domains:faces}
 B=B_b\sqcup B_w,\qquad
 H=H_B(\nu),\qquad P_B=\tr_BK,
\end{equation}
where $\nu$ points into $\Omega$.  Thus
\begin{equation}\label{n3:domains:face-expansions}
 H+P_B=c_b\quad\hbox{on }B_b,
 \qquad
 H-P_B=c_w\quad\hbox{on }B_w.
\end{equation}
The set $B_w$ may be empty, and $B_b$ need not contain every black
frontier component: a white component may cut a black face off from
the end.  We retain exactly the boundary of the component that reaches
infinity.  The compact, disjoint frontiers create no intersection corners.

The excluded side of $B$ contains $S_0$ and a collar of it.  Hence
every compact smooth enclosing cut $T$ in $\Omega$, with all its
components retained, is an enclosing competitor for the original
filled obstacle after adjoining the discarded regions to its inner
side.  It follows that
\begin{equation}\label{n3:domains:inherited-area}
 |T|_g\geq A_*.
\end{equation}
The same conclusion holds for perimeter competitors.  This is a
statement about inclusion of competitor classes; no minimizing
property of $B$ is used.

Every enclosing cut now carries the area lower bound $A_*$. The next
step is local: on each retained face we need a smooth outward family
whose expansion is strictly larger than its threshold. Stability
alone allows a zero principal eigenvalue, so that case requires the
maximality of the chosen region.

\subsection{Outward collars, including zero stability eigenvalue}

\begin{lemma}[Collars of a maximal frontier]\label{n3:domains:stable-collar}
Let $(N^d,g)$ be smooth, with $d\in\{3,4\}$ and a fixed smooth
symmetric tensor $Q$.  Let $\Sigma$ be a compact connected smooth
two-sided hypersurface without boundary, forming a component of
the frontier of a closed region $D$, with normal pointing out of
$D$.  Assume $\Theta_Q(\Sigma)=c$ and
stability for outward variations.  Assume also local
\emph{outward-replacement maximality}: in some collar of $\Sigma$
disjoint from the other frontier components and all designated
barrier faces, no positive outward normal graph can replace
$\Sigma$, leaving $D$ unchanged outside that collar, and have
expansion at most $c$ everywhere on the new component.
There is a smooth foliation $\Sigma_s$, $0\leq s<s_0$, of an outward
collar, with $\Sigma_0=\Sigma$ and positive outward normal velocity,
such that
\begin{equation}\label{n3:domains:collar-expansion}
 \Theta_Q(\Sigma_s)>c\quad(0<s<s_0),\qquad
 \Theta_Q(\Sigma_0)=c.
\end{equation}
In particular its leaf parameter $s$ is a smooth defining function
and $|\nabla s|$ is bounded above and below by positive constants.
\end{lemma}

\begin{proof}
Use normal graphs $\Sigma(v)$ over $\Sigma$ and pull their expansion
back to $\Sigma$.  The map
\[
 \Phi:C^{2,\alpha}(\Sigma)\longrightarrow C^{0,\alpha}(\Sigma),
 \qquad
 \Phi(v)=\Theta_Q(\Sigma(v))-c
\]
is smooth near zero.  Let $\nu$ be the outward normal and
$\mathrm{II}$ its second fundamental form.  The normal-variation
formulas give the linearization $L=D\Phi(0)$ explicitly as
\[
 Lv=-\Delta_\Sigma v+2Q(\nu,\grad_\Sigma v)
   +\bigl(\tr_\Sigma\nabla_\nu Q-|\mathrm{II}|^2
                          -\Ric(\nu,\nu)\bigr)v.
\]
The hypersurface dimension is $d-1$, so
$\Theta_{Q-cg/(d-1)}=\Theta_Q-c$.  This constant shift has zero
derivative under hypersurface variations; hence $L$ is also the
MOTS stability operator for $Q-cg/(d-1)$, with principal part
$-\Delta_\Sigma$.  Stability gives a principal eigenvalue
$\lambda\geq0$ and a positive smooth eigenfunction $\varphi$
\cite[Proposition~5.1]{AnderssonMarsSimon2008}.
If $\lambda>0$, take $v=s\varphi$.  Smooth dependence gives
$\Phi(s\varphi)=s\lambda\varphi+O(s^2)$ in $C^0$, which is
positive for all sufficiently small $s>0$.  The graph velocity is
positive as well.

Suppose $\lambda=0$.  The kernel of $L$ is spanned by $\varphi$,
and the adjoint kernel is spanned by a positive smooth function
$\varphi^*$.  These are the principal-eigenfunction facts for a
scalar elliptic operator on a compact connected closed manifold
\cite[Section~4, Lemma~4.1]{AnderssonMarsSimon2008}.  For
example, simplicity of the kernel also follows by writing a kernel
element as $v=\varphi w$ and applying the strong maximum principle
to the resulting equation for $w$, whose zeroth-order term vanishes.
Consider the augmented map
\[
 \mathcal G(v,a)
   =\left(\Phi(v)-a,\ \int_\Sigma v\,\dd A\right).
\]
Its derivative at $(0,0)$ is
\begin{equation}\label{n3:domains:augmented-linearization}
 (v,a)\longmapsto
 \left(Lv-a,\ \int_\Sigma v\,\dd A\right).
\end{equation}
This map is an isomorphism.  Indeed, to solve $Lv-a=f$ with
$\int v=b$, the Fredholm compatibility condition uniquely fixes
\[
 a=-\frac{\int_\Sigma\varphi^*f\,\dd A}
          {\int_\Sigma\varphi^*\,\dd A}.
\]
The equation for $v$ then has a solution; adding a multiple of
$\varphi$ imposes its prescribed integral uniquely.  Elliptic
estimates and this one-dimensional normalization give a bounded
inverse between the indicated H\"older spaces.

The implicit-function theorem consequently supplies smooth
$v(s),a(s)$ with $\mathcal G(v(s),a(s))=(0,s)$.  Differentiating at
zero and pairing with $\varphi^*$ gives $a'(0)=0$ and
\[
 v'(0)=\frac{\varphi}{\int_\Sigma\varphi\,\dd A}>0.
\]
After shortening the interval, all these graphs form an outward
foliation and have constant expansion $c+a(s)$.  If $a(s)\leq0$
at any positive $s$ in this interval, replace only this frontier
component by $\Sigma(v(s))$.  This graph strictly enlarges $D$
inside the chosen collar, leaves every other frontier component
and designated barrier fixed, and has expansion at most $c$ on
the replaced component.  This contradicts the assumed local
outward-replacement maximality.  Therefore $a(s)>0$ for every such
$s>0$.

In either case positive graph velocity and compactness give the
asserted bounds on $|\nabla s|$.  Smooth elliptic bootstrapping of
the graph equation gives smooth leaves in the zero-eigenvalue case.
\end{proof}

In the present application $d=3$.  The full black and white regions
are $\mathcal B_{c_b}$ and $\mathcal W_{c_w}$, with $Q=K$ and
$Q=-K$, respectively.  Theorem~\ref{n3:domains:region-theorem}
gives their smooth two-sided stable frontiers, whose retained
components have expansion $c_b$ and $c_w$.  Total-region maximality
verifies the local replacement hypothesis: a positive outward
replacement with expansion at most the threshold, leaving all
other frontier components and designated barriers fixed, would
be an admissible strictly larger trapped region.  For white faces
the entire replacement lies in the fixed black complement.
Thus the lemma applies separately to each retained face.
There are finitely many faces, so choose disjoint
collars in $\Omega$ and shorten them to have a common positive
parameter range when convenient.  Write $\nu_s=\nabla s/|\nabla s|$
for their outward leaf normal; it agrees with $\nu$ at $B$.

\subsection{From weak exterior supports to a smooth enclosure}

The sublevel limits used below need not have smooth frontiers.  We now
show how an expansion bound on their exterior supports produces a smooth
enclosure.  The argument applies to arbitrary closed sets, without
positive reach or a curvature bound on the supporting surfaces.

\begin{definition}\label{n3:domains:support-definition}
Let $(N^d,g)$ be a smooth Riemannian manifold, $d\in\{3,4\}$,
with a smooth symmetric tensor $Q$.  Let $D\subset N$ be closed,
and let $x\in\partial D$ be an interior point of $N$.
A smooth exterior support is a smooth function $\phi$ near $x$ with
$\phi(x)=0$, $\dd\phi(x)\ne0$, and
$D\subset\{\phi\leq0\}$ locally near $x$.  Its normal is
$n=\nabla\phi/|\nabla\phi|$.  Define
\begin{equation}\label{n3:domains:test-expansion}
 \mathcal E_Q[\phi](x)
   =\frac{(g^{ij}-n^in^j)\nabla_i\nabla_j\phi}{|\nabla\phi|}
       +\tr Q-Q(n,n).
\end{equation}
We say $D$ has exterior-support expansion at most $c$ if every such
support at every point of its frontier has
$\mathcal E_Q[\phi]\leq c$.
\end{definition}

For a smooth domain this is exactly the usual upper bound on its
outward expansion: at contact an exterior supporting graph has mean
curvature at most that of the enclosed smooth graph.  The definition
is unchanged by increasing smooth reparametrizations of $\phi$,
because the additional Hessian term is a multiple of
$\dd\phi\otimes\dd\phi$ and has zero tangential trace.

\begin{lemma}[Parallel neighborhoods and additive error]
\label{n3:domains:parallel-support}
Let $(N^d,g,Q)$ be smooth, with $d\in\{3,4\}$.  Suppose the
frontier of a closed set $D$ lies in a fixed compact interior
region and has exterior-support expansion at most $c$.  Put
$D_z=\{y:\operatorname{dist}(y,D)\leq z\}$.  For sufficiently small
$z>0$, all exterior supports of $D_z$ have expansion at most
\begin{equation}\label{n3:domains:parallel-estimate}
 c+Cz.
\end{equation}
The allowable radius and $C$ depend only on the compact ambient
geometry and $Q$, and not on the support, its second fundamental
form, or the regularity of $D$.
\end{lemma}

\begin{proof}
If $D=\varnothing$, use $\operatorname{dist}(y,D)=+\infty$;
then $D_z=\varnothing$ and the assertion is vacuous.  Otherwise,
take a support $\phi$ to $D_z$ at $x'$.  Since $x'\in\partial D_z$,
its distance from $D$ is $z$.  For sufficiently small radii the
relevant distance neighborhood of $\partial D$ lies in a fixed
compact interior neighborhood.  There is therefore a nearest point
$x\in D$ and a length-$z$ minimizing unit-speed geodesic
$\gamma:[0,z]\to N$ from $x$ to $x'$.  This nearest point lies on
$\partial D$: if it were interior to $D$, moving a short distance
along $\gamma$ toward $x'$ would decrease the distance.  The entire
segment lies in the fixed smooth compact neighborhood and is
shorter than its injectivity radius.  The ball
$\overline B_z(x)$ is contained in $D_z$ and touches its exterior
support at $x'$.  Its tangent plane there therefore agrees with that
of the support and
\begin{equation}\label{n3:domains:radial-normal}
 \frac{\nabla\phi(x')}{|\nabla\phi(x')|}=\dot\gamma(z).
\end{equation}
Let $P_t:T_xN\to T_{\gamma(t)}N$ denote parallel transport and put
$P=P_z$.

For every $X\in T_xN$ solve the Jacobi boundary problem
\begin{equation}\label{n3:domains:jacobi-endpoints}
 J_X''+\operatorname{Rm}(J_X,\dot\gamma)\dot\gamma=0,
 \qquad J_X(0)=X,\qquad J_X(z)=PX.
\end{equation}
Here primes denote covariant differentiation along $\gamma$.
The short interval and absence of conjugate points make this problem
uniquely solvable.  Its estimates can be seen without division by an
uncontrolled quantity.  In a parallel frame write
$j_X(t)=P_t^{-1}J_X(t)$ and $\mathcal R(t)$ for the curvature
coefficient.  The integral equation is
\[
 j_X(t)=X+tA_X-
       \int_0^t(t-s)\mathcal R(s)j_X(s)\,\dd s,
 \qquad
 A_X=\frac1z\int_0^z(z-s)\mathcal R(s)j_X(s)\,\dd s.
\]
It implies, by absorption for small $z$,
\begin{equation}\label{n3:domains:jacobi-estimates}
 \begin{aligned}
 \sup_{0\leq t\leq z}|j_X(t)|&\leq2|X|,
 \qquad |A_X|\leq Cz|X|,\\
 \sup_{0\leq t\leq z}|j_X(t)-X|&\leq Cz^2|X|.
 \end{aligned}
\end{equation}
The second line follows from the same integral equation and the
first-line bounds.
Also $\langle J_X,\dot\gamma\rangle$ is affine in $t$ and has
equal endpoint values.  Thus
$\langle A_X,\dot\gamma(0)\rangle=0$.

We may therefore prescribe the first jet of a smooth unit vector
field $V$ near $x$ by
\[
 V(x)=\dot\gamma(0),\qquad \nabla_XV(x)=A_X.
\]
To realize it with uniform second-jet bounds, choose a smooth local
orthonormal frame obtained by radial parallel transport from $x$.
In that frame take the affine coefficient vector with the prescribed
value and first derivatives, and normalize its length.  The
orthogonality just proved means normalization preserves its prescribed
first derivatives.  The frame has uniformly bounded derivatives on
a fixed small coordinate ball, and Equation~\eqref{n3:domains:jacobi-estimates}
therefore gives uniformly bounded second derivatives of $V$.
This construction uses the radial direction and the ambient
geometry, with no derivatives of the support $\phi$.

Define the local endpoint map
\[
 F_z(y)=\exp_y\bigl(zV(y)\bigr).
\]
Its differential at $x$ is the terminal value of the Jacobi field
with initial values $X,\nabla_XV(x)$.  By
Equation~\eqref{n3:domains:jacobi-endpoints}, it has the exact property
\begin{equation}\label{n3:domains:exact-endpoint-isometry}
 F_z(x)=x',\qquad (\dd F_z)_x=P,
 \qquad |(\nabla\dd F_z)_x|\leq Cz.
\end{equation}
The last bound follows by differentiating the
smooth map $(y,v,z)\mapsto\exp_y(zv)$ twice in $y$, including the
first and second derivatives of $V$.  At $z=0$ the map is the
identity and its covariant second derivative vanishes.  Its
$z$-derivative is uniformly bounded for the bounded jets just
constructed, giving the displayed $Cz$ bound.  The exact middle
identity, rather than a near-isometry estimate, is essential.

The map $F_z$ takes $D$ locally into $D_z$, since its defining
geodesic from any $y\in D$ has length $z$.  Consequently
$\psi=\phi\circ F_z$ is an exterior support to $D$ at $x$.
Its derivative is $P^*\dd\phi$, so its gradient length is exactly
$|\nabla\phi(x')|$ and, by
Equation~\eqref{n3:domains:radial-normal}, its normal is
$\dot\gamma(0)$.  The covariant chain rule gives
\[
 \Hess\psi(X,Y)
   =\Hess\phi(PX,PY)
        +\dd\phi\bigl((\nabla\dd F_z)(X,Y)\bigr)
 \quad\hbox{at }x.
\]
Trace over an orthonormal basis of $\dot\gamma(0)^\perp$ and divide
by the common gradient length.  The first term is exactly the
corresponding normalized tangential trace at $x'$, and the remaining
term has absolute value at most $Cz$.  Finally, the tensor trace
changes by at most $Cz$, by the $C^1$ bound on $Q$ and parallel
transport.  It follows that
\[
 \bigl|\mathcal E_Q[\psi](x)
             -\mathcal E_Q[\phi](x')\bigr|\leq Cz.
\]
The assumed support inequality at $x$ proves
Equation~\eqref{n3:domains:parallel-estimate}.  There is no term
containing $z|\Hess\phi|$.
\end{proof}

\begin{lemma}[A smooth enclosure from weak supports]
\label{n3:domains:weak-support}
Let $N$ be a smooth compact connected three-manifold with boundary
$\Gamma_-\sqcup\Gamma_+$, where $\Gamma_+$ is nonempty and
$\Gamma_-$ may be empty.  Let $D$ be a nonempty closed subset which
contains a relative collar of $\Gamma_-$ when that boundary is
prescribed, and stays a positive distance from $\Gamma_+$.  Suppose
its interior frontier has exterior-support expansion at most $c$
for a smooth tensor $Q$.

For every $c'>c$ such that
$\Theta_Q(\Gamma_+)>c'$ with its normal out of $N$, there is a
smooth bounding domain $E$ containing $D$ in its relative interior,
containing the required inner collar, and avoiding $\Gamma_+$,
whose free boundary satisfies
\begin{equation}\label{n3:domains:smooth-enclosure-expansion}
 \Theta_Q(\partial E\setminus\Gamma_-)=c'.
\end{equation}
Its boundary and the domain may be disconnected.  In particular $D$
is contained in the closed total trapped region at any such larger
threshold.  No regularity or positive-thickness assumption on $D$
is required.
\end{lemma}

\begin{proof}
All distance neighborhoods in this proof are relative to $N$.
Since $D$ contains a collar of $\Gamma_-$ and avoids $\Gamma_+$,
their new frontiers, at sufficiently small radii, lie in a compact
subset of the interior.  A minimizing segment ending on a new
frontier has its nearest point on the interior frontier of $D$;
it cannot start on $\Gamma_-$ because its covered collar has fixed
positive width.  Thus the proof of
Lemma~\ref{n3:domains:parallel-support} applies unchanged.  Choose
$\beta>0$ so small that its conclusion holds for $0<z\leq\beta$,
all these frontiers avoid the prescribed boundary, and
\begin{equation}\label{n3:domains:band-gap}
 c+C\beta<c'.
\end{equation}

We first construct a smooth inner enclosure without imposing any
expansion bound on it.  Smoothly approximate the continuous distance
function $d_D=\operatorname{dist}(\cdot,D)$ on $N$ uniformly, with
error less than $\beta/32$, using a finite coordinate cover,
mollification, and a partition of unity.  Choose a regular value in
$(\beta/3,\beta/2)$ of the resulting smooth function.  Its sublevel
set $U$ has smooth interior boundary and, for example, satisfies
\begin{equation}\label{n3:domains:initial-smooth-neighborhood}
 D_{\beta/4}\subset U\subset D_{3\beta/4}.
\end{equation}
The boundary $\Gamma_-$ is contained in its relative interior and
$\Gamma_+$ is excluded.  If a component of $N\setminus\overline U$
meets no component of $\Gamma_+$, fill it.  This only deletes
components of the smooth free boundary.  After this filling the
second inclusion in
Equation~\eqref{n3:domains:initial-smooth-neighborhood} need not hold
for its interior, but every remaining free boundary component still
lies in the indicated thin neighborhood, and the first inclusion
still holds.

Choose $\eta\in C^\infty(N)$, $0\leq\eta\leq1$, equal to one
on that free boundary and supported in $\{d_D<\beta\}$.  Its
existence uses the positive separation between this compact smooth
boundary and the closed set $\{d_D\geq\beta\}$; it requires no
smoothness of $d_D$.  For a sufficiently large finite constant $A$,
the tensor
\begin{equation}\label{n3:domains:modified-tensor}
 Q_A=Q-\tfrac12(c'+A\eta)g
\end{equation}
makes the free boundary of $U$ a strictly negative inner barrier:
\[
 \Theta_{Q_A}(\partial U)
       =\Theta_Q(\partial U)-c'-A<0.
\]
All prescribed outer faces remain strictly positive for $Q_A$,
because $\eta=0$ there and their original expansion exceeds $c'$.
Apply the barrier part of
Theorem~\ref{n3:domains:region-theorem} on each component of
$N\setminus\overline U$.  Each meets $\Gamma_+$ by the filling
step and has nonempty inner boundary by connectedness of $N$ and
$U\ne\varnothing$.  We obtain smooth MOTSs $\Sigma$ bounding a
domain $E$ which contains $\overline U$ and avoids $\Gamma_+$.
They satisfy, with their normal out of $E$,
\begin{equation}\label{n3:domains:modified-mots-expansion}
 \Theta_Q(\Sigma)=c'+A\eta\geq c'.
\end{equation}

We claim $\Sigma$ misses $\{d_D<\beta\}$.  Otherwise set
$z=\min_{\Sigma}d_D<\beta$.  The first inclusion in
Equation~\eqref{n3:domains:initial-smooth-neighborhood} gives
$z\geq\beta/4>0$.  Moreover $D_z$ lies on the inner side of
$\Sigma$.  To see this directly, if $d_D(y)<z$, a shortest path
from $y$ to $D$ of length less than $z$ cannot cross $\Sigma$:
such a crossing would have distance from $D$ less than $z$.
Since $D\subset E$, the whole path and $y$ lie in $E$.
Closure gives $D_z\subset\overline E$.

At a point attaining the minimum, a defining function for the
smooth boundary $\Sigma$ is therefore an exterior support to
$D_z$, with the same outward normal as in
Equation~\eqref{n3:domains:modified-mots-expansion}.  The parallel
support bound and Equation~\eqref{n3:domains:band-gap} imply
\[
 \Theta_Q(\Sigma)\leq c+Cz<c',
\]
contradicting Equation~\eqref{n3:domains:modified-mots-expansion}.
This excludes the entire modification range, so $\eta=0$ on
$\Sigma$ and proves Equation~\eqref{n3:domains:smooth-enclosure-expansion}.
The contained $D_{\beta/4}$ gives strict enclosure of $D$.
If desired, fill any remaining complementary components which meet
no designated outer face; this deletes smooth boundary components
and preserves all stated properties.
\end{proof}

For an exterior, truncate by a distant sphere that is strict at the
larger threshold.  In the black construction, $\mathcal B$ already
contains the original inner collar.  In the white construction there is
no required inner boundary, and the black faces belong to $\Gamma_+$.
The later applications therefore need compact avoidance of the foreign
faces; they need no boundary regularity of the weak set.  Empty weak sets
can be omitted.

The support argument has supplied the geometric compactness needed
for the later height estimate. We can now fix the small offsets in
the trace equation. Their size is measured against the strict
constraint margin, while their normal derivatives distinguish the
two boundary signs.

\subsection{Small offsets and the prepared domain}

\begin{proposition}[Prepared domain]\label{n3:domains:prepared-domain}
The thresholds and the domain above can be chosen with the following
properties.  The numbers $c_b,c_w<0$ are right-continuity points of
their respective full maximal-region families.  The boundary of
$\Omega$ consists of the faces in
Equation~\eqref{n3:domains:face-expansions}, has disjoint smooth
outward leaf collars as in
Lemma~\ref{n3:domains:stable-collar}, and preserves the enclosing-area
lower bound in Equation~\eqref{n3:domains:inherited-area}.

There are a smooth compactly supported function $C$ on $\Omega$,
constants $C_b\geq0$, $C_w\leq0$, and positive constants $k_B$
on its face collars such that
\begin{equation}\label{n3:domains:offset-collars}
 \begin{aligned}
  C_w\leq C\leq C_b,&\qquad
  C=C_b-k_Bs&&\text{in a black collar near }B_b,\\
  &\qquad C=C_w+k_Bs&&\text{in a white collar near }B_w.
 \end{aligned}
\end{equation}
For an absent type, its corresponding constant is simply a bound.
Set
\begin{equation}\label{n3:domains:assigned-heights}
 b_-=c_b-C_b<0,
 \qquad b_+=-c_w-C_w>0.
\end{equation}
There is a smooth positive function $\rho$, equal to a positive
constant times $r^{-4}$ on the far end and satisfying symbol
derivative bounds there, such that, for every $h\in[b_-,b_+]$,
\begin{equation}\label{n3:domains:offset-smallness}
 |(h+C)\tr K|+|\nabla C|+\rho
      \leq\tfrac12\mathfrak m
 \quad\hbox{everywhere on }\Omega.
\end{equation}
The thresholds, collars, offset, and $\rho$ are fixed before the
deformation parameters are chosen.  After fixing the thresholds and
collars, the offset and its first derivatives can be made arbitrarily
small.
\end{proposition}

\begin{proof}
The large-sphere comparison above confines all negative-threshold
regions to a fixed compact set, independently of how close the
thresholds are to zero.  Choose a slightly larger compact set
$\mathcal K$ containing this range and $\supp K$ in its interior.
Let
\[
 m_0=\min_{\mathcal K}\mathfrak m>0,
 \qquad T_0=\sup_{\mathcal K}|\tr K|.
\]
First restrict both negative thresholds to an interval so close to
zero that
\[
 \max\{|c_b|,|c_w|\}\,T_0<m_0/8,
\]
while retaining Equation~\eqref{n3:domains:black-interval}.  Lemma~\ref{n3:domains:right-continuity}
allows the required successive choices in that interval.  Construct
the two regions, select $\Omega$, and fix their disjoint collars
inside $\mathcal K$ as above.  All ensuing constants may depend on
these fixed choices.

Here is an explicit offset construction.  On each collar choose a
smooth nonincreasing cutoff $\chi(s)$, equal to one near $s=0$ and
zero near the collar's far end.  Choose $\kappa_B>0$ so that
$1-\kappa_Bs>0$ throughout that collar.  For a common amplitude
$a>0$, define the contribution of a black collar to be
$a\chi(s)(1-\kappa_Bs)$, and that of a white collar to be
$-a\chi(s)(1-\kappa_Bs)$.  Extend by zero outside the disjoint
collars and sum.  This produces a smooth function on the manifold
with boundary $\Omega$, supported in $\mathcal K$, with
$C_b=a$, $C_w=-a$, and $k_B=a\kappa_B$ on the respective collars.
These remain valid bounds when a type is absent.  Its $C^1$ norm
tends to zero with $a$, even though the collar widths have already
been fixed.

For $h\in[b_-,b_+]$ and $C\in[C_w,C_b]$,
\[
 |h+C|\leq\max\{|c_b|,|c_w|\}+C_b-C_w.
\]
Choose $a$ so small that
$(C_b-C_w)T_0+\|\nabla C\|_\infty<m_0/8$.
The first two terms of
Equation~\eqref{n3:domains:offset-smallness} are then at most
$m_0/4$ on $\mathcal K$ and vanish outside $\mathcal K$.
Finally take a fixed smooth positive weight equal to $r^{-4}$ far
out and multiply it by a sufficiently small positive constant.
The end lower bound for $\mathfrak m$ and $\delta<1$, together with
compact positivity, allow the choice $\rho\leq\mathfrak m/4$
everywhere.  This proves Equation~\eqref{n3:domains:offset-smallness}.
All remaining assertions were proved in constructing the regions
and collars.
\end{proof}

At the assigned face heights the quantity $h+C$ consequently has
the useful exact values
\begin{equation}\label{n3:domains:compatible-face-values}
 \begin{aligned}
  h+C&=c_b=P_B+H&&\text{on }B_b\text{ when }h=b_-,\\
  h+C&=-c_w=P_B-H&&\text{on }B_w\text{ when }h=b_+.
 \end{aligned}
\end{equation}
These identities determine the deformation's boundary signs.
Right-continuity and Lemma~\ref{n3:domains:weak-support} will exclude
limiting sublevel sets away from the corresponding full trapped regions.

\section{The elliptic deformation and exclusion of the floor}
\label{n3:deformation:section}

We define the coupled system, derive its scalar-curvature identity, and
choose a continuation path along which the conformal factor cannot touch
its lower bound.  We use the fixed data of
Proposition~\ref{n3:domains:prepared-domain}. Thus $\Omega$ has one controlled
asymptotically flat end, $K$ is compactly supported, and its compact boundary
$B=B_b\cup B_w$ consists of black and white faces. Either type may be absent.
The normal $\nu$ points into $\Omega$, and
\[
 H=\Div_B\nu,\qquad P_B=\tr_BK.
\]
The numbers $b_-<0<b_+$, the smooth compactly supported offset $C$, and
the positive function $\rho=O(r^{-4})$ are fixed as in that proposition.
In particular,
\begin{equation}\label{n3:deformation:face-thresholds}
 b_-+C=P_B+H\ \hbox{on }B_b,\qquad
 b_++C=P_B-H\ \hbox{on }B_w.
\end{equation}
All compact sets and geometric constants in this section are fixed after
this preparation. Extend the end radius to a smooth function on
$\overline\Omega$ bounded below by one. Write $\Omega_R$ for the large
spherical truncations and $S_R$ for their outer boundary.

A \emph{polynomial constant} is bounded by $C(1+N)^A$, with $C,A$
depending on the prepared data and fixed auxiliary profiles, but not on
$R$, the homotopy parameter, or the solution. Constants permitted to
depend arbitrarily on fixed $N$ will be identified explicitly.
First fix $0<\epsilon<1$, then choose $N$ sufficiently large, and finally
require $R\ge R_{\min}(N,\epsilon)$. We always take $N\ge4$ and
$N>2/\epsilon$.

\subsection{Variables and equations}
\label{n3:deformation:variables}

Choose a smooth nonincreasing $\vartheta:\R\to[0,1]$ equal to one on
$(-\infty,0]$ and zero on $[1,\infty)$. Define
\begin{equation}\label{n3:deformation:l-profile}
 \begin{gathered}
 p(t)=N\vartheta(Nt),\qquad
 l(t)=\exp\left(\int_0^t p(s)\dd s\right),\\
 L(t)=l(t)e^{2t},\qquad p_L=p+2.
 \end{gathered}
\end{equation}
Thus $l(0)=1$, $l(t)=e^{Nt}$ for $t\le0$, and $l$ is constant for
$t\ge1/N$. Derivatives of $p$ of every fixed order are polynomially bounded.
Set
\begin{equation}\label{n3:deformation:small-parameters}
 \ell=e^{-\epsilon N},\qquad \tau=\ell^{3/2}.
\end{equation}
Whenever $t\ge-\epsilon$, one has $\ell\le l(t)\le e$.

For a function $f$, set
\begin{equation}\label{n3:deformation:basic-variables}
 \begin{gathered}
 \sigma=|\nabla f|,\quad D=1+l^2\sigma^2,\quad d=D^{-1},\\
 a=l\sigma/\sqrt D,\quad v=a^2=1-d,\quad
 w=l\nabla f/\sqrt D,\quad \chi=4d/(4+pv),\\
 Z=t+\tfrac18\log D,\quad u=L\sqrt D,\quad U=l\sqrt D,\quad h=\tau f.
 \end{gathered}
\end{equation}
Derivatives and contractions refer to $g$ unless indicated otherwise.
At $\sigma>0$, write $e=\nabla f/\sigma$ and
$A_j=I+(j-1)e\otimes e$. The matrices used below have the direction-free
expressions
\begin{equation}\label{n3:deformation:axial-matrices}
 A_d=I-w\otimes w,\qquad
 A_\chi=I-\frac{p+4}{4+pv}w\otimes w.
\end{equation}
Vectors and covectors are identified with $g$.

The unknowns will be $(f,Z)$. For fixed $\nabla f$,
$\partial Z/\partial t=1+pv/4>0$, and its defining expression tends to
$-\infty$ and $+\infty$ at the two ends of the $t$-axis. It therefore
defines a unique smooth $t=t(Z,\nabla f)$ for all inputs, without a floor
assumption. The coefficients $t,d,v,w,\chi,u,U$ and the direction-free
matrices are smooth at $\nabla f=0$; occurrences of $\sigma$ and $a$
in the equations combine into smooth invariant expressions. One has
$0<\chi\le d\le1$ for finite inputs.
Define
\begin{equation}\label{n3:deformation:metrics}
 \bar g=g+l^2\dd f^2,\qquad \hat g=e^{4t}\bar g,\qquad
 H^f=\frac l{\sqrt D}\Hess f,\qquad S=K+H^f.
\end{equation}
The original trace equation is
\begin{equation}\label{n3:deformation:trace}
 \tr_{A_\chi}S=F,\qquad F=h+C,
\end{equation}
where $\tr_A S$ is contraction with the contravariant matrix $A$.
Along the continuation path we will specify modified right sides $F$.
The full flux is
\begin{equation}\label{n3:deformation:flux}
 V=uA_\chi\bigl(4\nabla Z+K(w,\cdot)\bigr).
\end{equation}
Here $f$ is the graph height, while $Z$ combines the conformal
variable $t$ with the logarithm of the graph slope. The numbers $d$
and $\chi$ record ellipticity in the gradient direction, and $u$
weights the divergence flux. The scaled height $h=\tau f$ makes the
prescribed trace strictly increasing in $f$.
At $S_R$ prescribe $f=Z=0$. Initially the inner Dirichlet values are
$h=b_-$ on $B_b$ and $h=b_+$ on $B_w$.

\subsection{The exact scalar-curvature identity}

The graph-curvature strategy originates in Jang's equation and the
Schoen--Yau positive-energy argument
\cite{Jang1978,SchoenYau1981}. Its generalized warped-graph form and
Penrose coupling were developed by Bray and Khuri
\cite{BrayKhuri2010,BrayKhuri2011}; see in particular Identity~9
of~\cite{BrayKhuri2011}. We use the independently proved algebraic identity of~\cite[Proposition~4.1]{NumericalCompanion},
with the normalization below. The prescribed trace $F=h+C$, nonzero source,
and floor penalty below define a different coupled system; the
identity does not import global solvability of a zero-divergence
coupling.

The next identity separates the original constraint densities, a
nonnegative quadratic term, and a divergence.  The second equation will
prescribe that divergence so that the deformed metric has nonnegative
scalar curvature.

\begin{proposition}[Deformation identity]\label{n3:deformation:identity}
Let $f,t$ be smooth and $F=\tr_{A_\chi}(K+H^f)$. At $\sigma>0$ decompose
$S$ relative to $\R e\oplus e^\perp$:
\[
 T=S|_{e^\perp\times e^\perp},\quad
 M=S(e,\cdot)|_{e^\perp},\quad q=S(e,e),\quad \tr T=F-\chi q.
\]
Then
\begin{equation}\label{n3:deformation:scalar-identity}
 \begin{split}
 \tfrac12e^{4t}\Scal_{\hat g}
 ={}&8\pi(\mu+J(w))+\mathcal T
       +w(F)-F\tr K-u^{-1}\Div V,
 \end{split}
\end{equation}
where
\begin{equation}\label{n3:deformation:quadratic-term}
 \begin{split}
 \mathcal T={}&\tfrac12|T^{\operatorname{tf}}|^2
 +(M+pa\nabla_\perp t)\cdot(M+(p+2)a\nabla_\perp t)\\
 &+4(p+1)|\dd t|_{A_d}^2
   +(\chi-\chi^2/4)q^2+2(p+1)\chi qa\,\partial_e t\\
 &-\chi Fq/2+3F^2/4.
 \end{split}
\end{equation}
After substituting $F=\tr_{A_\chi}S$, this expression extends smoothly
across $\sigma=0$.
\end{proposition}

\begin{proof}
Apply the scalar identity of~\cite[Proposition~4.1]{NumericalCompanion} in spatial
dimension three, so its surface dimension is $k=2$. To distinguish
the two conventions temporarily, add a subscript $\mathrm A$ to its
variables. The exact dictionary is
\begin{equation}\label{n3:deformation:normalization-dictionary}
 \begin{gathered}
 t_{\mathrm A}=2t,\quad Z_{\mathrm A}=2Z,\quad
 l_{\mathrm A}(s)=l(s/2),\quad p_{\mathrm A}=p/2,\\
 \mu_{\mathrm A}=8\pi\mu,\qquad J_{\mathrm A}=8\pi J,
 \qquad S_{\mathrm A}=S.
 \end{gathered}
\end{equation}
The quantities $D,d,w,a,v,H^f,F,U$ are identical. In particular,
\[
 \frac{2d}{2+p_{\mathrm A}v}=\frac{4d}{4+pv}=\chi,
 \quad u_{\mathrm A}=e^{t_{\mathrm A}}l\sqrt D=u,
 \quad 2\nabla Z_{\mathrm A}=4\nabla Z,
\]
and the companion metric $e^{2t_{\mathrm A}}\bar g$ is precisely
$\hat g=e^{4t}\bar g$. Thus the divergence term, trace prescription
and metric match without rescaling their flux. The physical density
conversion gives $8\pi(\mu+J(w))$; the ADM energy and momentum are
not changed by that conversion.

For the quadratic expression, the companion coefficient
$s_p=p_{\mathrm A}+(k-1)/2$ is $(p+1)/2$, and
$dt_{\mathrm A}=2dt$. Its transverse square is
$|M+(p+1)a\nabla_\perp t|^2$. Combining its residual
$-[a^2|\nabla_\perp t|^2]$ with that square gives
$(M+pa\nabla_\perp t)\cdot(M+(p+2)a\nabla_\perp t)$.
The remaining gradient contribution is
$4(p+1)|dt|_{A_d}^2$. Its normal terms become
\[
 (\chi-\chi^2/4)q^2+2(p+1)\chi qa\,\partial_e t
       -\chi Fq/2+3F^2/4.
\]
These substitutions give Equations~\eqref{n3:deformation:scalar-identity}
and~\eqref{n3:deformation:quadratic-term}. Smoothness at zero gradient
is part of the invariant identity: the temporary axis has not been
used to define its coefficients. The import is pointwise and
requires neither an energy condition nor a solved divergence equation.
\end{proof}

\begin{lemma}[Polynomial coercivity]\label{n3:deformation:coercivity}
The weighted coercivity theorem of~\cite[Proposition~4.2]{NumericalCompanion}, with
Equation~\eqref{n3:deformation:normalization-dictionary}, makes
$\mathcal T$ nonnegative. For $0\le p\le N$,
\begin{equation}\label{n3:deformation:coercive-estimate}
 |\dd t|_{A_d}^2+|T|^2+|M|^2+dq^2+F^2
 \le C(1+N)^A\mathcal T
\end{equation}
with universal $C,A$.
At $\dd t=0$ there is the stronger comparison, uniform in $p,d$,
\begin{equation}\label{n3:deformation:floor-norm}
 c\bigl(|T^{\operatorname{tf}}|^2+|M|^2+\chi q^2+F^2\bigr)
 \le\mathcal T
 \le C\bigl(|T^{\operatorname{tf}}|^2+|M|^2+\chi q^2+F^2\bigr).
\end{equation}
\end{lemma}

\begin{proof}
With $x=\partial_e t$ and $y=\nabla_\perp t$, exact completion gives
\begin{equation}\label{n3:deformation:square-completion}
 \begin{split}
 \mathcal T={}&\tfrac12|T^{\operatorname{tf}}|^2
 +|M+(p+1)ay|^2+[4(p+1)-v]|y|^2\\
 &+4(p+1)d\left(x+\frac{\chi aq}{4d}\right)^2\\
 &+\tfrac34\left(F-\frac{\chi q}{3}\right)^2
 +\frac{4d(9-d)}{3(4+pv)^2}q^2.
 \end{split}
\end{equation}
Completing first in $x$ leaves $3\chi q^2/(4+pv)$; completing next in
$F$ leaves $3\chi/(4+pv)-\chi^2/12=4d(9-d)/[3(4+pv)^2]$.
The latter coefficient is at least $32d/[3(4+N)^2]$.
It controls $dq^2$ polynomially. The $F$-square then controls $F^2$,
since $\chi^2q^2\le dq^2$.
The transverse gradient coefficient is at least three, and the shifted
$M$-square controls $M$ with an additional factor $(1+N)^2$.
The squared $d$-norm of the shift in the $x$-square is
$\chi^2vq^2/(16d)\le dq^2/16$, so it also controls $dx^2$.
Finally $|T|^2=|T^{\operatorname{tf}}|^2+(F-\chi q)^2/2$.
This proves \eqref{n3:deformation:coercive-estimate}.
If $\dd t=0$, the sole cross term is $-\chi Fq/2$.
Using $|\chi Fq|/2\le\chi q^2/4+\chi F^2/4$ and $0<\chi\le1$
proves \eqref{n3:deformation:floor-norm}.
\end{proof}

The scalar identity separates curvature from divergence. At the inner
boundary we must also relate this divergence flux to mean curvature.
Differentiating the definition of $Z$ gives, wherever $df\ne0$,
\begin{equation}\label{n3:deformation:Z-differential}
 4\,dt+\tfrac12\,d\log D
       =4\,dZ=(4+pv)\,dt+v\,d\log\sigma.
\end{equation}
This relation is smooth in its direction-free form and will be used
only in contractions that extend across $df=0$.

\begin{lemma}[Boundary identity]\label{n3:deformation:boundary-identity}
On a face where $f$ is constant,
\begin{equation}\label{n3:deformation:boundary}
 V_\nu/u=d(H+4\partial_\nu t)-H+w_\nu(F-P_B).
\end{equation}
Its mean curvature in $\hat g$, with normal into $\Omega$, is
$e^{-2t}\sqrt d\,(H+4\partial_\nu t)$.
\end{lemma}

\begin{proof}
The tangential Hessian is
$\Hess f|_{TB}=(\partial_\nu f)\operatorname{II}_g$.
Thus $\tr_\perp H^f=w_\nu H$ and
$\chi q=F-P_B-w_\nu H$.
The normal component of \eqref{n3:deformation:flux}, using
\eqref{n3:deformation:Z-differential}, is
\[
 V_\nu/u=4d\partial_\nu t+\chi w_\nu q
        =4d\partial_\nu t+w_\nu(F-P_B)-vH.
\]
The $\bar g$-normal is $\sqrt d\,\nu$. Its mean curvature is $\sqrt d\,H$,
because the tangential metric is unchanged and its normal first variation
is multiplied by $\sqrt d$. The conformal mean-curvature formula gives
the asserted mean curvature for $\hat g$.
\end{proof}

\subsection{The source, boundary condition, and continuation path}

We now prescribe the divergence and inner flux. For the original system,
the source is chosen to give the deformed metric nonnegative scalar
curvature. A penalty near the lower bound will exclude contact with that
bound along the continuation path.

Choose $3/4<\gamma<1$ and fixed positive smooth weights
\[
 \rho_0\asymp r^{-4},\qquad \rho_1\asymp r^{-2\gamma},
\]
with differentiated end bounds. Let $m:\R\to[0,1]$ be smooth, equal to one
on $(-\infty,0]$ and zero on $[1,\infty)$, and set
\[
 m_0(t)=m(N(t+\epsilon)),\qquad
 \mathcal P=C_N(\rho_0+v\rho_1).
\]
The positive $\delta_0$ is fixed independently of $N,R$ in
Proposition~\ref{n3:deformation:floor}; thereafter $C_N$ is a sufficiently
large polynomial constant. Define
\begin{equation}\label{n3:deformation:source}
 \Xi=\delta_0\mathcal T+\rho+\delta_0\tau a\sigma-m_0(t)\mathcal P.
\end{equation}
Choose a smooth face function $N_B>1+\sup_B|H|$.
The original second equation and its inner condition are
\begin{equation}\label{n3:deformation:original-second-equation}
 \Div V=u\Xi,\qquad V_\nu/u=-H+w_\nu(F-P_B)-N_Bd.
\end{equation}
Lemma~\ref{n3:deformation:boundary-identity} gives
$H+4\partial_\nu t=-N_B$.
Only $-\epsilon\le t\le-\epsilon+1/N$ is penalized in the admissible range.
Since $N>2/\epsilon$, this band has negative $t$, and
\begin{equation}\label{n3:deformation:penalty-band}
 \ell\le l\le e\ell,\qquad c\ell\le L\le C\ell
 \quad\hbox{on the admissible penalty band}.
\end{equation}
The latter constants can be uniform for $0<\epsilon<1$.

\begin{proposition}[Specified homotopy]\label{n3:deformation:homotopy}
There is a continuous piecewise smooth family of systems indexed by
$s\in[0,4]$, beginning at
\eqref{n3:deformation:trace}, \eqref{n3:deformation:original-second-equation}
and the original Dirichlet values, and ending at $f=0$ and the scalar
problem with zero source and zero inner conormal flux.
Throughout its first three stages the source is
\eqref{n3:deformation:source}, with $\mathcal T$ evaluated using the actual
trace $F$. Each trace equation has right-side derivative at least one
in $h$, with $x,Z,\nabla f$ fixed.
\end{proposition}

\begin{proof}
Fix a smooth $j:\R\to[0,1]$ which vanishes for $t\le0$ and equals one
for $t\ge1$. For $0\le\lambda\le1$ put
\[
 V_\lambda=uA_\chi(4\nabla Z+\lambda K(w,\cdot)),\qquad
 \mathcal B=-H+w_\nu(F-P_B)-N_Bd.
\]
Before the last stage use the boundary condition
\begin{equation}\label{n3:deformation:homotopy-boundary}
 \begin{split}
 (V_\lambda)_\nu/u
 ={}&[1-(1-\lambda)j(t)]\\
 &\cdot[\mathcal B-(1-\lambda)(A_\chi K(w,\cdot))_\nu].
 \end{split}
\end{equation}
Outer values remain $f=Z=0$ throughout.

We specify two smooth collar modifications. Let
$\nu_s=\nabla s/|\nabla s|$ be the outward unit normal to a collar leaf;
the leaf coordinate $s$ here is local and is separate from the homotopy
parameter. Choose disjoint collar cutoffs equal to one at the faces.
Take $\eta>0$ so small that $b_-+3\eta<b_+$.
A black contribution to $D_0(x,w,h)$ is a negative constant times a
collar cutoff, times a smooth directional cutoff supported where
$w\cdot\nu_s<-1/2$ and equal to one at $-1$, times a nonincreasing
height cutoff equal to one below $b_-+\eta$ and zero above $b_-+2\eta$.
A white contribution has positive sign, its directional cutoff supported
where $w\cdot\nu_s>1/2$ and equal to one at $1$, and a nondecreasing
height cutoff equal to zero below $b_+-2\eta$ and one above $b_+-\eta$.
Decrease $\eta$ to separate these ranges. Their sum $D_0$ is smooth,
bounded, supported in the collars, and nondecreasing in $h$.
It is nonpositive for $h\le b_-+\eta$, nonnegative for
$h\ge b_+-\eta$, and vanishes on black outward and white inward
directions, for every length $0\le|w|\le1$.

Choose the magnitudes larger than $2\sup_B|H|+1$.
Equation~\eqref{n3:deformation:face-thresholds} then gives, for
$F=h+C+D_0$ at the assigned boundary height $b=b_\pm$,
\begin{equation}\label{n3:deformation:directional-barriers}
 F(x,+\nu,b)\ge P_B+H,\qquad F(x,-\nu,b)\le P_B-H.
\end{equation}
These are limiting directional evaluations at $a=1,\chi=0$.
The compatible direction has equality; $D_0$ corrects the other one.
Its derivatives on bounded height ranges are independent of $N$.
Choose a smooth compactly supported extension $b(x)$, constant equal
to $b_\pm$ on each smaller collar.
On these collars let $D_1(x,w)=A_1w\cdot\nu_s$, extended by a collar
cutoff, with $A_1>\sup_B|H|+1$.
Then $D_1(x,0)=0$ and $D_1(x,\pm\nu)=\pm A_1$ on $B$.

The four stages are the following.
\begin{enumerate}[label=\textup{(\roman*)}]
 \item For $0\le s\le1$, set $\lambda=1-s$, keep $F=h+C$ and $h|_B=b$,
 and use $\Div V_\lambda=u\Xi$ and
 \eqref{n3:deformation:homotopy-boundary}.
 \item For $1\le s\le2$, set $\lambda=0$, $\alpha=s-1$, retain $h|_B=b$,
 and replace the trace right side by $F=h+C+\alpha D_0(x,w,h)$.
 Keep the same scalar equation and boundary prescription.
 \item For $2\le s\le3$, set $\lambda=0$, $\zeta=s-2$, prescribe
 $h|_B=(1-\zeta)b$, and use
 \begin{equation}\label{n3:deformation:interpolation-trace}
 \begin{split}
 F={}&h+(1-\zeta)[C+D_0(x,w,h+\zeta b(x))]\\
    &+\zeta[\tr_{A_\chi}K+D_1(x,w)].
 \end{split}
 \end{equation}
 Keep the same scalar equation and boundary prescription.
 At a face, $h+\zeta b=b$. The value at $w=\pm\nu$ is the convex
 combination of the preceding directional value and
 $P_B+D_1(x,\pm\nu)$, so \eqref{n3:deformation:directional-barriers} persists.
 \item For $3\le s\le4$, retain the trace problem at $\zeta=1$ and
 multiply both $\Xi$ and the right side of
 \eqref{n3:deformation:homotopy-boundary} by $\eta_s=4-s$.
\end{enumerate}
The prescriptions agree at the common endpoints.
Their $h$-derivatives are $1$, $1+\alpha\partial_hD_0$, or
$1+(1-\zeta)\partial_hD_0$, all at least one.

At $\zeta=1$ the trace equation is
$\tr_{A_\chi}H^f=h+D_1(x,w)$ with zero Dirichlet values.
At a positive interior maximum of $f$, one has $w=0$, $A_\chi=I$,
$D_1=0$, and $l\Delta f=\tau f>0$, a contradiction.
A negative minimum is excluded in the same way.
Thus $f=0$ and $t=Z$ in the last stage.
At $s=4$ the remaining equation is
$\Div(4L(Z)\nabla Z)=0$, with zero inner conormal flux and $Z=0$ on
$S_R$. Multiplication by $Z$ and integration give $Z=0$.
For the fixed-point construction in Section~\ref{n3:existence:section},
the final linear scalar problem has zero source and zero inner flux
for every input; consequently that last map has output identically zero.
\end{proof}

\subsection{Elementary height and end barriers}

\begin{lemma}[Height control]\label{n3:deformation:height}
Every smooth trace solution in Proposition~\ref{n3:deformation:homotopy}
satisfies $|h|\le C_0$, independently of $N,R$, the homotopy parameter,
and the input $Z$. Its actual right side satisfies $|F|\le C_F$
with the same independence.
\end{lemma}

\begin{proof}
The collar modifications vanish at $w=0$.
At an interior extremum the trace equation reads either
$\tr K+l\Delta f=h+C$ or
$\tr K+l\Delta f=h+(1-\zeta)C+\zeta\tr K$.
The maximum and minimum inequalities bound $h$ by fixed bounds for
$|\tr K|+|C|$. Boundary values range between $b_-,0,b_+$.
This proves the height bound without using the floor.
Boundedness of $D_0,D_1$ and $|\tr_{A_\chi}K|\le3|K|$ proves the
assertion for $F$.
\end{proof}

\begin{lemma}[End height and outer boundary]\label{n3:deformation:end-height}
There are $r_0,C>0$, independent of $N,R$ and the homotopy parameter, such
that for $R\ge2r_0$,
\begin{equation}\label{n3:deformation:end-height-estimate}
 |h|\le C(r^{-\gamma}-R^{-\gamma})\quad(r_0\le r\le R).
\end{equation}
In particular,
\begin{equation}\label{n3:deformation:outer-gradient}
 |\nabla f|\le C R^{-1-\gamma}/\tau\quad\hbox{on }S_R.
\end{equation}
For fixed $\epsilon,N$, a sufficiently large lower bound for $R$ gives
$t>-\epsilon$ on $S_R$.
\end{lemma}

\begin{proof}
Take $r_0$ outside the supports of $K,C,D_0,D_1$.
There $F=h$ in the first three stages, and $h=0$ in the last one.
For $\psi=r^{-\gamma}$, with axis $\nabla\psi/|\nabla\psi|$,
the controlled end gives, uniformly for $0<\chi\le1$,
\[
 \tr_{A_\chi}\Hess\psi
 =\gamma((\gamma+1)\chi-2)r^{-\gamma-2}
      +O(r^{-\gamma-3})<0
\]
after enlarging $r_0$.
At a comparison contact, coefficients use the given $Z$ and the common
gradient of test and solution. The positive factor $l/(\tau\sqrt D)$
preserves this sign. Height monotonicity therefore makes
$C(\psi-R^{-\gamma})$ an upper barrier and its negative a lower barrier.
Choose $C$ to dominate $C_0$ on $r=r_0$ for every $R\ge2r_0$.
Differentiating at the common outer value gives
\eqref{n3:deformation:outer-gradient}; tangential derivatives there vanish.

At $S_R$, $Z=0$. The monotonicity of its defining expression implies
$t>-\epsilon$ if $\ell^2|\nabla f|^2<e^{8\epsilon}-1$.
It suffices to require
$C^2\ell^2\tau^{-2}R^{-2-2\gamma}<e^{8\epsilon}-1$,
which is a permissible lower bound depending on fixed $N,\epsilon$.
\end{proof}

\subsection{Exclusion of a contact with the floor}

We exclude a proposed floor contact using the equations at that point.
These estimates require neither an upper bound for $Z$ nor a bound for
$|\nabla f|$.

\begin{lemma}[Curvature at a minimum of $t$]\label{n3:deformation:floor-curvature}
At an interior point where $\dd t=0$ and $\Hess t\ge0$,
\begin{equation}\label{n3:deformation:floor-gauss}
 \tfrac12e^{4t}\Scal_{\hat g}
 \le \tfrac12\Scal_g-v\Ric_g(e,e)+\mathcal S(H^f),
\end{equation}
where, for a symmetric tensor $A$,
\begin{equation}\label{n3:deformation:gauss-quadratic}
 \begin{split}
 \mathcal S(A)={}&\tfrac14(\tr_\perp A)^2
 -\tfrac12|A_{\perp\perp}^{\operatorname{tf}}|^2\\
 &+dA_{ee}\tr_\perp A-d|A_{e\perp}|^2.
 \end{split}
\end{equation}
There are universal $c_*,\eta_*>0$ and polynomial constants $C_N$ such
that, at $\dd t=0$,
\begin{equation}\label{n3:deformation:floor-algebra-general}
 \mathcal T-\mathcal S(H^f)
 \ge c_*\mathcal T-C_N(F^2+|K|^2).
\end{equation}
If $K=0$, the more useful uniform estimate is
\begin{equation}\label{n3:deformation:floor-algebra-end}
 \mathcal T-\mathcal S(H^f)\ge c_*\mathcal T-C_N vF^2.
\end{equation}
\end{lemma}

\begin{proof}
Consider the graph of $f$ in the Riemannian product with warped fiber
$(\Omega\times\R,g+l^2\dd b_0^2)$. Its induced metric is $\bar g$.
At $\dd t=0$ one also has $\dd l=0$, and its second fundamental form,
up to orientation, is $H^f$. Write $E=l^{-1}\partial_{b_0}$.
The graph unit normal is $\sqrt d\,E-ae$. The ambient curvature formulas
at this point are
\[
 \begin{gathered}
 \Scal_{\mathrm{amb}}=\Scal_g-2\Delta l/l,\\
 \Ric_{\mathrm{amb}}|_{T\Omega}=\Ric_g-\Hess l/l,\qquad
 \Ric_{\mathrm{amb}}(E,E)=-\Delta l/l,
 \end{gathered}
\]
with zero mixed Ricci terms. Hence the ambient contribution to half the
graph scalar curvature is
\[
 \tfrac12\Scal_g-v\Ric_g(e,e)
      -v\,\tr_\perp\Hess l/l.
\]
The last term is nonpositive: at $\dd t=0$,
$\Hess l=lp\,\Hess t$ and $p\ge0$.
The second-form contribution is
$[(\tr_{\bar g}H^f)^2-|H^f|_{\bar g}^2]/2$.
Since the inverse graph metric is $A_d$, its expansion is exactly
\eqref{n3:deformation:gauss-quadratic}.
Finally the conformal contribution at this point is
$-4\bar\Delta t\le0$. This proves \eqref{n3:deformation:floor-gauss}.

The pointwise algebraic clause of~\cite[Lemma~4.3]{NumericalCompanion} applies
under Equation~\eqref{n3:deformation:normalization-dictionary}. We also
record its three-dimensional calculation, because the boundary
homotopy below needs its precise degenerate normal coefficient. First evaluate $\mathcal S$ on
$S=K+H^f$, whose trace in the transverse plane is $F-\chi q$.
Subtracting \eqref{n3:deformation:gauss-quadratic} from
\eqref{n3:deformation:quadratic-term} at $\dd t=0$ gives
\begin{equation}\label{n3:deformation:floor-difference}
 \begin{split}
 \mathcal T-\mathcal S(S)
 ={}&|T^{\operatorname{tf}}|^2+(1+d)|M|^2\\
 &+\chi(1+d-\chi/2)q^2-dFq+F^2/2.
 \end{split}
\end{equation}
The normal coefficient is at least $\chi$, while
$|dFq|\le\chi q^2/2+d^2F^2/(2\chi)$.
Because
\[
 d^2/\chi=d(4+pv)/4\le1+N/4,
\]
Equation~\eqref{n3:deformation:floor-norm} implies
$\mathcal T-\mathcal S(S)\ge c\mathcal T-C(1+N)F^2$
with universal $c>0$.

To replace $S$ by $H^f=S-K$, polarize
\eqref{n3:deformation:gauss-quadratic}. The difference is bounded by
\[
 C|K|(|T^{\operatorname{tf}}|+|M|+|F|+d|q|)+C|K|^2.
\]
Indeed the only extra normal product is $dq\,\tr_\perp K$;
all remaining products involve transverse entries, $dM$, or
$\tr_\perp S=F-\chi q$.
By \eqref{n3:deformation:floor-norm} and $d/\sqrt\chi\le C\sqrt{1+N}$,
this bound is at most
$C\sqrt{1+N}|K|\sqrt{\mathcal T}+C|K|^2$.
Young's inequality absorbs a fixed small fraction of $\mathcal T$
and costs only $C(1+N)|K|^2$. This proves
\eqref{n3:deformation:floor-algebra-general}.

If $(1+p)v\le\eta_*$, then $d$ and $\chi$ are within $C\eta_*$
of one. At $d=\chi=1$ the normal block in
\eqref{n3:deformation:floor-difference} is
$3q^2/2-Fq+F^2/2$; its symmetric matrix has smallest eigenvalue
$(2-\sqrt2)/2>0$.
Choose $\eta_*>0$ small enough that its coefficients remain within
half this spectral gap. The transverse coefficients are already
bounded below. For $K=0$, Equation~\eqref{n3:deformation:floor-norm}
therefore gives $\mathcal T-\mathcal S(H^f)\ge c_*\mathcal T$
in this regime.
In the remaining regime $v>\eta_* /(1+p)$, use
\eqref{n3:deformation:floor-algebra-general} with $K=0$ and absorb
$(1+N)F^2$ into $C(1+N)^2vF^2$.
This proves \eqref{n3:deformation:floor-algebra-end}.
All formulas extend to $\sigma=0$ by their invariant expressions;
there $d=\chi=1$ and the full quadratic expressions are independent of
the temporary choice of axis.
\end{proof}

\begin{lemma}[Homotopy errors at a floor contact]
\label{n3:deformation:floor-errors}
There is a fixed compact set $\mathcal K\subset\overline\Omega$ containing
the supports of $K,C,b,D_0,D_1$ such that, at an interior minimum
$t=-\epsilon$ of a smooth solution in the first three stages,
\begin{equation}\label{n3:deformation:floor-divergence-lower}
 \Div V_\lambda/u
 \ge 2\delta_0\mathcal T+\tau a\sigma
      -C(1+N)^A
       [\mathbf1_{\mathcal K}+v(\rho_0+\rho_1)].
\end{equation}
Here $\delta_0>0$ can be fixed sufficiently small independently of
$N,R$ and the homotopy parameter.
\end{lemma}

\begin{proof}
At this contact $\dd t=0$ and $p=N$ is locally constant as a function
of $t$ near its value $-\epsilon$. Differentiating $w$ in an orthonormal
frame gives
\begin{equation}\label{n3:deformation:floor-w-derivative}
 \nabla_iw_j=H^f_{ij}-w_jH^f_{ik}w^k.
\end{equation}
In the $e$-frame the normal vector slot is multiplied by $d$.
Equations \eqref{n3:deformation:floor-norm} and $d/\sqrt\chi\le C\sqrt{1+N}$
therefore show
\begin{equation}\label{n3:deformation:floor-w-control}
 |\nabla w|\le C(1+N)^A(|K|+\sqrt{\mathcal T}).
\end{equation}
No bound for the uncontracted normal Hessian is asserted.

Each first-three-stage right side can be written as a smooth function
$F(x,w,h,p)$. Its $h$-derivative is at least one.
On the bounded height range supplied by Lemma~\ref{n3:deformation:height},
its explicit $x$- and $w$-derivatives are polynomially bounded and
compactly supported, apart from the derivative of the term $h$.
This follows directly from the fixed collar cutoffs and
$A_\chi=I-(p+4)w\otimes w/(4+pv)$, whose denominator is at least four.
For example its first $w$-derivatives on $|w|\le1$ are bounded by
$C(1+N)^2$, and the other explicit modifications have bounds independent
of $N$. Combining these bounds with
\eqref{n3:deformation:floor-w-control} and Young's inequality therefore
produces a fixed finite power of $1+N$.
Since $\nabla p=0$ at the contact and $w(h)=\tau a\sigma$, the chain rule
and \eqref{n3:deformation:floor-w-control} give, for every fixed $\eta>0$,
\begin{equation}\label{n3:deformation:floor-F-control}
 w(F)\ge\tau a\sigma-\eta\mathcal T
                 -C_\eta(1+N)^A\mathbf1_{\mathcal K}.
\end{equation}
The dependence $D_0(x,w,h+\zeta b(x))$ only adds the fixed compact
derivative of $b$ and does not change this conclusion.

The omitted flux is $(1-\lambda)uA_\chi K(w,\cdot)$.
Let $B_K=A_\chi K(w,\cdot)$. Its derivative is bounded by a polynomial
constant times $1+\sqrt{\mathcal T}$ on $\mathcal K$, by
\eqref{n3:deformation:floor-w-control}; derivatives of $p$ vanish.
The weight term requires the contraction rather than a bound for
$|\nabla\log u|$. At the contact,
\[
 \nabla\log u=H^f(w,\cdot),\qquad
 \langle\nabla\log u,B_K\rangle
 =a^2\bigl(\chi H^f_{ee}K_{ee}
                    +H^f_{e\perp}\cdot K_{e\perp}\bigr).
\]
This involves only $\chi H^f_{ee}$ and $H^f_{e\perp}$ and is controlled
by \eqref{n3:deformation:floor-norm}. It follows that
\begin{equation}\label{n3:deformation:floor-flux-control}
 |\Div(uB_K)|/u
 \le\eta\mathcal T+C_\eta(1+N)^A\mathbf1_{\mathcal K}.
\end{equation}

Combine the scalar identity with \eqref{n3:deformation:floor-gauss}.
The difference $8\pi\mu-\Scal_g/2$ equals
$((\tr K)^2-|K|^2)/2$ and is compactly supported; so is $J$.
The term $F\tr K$ is compactly bounded by Lemma~\ref{n3:deformation:height}.
On $\mathcal K$, Equation~\eqref{n3:deformation:floor-algebra-general}
therefore contributes $c_*\mathcal T-C(1+N)^A$.
Outside $\mathcal K$, $K=0$, $F=h$, and
\eqref{n3:deformation:floor-algebra-end} together with
Lemma~\ref{n3:deformation:end-height} gives
$c_*\mathcal T-C(1+N)^Av r^{-2\gamma}$.
The remaining Ricci term is bounded below by
$-C\mathbf1_{\mathcal K}-Cv r^{-3}$; enlarge $\mathcal K$ if necessary.
Both end weights are bounded by a multiple of $\rho_0+\rho_1$.
Finally subtract the omitted flux and use
\eqref{n3:deformation:floor-F-control}--\eqref{n3:deformation:floor-flux-control},
choosing their two $\eta$'s with sum at most $c_*/2$.
Fix $0<\delta_0<\min(c_*/4,1/4)$. Decreasing it slightly if necessary
gives \eqref{n3:deformation:floor-divergence-lower}.
\end{proof}

\begin{proposition}[Floor exclusion]\label{n3:deformation:floor}
Choose $\delta_0$ as in Lemma~\ref{n3:deformation:floor-errors}.
There is a polynomial choice $C_N=C(1+N)^A$ in the penalty such that,
for all sufficiently large $N$ and all
$R\ge R_{\mathrm{floor}}(N,\epsilon)$, no smooth solution anywhere along
Proposition~\ref{n3:deformation:homotopy} can satisfy $t\ge-\epsilon$
and attain $t=-\epsilon$.
In particular every such admissible solution has $t>-\epsilon$ on
$\overline{\Omega_R}$.
\end{proposition}

\begin{proof}
Choose $R_{\mathrm{floor}}$ to include the requirement of
Lemma~\ref{n3:deformation:end-height}. A floor contact cannot occur on $S_R$.
In the first three stages, an inner contact has $j(t)=0$.
Equation~\eqref{n3:deformation:homotopy-boundary} then cancels the omitted
flux on both sides and gives $V_\nu/u=\mathcal B$.
Equation~\eqref{n3:deformation:boundary} forces
$H+4\partial_\nu t=-N_B$, so $\partial_\nu t<0$.
This contradicts the nonnegative derivative into $\Omega_R$ at a
boundary minimum.

Consider an interior contact. Let
$E_N=C(1+N)^A[\mathbf1_{\mathcal K}+v(\rho_0+\rho_1)]$
denote the error in \eqref{n3:deformation:floor-divergence-lower}.
Because $\rho_0$ has a positive minimum on $\mathcal K$,
$\mathbf1_{\mathcal K}\le C_{\mathcal K}\rho_0$.
Also $v\rho_0\le\rho_0$, and $\rho/\rho_0$ is globally bounded.
We can therefore choose a polynomial $C_N$ so large that, pointwise
for all $0\le v\le1$,
\begin{equation}\label{n3:deformation:penalty-choice}
 \mathcal P\ge E_N+2\rho+\rho_0.
\end{equation}
At the floor $m_0=1$, so the scalar equation and the lower estimate give
\[
 \begin{split}
 0
 &\ge \delta_0\mathcal T+(1-\delta_0)\tau a\sigma
                     +\mathcal P-E_N-\rho\\
 &\ge \rho+\rho_0>0,
 \end{split}
\]
a contradiction. Enlarging $C_N$ below preserves this inequality and
preserves its polynomial dependence.

In the last stage $f=0$, $t=Z$, $u=L$, and $v=0$.
At an interior minimum,
$\mathcal T=(|K|^2+(\tr K)^2)/2$.
Increase $C_N$ polynomially if necessary so that
$\delta_0\mathcal T+\rho-C_N\rho_0<0$ at all such points.
For $\eta_s>0$, the scalar equation at a floor minimum has left side
$4L\Delta t\ge0$ and strictly negative right side.
At an inner floor minimum its flux is
$4\partial_\nu t=\eta_s(-H-N_B)<0$, also impossible.
If $\eta_s=0$, multiplication of
$\Div(4L(t)\nabla t)=0$ by $t$, with zero outer value and zero inner
flux, gives $t=0$.
Thus the last stage has no floor contact either.
\end{proof}

We have excluded floor contact for every smooth admissible solution
along the specified homotopy. Section~\ref{n3:apriori:section} supplies the
remaining a priori bounds, and Section~\ref{n3:existence:section} proves
existence. The choice of $\mathcal P$ uses only polynomial losses in $N$;
it is independent of all fixed-$N$ Schauder constants.

\section{A priori estimates for the deformation}
\label{n3:apriori:section}

We bound the height, boundary slopes, and conformal variable before
using classical regularity. Fix the prepared exterior of
Proposition~\ref{n3:domains:prepared-domain}, its thresholds and collars,
and $\epsilon>0$. We use the variables and four stages of
Proposition~\ref{n3:deformation:homotopy}. A solution is smooth on a finite
truncation and satisfies $t\ge-\epsilon$, so the estimates include
putative floor-contact solutions. No Hessian bound or upper bound for
$Z$ is assumed.

A quantity denoted by $\Pi_N$
is bounded by $C(1+N)^A$, with $C,A$ independent of the truncation,
the homotopy parameter, and $N$.  It can increase at successive uses.
The fixed data, collars, and $\epsilon$ can enter $C,A$.  A constant
denoted by $C(N)$ can depend arbitrarily on these fixed data and on
$N,\ell,\tau$; it is still independent of the truncation and the
homotopy parameter.  Constants explicitly described as fixed do not
depend on $N$.  Recall that
\begin{equation}
 \ell=e^{-\epsilon N},\qquad \tau=\ell^{3/2},\qquad
 \ell\le l\le e,\qquad 0\le p\le N,
 \qquad \frac{4d}{4+N}\le\chi\le d.
 \label{n3:apriori:parameter-comparisons}
\end{equation}
In particular $\Pi_N\tau/\ell\to0$. This small factor will absorb
only terms whose constants are polynomial.

\subsection{Sublevels and exclusion of boundary layers}

The limiting height will be only lower semicontinuous.  We first show
how an expansion bound on its lower tests controls every exterior
support of a closed sublevel, including a plateau.  This local step
will also apply to the four-dimensional constructions.

\begin{lemma}[Exterior supports of a closed sublevel]
\label{n3:apriori:sublevel-support}
Let $U$ be an open subset of a smooth Riemannian manifold of dimension
$d\in\{3,4\}$, let $Q$ be a smooth symmetric tensor on $U$, and let
$u:U\to\R$ be finite and lower semicontinuous.  Fix $k<k'$ and
$q_0\in\R$.  A smooth lower test of $u$ at $x$ is a smooth function
$\phi$ with $\phi(x)=u(x)$ and $\phi\le u$ near $x$.
Suppose every such test with $u(x)<k'$ and $\dd\phi(x)\ne0$ satisfies
\[
 \mathcal E_Q[\phi](x)\le q_0,
\]
where $\mathcal E_Q$ is the expansion in
Equation~\eqref{n3:domains:test-expansion} of the regular level through
$x$, oriented toward increasing $\phi$.  Equivalently, subtract
$\phi(x)$ before applying that formula.
Then every smooth exterior support of the relatively closed set
$D_k=\{x\in U:u(x)\le k\}$, at a frontier point in $U$, has
expansion at most $q_0$.
\end{lemma}

\begin{proof}
For integers $j\ge1$ define
\begin{equation}\label{n3:apriori:sigmoid}
 S_j(s)=\bigl(1+\exp[-j^2(s-k-j^{-1})]\bigr)^{-1},
 \qquad v_j=S_j\circ u.
\end{equation}
These functions are lower semicontinuous and take values in $(0,1)$.
Their lower relaxed limit $\liminf_{j\to\infty,\,y\to x}v_j(y)$
is the function $I_k(x)$ equal to zero on $D_k$ and one on its
complement.  Indeed, at $x\in D_k$ the constant
sequence gives $v_j(x)\le S_j(k)\to0$, while $v_j\ge0$ everywhere.
If $u(x)>k$, lower semicontinuity gives a neighborhood on which
$u\ge k+\delta$ for some $\delta>0$, and $v_j\to1$ uniformly there.

Let $\phi$ define an exterior support of $D_k$ at $x$:
$\phi(x)=0$, $\dd\phi(x)\ne0$, and $\phi\le0$ on $D_k$ near $x$.
Shrink to a small normal-coordinate ball compactly contained in $U$ on which
$|\phi|<1/2$.  Then $\phi$ is a lower test of $I_k$ at $x$.
Subtract a small positive multiple of the fourth power of the distance
from $x$ to make this contact strict without changing its two-jet;
write the resulting function as $\phi_0$.
Minimize $v_j-\phi_0$ on the closed ball.  The preceding relaxed-limit
description, the strict contact, and $v_j(x)\to0$ show that the
minimizers $x_j$ lie in its interior for large $j$, satisfy $x_j\to x$,
and have contact values $v_j(x_j)\to0$.  Thus $\phi_0+a_j$ is a
lower test of $v_j$ at $x_j$, where $a_j\to0$.

For each fixed large $j$, restrict to a neighborhood of $x_j$ on
which $\phi_0+a_j$ takes values in $(0,1)$.  The increasing function
$S_j^{-1}$ makes $S_j^{-1}(\phi_0+a_j)$ a smooth lower test of $u$.
Its contact height is
\[
 u(x_j)=k+j^{-1}
       +j^{-2}\log\frac{v_j(x_j)}{1-v_j(x_j)}
       \le k+j^{-1}<k'
\]
for all sufficiently large $j$.  Its gradient is nonzero, since
$\dd\phi_0(x_j)\to\dd\phi(x)\ne0$.
An increasing smooth composition preserves the oriented level expansion:
it preserves the unit normal, and for $\theta'>0$,
\[
 \Hess(\theta\circ\phi_0)
  =\theta'\Hess\phi_0
       +\theta''\dd\phi_0\otimes\dd\phi_0,
 \qquad
 \frac{\tr_{(\dd\phi_0)^\perp}\Hess(\theta\circ\phi_0)}
 {|\dd(\theta\circ\phi_0)|}
 =\frac{\tr_{(\dd\phi_0)^\perp}\Hess\phi_0}{|\dd\phi_0|}.
\]
The hypothesis therefore gives
$\mathcal E_Q[\phi_0+a_j](x_j)\le q_0$.
Letting $j\to\infty$ yields $\mathcal E_Q[\phi](x)\le q_0$,
because $\phi_0$ and $\phi$ have the same two-jet at $x$.
\end{proof}

For notation in this section, let $\mathcal B_c$ be the full closed
maximal black region at threshold $c$, and let $\mathcal W_c$ be the
corresponding white region with $\mathcal B_{c_b}$ held fixed.  These
are the families used in Proposition~\ref{n3:domains:prepared-domain}.

\begin{lemma}[Uniform separation from the wrong height]
\label{n3:apriori:sublevels}
In the first two homotopy stages, the following holds.  If
$Q\subset\bar\Omega$ is compact and disjoint from $\mathcal B_{c_b}$,
there are $\eta_Q>0$ and $N_Q$ such that
\begin{equation}
 h\ge b_-+\eta_Q\quad\hbox{on }Q
 \label{n3:apriori:lower-height}
\end{equation}
for every $N\ge N_Q$, every sufficiently large allowed truncation,
and every admissible or floor-contact solution in those stages.
If $Q$ is disjoint from $\mathcal W_{c_w}$, the analogous conclusion is
$h\le b_+-\eta_Q$.  The compact sets are relative to the full closure:
they may contain faces of the opposite color.
\end{lemma}

\begin{proof}
We prove the lower assertion.  It is enough to consider an arbitrary
sequence $N_i\to\infty$, $R_i\to\infty$ and solutions in either of the
first two stages.  The elementary height and end estimates in
Lemmas~\ref{n3:deformation:height} and~\ref{n3:deformation:end-height} give
\begin{equation}
 |h_i|\le C_0,\qquad
 |h_i|\le C(r^{-\gamma}-R_i^{-\gamma})\quad\hbox{on the end},
 \qquad 3/4<\gamma<1.
 \label{n3:apriori:elementary-height}
\end{equation}
Near each white face extend $h_i$ to the other side by the constant
$b_+$.  Extend the fixed geometric coefficients smoothly there.
The extended functions are continuous.  Define the lower relaxed limit
\[
 \underline h(x)=\lim_{j\to\infty}
       \inf\{h_i(y):i\ge j,\ \operatorname{dist}(x,y)<1/j\}.
\]
It is lower semicontinuous, bounded, and satisfies the limiting end
bound.  No equicontinuity of $h_i$ is asserted.

Choose $b_-<k'<\min\{0,b_-+\eta\}$, where $\eta$ is the low-height
sign range of the added term $D_0$.  Suppose that a smooth function
$\phi$ touches $\underline h$ from below at $x$, with
$\underline h(x)\le k'$ and $d\phi(x)\ne0$.  Exclude for the moment
the black faces.  Subtracting a small fourth-order localization term
from $\phi$ makes the contact strict without changing its two-jet.
Minimization on a small closed ball then produces, after passing to a
subsequence, points $x_i\to x$ and constants $c_i\to0$ such that
$\phi+c_i$ touches $h_i$ from below
at $x_i$, and $h_i(x_i)\to\underline h(x)$.  The boundary of this
ball cannot contain $x_i$ by strictness.  Nor can $x_i$ lie on a white
face or on the extended side: their values there are $b_+>k'$.
Thus every sufficiently large $i$ uses the interior trace equation.

At these contacts,
\[
 \nabla f_i=\tau_i^{-1}\nabla\phi,
 \qquad
 d_i\le\frac{\tau_i^2}{\ell_i^2|d\phi(x_i)|^2}
          \longrightarrow0,
 \qquad a_i\longrightarrow1,
 \qquad \chi_i\longrightarrow0.
\]
Here $a_i=l_i|\nabla f_i|/\sqrt{D_i}$ is the deformation variable.
Positivity of $A_{\chi_i}$ and the Hessian
inequality at a lower contact imply
\[
 \tr_{A_{\chi_i}}K+
 \frac{l_i}{\tau_i\sqrt{D_i}}
                 \tr_{A_{\chi_i}}\Hess\phi
 \le F_i(x_i).
\]
In the second stage $D_0\le0$ at these values; in the first stage it
is absent.  Consequently $\limsup F_i(x_i)\le k'+C_b$.
Passing to the limit gives
\begin{equation}
 \frac{\tr_{(d\phi)^\perp}\Hess\phi}{|d\phi|}
           +\tr_{(d\phi)^\perp}K\le k'+C_b.
 \label{n3:apriori:relaxed-test}
\end{equation}
This is the expansion of the level surface of $\phi$, oriented toward
increasing $\phi$.  The only Hessian used in this passage is the
fixed test Hessian.  Removing the localization proves the assertion
for non-strict contacts as well.

Fix $b_-<k<k'$ and put $E_k=\{\underline h\le k\}$.
Apply Lemma~\ref{n3:apriori:sublevel-support} with
$u=\underline h$, tensor $Q=K$, and $q_0=k'+C_b$.
The function is finite and lower semicontinuous, and
Equation~\eqref{n3:apriori:relaxed-test} supplies the required bound
for every nonzero-gradient lower test below $k'$.  We apply the lemma
locally away from black faces, using the constant extensions to obtain
open neighborhoods across white faces.  Every exterior support of
$E_k$ there consequently has expansion at most $k'+C_b$, with no
regularity assumption on the sublevel.

This argument also rules out a hidden layer at a white face.
At every point strictly on the added side, the relaxed limit equals
$b_+>k$.  Thus $E_k$ lies locally in the original closed exterior,
and the reversed white face is an exterior support of $E_k$ at any
contact there, with expansion
\[
 -H_B+P_B=-c_w>0.
\]
For $k'$ sufficiently close to $b_-$, this is strictly greater than
$k'+C_b$.  The high extension ensured that all preceding low-value
tests came from the interior equation, so the support contradiction
applies at the face itself.  Hence $E_k$ avoids every white face.
The end estimate makes $E_k$ compact.

Adjoin the entire closed region $\mathcal B_{c_b}$.
At a point on its smooth black frontier,
every exterior support of the union also contains the black region
locally.  The tangency comparison therefore bounds its expansion by
$c_b<k'+C_b$.  At all other frontier points the preceding sublevel
argument applies.  The resulting compact closed set
contains the required original inner barrier with a collar, because
the black region already does.  It avoids the outer barriers.
Lemma~\ref{n3:domains:weak-support} places it in a smooth black trapped
region at any threshold $c''>k'+C_b$ sufficiently close to $c_b$.
It is consequently contained in $\mathcal B_{c''}$.

Finally let $Q$ be as in the lemma.  Hausdorff right continuity at
$c_b$ gives a $\delta>0$ such that
$\mathcal B_{c_b+\delta}\cap Q=\varnothing$; decrease $\delta$ so
that the enlarged threshold is still negative and all barriers
remain strict.  Choose $k,k',c''$ with
\[
 b_-<k<k',\qquad k'+C_b<c''<c_b+\delta.
\]
The construction implies $E_k\cap Q=\varnothing$ for every sequence
under consideration.  If the uniform estimate $h>k$ on $Q$ failed,
choose violating solutions with $N_i\to\infty$, $R_i\to\infty$ and
$x_i\in Q$, $h_i(x_i)\le k$.  A subsequence has $x_i\to x\in Q$
and $\underline h(x)\le k$, a contradiction.  This gives
\eqref{n3:apriori:lower-height}, after decreasing its positive margin.

For the upper assertion apply the entire argument to
$(-h,-K,-C)$, interchanging colors.  The reversed black faces have
strictly positive white expansion $-c_b$.  Their exclusion and
compactness give positive distance from the black region before
applying the white version of Lemma~\ref{n3:domains:weak-support}.
Adjoining the full white region is therefore legitimate in the
fixed black complement.  Empty white regions and missing face types
cause no change: only an actually nonempty sublevel would supply a
trapped seed and hence a contradiction.
\end{proof}

For the truncation quantifier, start with any $R_{\min}(N)\to\infty$
that satisfies the outer estimates of
Section~\ref{n3:deformation:section}, and enlarge it when necessary.
Failure of any asserted estimate for arbitrarily large $N$ and
$R\ge R_{\min}(N)$ supplies exactly the sequence excluded above.
Only finitely many compact sets will be needed below, so their
thresholds can be chosen simultaneously.

The weak-set argument has kept low heights near black faces and high
heights near white faces. We now sharpen this separation into a
nonzero signed boundary slope. That slope makes the otherwise
uncontrolled inner flux small enough for the mass estimate.

\subsection{Signed collar barriers}

\begin{lemma}[Boundary slopes]
\label{n3:apriori:collars}
There are fixed constants $c,C_*>0$ such that, for sufficiently large
$N$ and the allowed truncations, the first two stages satisfy
\begin{equation}
 \begin{array}{ll}
 c/\tau\le\partial_\nu f\le C_*/\tau&\text{on black faces},\\[2pt]
 c/\tau\le-\partial_\nu f\le C_*/\tau&\text{on white faces}.
 \end{array}
 \label{n3:apriori:signed-slopes}
\end{equation}
During the third stage one still has
$|\partial_\nu f|\le C_*/\tau$.  On every fixed compact region needed
for the final scalar-curvature estimate, the first two stages satisfy
$b_-\le h\le b_+$.
\end{lemma}

\begin{proof}
Write $s$ for the smooth leaf parameter in an outward collar and
$\nu_s=\nabla s/|ds|$.  All geometric coefficients of this fixed collar
are bounded, with $|ds|$ bounded above and below.  At a comparison
contact with $h_0=b_-+\psi(s)$, $\psi'>0$, the shared gradient gives
\begin{equation}
 d\le C\frac{\tau^2}{\ell^2(\psi')^2},\qquad
 1-a\le d,\qquad N\chi\ge c_1d\quad(N\ge4).
 \label{n3:apriori:contact-degeneracy}
\end{equation}
If the slopes are in a fixed positive interval, the first bound is
$d\le C\ell$.  None of these comparisons uses an upper bound for $Z$.

The trace operator on this test is
\begin{align}
 &\tr_{A_\chi}\left(K+
       \frac{l}{\tau\sqrt D}\Hess h_0\right)\notag\\
 &\quad=P_s+aH_s+\chi K(\nu_s,\nu_s)
       +a\chi\left(
           \frac{\Hess s(\nu_s,\nu_s)}{|ds|}
                      +|ds|\frac{\psi''}{\psi'}\right),
 \label{n3:apriori:collar-trace}
\end{align}
where $P_s$ is the tangential trace of $K$ on the leaf.  Thus it is
the leaf expansion $H_s+P_s$, with an error of absolute value at most
$C d$, plus the displayed logarithmic-slope term.  The same formula
for a negative slope has $a$ replaced by $-a$ in the Hessian terms.
This is an evaluation at the contact variables; it does not
differentiate the unknown $t$.

On a black collar $H_s+P_s\ge c_b$ and $C=C_b-k_Bs$.
Choose small positive slopes so that $\psi(0)=0$ and
$\psi(s)\le k_Bs/2$.  At the outer end of a sufficiently short fixed
collar, Lemma~\ref{n3:apriori:sublevels} lets this test lie below the
solution.  In the compatible direction the add-on $D_0$ vanishes.
The residual of the trace equation on the test is therefore at least
\begin{equation}
 k_Bs/2-Cd+a\chi|ds|\frac{\psi''}{\psi'}.
 \label{n3:apriori:lower-collar-residual}
\end{equation}
Choose a smooth nonnegative $\beta_N(s)$ equal to $AN$ for
$0\le s\le1/N$, supported in $0\le s\le2/N$, and satisfying
$\int\beta_N\le2A$.  Define
\[
 \psi'(s)=m\exp\left(\int_0^s\beta_N(q)\dd q\right).
\]
Choose the fixed $A$ large enough that the final term of
\eqref{n3:apriori:lower-collar-residual} dominates $Cd$ on $s\le1/N$,
using \eqref{n3:apriori:contact-degeneracy} and $a\ge1/2$ there.
On $s\ge1/N$, the first term dominates $Cd\le C\ell$ for large $N$;
the taper contributes with the good sign.  Since the slope multiplier
is at most $e^{2A}$, a fixed sufficiently small $m>0$ enforces all
the previous small-slope requirements.  This constructs a strict
lower barrier with slopes bounded above and below independently of
$N$.

An upper barrier from $b_-$ is constructed in the same collar with
\[
 \psi'(s)=M\exp\left(-\int_0^s\beta_N(q)\dd q\right).
\]
The smooth leaf expansion differs from its boundary value by $O(s)$.
Choose the fixed $M$ so large that
$\psi(s)\ge C_2s$ dominates this variation, the offset variation,
and the height bound at the outer collar end.  The logarithmic-slope
term is now negative.  On $s\le1/N$ it dominates $Cd$; on the rest
of the collar the negative margin proportional to $s$ does so.
The upper barrier is strict at every possible positive contact.
Its slopes lie in $[Me^{-2A},M]$, an interval independent of $N$.

To check the comparison signs, at a negative interior minimum of
$h-h_0$ the Hessian of $h$ dominates that of $h_0$, the gradients
and $Z$ agree, and hence so do $t,l,A_\chi$.  The right side is
strictly increasing in the height.  Thus a strict positive residual
for a lower barrier contradicts the equation.  At a positive maximum
the inequalities reverse, proving the upper-barrier assertion.
Taking the one-sided derivative at $s=0$ proves the black part of
\eqref{n3:apriori:signed-slopes}.  Replacing $(h,K,C)$ by
$(-h,-K,-C)$ proves the white part.

In the third stage the boundary height is the assigned interpolated
constant, say $b_\zeta$.  Construct tests
$b_\zeta+\psi(s)$ and $b_\zeta-\psi(s)$ with large positive
$\psi'$ and $(\log\psi')'=-\beta_N$.  At $s=0$, the two directional
inequalities of Proposition~\ref{n3:deformation:homotopy} say precisely
that the positive-slope test has nonpositive limiting residual and
the negative-slope test nonnegative limiting residual.  At every
comparison contact, the fixed positive barrier slope and
Equation~\eqref{n3:apriori:contact-degeneracy} give $d\le C\ell$, so
$a\ge1/2$ for large $N$.  On $1/2\le|w|\le1$, the first
$w$-derivatives of $A_\chi$ in
Equation~\eqref{n3:deformation:axial-matrices} are bounded independently
of $0\le p\le N$.  The other collar terms have fixed smooth
coefficients on the bounded height range.  Thus, at fixed boundary
height the coefficient and geometric errors are $O(s+d+\chi)$,
since $|w-(\pm\nu_s)|=1-a\le d$ and the radial comparison segment stays
in this range.  Changing the test height supplies, by height
monotonicity, a residual of magnitude at least $\psi(s)$ with the
required sign.  Large fixed slopes dominate the $O(s)$ error and
the outer-end height bound.  The logarithmic-slope term has the
required sign for both tests, as follows from
\eqref{n3:apriori:collar-trace}; it dominates the $O(d)$ error on the
initial $1/N$ interval.  The remaining margin dominates it beyond
that interval.  This proves the two-sided upper derivative bound,
without using sublevel exclusion in the third stage.

Finally, on each own-color collar the lower barrier just constructed
gives the required one-sided height bound.  On the rest of a fixed
compact region use Lemma~\ref{n3:apriori:sublevels}; the opposite height
bound on an own-color collar follows from that lemma as well, since
the closed collar is disjoint from the other region.  A finite cover
proves $b_-\le h\le b_+$ on the compact regions in the statement.
\end{proof}

\subsection{A trace estimate with polynomial constants}

The collar trace estimate must have polynomial constants: it will
control the boundary flux both here and in the mass estimate.
We therefore prove it before comparing with the ordinary product graph.

\begin{lemma}[Weighted and unweighted collar traces]
\label{n3:apriori:trace}
In the first two stages let $\mathcal C$ be a fixed union of disjoint
collars of the inner faces, shortened if necessary.  For every
nonnegative smooth function $\varphi$ supported in those collars,
\begin{align}
 \int_B L\varphi\dd A_g
 &\le\Pi_N\int_{\mathcal C}u
       \bigl[(1+\sqrt{\mathcal T})\varphi
                          +|d\varphi|_{A_\chi}\bigr]\dd V_g,
 \label{n3:apriori:weighted-trace}\\
 \int_B\varphi\dd A_g
 &\le\Pi_N\int_{\mathcal C}\sqrt D
       \bigl[(1+\sqrt{\mathcal T})\varphi
                          +|d\varphi|_{A_\chi}\bigr]\dd V_g.
 \label{n3:apriori:unweighted-trace}
\end{align}
The estimates also hold without a support condition on $\varphi$
if the right sides include a fixed collar cutoff equal to one on
$B$, with its derivative included in the coefficient of $\varphi$.
\end{lemma}

\begin{proof}
For $\sigma>0$, write $e_f=\nabla f/\sigma$, the axis denoted by
$e$ in Section~\ref{n3:deformation:section}. Set
\[
 W=l\bar\nabla f=\frac{l\nabla f}{D}=\sqrt d\,a e_f.
\]
Its norm in $\bar g$ is $a\le1$.  For any covector $\xi$,
\begin{equation}
 |\xi(W)|\le |\xi|_{A_d}
                  \le\sqrt{1+N/4}\,|\xi|_{A_\chi}.
 \label{n3:apriori:W-gradient}
\end{equation}
Direct differentiation, using $uW=Lw$, gives
\begin{align}
 \frac{\Div(uW)}u
 &=\sqrt d\left[
       \tr_{A_d}H^f+(dp+p_L)a\partial_{e_f}t\right],
 \label{n3:apriori:W-divergence}\\
 \frac{\Div(\sqrt D\,W)}{\sqrt D}
 &=\sqrt d\left[
       \tr_{A_d}H^f+dp\,a\partial_{e_f}t\right].
 \label{n3:apriori:W-divergence-unweighted}
\end{align}
For example,
$\Div w=\tr_{A_d}H^f+dp\,a\partial_{e_f}t$,
while differentiating $L$ adds $p_La\partial_{e_f}t$ to the first
formula.  Coercivity in Lemma~\ref{n3:deformation:coercivity}
bounds both right sides by $\Pi_N(1+\sqrt{\mathcal T})$.
More explicitly, writing $q=(K+H^f)_{e_fe_f}$, the normal Hessian
term has coefficient $d^{3/2}$, so it is bounded by $\sqrt d|q|+C$.
The other Hessian terms are transverse traces, and the gradient
term is at most $(2N+2)\sqrt d|\partial_{e_f}t|$.
All are controlled by the stated coercivity estimate.  This also
establishes the formulas and bounds at zero gradient by their smooth
direction-free extensions.

Let $\varepsilon_B=1$ on a black collar and $-1$ on a white collar.
The signed slopes give $w_\nu=\varepsilon_B a$ on $B$ and
$a\ge1/2$ for large $N$.  Thus
\[
 uW_\nu=\varepsilon_B La,
 \qquad \sqrt D\,W_\nu=\varepsilon_B a.
\]
The outward integration normal is $-\nu$.
Apply the divergence theorem to $-\varepsilon_B uW\varphi$, with
a fixed collar cutoff.  Its boundary flux is $La\varphi$.
Equations~\eqref{n3:apriori:W-gradient} and
\eqref{n3:apriori:W-divergence}, and the fixed cutoff derivative, give
\eqref{n3:apriori:weighted-trace}.  Replace $u$ by $\sqrt D$ and use
\eqref{n3:apriori:W-divergence-unweighted} for
\eqref{n3:apriori:unweighted-trace}.  No factor $\ell^{-1}$ was used.
\end{proof}

\subsection{High-level energy and the first upper bound}

The signed slopes and polynomial trace estimate are now available.
We use them before allowing any constants that depend arbitrarily on
$N$. The first stage contains the drift in the divergence equation. We first
bound this stage uniformly for $\lambda\in[0,1]$.
There is a number $M_0=M_0(N)>0$ such that
\begin{equation}
 \Xi\ge\delta_0\mathcal T+\rho+
                      \tfrac12\delta_0\tau a\sigma
          \quad\hbox{on }\{Z>M_0\}.
 \label{n3:apriori:high-source}
\end{equation}
Indeed, outside the support of $m_0$ this is immediate.  On its
support, $-\epsilon\le t\le-\epsilon+1/N$, so $l$ is comparable
to $\ell$.  The identity $D=e^{8(Z-t)}$ shows that $a\sigma$ tends
uniformly to infinity as $Z\to\infty$ in that band.  The penalty
$\mathcal P$ is bounded on the whole exterior by $\Pi_N$.
It is therefore absorbed by $\delta_0\tau a\sigma/2$ above a fixed
$M_0(N)$.  This choice is independent of $R$ and $\lambda$.

Choose a smooth nondecreasing $k_0$ equal to zero on $Z\le M_0$,
equal to one on $Z\ge M_0+1$, and with bounded derivative.  Test
$\Div V_\lambda=u\Xi$ by $k_0(Z)$.  Its outer boundary term
vanishes since $Z=0$.  The drift cross term is bounded by
\[
 u k_0'\left|\langle dZ,K(w,\cdot)\rangle_{A_\chi}\right|
 \le 2u k_0'|dZ|_{A_\chi}^2+C u k_0'|K|^2.
\]
The last term has bounded integral: it is supported on a fixed
compact set and on $M_0\le Z\le M_0+1$, where the floor and
$D=e^{8(Z-t)}$ bound $u$ and $\sigma$ by $C(N)$.

At a black face $F-P_B=H$, $w_\nu=a$; at a white face
$F-P_B=-H$, $w_\nu=-a$.  In both cases
\[
 -H+w_\nu(F-P_B)=(a-1)H.
\]
The additional omitted-drift term in the homotopy boundary condition
has absolute value at most $C\chi\le Cd$, since
$A_\chi K(w,\cdot)$ has normal component
$\chi w_\nu K(\nu,\nu)$.  The remaining term is $-N_Bd$, and the
boundary multiplier lies between zero and one.  Hence
\begin{equation}
 |(V_\lambda)_\nu|\le Cud=CL\sqrt d
                       \le C\frac\tau\ell L\quad\hbox{on }B,
 \label{n3:apriori:small-boundary-flux}
\end{equation}
where the last inequality uses the lower slope in
\eqref{n3:apriori:signed-slopes}.  Integration now gives
\begin{align}
 &\int_{\Omega_R}u k_0
       \left(\delta_0\mathcal T+\rho+
                      \tfrac12\delta_0\tau a\sigma\right)\dd V_g
       +2\int_{\Omega_R}u k_0'|dZ|_{A_\chi}^2\dd V_g\notag\\
 &\hspace{18mm}\le C(N)+C\frac\tau\ell
                                  \int_B Lk_0\dd A_g.
 \label{n3:apriori:high-energy-before-trace}
\end{align}

Apply \eqref{n3:apriori:weighted-trace} with the fixed collar cutoff.
The boundary cost is at most
\[
 \Pi_N\frac\tau\ell\int_{\mathcal C}u
       \left[k_0(1+\sqrt{\mathcal T})
                              +k_0'|dZ|_{A_\chi}\right]\dd V_g.
\]
Let $\rho_{\mathcal C}>0$ be the minimum of $\rho$ on the closed
collars.  Pointwise,
$1+\sqrt{\mathcal T}\le C(\rho+\delta_0\mathcal T)$ there, with
$C$ fixed.  Since $\Pi_N\tau/\ell\to0$, the first summand is
absorbed in the first integral of
\eqref{n3:apriori:high-energy-before-trace}.  Young's inequality bounds
the derivative summand by
\[
 u k_0'|dZ|_{A_\chi}^2+
          C(\Pi_N\tau/\ell)^2u k_0'.
\]
The second term again has bounded compact integral because it lies
in the transition strip of $k_0$.  We obtain
\begin{equation}
 \int u k_0(\mathcal T+\rho+\tau a\sigma)\dd V_g
       +\int u k_0'|dZ|_{A_\chi}^2\dd V_g\le C(N).
 \label{n3:apriori:high-energy}
\end{equation}
The smallness requirement on $N$ in this absorption is independent
of $M_0$; only its resulting right-hand constant depends on $M_0$.

The polynomial boundary absorption is complete. We may now use
arbitrary fixed-$N$ comparison constants. Write
$\Gamma$ for the ordinary product graph of $f$ in $g+(\dd b_0)^2$,
and $\dd\Gamma=\sqrt{1+\sigma^2}\dd V_g$.  The measures
$\dd\Gamma$ and $\sqrt D\dd V_g$ are comparable by constants
depending only on $\ell$.  On every fixed compact set $Q$,
\begin{equation}
 \int_{\Gamma\cap\pi^{-1}(Q)}e^{2Z}\dd\Gamma\le C_Q(N),
 \label{n3:apriori:L2}
\end{equation}
where $\pi$ is projection to the base.  To check this without an
upper bound for $t$, observe that
\[
 e^{2Z}=e^{2t}D^{1/4},\qquad u=le^{2t}\sqrt D,
 \qquad D^{1/4}\le C(N)(1+\tau a\sigma).
\]
The last inequality follows by splitting $l\sigma\le1$ and
$l\sigma>1$, using $\ell\le l\le e$; on the second set
$a\ge1/\sqrt2$ and the right side grows linearly in $\sigma$,
whereas $D^{1/4}$ grows at most as $C\sqrt\sigma$.
Thus the high-level part of \eqref{n3:apriori:L2} follows from
\eqref{n3:apriori:high-energy} and the positive minimum of $\rho$ on
$Q$.  On $Z\le M_0+1$, both $\sigma$ and $e^{2Z}$ are bounded
by the floor identity, which controls the remaining compact integral.

\subsection{Sobolev inequality and iteration on the ordinary graph}

The high-level estimate gives integral control, but degree requires a
pointwise upper bound for $Z$. The graph Sobolev inequality and
iteration supply that bound. The earlier polynomial trace estimate,
not the constants in this iteration, will control the mass limit.

The comparisons in this subsection are uniform in the solution,
$R$, and the homotopy parameter, but may depend arbitrarily on $N$.
If $d_0=(1+\sigma^2)^{-1}$, then
\[
 |d\varphi|_{A_{d_0}}=|\nabla_\Gamma\varphi|,
 \qquad
 \frac{d}{d_0}=\frac{1+\sigma^2}{1+l^2\sigma^2}.
\]
The last ratio is bounded above and below by constants depending
only on $\ell$; Equation~\eqref{n3:apriori:parameter-comparisons}
therefore compares $A_{d_0},A_d,A_\chi$ without any bound for $Z$.

\begin{lemma}[Local graph Sobolev inequality]
\label{n3:apriori:sobolev}
For a smooth function $\omega$ supported over a fixed compact base
patch, which may meet an inner face,
\begin{equation}
 \|\omega\|_{L^6(\dd\Gamma)}^2
 \le C(N)\int_\Gamma
       \left[|\nabla_\Gamma\omega|^2+
                            (1+\mathcal T)\omega^2\right]\dd\Gamma.
 \label{n3:apriori:sobolev-equation}
\end{equation}
\end{lemma}

\begin{proof}
Extend the fixed base patch smoothly and embed a compact enlargement
isometrically in Euclidean space by Nash's embedding theorem
\cite[Theorem~2, p.~59]{Nash1956}.  Its product with the vertical
line gives an isometric Euclidean immersion of the graph.  The
second fundamental form of the base embedding is bounded, so its
contribution to the graph mean curvature vector is bounded.
The scalar mean curvature in the product is
\[
 H_\Gamma=
 \frac{\sqrt D}{l\sqrt{1+\sigma^2}}
       \left(\tr_{e_f^\perp}H^f+d_0H^f(e_f,e_f)\right).
\]
The prefactor is bounded by $C(N)$, and
$d_0/\sqrt d=\sqrt{1+l^2\sigma^2}/(1+\sigma^2)\le e$.
Coercivity therefore gives
$|\boldsymbol H_\Gamma|\le C(N)(1+\sqrt{\mathcal T})$.
This estimate involves controlled quadratic energy, not a pointwise
bound for the Hessian.

The ordinary Michael--Simon inequality \cite{MichaelSimon1973}, in
the compact-support form stated in
\cite[Chapter~4, Section~5, Theorem~5.7, p.~98]{Simon2014}, gives after the
boundary extension described below, for nonnegative $\varphi$,
\[
 \left(\int_\Gamma\varphi^{3/2}\dd\Gamma\right)^{2/3}
 \le C\left[
       \int_\Gamma(|\nabla_\Gamma\varphi|
                          +|\boldsymbol H_\Gamma|\varphi)\dd\Gamma
       +\int_{\partial\Gamma}\varphi\dd A\right].
\]
One can obtain the boundary form directly from the compact-support
version by smoothly extending each smooth graph across its boundary
and cutting off on a collar tending to zero: the cutoff derivative
integral converges to the boundary integral, and the added curvature
integral tends to zero.  This argument is applied to each smooth
graph before the uniform estimate is taken.  The only boundary in
the support is an inner face, where $f$ is constant and
$\dd A=\dd A_g$.  Equation~\eqref{n3:apriori:unweighted-trace} and
the fixed-$N$ graph comparisons bound its contribution by
\[
 C(N)\int_\Gamma
       [(1+\sqrt{\mathcal T})\varphi
                               +|\nabla_\Gamma\varphi|]\dd\Gamma.
\]
Apply the resulting inequality to $\varphi=|\omega|^4$.
Cauchy--Schwarz bounds its right side by
\[
 C(N)\|\omega\|_6^3
 \left(\int_\Gamma
       [|\nabla_\Gamma\omega|^2+(1+\mathcal T)\omega^2]
                                      \dd\Gamma\right)^{1/2}.
\]
The left side is $\|\omega\|_6^4$; division, with the zero case
separate, and squaring prove the assertion.
\end{proof}

To turn the integral bound into an upper bound, let $\eta$ be a
nonnegative cutoff over a fixed compact patch.  For $k$ sufficiently
large, independently of $R$ and the solution, test the first-stage
divergence equation by
\[
 q=\eta^2e^{2kZ}/L.
\]
The positive measure after cancellation of $L$ is
$\sqrt D\dd V_g$.  The differentiated test is
\[
 dq=\frac{e^{2kZ}}L
       [2\eta\,d\eta+\eta^2(2k\,dZ-p_L\,dt)].
\]
Its principal exponential contribution is
$8k\eta^2e^{2kZ}|dZ|_{A_\chi}^2$.  The source is bounded below
by $\delta_0\mathcal T-\Pi_N$ everywhere.
Coercivity and $A_\chi\le A_d$ give
\[
 4p_L|\langle dt,dZ\rangle_{A_\chi}|
 \le\tfrac14\delta_0\mathcal T+C(N)|dZ|_{A_\chi}^2.
\]
Choose $k\ge k_*(N)$ so that the second term is absorbed by the
principal contribution.  The drift terms, with $|K(w,\cdot)|_{A_\chi}
\le |K|$, obey
\begin{align*}
 2k|\langle K(w,\cdot),dZ\rangle_{A_\chi}|
   &\le k|dZ|_{A_\chi}^2+Ck|K|^2,\\
 p_L|\langle K(w,\cdot),dt\rangle_{A_\chi}|
   &\le\tfrac14\delta_0\mathcal T+C(N)|K|^2.
\end{align*}
The cutoff cross terms are handled by the same Cauchy inequalities,
leaving $C(N)k e^{2kZ}(|d\eta|_g^2+\eta^2)$.
The penalty is bounded and contributes to this latter expression.
Finally, Equation~\eqref{n3:apriori:small-boundary-flux} before its last
inequality gives
$|q(V_\lambda)_\nu|\le C\eta^2e^{2kZ}\sqrt d
\le C\eta^2e^{2kZ}$.
Converting to the ordinary graph yields
\begin{align}
 &\int_\Gamma e^{2kZ}\eta^2
       (\mathcal T+k|\nabla_\Gamma Z|^2)\dd\Gamma\notag\\
 &\quad\le C(N)k\int_\Gamma e^{2kZ}
                (\eta^2+|d\eta|_g^2)\dd\Gamma
               +C(N)\int_B e^{2kZ}\eta^2\dd A_g.
 \label{n3:apriori:iteration-energy-boundary}
\end{align}

Use the unweighted trace with $\varphi=e^{2kZ}\eta^2$.  The boundary
term is at most
\[
 C(N)\int_\Gamma e^{2kZ}
   \left[\eta^2(1+\sqrt{\mathcal T}
                   +k|\nabla_\Gamma Z|)+\eta|d\eta|_g\right]
                                                   \dd\Gamma.
\]
Young's inequality absorbs its $\sqrt{\mathcal T}$ term into the
$\mathcal T$ integral and its $k|\nabla_\Gamma Z|$ term into the
$k|\nabla_\Gamma Z|^2$ integral.  The latter absorption costs
$C(N)k\eta^2e^{2kZ}$, which has the allowed size.  Hence
\begin{equation}
 \int_\Gamma e^{2kZ}\eta^2
       (\mathcal T+k|\nabla_\Gamma Z|^2)\dd\Gamma
 \le C(N)k\int_\Gamma e^{2kZ}
                   (\eta^2+|d\eta|_g^2)\dd\Gamma,
 \label{n3:apriori:iteration-energy}
\end{equation}
with constants independent of $k\ge k_*$.

Apply Lemma~\ref{n3:apriori:sobolev} to $\omega=\eta e^{kZ}$ and use
\eqref{n3:apriori:iteration-energy}.  For nested base patches
$U_r\subset U_{r'}$ and a cutoff satisfying
$|d\eta|_g\le C/(r'-r)$, this gives
\begin{equation}
 \|e^Z\|_{L^{6k}(\Gamma|_{U_r})}
 \le\left[C(N)k^2(1+(r'-r)^{-2})\right]^{1/(2k)}
                  \|e^Z\|_{L^{2k}(\Gamma|_{U_{r'}})}.
 \label{n3:apriori:moser-step}
\end{equation}
Here $\Gamma|_U$ denotes the part of the graph projecting to $U$.
Iterate with $k_j=3^jk_*$ and radii decreasing geometrically from
$r'$ to $r$.  The sums $\sum k_j^{-1}$ and $\sum j/k_j$ are finite.
Taking the product in \eqref{n3:apriori:moser-step} therefore gives
\begin{equation}
 S(r):=\sup_{U_r}e^Z
 \le C(N)(r'-r)^{-A(N)}
             \|e^Z\|_{L^{2k_*}(\Gamma|_{U_{r'}})}.
 \label{n3:apriori:moser-high-norm}
\end{equation}
Take all these patches inside a slightly larger fixed patch on
which \eqref{n3:apriori:L2} holds.  Increasing $k_*$ to at least two,
interpolation gives
\[
 \|e^Z\|_{2k_*}
 \le S(r')^{1-1/k_*}
                  \left(\int e^{2Z}\dd\Gamma\right)^{1/(2k_*)}.
\]
Thus $S(r)\le C(N)(r'-r)^{-A(N)}S(r')^{1-1/k_*}$.
For every $\alpha>0$, Young's inequality makes this
\begin{equation}
 S(r)\le\alpha S(r')+
                   C(N,\alpha)(r'-r)^{-A(N)k_*}.
 \label{n3:apriori:moser-interpolation}
\end{equation}
Choose radii increasing geometrically to a fixed outer radius $r_1$,
and then choose $\alpha<2^{-A(N)k_*-1}$.  Iterating
\eqref{n3:apriori:moser-interpolation} has a convergent geometric sum.
Its final remainder $\alpha^jS(r_j)$ tends to zero, because each
individual smooth solution is bounded on the closed outer patch.
The resulting smaller-patch bound is independent of that individual
supremum.  A finite cover proves a uniform compact upper bound for
$Z$, including at every inner face.

Beyond a fixed sphere containing the drift support, the equation is
$\Div(4uA_\chi\nabla Z)=u\Xi$, and $\Xi>0$ above $M_0$.
At an interior maximum above $M_0$ its left side is nonpositive,
because $dZ=0$ and $\Hess Z\le0$.  The outer boundary value is zero.
The compact bound on that fixed sphere consequently extends the
first-stage upper bound to the whole truncation.

\subsection{The remaining stages and the bounded-variable conclusion}

In the second and third stages the drift is zero.  Their source has
the same high-$Z$ positivity as \eqref{n3:apriori:high-source}, with a
uniform threshold for the bounded homotopy range.  By
Lemma~\ref{n3:apriori:collars}, on an inner face
\begin{equation}
 Z=t+\tfrac18\log(1+l^2\sigma^2)
       \le t+\tfrac18\log(1+e^2C_*^2/\tau^2).
 \label{n3:apriori:high-boundary-t}
\end{equation}
Thus sufficiently high boundary $Z$ forces $j(t)=1$.
The prescribed boundary flux is then zero, since $\lambda=0$.
At the face $A_\chi\nu=\chi\nu$, so this is exactly
$\partial_\nu Z=0$.  The maximum principle excludes a high interior
maximum, and the boundary point principle excludes a high maximum
on a smooth inner face with this zero conormal derivative.  Both
principles are applied to the smooth individual solution on its
finite domain; its coefficients are positive definite there.
Locating the maximum requires no uniform ellipticity constant.
The outer value is zero, so these two stages also have a uniform upper
bound.

In the fourth stage $f=0$ and $t=Z$.  When the common source/flux
scale is positive, the same maximum and boundary argument works
with its positive high-level source.  At scale zero the homogeneous
mixed problem and the outer zero Dirichlet value give $Z=0$.
The upper bound is consequently uniform also as the scale tends
to zero.

\begin{proposition}[Bounds before regularity]
\label{n3:apriori:bounds}
Fix the prepared data and $\epsilon>0$.  There are $N_0$ and
$R_{\min}(N)\to\infty$ such that, for each fixed $N\ge N_0$ and
every $R\ge R_{\min}(N)$, all smooth solutions of the four-stage
homotopy with $t\ge-\epsilon$ satisfy
\begin{equation}
 |h|\le C_0,\qquad |f|\le C_0/\tau,\qquad
 -\epsilon\le t\le Z\le C(N),\qquad
 |\nabla f|\le C(N).
 \label{n3:apriori:bounded-variables}
\end{equation}
The constants are uniform in $R$ and in the homotopy parameter.
They include putative floor-contact solutions, which are then
excluded by Proposition~\ref{n3:deformation:floor}.
The fixed signed slopes, the third-stage upper slopes, and the
polynomial trace estimates are those of
Lemmas~\ref{n3:apriori:collars} and~\ref{n3:apriori:trace}.
\end{proposition}

\begin{proof}
It remains to extract the bounds for the conformal variable and the
gradient. Since $Z=t+\log D/8$ and $D\ge1$, the floor gives
$-\epsilon\le t\le Z$.  The upper bound for $Z$ gives
\[
 \sigma^2=\frac{e^{8(Z-t)}-1}{l^2}
       \le\ell^{-2}\bigl(e^{8(C(N)+\epsilon)}-1\bigr).
\]
This proves the asserted fixed-$N$ gradient bound.  The height bounds
are Lemma~\ref{n3:deformation:height}.  The finite collection of
sublevel and collar estimates and the polynomial absorption choose
$N_0$ once.  Enlarge $R_{\min}(N)$ to include the outer estimates and
all fixed compact patches used above. The next section derives Hessian
and continuity estimates from these bounds and uses them to prove
existence.
\end{proof}

\section{Solving the three-dimensional boundary deformation}
\label{n3:existence:section}

The boundary estimates have bounded $f$, $Z$ and $Df$ for every
admissible solution of the continuation equations. We now turn those
bounds into existence. Two points need attention. The divergence
equation for $Z$ contains $|D^2f|^2$, so the bounded variables do not
immediately give a classical elliptic estimate. Also, a compact
solution map must be defined on an open set of arbitrary inputs,
before imposing the admissibility inequality $t>-\epsilon$.

We address these points in order. A gradient-aligned scalar equation
gives both continuity of $Df$ and a local integral bound for $D^2f$.
An oscillation argument then gives continuity of $Z$, after which
ordinary boundary regularity applies. Separately, a scalar Dirichlet
problem can be solved for every bounded input $Z$, even when its
associated $t$ lies below the floor. Freezing that scalar solution in
the divergence equation produces the compact map. The floor exclusion
and the preceding estimates keep its fixed points inside the
admissible set along the homotopy of
Proposition~\ref{n3:deformation:homotopy}.

All regularity constants in this section may depend on the fixed
deformation parameters $N$, $\epsilon$ and $\tau$. They are not used
in any earlier argument requiring polynomial dependence on $N$.
The estimates needed for exhaustion will be uniform in the outer
radius $R$ at fixed $N$.

\subsection{From a bounded gradient to a Hessian integral bound}

The first equation has two equal transverse eigenvalues; its remaining
eigenvector is the gradient of its solution. This structure gives a
gradient estimate before the axial eigenvalue is continuous. The
accompanying Hessian estimate will control the source in the second
equation.

In the next two lemmas, $D$ denotes coordinate derivatives; it is not
the scalar $1+l^2|\nabla f|^2$ from the deformation.  Balls
are Euclidean coordinate balls, and a boundary ball is their
intersection with one side of a smooth boundary.  A family of boundary
patches has bounded geometry here if normal-coordinate charts, their
inverses, the metric and its inverse, and the derivatives through the
orders used below have fixed bounds.  We only apply the lemmas to
smooth metrics and smooth boundaries.  Their doubled metrics will be
Lipschitz and smooth on each side of a single plane.

\begin{lemma}[Axial gradient estimate]\label{n3:existence:gradient}
Let $U'$ be compactly contained in an interior or boundary coordinate
patch $U$ in dimension three.  In a boundary patch assume that $z$ is
constant on the boundary face.  Suppose $z$ is $C^3$, up to that face when present,
$|\nabla z|\le L$, and, almost everywhere,
\begin{equation}\label{n3:existence:axial-equation}
 A_\chi:\Hess z=G,\qquad
 A_\chi=I-(1-\chi)e\otimes e,\qquad
 e=\frac{\nabla z}{|\nabla z|},\qquad
 0<\chi_0\le\chi\le1,
\end{equation}
where $\|G\|_\infty\le Q$.  At zero gradients any measurable unit
axis for which the equation holds is allowed.  Then there are
$\alpha\in(0,1)$ and $C<\infty$, depending only on $\chi_0$, the
patch geometry and the separation of $U'$ from the other patch
boundaries, such that
\begin{equation}\label{n3:existence:gradient-estimate}
 \|\nabla z\|_{C^{0,\alpha}(U')}
 \le C(L+Q).
\end{equation}
For all sufficiently small coordinate balls centered in $U'$,
with intersection with the domain understood at a boundary, one also
has
\begin{equation}\label{n3:existence:hessian-morrey}
 \int_{B_r\cap U}|\Hess z|^2\dd V_g
 \le C(L+Q)^2r^{1+2\alpha}.
\end{equation}
In particular, the constants are independent of continuity moduli of
$\chi$ and $G$.
\end{lemma}

\begin{proof}
We first obtain an $L^p$ Hessian estimate for some $p>2$.
A divergence identity then shows that a gradient nearly attaining its
upper bound forces the solution to be nearly affine. The equation
with a fixed gradient direction improves that affine approximation.
The same steps survive reflection at a constant Dirichlet face,
and iteration gives the stated estimate.

\paragraph{A Cordes estimate above exponent two.}
For a Euclidean unit vector $e$,
\begin{equation}\label{n3:existence:cordes-defect}
 \|I-A_\chi\|_F=1-\chi\le1-\chi_0.
\end{equation}
For a compactly supported smooth function $v$, Plancherel's theorem
gives $\|D^2v\|_2=\|\Delta v\|_2$.  The Calder\'on--Zygmund
operator taking $\Delta v$ to the full array $D^2v$ is bounded on
$L^p$, and interpolation makes its norm arbitrarily close to one as
$p\to2$.  Choose once and for all
\begin{equation}\label{n3:existence:p-zero}
 2<p_0<3,\qquad
 C_{p_0}(1-\chi_0/2)<1.
\end{equation}
Rescale a metric-normalized patch until the coordinate principal
matrix differs from its Euclidean axial matrix by less than
$\chi_0/2$ in Frobenius norm.  Absorbing this difference in the
Laplacian estimate, applying a cutoff, and using
$\|Dv\|_p\le\eta\|D^2v\|_p+C_\eta\|v\|_p$ gives
\begin{equation}\label{n3:existence:local-cordes}
 \|D^2v\|_{L^{p_0}(B_{1/2})}
 +\|Dv\|_{L^{p_0}(B_{1/2})}
 \le C\bigl(\|v\|_{L^{p_0}(B_1)}
       +\|A_\chi:\Hess v\|_{L^{p_0}(B_1)}\bigr).
\end{equation}
The Christoffel terms are bounded first-order terms in this estimate.
The same argument applies when $v=z-L_0$ and $L_0$ is coordinate
affine: its extra right side is bounded by
$C\|Dg\|_\infty|DL_0|$.
The estimate uses bounded measurable principal coefficients only;
the assumed $C^3$ regularity allows all of its applications
without an approximation issue.  The linear estimates just used are
the ordinary Laplacian estimates; see
\cite[Chapters 7 and 9]{GilbargTrudinger2001}.

\paragraph{A nearly maximal gradient forces small Hessian energy.}
The preceding estimate does not yet give continuity of $Dz$. To gain
continuity we use the fact that the axial direction is $\nabla z$
itself. Put $H=\Hess z$ and $P=\nabla z$.  Equation~\eqref{n3:existence:axial-equation}
gives the pointwise vector identity
\begin{equation}\label{n3:existence:gradient-flux}
 A_\chi\nabla\frac{|P|^2}{2}-GP
       =(H-\Delta z\,g)P.
\end{equation}
Indeed $HP=|P|H(e,\cdot)$ and
$G=\Delta z-(1-\chi)H(e,e)$, so the two axial terms cancel.
The identity is also valid at $P=0$.  Taking the divergence of its
right side, rather than differentiating $\chi$, yields
\begin{equation}\label{n3:existence:bochner-flux}
 \Div\left(A_\chi\nabla\frac{|P|^2}{2}-GP\right)
 =|H|^2-(\Delta z)^2+\Ric(P,P).
\end{equation}
Choose $a_0>0$ so small that
$(1+a_0)(1-\chi_0)^2<1$.  Young's inequality and the equation show
\begin{equation}\label{n3:existence:positive-hessian}
 |H|^2-(\Delta z)^2
 \ge c_0|H|^2-C_0G^2,
\end{equation}
where $c_0>0$ depends only on $\chi_0$.

Consider solutions normalized by $|\nabla z|\le1$ on $B_1$.
After making the scale small, let the forcing and the rescaled
geometric errors have size at most $\delta$.  The nonnegative function
$s=1-|\nabla z|^2$ then satisfies
\begin{equation}\label{n3:existence:s-supersolution}
 \Div(A_\chi\nabla s+2G\nabla z)
 \le -2c_0|\Hess z|^2+C\delta
\end{equation}
in a weak sense.  In the Euclidean case the last term can be replaced
by $C\delta^2$.  We allow $C\delta$ to include the geometric errors.

For the compactness step, suppose
the normalized errors tend to zero in a sequence and
$\inf_{B_{1/2}}s\to0$.  The inhomogeneous weak Harnack inequality,
applied to Equation~\eqref{n3:existence:s-supersolution} after dropping
its favorable Hessian term, gives $s\to0$ in measure on $B_{1/2}$.
The boundedness $0\le s\le1$ also gives convergence in every finite
$L^p$ norm there.  To recover the Hessian information, take a cutoff
$\eta$ supported inside $B_{1/2}$ and use
$\eta^2(b-s)_+$ in the weak supersolution inequality.  Ellipticity
and Young's inequality give
\begin{equation}\label{n3:existence:low-truncation}
 \int_{\{s<b\}}\eta^2|Ds|^2
 \le Cb^2\int|D\eta|^2+o(1)
\end{equation}
for each fixed $b>0$ as the errors vanish.  The error fields in
Equation~\eqref{n3:existence:s-supersolution} are bounded and tend to
zero, so their products with $Ds$ are absorbed in deriving this
inequality.  Equation~\eqref{n3:existence:local-cordes} supplies a
uniform $L^2$ bound for $Ds$.  Its integral over $\{s\ge b\}$ tends
to zero in $L^1$ by Cauchy's inequality and convergence in measure.
On $\{s<b\}$ use Equation~\eqref{n3:existence:low-truncation}, then let
$b\downarrow0$.  Thus $Ds\to0$ locally in $L^1$.
Testing Equation~\eqref{n3:existence:bochner-flux} with a nonnegative
compact cutoff and using Equations~\eqref{n3:existence:positive-hessian}
and \eqref{n3:existence:gradient-flux} now gives
\begin{equation}\label{n3:existence:hessian-flat-limit}
 \int_{B_{1/4}}|\Hess z|^2\longrightarrow0.
\end{equation}
Uniform Lipschitz compactness implies that, after subtracting a
constant and taking a subsequence, $z$ tends uniformly on a smaller
ball to an affine function.  Its slope has length one because
$|\nabla z|\to1$ in measure and the Hessians tend to zero.

Consequently, for any prescribed sufficiently small $b_*>0$ there
are fixed $r_*\in(0,1/4)$, $k_*\in(r_*,1)$ and $\delta_*>0$ such
that every normalized solution with errors at most $\delta_*$ has
one of the following properties:
\begin{equation}\label{n3:existence:drop-or-flat}
 \begin{split}
 &\sup_{B_{r_*}}|\nabla z|\le k_*;\qquad\text{or}\\
 &\sup_{B_{r_*}}|z-L_0|\le b_*r_*,
       \qquad \tfrac12\le|DL_0|\le2.
 \end{split}
\end{equation}
If this conclusion failed, solutions whose gradient suprema approach
one and whose errors approach zero would contradict
Equation~\eqref{n3:existence:hessian-flat-limit}.  A gradient-drop
constant can always be weakened toward one to ensure $k_*>r_*$.

\paragraph{Regularity when the direction is fixed.}
The second alternative in \eqref{n3:existence:drop-or-flat} supplies
a nonzero affine slope. The comparison equation for that slope has
a fixed direction, although its axial coefficient remains merely
measurable. Let $e_0$ be a Euclidean unit vector and suppose
\begin{equation}\label{n3:existence:fixed-axis}
 \Delta_{e_0^\perp}Y_0+\chi(x)\partial_{e_0e_0}Y_0=0
 \quad\text{in }B_{3/4},\qquad
 \|Y_0\|_{W^{2,2}(B_{3/4})}\le C_1.
\end{equation}
The derivative $v=\partial_{e_0}Y_0$ belongs to $W^{1,2}$ and,
in distributions, satisfies
\begin{equation}\label{n3:existence:axial-derivative-divergence}
 \Delta_{e_0^\perp}v+
 \partial_{e_0}(\chi\,\partial_{e_0}v)=0.
\end{equation}
This follows by differentiating the distribution $\chi\partial_{e_0}v$
as a whole.  It does not assume a derivative of $\chi$ exists as a
function.  De Giorgi's estimate for this uniformly elliptic
divergence equation, followed by Caccioppoli applied to
$v-v(x_0)$, gives an exponent $\alpha_1\in(0,1)$ and
\begin{equation}\label{n3:existence:axial-derivative-morrey}
 [v]_{C^{0,\alpha_1}(B_{1/2})}\le C,
 \qquad
 \int_{B_r(x_0)}|Dv|^2\le Cr^{1+2\alpha_1}
\end{equation}
for balls in a smaller fixed interior region.  These are the scalar
divergence estimates with ellipticity $\chi_0$; see
\cite[Chapter 8]{GilbargTrudinger2001}.

All components of $DY_0$, not only $v$, are controlled by this
estimate.  In fact,
\begin{equation}\label{n3:existence:poisson-source}
 \Delta Y_0=f_0,\qquad
 f_0=(1-\chi)\partial_{e_0}v,
 \qquad
 \int_{B_r(x_0)}|f_0|\le Cr^{2+\alpha_1}.
\end{equation}
The last inequality follows from Cauchy's inequality and
Equation~\eqref{n3:existence:axial-derivative-morrey}.  Take the Newton
potential of $f_0$ restricted to a slightly larger interior ball;
its difference from $Y_0$ is harmonic in that ball.  For two points
at distance $h$, the contribution to the difference of gradient
potentials from annuli of radius $r\le2h$ is bounded by
$C\sum r^{\alpha_1}\le Ch^{\alpha_1}$.  On the remaining annuli
the difference of gradient kernels is at most $Ch r^{-3}$, so their
contribution is
\begin{equation}\label{n3:existence:potential-annuli}
 C h\sum_{r\ge2h}r^{\alpha_1-1}
 \le Ch^{\alpha_1}.
\end{equation}
The harmonic remainder has its usual interior derivative estimates.
Therefore, decreasing the interior ball if needed,
\begin{equation}\label{n3:existence:model-gradient}
 \|Y_0\|_{C^{1,\alpha_1}(B_{1/4})}\le C.
\end{equation}

\paragraph{Improving the affine approximation.}
We now compare the original equation with the fixed-direction equation.
Assume that, for a coordinate-affine function $L_0$,
\begin{equation}\label{n3:existence:flat-input}
 \|z-L_0\|_{L^\infty(B_1)}\le b,
 \qquad \tfrac14\le|DL_0|\le4,
 \qquad |\nabla z|\le8.
\end{equation}
Let the forcing and the first-derivative metric errors be at most
$b\delta$.  For $Y=(z-L_0)/b$, Equation~\eqref{n3:existence:local-cordes}
on fixed nested balls gives
\begin{equation}\label{n3:existence:normalized-residual}
 \|Y\|_{W^{2,p_0}(B_{7/8})}\le C.
\end{equation}
In particular, this estimate does not require a bound for
$DY$ in the supremum norm.  Since $Dz=DL_0+bDY$, the variable axis
converges in measure to $e_0=DL_0/|DL_0|$ as $b\to0$.  More
quantitatively, the coefficient error $E$ caused by replacing this
axis by $e_0$ is bounded and satisfies
\begin{equation}\label{n3:existence:axis-error}
 \|E\|_{L^q(B_{7/8})}\le Cb^{\vartheta}+Cb\delta,
 \qquad q=\frac{2p_0}{p_0-2},
 \qquad \vartheta=\frac{p_0-2}{2}>0.
\end{equation}
Indeed the unit-direction map is Lipschitz where
$|bDY|<|DL_0|/2$, and elsewhere its difference is bounded; the
$L^{p_0}$ bound in Equation~\eqref{n3:existence:normalized-residual}
and interpolation give Equation~\eqref{n3:existence:axis-error}.
The same estimate includes the metric principal-part error.
H\"older's inequality, with $1/2=1/q+1/p_0$, shows that
\begin{equation}\label{n3:existence:small-fixed-axis-source}
 \bigl\|[I-(1-\chi)e_0\otimes e_0]:D^2Y\bigr\|_{L^2(B_{3/4})}
 \le C(b^{\vartheta}+\delta).
\end{equation}

Solve for a zero-Dirichlet correction $w$ on $B_{3/4}$ with the
left side of Equation~\eqref{n3:existence:small-fixed-axis-source} as
its source.  This does not require a nondivergence Green function.
The convex Dirichlet Hessian inequality
$\|D^2w\|_2\le\|\Delta w\|_2$ and
Equation~\eqref{n3:existence:cordes-defect} make the Laplace solution
operator a contraction on $W^{2,2}\cap W^{1,2}_0$.  Hence
\begin{equation}\label{n3:existence:fixed-axis-correction}
 \|w\|_{W^{2,2}(B_{3/4})}
 \le C\chi_0^{-1}(b^{\vartheta}+\delta),\qquad
 \|w\|_\infty\le C(b^{\vartheta}+\delta).
\end{equation}
The last implication uses $W^{2,2}\hookrightarrow C^{0,1/2}$ in
dimension three.  The convex Hessian inequality is the
Miranda--Talenti estimate
\cite[Section~2, Theorem~3, p.~297]{Talenti1965}.
Now $Y_0=Y-w$ satisfies Equation~\eqref{n3:existence:fixed-axis} with
a uniform $W^{2,2}$ bound.

Choose $0<\alpha_2<\alpha_1$ and then a fixed
$\lambda_0\in(0,1/4)$ such that
$C\lambda_0^{1+\alpha_1}\le\lambda_0^{1+\alpha_2}/4$ in
Equation~\eqref{n3:existence:model-gradient}.  By first taking
$\delta$ small and then $b\le b_*$ small,
Equation~\eqref{n3:existence:fixed-axis-correction} has supremum norm
at most $\lambda_0^{1+\alpha_2}/4$.  Taylor expansion of $Y_0$
at the center therefore gives an affine function $L_1$ with
\begin{equation}\label{n3:existence:improved-flatness}
 \sup_{B_{\lambda_0}}|z-L_1|
 \le b\lambda_0^{1+\alpha_2},\qquad
 |DL_1-DL_0|\le Cb.
\end{equation}
All constants are uniform in the slope range in
Equation~\eqref{n3:existence:flat-input}.

\paragraph{Reflection at a constant Dirichlet face.}
Subtract the boundary value and use normal coordinates with the face
given by $x_3=0$.  Reflect the metric across this plane and reflect
$z$ oddly.  Tangential derivatives of $z$ vanish on the face, so
the extension is $C^1$.  There is no distributional Hessian atom.
The doubled metric is Lipschitz; the coefficient and lower-order
estimates in Equation~\eqref{n3:existence:local-cordes} hold with its
essential first-derivative bound.

The flux in Equation~\eqref{n3:existence:gradient-flux} may have a
normal jump.  On either side of the face its normal trace is
\begin{equation}\label{n3:existence:reflected-flux}
 \bigl((\Hess z-\Delta z\,g)\nabla z\bigr)_\nu
   =-(\tr_T\Hess z)\,\partial_\nu z.
\end{equation}
Because the boundary value is constant, $\tr_T\Hess z$ is its
mean-curvature coefficient times $\partial_\nu z$.  The jump is
therefore bounded by $C|\mathrm{II}|\,|\nabla z|^2$.  At scale
$r$ it is a bounded plane density of size $O(r)$.  Write the plane
density and flux after multiplication by the reflected coordinate
volume density, so that the following divergence is Euclidean.
Such a signed density $a(x')\,\delta_{\{x_3=0\}}$ is the divergence of the
bounded field $a(x')\mathbf 1_{\{x_3>0\}}\partial_3$.
Thus its negative part is another small divergence error in
Equation~\eqref{n3:existence:s-supersolution}.  Smooth-side curvature
errors are $O(r^2)$.  Weak Harnack and the low-truncation argument
remain valid, and the flatness comparison uses the same Lipschitz
metric estimates.  This proves all the preceding alternatives for
balls meeting the reflection plane as well.

\paragraph{Iteration and the required Hessian bound.}
There are now two ways to improve the solution on a smaller ball:
its gradient bound decreases, or its approximation by a nonconstant
affine function improves. We choose the constants once so that either
route can be repeated.
Keep $\alpha_2,\lambda_0$ from the affine improvement fixed.  Choose
$b_*$ sufficiently small that
$Cb_*/(1-\lambda_0^{\alpha_2})<1/4$, with $C$ the slope-increment
constant in Equation~\eqref{n3:existence:improved-flatness}.
Then choose $r_*,k_*,\delta_*$ in
Equation~\eqref{n3:existence:drop-or-flat} for this fixed $b_*$.
Normalize initially by a fixed multiple of $L+Q$ and choose a
uniformly small initial radius so that the metric and forcing errors
satisfy all the preceding smallness conditions.  If the first case
of Equation~\eqref{n3:existence:drop-or-flat} holds, rescale the ball
by $r_*$ and the gradient bound by $k_*$.  The normalized forcing
decreases by $r_*/k_*<1$.  Repeating this step either continues
forever or reaches the second case.  In the former case the gradient
decays with exponent $\log k_*/\log r_*$.
In the latter case apply Equation~\eqref{n3:existence:improved-flatness}
successively, retaining exponent $\alpha_2$.  At each affine step,
rescale $z$ by dividing by the spatial scale, subtracting only an
additive constant; this preserves its gradient bound.
Its slope changes have total size at most
$Cb_*/(1-\lambda_0^{\alpha_2})$, so the slopes remain in $[1/4,4]$.
At the $j$th step the error size is $b_*\lambda_0^{j\alpha_2}$,
while the normalized forcing and metric first-derivative errors have
size $O(\lambda_0^j)$.  Since $\alpha_2<1$, their ratio to the error
size only decreases.  This justifies every subsequent application
of the improvement lemma.  To combine the two branches, now choose
\begin{equation}\label{n3:existence:iteration-exponent}
 0<\alpha\le\min\left\{\alpha_2,
              \frac{\log k_*}{\log r_*}\right\},\qquad \alpha<1.
\end{equation}

At every center we have consequently obtained affine functions
$L_{x,r}$ satisfying
\begin{equation}\label{n3:existence:campanato-affine}
 \sup_{B_r(x)}|z-L_{x,r}|
 \le C(L+Q)r^{1+\alpha}.
\end{equation}
The slope differences at consecutive radii are bounded by
$C(L+Q)r^\alpha$; comparing the affine approximations on
overlapping balls gives Equation~\eqref{n3:existence:gradient-estimate}.
This also identifies their limiting slope with $Dz$, since $z$ is $C^1$.  The reflected version gives the estimate up to a boundary.

Finally apply the $L^2$ version of
Equation~\eqref{n3:existence:local-cordes} to $z-L_{x,2r}$ after
rescaling.  The affine error contributes
$C(L+Q)r^{1/2+\alpha}$ to the Hessian $L^2$ norm; the bounded
forcing and affine Christoffel term contribute $C(L+Q)r^{3/2}$.
Squaring and using $\alpha<1$ proves
Equation~\eqref{n3:existence:hessian-morrey}.
\end{proof}

\subsection{Continuity in the presence of a quadratic source}

The preceding lemma controls $|D^2f|^2$ by its mass on small balls.
We next show why that control suffices for the $Z$-equation.
An exponential change of variable absorbs its quadratic gradient
term. The remaining measure has a small Green potential, and weak
Harnack then decreases the oscillation. We allow a measure source
so that the argument also covers the inner boundary after reflection.

\begin{lemma}[Natural growth with a measure source]
\label{n3:existence:natural-growth}\label{b:lem:natural-growth}
Let $n\in\{3,4\}$ and let $Z\in L^\infty\cap W^{1,2}_{\mathrm{loc}}$
on a ball in $\R^n$. Fix a Borel representative with $|Z|\le M$ for
products with measures. Suppose, in distributions, that
\begin{equation}\label{n3:existence:natural-equation}
 \operatorname{div}(aDZ+b)=e,
\end{equation}
where $a$ is measurable and symmetric, $\lambda I\le a\le\Lambda I$
for $0<\lambda\le\Lambda$, $|b|\le B$, and $e$ is a signed measure
satisfying
\begin{equation}\label{n3:existence:natural-source}
 |e|\le C_0(1+|DZ|^2)\dd x+\mu,\qquad
 \mu(B_r(x))\le C_\mu r^{n-2+\beta},\qquad \beta>0.
\end{equation}
Here $\mu$ is a positive measure, and its bound holds at every center
and all sufficiently small radii in a larger patch. Assume the weak
exponential product rule \eqref{n3:existence:exponential-identity} below
is valid with the specified representative. This is true, in particular,
for a function continuous across a flat interface and smooth on its
two closed sides, when the equation includes the normal-flux jump.
Then the Sobolev class of $Z$ has a H\"older representative on smaller
balls, with constants depending only on the displayed bounds and patch
separation. If $Z$ is already continuous, this is a uniform estimate
for that function. The conclusion applies up to a smooth face whenever
its reflected equation satisfies these hypotheses.
\end{lemma}

\begin{proof}
The proof has two steps. An exponential change absorbs the quadratic
term in $DZ$; a small potential correction then permits weak Harnack
to decrease the essential oscillation. Both steps work in the two
dimensions stated in the lemma. Restricting all radii to at most one,
we may replace $\beta$ by $\min\{\beta,1\}$ in the measure bound.
We retain the symbol $\beta$ for this positive reduced exponent;
this also makes the small-scale mollification bound below immediate.

Write $\mathcal L=-\operatorname{div}(aD)$ and, for $s\in\{-1,1\}$,
set $V_s=e^{skZ}$. The assumed product rule is
\begin{equation}\label{n3:existence:exponential-identity}
 \mathcal L V_s
 =\operatorname{div}(skV_sb)-skV_se
       -k^2V_s\bigl(aDZ\cdot DZ+b\cdot DZ\bigr).
\end{equation}
The term $V_se$ uses the fixed bounded representative. For a continuous
piecewise smooth function the identity follows by integration on the
two sides; its common trace multiplies the flux jump. Choose
$k\ge\max\{1,4C_0/\lambda\}$ and use
$B|DZ|\le(\lambda/2)|DZ|^2+B^2/(2\lambda)$.
The gradient-square term from $e$ is absorbed by the negative quadratic
term for both signs. Since $e^{-kM}\le V_s\le e^{kM}$, this gives
\begin{equation}\label{n3:existence:exponential-subsolution}
 \mathcal L V_s\le\operatorname{div}F_s+\nu,
 \qquad |F_s|\le C,
 \qquad \nu=C(\dd x+\mu)\ge0.
\end{equation}
Thus no small-ball estimate for $|DZ|^2$ is needed.

On a ball $B_R$ in the smaller patch solve the zero-Dirichlet problem
$\mathcal L w_s=\operatorname{div}F_s+\nu$.
For the bounded vector source, the scalar bounded-solution estimate
with any fixed $p>n$ gives
\[
 \|w_{F_s}\|_\infty
 \le CR^{1-n/p}\|F_s\|_{L^p(B_R)}\le CR;
\]
this is the scaled divergence-source estimate of
\cite[Theorem~2.6, p.~50]{LittmanStampacchiaWeinberger1963}.
For the positive measure source, extend $a$ elliptically to $B_{2R}$.
The scalar Green comparison theorem
\cite[Theorem~7.1 and Remark~2, p.~66]{LittmanStampacchiaWeinberger1963},
and domain monotonicity, give
\[
 0\le\mathcal G_{B_R}(x,y)\le\mathcal G_{B_{2R}}(x,y)
                    \le C|x-y|^{2-n},\qquad x,y\in B_R.
\]
Using $B_{2R}$ keeps both points in a fixed compact interior region
of the comparison ball after rescaling. Dyadic annuli and
\eqref{n3:existence:natural-source} imply
\begin{equation}\label{n3:existence:green-morrey-correction}
 \begin{split}
 \sup_{x\in B_R}\int_{B_R}\mathcal G_{B_R}(x,y)\dd\nu(y)
 &\le C\sum_{j\ge0}(2^{-j}R)^{2-n}
                       \nu(B_{2^{1-j}R}(x))\\
 &\le C(R^2+R^\beta).
 \end{split}
\end{equation}
All balls used in this estimate lie in the chosen larger patch.

This potential also gives an energy solution. Mollify the restriction
of $\nu$ to $B_R$ and restrict the resulting positive density again to
$B_R$. Its ball-growth bounds are uniform: above the smoothing scale
use an enlarged ball, and below it use the mollified density bound.
The variational solutions have the preceding supremum bound and obey
\[
 \lambda\int|Dw_j|^2\le\int w_j\dd\nu_j
              \le\|w_j\|_\infty\nu_j(B_R).
\]
Weak compactness in $W^{1,2}_0$ and distributional convergence produce
the desired measure-source solution. Adding the vector-source solution,
and taking $R\le1$, yields
\begin{equation}\label{n3:existence:correction-size}
 \|w_s\|_\infty\le\varepsilon_R,
 \qquad\varepsilon_R=C(R+R^\beta).
\end{equation}

All extrema in the rest of this proof are essential extrema. Put
$M_R=\operatorname*{ess\,sup}_{B_R}Z$ and
$m_R=\operatorname*{ess\,inf}_{B_R}Z$. The functions
\begin{equation}\label{n3:existence:exponential-deficits}
 H_+=e^{kM_R}-e^{kZ}+w_++\varepsilon_R,
 \qquad
 H_-=e^{-km_R}-e^{-kZ}+w_-+\varepsilon_R
\end{equation}
are nonnegative almost everywhere and satisfy $\mathcal L H_\pm\ge0$.
At least half of $B_{R/2}$ satisfies either
$Z\le(M_R+m_R)/2$ or $Z\ge(M_R+m_R)/2$.
On the first set the positive exponential deficit is at least
$c(M_R-m_R)$; on the second set the negative exponential deficit
has that lower bound. Their derivatives are bounded above and away
from zero on $[-M,M]$.
Weak Harnack for the corresponding $H_\pm$ therefore either lowers
the essential supremum or raises the essential infimum on $B_{R/4}$
by at least $c'(M_R-m_R)-C\varepsilon_R$. Hence
\begin{equation}\label{n3:existence:oscillation-contraction}
 \operatorname*{ess\,osc}_{B_{R/4}}Z
 \le(1-c')\operatorname*{ess\,osc}_{B_R}Z+C(R+R^\beta).
\end{equation}
Choose a positive exponent smaller than $\min\{1,\beta\}$ and than
$-\log(1-c')/\log4$. Iteration gives an essential oscillation bound
$Cr^\gamma$ at every center in a smaller patch.
The averages of $Z$ on balls centered at $x$ are consequently Cauchy
as their radii decrease to zero. Their limit defines $Z^*(x)$ at
every center; comparison on overlapping balls gives
$|Z^*(x)-Z^*(y)|\le C|x-y|^\gamma$.
Lebesgue differentiation identifies $Z^*=Z$ almost everywhere.
This proves the representative assertion. For continuous $Z$ the
representatives agree everywhere. No conclusion about arbitrary
values of the initially chosen Borel representative on a null set
is required.

For the boundary application, flatten a face and set
$S=\operatorname{diag}(1,\ldots,1,-1)$. An even extension uses
$a_-=Sa_+S$ and $b_-=Sb_+$; its plane source is twice the
prescribed total normal flux. An odd extension after subtracting a
constant Dirichlet value uses $a_-=Sa_+S$, $b_-=-Sb_+$ and the
sign-reflected bulk source; its normal flux is continuous.
The coordinate volume density is included in these coefficients.
For the smooth reflected solutions used here the product rule holds,
and a bounded plane density has ball mass $O(r^{n-1})$, giving
$\beta=1$ for that part of the measure. This verifies the boundary
clause whenever the remaining reflected measure has the stated growth.
\end{proof}

\subsection{Boundary regularity for the coupled equations}

We return to the deformation. Its regularity must be obtained in
the order $Df$, then $Z$, then $D^2f$, and finally $DZ$. In the
last step, small oscillation of $Z$ absorbs its quadratic gradient
term. This order avoids assuming the Hessian estimate that the
second equation still needs.

\begin{proposition}[Classical bounds for bounded-variable solutions]
\label{n3:existence:classical-bounds}
Fix the deformation parameters and assume the bounds in
Proposition~\ref{n3:apriori:bounds}.  Every smooth solution in the
homotopy of Proposition~\ref{n3:deformation:homotopy}, with
$t\ge-\epsilon$, has uniform local classical estimates of every
finite order, up to the inner and outer faces.  At fixed $N$ the
constants are independent of the homotopy parameter and of all
sufficiently large outer radii $R$.  The local patches can be chosen
with a uniform positive radius throughout the end.
\end{proposition}

\begin{proof}
\emph{First, control $Df$ and the Hessian measure.}
Proposition~\ref{n3:apriori:bounds} bounds $f,Z,t$ and
$\sigma=|\nabla f|$. At fixed parameters this bounds $u,l,D$
above and away from zero and gives $\chi\ge\chi_0>0$.  After dividing the scalar trace
equation by $l/\sqrt D$, its right side is
\begin{equation}\label{n3:existence:normalized-f-source}
 G_f=\frac{\sqrt D}{l}
         \bigl(F_s-\tr_{A_\chi}K\bigr).
\end{equation}
It is bounded throughout the homotopy, including its added collar
terms.  Lemma~\ref{n3:existence:gradient} therefore gives a uniform
$C^{1,\alpha}$ estimate for $f$ and
\begin{equation}\label{n3:existence:f-morrey-application}
 \int_{B_r}|\Hess f|^2\le Cr^{1+2\alpha}
\end{equation}
on the uniform patches, including boundary patches where $f$ is
constant.

\emph{Second, control the oscillation of $Z$.}
After multiplication by the Riemannian volume density, the second
equation has coordinate principal coefficients
$a^{ij}=4\sqrt{\det g}\,uA_\chi^{ij}$, a bounded flux vector $b$, and a source
bounded in absolute value by
\begin{equation}\label{n3:existence:Z-quadratic-source}
 C\bigl(1+|DZ|^2+|D^2f|^2\bigr).
\end{equation}
Indeed $t$ is a smooth function of $x,Z,Df$ on the bounded variable
range.  Thus $Dt$ is a bounded linear combination of $DZ,D^2f$
and fixed smooth terms; the expression $\mathcal T$ is quadratic
in these quantities.  All other source terms have lower order.
The same volume density is included in $b,e$; it is smooth and bounded
above and away from zero on the chosen patches, so it preserves these
bounds and uniform ellipticity.  The prescribed conormal flux at the inner boundary is
bounded.  At the outer face $Z=0$.
Apply Lemma~\ref{n3:existence:natural-growth}, with
$\mu=C|D^2f|^2\dd x$ and the bounded reflected plane density.
Equation~\eqref{n3:existence:f-morrey-application} supplies its measure
growth hypothesis.  We obtain a uniform positive H\"older exponent
for $Z$.

\emph{Third, gain two derivatives of $f$.}
The coefficients and source of the scalar equation are smooth
functions of $x,Z,Df$ on the bounded range. They are therefore
H\"older now that $Z$ and $Df$ are H\"older.  Interior and constant-Dirichlet boundary Schauder
estimates give $f\in C^{2,\theta}$ for a uniform $\theta>0$.
Expand the $Z$-equation in nondivergence form.  Its principal
coefficients are uniformly elliptic and $C^{0,\theta}$, while its
right side is bounded by $C(1+|DZ|^2)$.
On an inner face $Df$ is normal.  Therefore $A_\chi$ has zero
normal--tangential components there and normal eigenvalue $\chi$.
Its boundary equation solves for
\begin{equation}\label{n3:existence:ordinary-normal-condition}
 \partial_\nu Z=\mathcal H_s(x,Z,f,Df),
\end{equation}
where $\mathcal H_s$ is smooth on a fixed bounded variable range.
Its derivatives with respect to $x$ and its displayed scalar
arguments are bounded at fixed parameters.  The outer condition
remains constant Dirichlet.  These are different, disjoint smooth
boundary components; there is no change of boundary condition along
an edge or a corner.

\emph{Fourth, absorb the quadratic growth and control $DZ$.}
Continuity of $Z$ supplies the small factor needed for this step.
Fix $p>3$.  The linear local $W^{2,p}$ estimate for the leading
nondivergence operator with Equation~\eqref{n3:existence:ordinary-normal-condition}
or outer Dirichlet data has the form, on fixed nested patches
$U_x\Subset V_x$,
\begin{equation}\label{n3:existence:Wp-pre-absorption}
 \|Z\|_{W^{2,p}(U_x)}
 \le C\left(1+\|DZ\|_{L^{2p}(V_x)}^2
                 +\|Z\|_{W^{1,p}(V_x)}\right).
\end{equation}
Here the boundary norm is controlled by
\begin{equation}\label{n3:existence:boundary-trace-composition}
 \|\mathcal H_s(x,Z,f,Df)\|_{W^{1-1/p,p}(\partial V_x)}
 \le C\bigl(1+\|Z\|_{W^{1,p}(V_x^+)}\bigr),
\end{equation}
using trace and composition after extending its smooth expression
through the collar.  The derivatives of $f,Df$ entering this
extension are already bounded.  The leading matrix has a fixed
H\"older modulus, so its small oscillation on sufficiently small
patches permits the usual freezing and absorption in the linear
estimate.  The normal boundary operator is uniformly oblique and
satisfies the complementing condition. The general boundary theory is
developed in the two papers of Agmon, Douglis and Nirenberg
\cite{AgmonDouglisNirenberg1959,AgmonDouglisNirenberg1964}. The scalar
estimate used here is
\cite[Theorem~15.2, estimate~(15.5), p.~704]{AgmonDouglisNirenberg1959}.

The outer patches $U_x,V_x,V_x^+$ and the constant in
\eqref{n3:existence:Wp-pre-absorption} are fixed before the next radius
$h$ is chosen; only the interpolation is localized further.
Let $[Z]_{C^{0,\theta}}\le H_0$. On balls or half-balls of
radius $h$, the scaled Gagliardo--Nirenberg inequality, applied after
subtracting a constant, gives
\begin{equation}\label{n3:existence:small-oscillation-GN}
 \begin{split}
 \|DZ\|_{L^{2p}(B_h)}^2
 &\le C\operatorname{osc}_{B_h}Z\,
                    \|D^2Z\|_{L^p(B_h)}\\
 &\quad+Ch^{3/p-2}(\operatorname{osc}_{B_h}Z)^2.
 \end{split}
\end{equation}
One can use a bounded extension operator on a half-ball; its constants
are uniform in the boundary charts.  Here the interpolation exponents are spatial dimension $3$,
derivative orders $2$ and $1$, top-derivative exponent $p>3$,
height exponent $\infty$, and interpolation parameter $1/2$;
the resulting gradient exponent is $2p$. This is
\cite[Lecture~II, Equation~(2.2), p.~125, and Remark~5, p.~126]{Nirenberg1959},
including the lower-order term on a bounded patch. Subtracting a
constant gives the oscillation, rather than the absolute height,
as the coefficient of $\|D^2Z\|_p$.
Cover $V_x$ by radius-$h$ patches of bounded overlap and take the
$\ell^p$ sum of Equation~\eqref{n3:existence:small-oscillation-GN}.
It follows that
\begin{equation}\label{n3:existence:global-local-GN}
 \|DZ\|_{L^{2p}(V_x)}^2
 \le CH_0h^\theta\|D^2Z\|_{L^p(V_x^+)}
           +CH_0^2h^{2\theta-2}.
\end{equation}
Also, for any $\eta>0$,
$\|Z\|_{W^{1,p}(V_x)}
 \le\eta\|D^2Z\|_{L^p(V_x^+)}+C_\eta$.
Each enlarged patch $V_x^+$ is covered by a bounded number of the
patches $U_y$, independently of $R$.  Set
$S=\sup_y\|Z\|_{W^{2,p}(U_y)}$, finite for each smooth finite-domain
solution.  Equations~\eqref{n3:existence:Wp-pre-absorption} and
\eqref{n3:existence:global-local-GN} give
\begin{equation}\label{n3:existence:uniform-local-absorption}
 S\le C(H_0h^\theta+\eta)S+C(h,\eta).
\end{equation}
Choose $h$ and then $\eta$ so the first coefficient is less than
$1/2$.  We obtain a uniform $W^{2,p}$ bound and hence a
$C^{1,1-3/p}$ bound for $Z$.

\emph{Finally, obtain all remaining derivatives.}
The bounds $f\in C^{2,\theta}$ and $Z\in C^{1,1-3/p}$ make the
expanded $Z$ source H\"older and its normal datum
$\mathcal H_s(x,Z,f,Df)$ of class $C^{1,\theta'}$, after reducing
the positive exponent. Its principal coefficients are H\"older as
well, so both the interior and normal-boundary Schauder estimates apply.
First the nondivergence $Z$-equation has H\"older source, giving
$Z\in C^{2,\theta'}$; the scalar $f$-equation then gains another
derivative.  Repetition gives every finite classical derivative
supported by the smooth data.  The same procedure holds on the
uniform unit patches of the asymptotic end and on the large outer
spheres, whose normal-coordinate constants are uniform.  This
proves the asserted independence of $R$.  The prepared background has end symbol bounds of every order by
Definition~\ref{n3:prepared:data}; no regularity constant here
requires additional derivatives of the original weak-decay data.
\end{proof}

\subsection{Solving the scalar equation for arbitrary inputs}

The previous proposition estimates solutions of the coupled system.
To construct such a solution by degree, we first need an operator
defined on all of a Banach space of trial functions $Z$. In
particular, $t\ge-\epsilon$ cannot be assumed at this stage.
The scalar equation remains solvable because its loss of axial
ellipticity at large slope is at most reciprocal-linear.

\begin{lemma}[Scalar solve for unrestricted bounded inputs]
\label{n3:existence:scalar-solve}
Fix $N\ge4$, $\tau>0$, a finite smooth truncation $\Omega_R$,
and a homotopy parameter $s\in[0,4]$.  For every
$Z\in C^{1,b}(\overline{\Omega_R})$, $0<b<1$, the scalar trace
equation of Proposition~\ref{n3:deformation:homotopy}, with its
prescribed componentwise constant Dirichlet values for $f$, has a
unique solution.  It lies in $C^{3,b}$ and is smooth to the extent
allowed by the input.  The solution operator
\begin{equation}\label{n3:existence:scalar-operator}
 f=\mathcal S_s(Z)
\end{equation}
is continuous in $(s,Z)$, including across the concatenation points
of the homotopy.  On bounded subsets of $C^{1,b}$ its $C^{3,b}$
norm is uniformly bounded at these fixed parameters.  No lower
bound for $t$ is required.
\end{lemma}

\begin{proof}
For the scalar continuation below, use the affine space of
$C^{3,b}$ functions with the assigned Dirichlet values and the equation
space $C^{1,b}$. Its individual solutions meet the $C^3$ hypothesis of
Lemma~\ref{n3:existence:gradient}; no uniform third-derivative bound is
assumed before that lemma is applied.
The monotonicity $\partial Z/\partial t=1+pv/4>0$ defines
$t=t(x,Z,Df)$ smoothly for every finite input, including $Df=0$.
The principal matrix has no dependence on the value of $f$.
The right side is increasing in $h=\tau f$, with derivative
at least one. At a height extremum the collar terms vanish and
$Df=0$, so the equation gives fixed constant upper and lower
barriers. Including all prescribed boundary heights, we obtain
\begin{equation}\label{n3:existence:scalar-height}
 |h|\le C_0,
\end{equation}
where the constant can be chosen independently of $Z$ and the
homotopy parameter.  We include the fixed boundary values in its
choice.

\emph{Large slopes.}
Only $|Z|\le M$ is assumed here. As $\sigma\to\infty$,
the implicit equation forces
$t\to-\infty$ uniformly in $Z$.  There $p=N$ and $l=e^{Nt}$,
so its exact form is
\begin{equation}\label{n3:existence:implicit-large-slope}
 e^{8Z}=e^{8t}+e^{(2N+8)t}\sigma^2.
\end{equation}
Consequently
\begin{equation}\label{n3:existence:large-slope-asymptotics}
 \begin{split}
 t&=-\frac{\log\sigma}{N+4}+O(1),\qquad
 l\asymp\sigma^{-N/(N+4)},\\
 d&\asymp\sigma^{-8/(N+4)},\qquad
 \chi\asymp\sigma^{-8/(N+4)},\qquad
 \frac{\sqrt D}{l}\asymp\sigma.
 \end{split}
\end{equation}
The constants are uniform for $|Z|\le M$ at fixed $N$.
Since $N\ge4$, the exponent $8/(N+4)$ is at most one.  Combining
the large-slope estimates with compactness of the remaining variable
range yields
\begin{equation}\label{n3:existence:scalar-growth}
 \chi\ge\frac{c}{1+\sigma},\qquad
 |G_f|\le C(1+\sigma)
\end{equation}
whenever Equation~\eqref{n3:existence:scalar-height} holds.

\emph{Boundary slopes.}
The estimate $\chi\ge c/(1+\sigma)$ suggests a concave logarithmic
barrier: its axial second derivative can dominate terms of linear
growth in its slope.
Let $\delta$ be smaller than a fixed collar radius, half the distance
between distinct boundary components, and a normal injectivity
radius of a smooth extension of the finite domain.  On a boundary
distance collar use
\begin{equation}\label{n3:existence:logarithmic-modulus}
 \psi(r)=k^{-1}\log(1+Br),\qquad
 \psi''=-k(\psi')^2.
\end{equation}
The upper and lower barriers are the assigned face value plus or
minus $\psi$.  On these tests the axis is normal to the distance
leaves, so their principal trace is
$\chi\psi''+\psi'\Delta r$.  At slopes at least one,
Equation~\eqref{n3:existence:scalar-growth} gives
\begin{equation}\label{n3:existence:logarithmic-dominance}
 -k\chi(\psi')^2\le-\tfrac12ck\psi'.
\end{equation}
Choose $k$ large enough to dominate all the bounded geometry terms
and the linear growth in Equation~\eqref{n3:existence:scalar-growth}.
Next shrink $\delta$ so that $1/(2k\delta)$ exceeds any fixed
large-slope threshold used here.  Finally choose $B$ large enough
that $B\delta\ge1$ and $\psi(\delta)>4C_0/\tau$.
The latter exceeds twice the bound for $\operatorname{osc}f$
from \eqref{n3:existence:scalar-height}.  Then $\psi'\ge1/(2k\delta)$ on this collar and
the two barriers compare by the height monotonicity.  This gives
a uniform boundary gradient bound.

\emph{Interior slopes.}
Use the same modulus for differences at two points. After increasing
$B$ if needed, consider
\begin{equation}\label{n3:existence:doubling-function}
 f(x)-f(y)-\psi(\operatorname{dist}(x,y)),
 \qquad 0<\operatorname{dist}(x,y)<\delta.
\end{equation}
The distance is taken in the fixed smooth extension.  The boundary barriers exclude a positive maximum with one endpoint
on a face: the other point lies in that face's collar, its distance
to the face is at most the distance between the pair, and $\psi$ is
increasing. The separation choice prevents endpoints on distinct faces.  The inequality $\psi(\delta)>2\operatorname{osc}f$
excludes the other boundary of the two-point domain.  At a positive
interior maximum let $r=\operatorname{dist}(x,y)$ and
$q=\psi'(r)$.  The two gradients have length $q$ and are parallel
transports along the connecting geodesic.  Vary both endpoints
simultaneously in parallel transverse directions.  The second
variation formula for the short geodesic yields
\begin{equation}\label{n3:existence:two-point-transverse}
 \tr_{e_x^\perp}\Hess f(x)
 -\tr_{e_y^\perp}\Hess f(y)\le Cqr.
\end{equation}
Varying each endpoint separately along that geodesic gives
$f_{;e_xe_x}(x)\le\psi''(r)$ and
$f_{;e_ye_y}(y)\ge-\psi''(r)$.  Subtracting the two scalar
equations therefore gives
\begin{equation}\label{n3:existence:two-point-contradiction}
 -2C(1+q)
 \le G_f(x)-G_f(y)
 \le Cqr+(\chi_x+\chi_y)\psi''(r).
\end{equation}
By Equations~\eqref{n3:existence:scalar-growth} and
\eqref{n3:existence:logarithmic-modulus}, the last term is at most
$-ckq$ at the large comparison slopes.  Our choice of $k$ gives
a contradiction.  Thus Equation~\eqref{n3:existence:doubling-function}
is nonpositive.  Reversing $x,y$ proves the same bound for the
absolute difference and, on letting $y\to x$, a global Lipschitz
bound for $f$.  The endpoint variations are interior variations in the smooth
extension; the connecting short geodesic need not stay inside
$\Omega_R$. This argument used no continuity modulus for
$\chi_x-\chi_y$; only the common transverse eigenvalues and the
radial lower bound were needed.

\emph{Continuation and uniqueness.}
The height and gradient estimates supply compactness for a scalar
continuation. Interpolate the normalized equation with
$\Delta f=\tau f$, retaining its prescribed boundary values.  Explicitly,
for $a\in[0,1]$ take principal matrix
$(1-a)I+aA_\chi$ and right side
$(1-a)\tau f+aG_f$.  The height sign, the linear gradient growth,
and Equation~\eqref{n3:existence:scalar-growth} persist uniformly in
$a$.  The height and logarithmic gradient bounds just proved are
therefore uniform.  After these bounds the interpolating equations
are uniformly elliptic.  Their matrices still have the axial form,
with axial eigenvalue $(1-a)+a\chi$, so
Lemma~\ref{n3:existence:gradient} applies.  Schauder estimates then
give a uniform $C^{2,\theta}$ bound at fixed input $Z$.
The coefficients and right side are now Lipschitz in position on
bounded input sets, which permits their $C^{0,b}$ estimates and
an upgrade to $C^{2,b}$.  Differentiating once, or using the next
Schauder estimate, yields $f\in C^{3,b}$ because $Z\in C^{1,b}$.

The scalar linearization has the form
$a^{ij}D_{ij}\varphi+b^iD_i\varphi-c\varphi$ with
$c\ge c_1\tau>0$ on the bounded variable range.  Indeed the
principal matrix is independent of the height, while
$\partial_fG_f\ge\tau\sqrt D/l>0$.
The Dirichlet maximum principle gives uniqueness and a zero kernel;
the linear elliptic Dirichlet theory gives an inverse from
$C^{1,b}$ to $C^{3,b}_0$, where the subscript denotes zero data on
the entire boundary.  The
implicit-function theorem gives openness of the interpolation
parameter set.  The uniform classical bounds and compactness give
closedness, starting from the invertible Laplace endpoint.  This
proves existence.  The same maximum principle applied to the
difference of two solutions proves uniqueness of the nonlinear
solution.  The implicit-function theorem, or compactness and
uniqueness, gives continuous dependence on $Z$ and on the
homotopy parameter.  At the finitely many concatenation points the
scalar equations agree, so the same continuity holds there.
\end{proof}


\subsection{Degree on the admissible set and exhaustion}

We can now construct the compact map. The homotopy parameter
$s\in[0,4]$ has a separate role from the deformation integer $N$:
in Proposition~\ref{n3:deformation:homotopy}, $s=0$ is the original
system, $s=1$ removes the drift, $s=2$ completes the collar
modification, and $s=3$ makes the scalar solution $f$ identically
zero. On $[3,4]$ the divergence source and the inner flux are
multiplied together by $4-s$. The final compact map is consequently
the zero map. The remaining task is to show that no fixed point
can cross the boundary of the admissible set along this path.

\begin{lemma}[Degree on varying open sets]\label{b:lem:moving-degree}
Let $I=[a,b]$, let $X$ be a real Banach space, and let
$T:I\times X\to X$ be continuous and map bounded sets to relatively
compact sets. Let $\mathcal U\subset I\times X$ be bounded and
relatively open. Suppose no fixed point $Z=T(s,Z)$ belongs to the
relative boundary of $\mathcal U$. Then, for
$\mathcal U_s=\{Z:(s,Z)\in\mathcal U\}$, the Leray--Schauder degree
\[
 \deg(I_X-T_s,\mathcal U_s,0)
\]
is well defined and independent of $s$. Empty sections have degree zero.
If $T_b$ is the zero map and $0\in\mathcal U_b$, the common degree
is one, and every section contains a fixed point.
\end{lemma}

\begin{proof}
The fixed points in $\overline{\mathcal U}$ form a compact set
$\mathfrak F$. Indeed, boundedness of $\mathcal U$, compactness of $I$
and of the images of $T$ give a convergent subsequence of any sequence
of such fixed points; continuity preserves the fixed-point equation.
The hypothesis places $\mathfrak F$ inside $\mathcal U$.
Also $\{s\}\times\partial\mathcal U_s$ lies in the relative boundary
of $\mathcal U$, so each section degree is defined.

Fix $s_0\in I$. Cover its compact fixed-point slice by finitely many
open balls in $X$ whose closed balls remain in $\mathcal U_s$ for
all $s$ sufficiently close to $s_0$. Relative openness and the finite
cover permit this common parameter interval. Let $U$ be their union.
All fixed points in the closed sections for these nearby parameters
lie in $U$ after shortening the interval: otherwise compactness of
$\mathfrak F$ would give an omitted fixed point in the slice at $s_0$.
Excision identifies the degree on each $\mathcal U_s$ with that on
$U$, and fixed-domain homotopy invariance makes the latter constant.
If the slice is empty, compactness instead gives no nearby fixed
points in the closed sections, so those degrees are all zero.
Thus the section degree is locally constant, and connectedness of $I$
makes it constant. These are the excision and homotopy properties of
Leray--Schauder degree
\cite[Sections~12--14, pp.~59--61]{LeraySchauder1934}; see also \cite{Deimling1985}.
At the stated terminal map, normalization gives degree one; the
existence property of nonzero degree proves the last assertion.
\end{proof}

\begin{theorem}[Existence on finite truncations]
\label{n3:existence:deformation}
For the prepared domain of Proposition~\ref{n3:domains:prepared-domain}
and every sufficiently small fixed $\epsilon>0$, there is an
integer $N_0\ge4$ and, for each $N\ge N_0$, a radius
$R_{\min}(N)$, with $R_{\min}(N)\to\infty$, such that the original
deformation system has a smooth solution on every $\Omega_R$ with
$R\ge R_{\min}(N)$.  It satisfies $t>-\epsilon$ on the closure.
At fixed $N$, these solutions have the local bounds of
Proposition~\ref{n3:existence:classical-bounds}, uniformly as
$R\to\infty$.  In particular, any sequence $R_j\to\infty$ of such radii has a
subsequence converging smoothly on compact subsets of
$\overline\Omega$ to a solution of the original equations with
$t\ge-\epsilon$.
\end{theorem}

\begin{proof}
\emph{Choose the parameters before defining degree.}
Fix the prepared data, thresholds, collars and $\epsilon$.
Choose an integer $N_0$ so that
\begin{equation}\label{n3:existence:n-choice}
 N_0\ge\max\{4,\lfloor2/\epsilon\rfloor+1\}
\end{equation}
and so that the floor test and all estimates of
Proposition~\ref{n3:apriori:bounds} apply for $N\ge N_0$.
Let $R_{\mathrm{ap}}(N)$ dominate every lower-radius requirement
in the a priori estimates and in the floor exclusion; let
$R_{\mathrm{geo}}$ contain all compact boundary faces, collars and
supports of the construction.  The end barrier of
that proposition gives a constant $C_{\mathrm{out}}$, independent
of $N,R,s$, such that every admissible or floor-contact homotopy
solution has
\begin{equation}\label{n3:existence:outer-gradient}
 |\nabla f|\le
       C_{\mathrm{out}}\tau^{-1}R^{-1-\gamma}
       \quad\text{on the outer sphere}.
\end{equation}
The scalar last stage has $f=0$ and satisfies the same bound.
Choose, enlarging to a radius of a smooth truncation if necessary,
\begin{equation}\label{n3:existence:R-choice}
 \begin{split}
 R_{\min}(N)=\max\biggl\{&N,R_{\mathrm{geo}},R_{\mathrm{ap}}(N),\\
 &\left(\frac{2C_{\mathrm{out}}\ell}
               {\tau\sqrt{e^{8\epsilon}-1}}
        \right)^{1/(1+\gamma)}\biggr\}.
 \end{split}
\end{equation}
At $Z=0$, the value of its implicit expression at $t=-\epsilon$
is $-\epsilon+\frac18\log(1+\ell^2\sigma^2)$.
Equations~\eqref{n3:existence:outer-gradient} and
\eqref{n3:existence:R-choice} make this strictly negative.  Strict
monotonicity of that expression in $t$ therefore excludes an outer
floor contact.  The term $N$ in Equation~\eqref{n3:existence:R-choice}
ensures $R_{\min}(N)\to\infty$.

The sublevel and collar estimates apply to arbitrary smooth
admissible or floor-contact solution sequences with $N\to\infty$
and expanding outer cutoffs, uniformly over the compact ranges of
the homotopy parameters.  If the needed conclusion failed for
arbitrarily large $N$ with $R\ge R_{\min}(N)$, one could select
such a violating sequence.  The chosen radius convention would
make it precisely an expanding sequence excluded by those proofs.
Thus a single $N_0$, followed by the radius choices above, suffices
for every homotopy parameter.  No upper bound for $Z$ is inserted
in the earlier contact estimates: there $\sigma\asymp\tau^{-1}$
and $l\ge\ell$ alone give $d\lesssim(\tau/\ell)^2$.

\emph{Define the compact map on all bounded input sets.}
For the rest of the degree argument fix $N\ge N_0$ and
$R\ge R_{\min}(N)$. Work in the Banach space
\begin{equation}\label{n3:existence:banach-space}
 X_b=\{Z\in C^{1,b}(\overline{\Omega_R}):
                       Z=0\text{ on the outer face}\},
 \qquad 0<b<1.
\end{equation}
For an input $Z\in X_b$, first solve
$f=\mathcal S_s(Z)$ by Lemma~\ref{n3:existence:scalar-solve}.
Compute $t,u,A_\chi,\mathcal T$ and the full homotopy source
from this input pair.  Write $B_s$ for the frozen drift summand in
the flux, $E_s$ for the full right side of its divergence equation,
and $q_s$ for the prescribed inner flux.  Thus before the final
scaling $B_s=\lambda_suA_\chi K(w,\cdot)$ and $E_s=u\Xi$;
on the final interval $E_s$ and $q_s$ have the simultaneous scalar
factor prescribed in Proposition~\ref{n3:deformation:homotopy}.
Define $\mathcal K_s(Z)=\widetilde Z$ by the linear mixed problem
\begin{equation}\label{n3:existence:compact-map-equation}
 \begin{cases}
  \Div(4uA_\chi\nabla\widetilde Z+B_s)=E_s
                           &\text{in }\Omega_R,\\
  (4uA_\chi\nabla\widetilde Z+B_s)_\nu=q_s
                           &\text{on }B,\\
  \widetilde Z=0           &\text{on the outer face}.
 \end{cases}
\end{equation}
All coefficients in this problem are computed from the input
$(\mathcal S_s(Z),Z)$; the only unknown is $\widetilde Z$.
The normal $\nu$ on the inner faces points into the domain.
For a test function $\varphi$ vanishing on $S_R$, the weak form is
\[
 \int_{\Omega_R}4uA_\chi\nabla\widetilde Z\cdot\nabla\varphi
 =-\int_{\Omega_R}E_s\varphi
  -\int_{\Omega_R}B_s\cdot\nabla\varphi
  -\int_B q_s\varphi.
\]
The sign of the boundary term comes from using $-\nu$ as the
integration normal. The scalar estimate bounds the leading
coefficient above and away from zero on each bounded input set.
The nonempty outer Dirichlet face supplies Poincar\'e's inequality,
so the bilinear form is coercive and has a unique solution.

On bounded input sets in $C^{1,b}$,
Lemma~\ref{n3:existence:scalar-solve} bounds $f$ in $C^{3,b}$.
Hence the leading coefficient and drift in
Equation~\eqref{n3:existence:compact-map-equation} are $C^{1,b}$,
the bulk source is $C^{0,b}$, and the prescribed boundary data are
$C^{1,b}$.  The derivatives of $t(x,Z,Df)$ involve only $DZ$
and $D^2f$, so this assertion includes the quadratic expression
$\mathcal T$.  Linear regularity gives a bounded output set in
$C^{2,b}$.  The inclusion $C^{2,b}\hookrightarrow C^{1,b}$ is
compact on the finite domain.  Continuous dependence of the scalar
solve and the linear problem makes $(s,Z)\mapsto\mathcal K_s(Z)$
a continuous compact homotopy on bounded sets.  These assertions
hold also at the concatenation parameters because the defining
equations coincide there.  A fixed point solves the coupled system;
repeated Schauder estimates make it smooth.

\emph{Exclude fixed points from the boundary of admissibility.}
Write $t_s[Z]=t(x,Z,D\mathcal S_s(Z))$. The scalar solution must
be computed before this admissibility test. Define
\begin{equation}\label{n3:existence:admissible-open-set}
 \mathcal O=\{(s,Z)\in[0,4]\times X_b:
       \min_{\overline{\Omega_R}}t_s[Z]>-\epsilon,
       \ \|Z\|_{C^{1,b}}<M_1\}.
\end{equation}
It is relatively open because the scalar solve and the implicit
function defining $t$ depend continuously on $(s,Z)$.
Propositions~\ref{n3:apriori:bounds} and
\ref{n3:existence:classical-bounds} give a uniform bound in $C^{1,b}$
for all fixed points satisfying $t\ge-\epsilon$.  If initially
needed, use their higher classical bounds to reach the chosen
exponent $b$.  Choose $M_1$ larger than this bound.  No fixed point
lies on the norm boundary of $\mathcal O$.  No fixed point lies on
its floor boundary either: outer contacts were excluded above and
all remaining contacts are excluded by
Proposition~\ref{n3:deformation:floor}.

\emph{Apply degree on the varying admissible sections.}
The sections $\mathcal O_s$ need not be constant in $s$.
Lemma~\ref{b:lem:moving-degree} applies with $I=[0,4]$, $X=X_b$,
$T_s=\mathcal K_s$ and $\mathcal U=\mathcal O$.
We have proved continuity and compactness on bounded input sets;
the norm condition makes $\mathcal O$ bounded and the scalar solve
makes it relatively open. Its relative boundary can contain only
norm contact or floor contact, and neither can be a fixed point by
the preceding estimates. Thus its section degree is constant.

At $s=4$, the scalar equation has the unique solution $f=0$.
The drift vanishes, and both $E_4$ and $q_4$ vanish for every input.
Equation~\eqref{n3:existence:compact-map-equation} then has output
identically zero.  Moreover, $t_4[0]=0>-\epsilon$, so
$0\in\mathcal O_4$ and
\begin{equation}\label{n3:existence:degree-one}
 \deg(I-\mathcal K_4,\mathcal O_4,0)=1.
\end{equation}
Homotopy invariance gives the same degree at $s=0$.  There is
therefore a fixed point in $\mathcal O_0$, which is the desired
smooth solution with $t>-\epsilon$.

\emph{Remove only the outer truncation.}
Keep $N$ and all other deformation parameters fixed.
Proposition~\ref{n3:existence:classical-bounds} gives estimates
independent of $R$ on an exhaustion by compact sets, using boundary
charts at the fixed inner faces. A diagonal subsequence converges
smoothly on compact subsets of $\overline\Omega$.
Both equations and their inner boundary conditions pass to the
limit. The strict finite-domain inequality gives the non-strict
limit inequality $t\ge-\epsilon$. The following section establishes
end normalization, decay and mass flux; those conclusions do not
follow from compact exhaustion alone.
\end{proof}

\section{The complete exterior and its energy}\label{n3:limit:section}

We now obtain the three geometric outputs stated at the beginning:
nonnegative scalar curvature, control of every cut, and an arbitrarily
small upper energy error. Their combination will also yield the
numerical comparison. Fix the prepared initial data of
Definition~\ref{n3:prepared:data}, and the prepared exterior
\(\Omega\) of Proposition~\ref{n3:domains:prepared-domain}.
Its boundary \(B\) separates the original inner obstacle from the end.
Let \(A_*\) be the enclosing-area infimum of that original obstacle,
measured in the reduced metric \(g\).
In particular every cut enclosing \(B\) has \(g\)-area at least \(A_*\).

We use the notation of Section~\ref{n3:deformation:section} and take the
physical endpoint of the homotopy:
\[
 F=h+C,\qquad h=\tau f,\qquad
 \Div V=u\Xi,\qquad
 \Xi=\delta_0\mathcal T+\rho+\delta_0\tau a\sigma-m_0(t)\mathcal P.
\]
Here \(0<\delta_0<1\), \(\ell=e^{-\epsilon N}\),
\(\tau=\ell^{3/2}\), and all constants in the prepared data are fixed
before \(N\) is taken large. Constants marked \(C_N\) in this section
are bounded by \(C(1+N)^C\), with fixed exponents. Other constants
explicitly allowed to depend on fixed \(N\) need not have such a bound.

\subsection{The limiting solution and its end}

\begin{lemma}[End control]\label{n3:limit:end}
Fix \(\epsilon>0\) and a sufficiently large
\(N\ge\max\{4,\lfloor2/\epsilon\rfloor+1\}\). The truncation solutions of
Theorem~\ref{n3:existence:deformation} have a subsequence converging
smoothly on compact subsets of \(\overline\Omega\) to a solution of
the physical system with \(t\ge-\epsilon\).
On the asymptotically flat end, for some \(c_N>0\) and fixed-\(N\)
constants \(C_j(N)\),
\[
 |\nabla^j f|\le C_j(N)e^{-c_N r}
 \quad\text{for each fixed }j,\qquad
 Z,t=O_2(r^{-1}).
\]
The differences \(t-Z\), \(\bar g-g\), and their coordinate derivatives
of any fixed order decay exponentially. The vector field \(V\)
satisfies
\begin{equation}\label{n3:limit:V-asymptote}
                    V=4\nabla_g t+O(r^{-3})
\end{equation}
and has a finite outward flux limit.
\end{lemma}

\begin{proof}
The compact convergence, including at the fixed inner boundary, follows
from Theorem~\ref{n3:existence:deformation} and its uniform fixed-\(N\)
classical estimates. All equations and inner boundary conditions pass
to this limit. To preserve the end normalization, we prove the decay
estimates on the truncations before passing to infinity.

Outside a fixed large sphere, \(K=C=0\), and the first equation is
\[
 \tr_{A_\chi}\Hess_g f=\frac{\tau\sqrt D}{l}\,f.
\]
At fixed \(N\) the preceding estimates make the left side uniformly
elliptic, and its positive height coefficient is bounded below by a
positive constant depending on \(N\). It is bounded above as well.
For sufficiently small \(c_N>0\), the functions
\(\exp[-c_N(r-r_0)]\) are positive supersolutions of the corresponding
linear equation outside a sufficiently large \(r_0\): their radial
second derivative has size \(c_N^2\), while their tangential second
derivatives have size \(c_N/r\), and the fixed positive height term
dominates both. Multiplying by a constant bounds the solution on
\(S_{r_0}\); the zero outer Dirichlet value is also bounded.
The maximum principle applied to both signs of \(f\) gives the
uniform exponential estimate. The fixed-\(N\) unit-patch estimates,
followed by the local homogeneous linear estimates for this equation,
give the same decay for derivatives. The corresponding boundary
estimates apply on the large outer spheres. Iterating the smooth
equations gives each required finite derivative order.

The relation \(Z=t+\tfrac18\log(1+l^2|\nabla f|^2)\) now shows that
\(t-Z\) is exponentially small, with derivatives. The graph errors
have the same property. On the end the second equation becomes
\begin{equation}\label{n3:limit:Z-end}
 \Delta_gZ+b_N(Z)|dZ|_g^2=O(r^{-4}),\qquad
 b_N(s)=p_L(s)-\delta_0\bigl(p_L(s)-1\bigr).
\end{equation}
Indeed \(A_\chi=I\) and \(u=L(Z)\) up to exponentially small errors,
whereas
\(\mathcal T=4(p_L(Z)-1)|dZ|^2\) up to such errors.
The \(\rho_0\) penalty weight is \(O(r^{-4})\); the term with
\(\rho_1\) is multiplied by the exponentially small \(v\).
All coefficients in this statement are bounded on the fixed-\(N\)
solution range. Differentiably controlled exponential errors can
absorb any fixed polynomial power of \(r\).

Define an increasing smooth function on that range by
\[
 \Phi_N(0)=0,\qquad
 \Phi_N'(s)=\exp\!\left(\int_0^s b_N(z)\,\mathrm dz\right).
\]
Its derivative is bounded above and below by positive fixed-\(N\)
constants. The chain rule in \eqref{n3:limit:Z-end} gives
\[
                         \Delta_g\Phi_N(Z)=O(r^{-4}).
\]
The function is bounded and vanishes at the outer cutoff. To obtain
a bound independent of that cutoff, choose \(0<\eta<1\).
For large \(r\), a positive multiple of
\(r^{-1}-r^{-1-\eta}\) satisfies
\[
 -\Delta_g(r^{-1}-r^{-1-\eta})\ge c r^{-3-\eta}.
\]
Here the Euclidean leading term is
\(\eta(1+\eta)r^{-3-\eta}\), and the metric error is \(O(r^{-4})\).
It dominates the right side after increasing the multiple and the
fixed inner radius. Comparison for both signs yields
\(|\Phi_N(Z)|\le C'(N) r^{-1}\), where \(C'(N)\) is a fixed-\(N\)
constant with no polynomial claim. Consequently \(Z=O(r^{-1})\).

At sufficiently large radius this estimate and the exponential graph
error give \(t>-\epsilon+1/N\); hence the floor penalty vanishes there.
Rescale each annulus of radii comparable to \(r\) to unit size.
The equation for \(\Phi_N(Z)\) has a bounded rescaled source and
uniformly controlled metric coefficients. Local \(W^{2,p}\)
estimates with \(p>3\) first give the gradient decay. Its source is
then a smooth weight times a smooth function of the bounded variables,
plus the already controlled exponential errors. Schauder estimates,
or a further Sobolev bootstrap followed by them, give
\(Z=O_2(r^{-1})\). Inverting \(\Phi_N\) preserves these bounds, and
the same holds for \(t\).

Since \(L(0)=1\), we have \(u=1+O(r^{-1})\) and
\(A_\chi=I+O(e^{-c_N r})\). With \(K=0\) on the end, the definition
of \(V\) gives \eqref{n3:limit:V-asymptote}.
Moreover \(\Div V=u\Xi\) is integrable there:
\(\mathcal T=O(r^{-4})\), \(\rho=O(r^{-4})\), the penalty has vanished,
and \(\tau a\sigma\) decays exponentially. The divergence theorem on
annuli proves the existence of the outward flux limit.
\end{proof}

\begin{lemma}[Curvature, completeness, and ADM charge]\label{n3:limit:geometry}
For the solution in Lemma~\ref{n3:limit:end},
\[
                         \hat g=e^{4t}(g+l^2df^2)
\]
is a smooth complete metric on \(\overline\Omega\), with nonnegative
scalar curvature and strictly negative mean curvature on \(B\) for
the normal pointing into \(\Omega\).
Its end is asymptotically flat with
\[
 \hat g-\delta=O_2(r^{-1}),\qquad
 R_{\hat g}=O(r^{-3-\delta'}),\qquad
 \delta'=\min\{\delta,1\}>0,
\]
where the reduced end has \(R_g=O(r^{-3-\delta})\).
If
\[
                 \mathfrak F_N=\lim_{r\to\infty}
                    \int_{S_r}V_\nu\dd A_g,
\]
then
\begin{equation}\label{n3:limit:adm}
                         E_{\hat g}=E-\frac{\mathfrak F_N}{8\pi}.
\end{equation}
\end{lemma}

\begin{proof}
At the physical endpoint the curvature identity gives
\begin{align*}
 \tfrac12e^{4t}R_{\hat g}
 ={}&8\pi(\mu+J(w))+(1-\delta_0)\mathcal T
       +(1-\delta_0)\tau a\sigma\\
    &+w\cdot\nabla C-F\tr K-\rho+m_0(t)\mathcal P.
\end{align*}
The height bounds on the compact supports of \(K\) and \(C\) give
\(h\in[b_-,b_+]\). Since \(|w|\le1\), the terms
\(|F\tr K|+|\nabla C|+\rho\) are dominated by
\(8\pi(\mu-|J|)\), by the choices in
Proposition~\ref{n3:domains:prepared-domain} and the choice of \(\rho\).
Outside those supports the same conclusion uses
\(8\pi(\mu-|J|)-\rho\ge0\).
All remaining terms are nonnegative. This proves
\(R_{\hat g}\ge0\).

The boundary identity and the prescribed flux give
\(H+4\partial_\nu t=-N_B<0\). Since \(f\) is constant on each face,
this is exactly the mean-curvature sign for \(\hat g\).
Also \(\hat g\ge e^{-4\epsilon}g\), so any escaping curve has infinite
length; the compact boundary is included. Smoothness and completeness
therefore follow.

Lemma~\ref{n3:limit:end} gives the stated metric decay. The prepared-data hypothesis supplies \(R_g=O(r^{-3-\delta})\) on the end. Since graph errors are exponentially small, the
conformal formula gives
\[
 R_{\hat g}
 =e^{-4t}\bigl(R_g-8\Delta_gt-8|dt|_g^2\bigr)
      +O(e^{-c_N r}).
\]
Equation~\eqref{n3:limit:Z-end}, the vanishing far-end penalty, and the
gradient estimate give \(\Delta_gt=O(r^{-4})\).
This proves the asserted scalar-curvature decay, including the
pointwise falloff required by the AF convention of
\cite[Definition~21 and Theorem~19]{Bray2001}.

Exponential graph errors contribute no ADM flux. The conformal change
and \(t=O_2(r^{-1})\) give
\[
 E_{\hat g}
 =E-\frac1{2\pi}\lim_{r\to\infty}
                 \int_{S_r}\partial_\nu t\dd A_g.
\]
Replacing the metric normal and area form by their Euclidean versions
costs \(o(1)\), and the \(O(r^{-3})\) error in
\eqref{n3:limit:V-asymptote} has vanishing sphere integral.
This proves \eqref{n3:limit:adm}.
\end{proof}

\subsection{The boundary and floor losses}

The mass formula reduces the remaining estimate to a lower bound for
the flux. We absorb the inner boundary term into the positive bulk
source and then estimate the penalty near the floor. Both steps use
only polynomial constants.

\begin{lemma}[Flux lower bound]\label{n3:limit:flux}
For fixed \(\epsilon>0\), the physical solutions satisfy
\[
 \mathfrak F_N\ge-\varepsilon_N,\qquad
 0\le\varepsilon_N\le
      C(1+N)^C\bigl(\ell+\ell^2/\tau\bigr)\longrightarrow0.
\]
All constants in this estimate are independent of the fixed-\(N\)
classical regularity constants.
\end{lemma}

\begin{proof}
The normal \(\nu\) on \(B\) points into \(\Omega\), so the divergence
theorem has the sign
\begin{equation}\label{n3:limit:flux-split}
 \mathfrak F_N
     =\int_\Omega u\Xi\dd V_g+\int_BV_\nu\dd A_g.
\end{equation}
Each integral is finite for a fixed \(N\), by
Lemma~\ref{n3:limit:end} and compact smoothness.

At a black face, \(F-P_B=H\) and \(w_\nu=a\).
At a white face, \(F-P_B=-H\) and \(w_\nu=-a\).
Thus at either type of face,
\[
                 V_\nu/u=-(1-a)H-N_Bd\ge-Cd.
\]
The signed boundary gradient bounds give
\(\sigma\ge c/\tau\), so
\[
 \sqrt d\le\frac1{l\sigma}\le C\tau/\ell,\qquad
 \int_BV_\nu\dd A_g\ge
             -C\frac{\tau}{\ell}\int_BL\dd A_g.
\]
Use the polynomial collar trace estimate from
Section~\ref{n3:apriori:section} with the constant test function:
\[
 \int_B L\dd A_g
    \le C(1+N)^C
       \int_{\mathcal C}u(1+\sqrt{\mathcal T})\dd V_g,
\]
where \(\mathcal C\) is a fixed compact union of collars.
Its polynomial constant is independent of the upper \(Z\) bound.
Since \(\rho\) has a fixed positive minimum on \(\mathcal C\),
\[
 \int_{\mathcal C}u(1+\sqrt{\mathcal T})\dd V_g
 \le C\int_\Omega u(\delta_0\mathcal T+\rho)\dd V_g.
\]
The factor \(C(1+N)^C\tau/\ell
=C(1+N)^C\ell^{1/2}\) tends to zero.
For large \(N\), the entire negative boundary term in
\eqref{n3:limit:flux-split} is therefore at most one quarter of
\(\int_\Omega u(\delta_0\mathcal T+\rho)\dd V_g\).

It remains to estimate the floor penalty. On its support,
\(-\epsilon\le t\le-\epsilon+1/N\), so \(l\) and \(L\) are comparable
to \(\ell\), with constants independent of \(N\).
Directly from their definitions,
\[
 \begin{split}
 u&=L\sqrt{1+l^2\sigma^2}
                      \le C(\ell+\ell^2\sigma),\\
 uv&=\frac{Ll^2\sigma^2}{\sqrt{1+l^2\sigma^2}}
                      \le C\ell^2\sigma,\\
 ua\sigma&=Ll\sigma^2\ge c\ell^2\sigma^2 .
 \end{split}
\]
Consequently
\[
 um_0\mathcal P
 \le C_N\bigl[\ell\rho_0+\ell^2\sigma(\rho_0+\rho_1)\bigr].
\]
Young's inequality, with the square term chosen to be a quarter of
\(\delta_0\tau ua\sigma\), yields
\[
 um_0\mathcal P
 \le \tfrac14\delta_0\tau ua\sigma+
 C_N\left[\ell\rho_0+
              \frac{\ell^2}{\tau}(\rho_0^2+\rho_1^2)\right].
\]
Squaring the constant only changes the polynomial
\(C_N\). The integrals of \(\rho_0\), \(\rho_0^2\), and \(\rho_1^2\)
are finite: on the end their orders are respectively \(r^{-4}\),
\(r^{-8}\), and \(r^{-4\gamma}\), with \(4\gamma>3\).
Integrating leaves at most
\(C(1+N)^C(\ell+\ell^2/\tau)\), after absorption.

Substitute both bounds in \eqref{n3:limit:flux-split}.
Positive fractions of the three bulk terms
\(\delta_0\mathcal T,\rho,\delta_0\tau a\sigma\) remain.
Dropping them proves the claim. Finally
\(\ell^2/\tau=\ell^{1/2}\), so the error tends to zero
exponentially against every polynomial loss.
\end{proof}

\begin{proof}[Proof of Theorem~\ref{n3:prepared:output}]
The fixed subexterior and preservation of every full cut are supplied
by Proposition~\ref{n3:domains:prepared-domain}. Fix $\epsilon$ and
take the global solutions of Lemma~\ref{n3:limit:end}, obtained by
first exhausting the end at fixed $N$. Lemma~\ref{n3:limit:geometry}
gives completeness, the two curvature signs, and the end decay. The
pointwise inequality $\widehat g\ge e^{-4\epsilon}g$ gives the
two-dimensional area inequality on every full cut, including every
component and every coincident portion. The energy formula and
Lemma~\ref{n3:limit:flux} give the asserted upper error with
$\eta_{\epsilon,N}=\varepsilon_N/(8\pi)$. This completes the
geometric construction before any Riemannian Penrose comparison.
\end{proof}

\subsection{The Riemannian horizon and the inequality}

The deformed metric has nonnegative scalar curvature, strictly negative
inner mean curvature toward the end, and the required asymptotic decay.
We replace its inner boundary by an outer area-minimizing minimal
enclosure, then compare its area with the original enclosing infimum.

\begin{theorem}[Reduced energy inequality]\label{n3:limit:energy}
The reduced data satisfy
\[
                              E\ge\sqrt{\frac{A_*}{16\pi}}.
\]
Consequently Theorem~\ref{n3:intro:exterior-inequality} holds in its
entire stated exterior class.
\end{theorem}

\begin{proof}
For each fixed sufficiently large \(N\), minimize \(\hat g\)-area
among filled cuts enclosing \(B\), using a large outer sphere as an
auxiliary barrier. The independent enclosure results of~\cite[Lemmas~2.1--2.3]{RiemannianCompanion}
apply to this smooth AF metric; we check their hypotheses and the
strict detachment here.
A fixed enclosing competitor bounds the area, and local perimeter
density estimates prevent minimizers from escaping to infinity.
The positive mean curvature of sufficiently large outer spheres
excludes contact there, as in that enclosure theorem.
For the inner barrier, let \(U_s\) be a short outward normal collar
of the filled obstacle and \(Y\) its outward unit foliation field.
Strict negativity gives \(\Div_{\hat g}Y<0\) throughout the collar.
If a perimeter minimizer \(E\) omitted a positive-volume portion
\(D=U_s\setminus E\), Gauss--Green and \(|Y|_{\hat g}\le1\) would give
\[
 P_{\hat g}(E\cup U_s)-P_{\hat g}(E)
       \le\int_D\Div_{\hat g}Y\dd V_{\hat g}<0.
\]
Thus \(E\) contains a whole inner collar; obstacle regularity upgrades
this inclusion to its regular representative. Its free boundary is
disjoint from the obstacle. The resulting cut \(\Gamma_N\) is compact,
smooth, and minimal in dimension three; it may be disconnected.
It is outer area-minimizing by its minimizing property.
Equivalently one may first use the ordinary trapped-region theorem
for \((\hat g,0)\) and then take the outer minimizing enclosure.
The boundary version of the Riemannian Penrose theorem applies to the
exterior of \(\Gamma_N\), by Lemma~\ref{n3:limit:geometry}
and \cite[Theorem~19]{Bray2001}.

For every two-plane the restriction of \(\bar g=g+l^2df^2\) dominates
that of \(g\). The floor therefore gives
\[
                  |\Gamma_N|_{\hat g}
          \ge e^{-4\epsilon}|\Gamma_N|_g
          \ge e^{-4\epsilon}A_*.
\]
The last inequality uses that \(B\), and hence \(\Gamma_N\), cuts off
the original filled obstacle. Lemmas~\ref{n3:limit:geometry}
and~\ref{n3:limit:flux} now give
\[
 E+\frac{\varepsilon_N}{8\pi}
       \ge E_{\hat g}
       \ge\sqrt{\frac{|\Gamma_N|_{\hat g}}{16\pi}}
       \ge e^{-2\epsilon}\sqrt{\frac{A_*}{16\pi}}.
\]
Let \(N\to\infty\) with \(\epsilon\) fixed, and then
\(\epsilon\downarrow0\). This proves the reduced inequality.

Finally apply the weak three-dimensional rest-preparation theorem
of~\cite[Proposition~7.13]{NumericalCompanion} to an exterior satisfying the hypotheses
of Theorem~\ref{n3:intro:exterior-inequality}.
Write $(g_j,K_j)$ for the resulting data. Their tensor fields have
compact support, their momenta vanish, and their metrics have all-order
$r^{-1}$ decay and the stipulated scalar falloff. Discard a finite
initial segment so that $E_j>0$. Strict trapping and strict DEC hold
on the compact core. The positive minimum there of
$8\pi(\mu_j-|J_j|_{g_j})\varrho^{3+\delta}$, together with the
end margin, gives the global weighted bound after decreasing $c_j$.
The full enclosing area is positive by the geometric enclosure
lemma. Thus each fixed $j$ satisfies Definition~\ref{n3:prepared:data};
no constant uniform in $j$ has been used in its deformation solve.
The reduced energies converge to its invariant ADM mass and the
full enclosing infima converge to the original enclosing infimum.
Applying the reduced bound to each member of that sequence and passing
to the limit gives Equation~\eqref{n3:intro:penrose}.
\end{proof}


\input{four-opening}
\input{four-geometry}
\input{four-system}
\input{four-bounds}
\input{four-regularity}
\input{four-existence}

\section{Returning to weakly decaying four-dimensional data}
\label{b:four-transfer}

The prepared boundary theorem can be used without imposing its improved
end decay on the original data.  The necessary reduction is the
independent rest-end replacement and strictification proved in
\cite[Proposition~8.8 and Lemma~8.9]{NumericalCompanion}.  We record
the application with its order of limits, since its area comparison
is one-sided.

\begin{corollary}[The one-ended weak-decay consequence]
\label{b:four-weak}
Let $(M^4,g,K)$ be a smooth connected orientable exterior, complete with
its nonempty compact smooth boundary included, and with exactly one
Euclidean coordinate end and compact complement.  Assume
\[
 g-\delta=O_2(r^{-q}),\qquad K=O_1(r^{-1-q}),\qquad q>1.
\]
Use the four-dimensional densities and ADM normalizations of
Section~\ref{four:sec:output}.  Suppose the densities are integrable,
the finite ADM charges satisfy $E>|P|$, the dominant energy condition
holds, and the boundary satisfies $H+\tr_{\mathrm{tan}}K\le0$ with
normal into $M$.  Then
\[
 \sqrt{E^2-|P|^2}\ge
       \frac12\left(\frac{A_*(g)}{\omega_3}\right)^{2/3},
 \qquad \omega_3=2\pi^2,
\]
where $A_*$ uses every full cut of Definition~\ref{four:def:cuts}.
\end{corollary}

\begin{proof}
Set $m=\sqrt{E^2-|P|^2}$.  The rest-end replacement
\cite[Proposition~8.8]{NumericalCompanion} applies with exactly these
one-ended weak-decay hypotheses.  It produces data $(g_R,K_R)$ on the
same exterior, satisfying the dominant energy condition and the weak
future-trapping inequality, with compactly supported $K_R$ and an exact
static Schwarzschild--Tangherlini tail.  Their energy is
$m_R\to m$, their momentum is zero, and their full enclosing infima
satisfy
\begin{equation}\label{b:four-transfer:one-sided}
 A_*(g_R)\ge(1-o_R(1))A_*(g).
\end{equation}
This is the lower bound required below; no equality or convergence of
these infima is used.

Fix $R$ first.  The strictification of
\cite[Lemma~8.9]{NumericalCompanion} gives
$(g_{R,s},K_{R,s})=(e^{2s\phi_R}g_R,e^{s\phi_R}K_R)$, for all
sufficiently small $s>0$.  It preserves completeness and compact tensor
support, makes the boundary strictly future trapped, and gives, for a
fixed $0<\delta_1<1$,
\[
 \mu_{R,s}-|J_{R,s}|\ge c_{R,s}\,r^{-4-\delta_1}>0,
 \quad g_{R,s}-\delta=O_j(r^{-2})\ (j\ge0),
 \quad R_{g_{R,s}}=O(r^{-4-\delta_1}).
\]
The densities are integrable and the ADM energy is finite.  Thus every
hypothesis of Theorem~\ref{four:thm:boundary-deformation} is met, with
all data and their constants now fixed.  Its complete boundary solve
and Corollary~\ref{four:thm:prepared-energy} give
\[
 E_{g_{R,s}}\ge
       \frac12\left(\frac{A_*(g_{R,s})}{\omega_3}\right)^{2/3}.
\]
The latter corollary already exhausts the domain, takes $N\to\infty$,
and removes the conformal floor for this fixed pair $(R,s)$.

Still at fixed $R$, strictification supplies
$E_{g_{R,s}}\to m_R$ and $A_*(g_{R,s})\to A_*(g_R)$ as $s\downarrow0$.
Taking this limit and then using \eqref{b:four-transfer:one-sided} gives
\[
 m_R\ge\frac12\left(\frac{A_*(g_R)}{\omega_3}\right)^{2/3}
 \ge\frac12\left(\frac{(1-o_R(1))A_*(g)}{\omega_3}\right)^{2/3}.
\]
Finally take $R\to\infty$.  The mass and density conventions coincide
with those of both preparation results, so no normalization factor is
introduced.  The proof has used their preparations and the boundary
construction, without importing a numerical spacetime theorem.
\end{proof}

\section{What changes when the tensor reaches infinity}
\label{b:transition}

The boundary construction controls two different interfaces of the same
deformation.  At the inner boundary it prescribes the normal flux so that
the new mean curvature is negative.  At infinity it estimates the total
flux so that the new energy is no larger than the old energy, up to an
error tending to zero.  Compact support of the second fundamental form
simplifies the second task: sufficiently far out, the trace equation
contains only the graph.  We now explain precisely how maximality replaces
that simplification.  This calculation is also the reason to retain a
separate maximal construction: it carries the original decaying tensor
through the deformation instead of removing it from the end.

We work in spatial dimension four.  Write $g$ for the background metric,
$K$ for its second fundamental form, and assume $\tr_gK=0$.
Let $N\ge4$ be a fixed deformation parameter, let $0<\varepsilon<1$,
and choose a smooth nonincreasing function $\vartheta$ equal to one on
$(-\infty,0]$ and zero on $[1,\infty)$.  Set
\[
 p(t)=N\vartheta(Nt),\qquad
 l(t)=\exp\!\left(\int_0^t p(s)\,\dd s\right),\qquad
 L_0(t)=e^{2t}l(t),\qquad
 \ell=e^{-\varepsilon N},\quad \eta_N=\ell^{3/2}.
\]
For functions $f,Z$, the equation
\[
 Z=t+\tfrac16\log D,\qquad D=1+l(t)^2|\nabla f|_g^2
\]
defines $t$ uniquely: its right side has positive $t$ derivative
$1+p(1-D^{-1})/3$, and tends to the two ends of the real line with $t$.
Put
\[
 d=D^{-1},\quad w=\frac{l\nabla f}{\sqrt D},\quad v=|w|^2=1-d,
 \quad\chi=\frac{3d}{3+pv},\quad
 A_\chi=I-\frac{3+p}{3+pv}w\otimes w,
 \quad u=L_0\sqrt D.
\]
The deformed metric and flux are
\[
 \widehat g=e^{2t}(g+l^2\dd f^2),\qquad
 V=uA_\chi\bigl(3\nabla Z+K(w,\cdot)\bigr).
\]
These are exactly the four-dimensional boundary variables.  Thus the
inner boundary calculation remains unchanged.  On a face where $f$ is
constant, writing $H=\operatorname{div}_B\nu$ and $P_B=\tr_BK$ for
the normal into the exterior gives
\begin{equation}\label{b:transition:boundary}
 \frac{V_\nu}{u}
 =d(H+3\partial_\nu t)-H+w_\nu(F-P_B),
 \qquad
 \widehat H=e^{-t}\sqrt d\,(H+3\partial_\nu t).
\end{equation}
These are the local identities of
Lemma~\ref{four:sys:boundary-lemma}, with
$H^f=lD^{-1/2}\nabla^2f$ and $F=\tr_{A_\chi}(K+H^f)$.
All graph variables and the normal convention agree with that lemma.
Its proof uses only the constant trace of $f$ on the face and the
four-dimensional conformal boundary formula; it imposes no support
or energy hypothesis on $K$.
In particular the prescribed flux
$V_\nu/u=-H+w_\nu(F-P_B)-N_Bd$, with $N_B>0$, still makes
$\widehat H=-e^{-t}\sqrt d\,N_B<0$.

The new issue lies at the other end of the domain.  There the compactly
supported offset $C$ vanishes and the trace equation is
\begin{equation}\label{b:transition:trace}
 \tr_{A_\chi}(K+H^f)=\eta_N f.
\end{equation}
Maximality is useful here for an exact algebraic reason:
\[
 \tr_{A_\chi}K
 =-\frac{3+p}{3+pv}\frac{l^2}{D}K(\nabla f,\nabla f).
\]
Consequently \eqref{b:transition:trace} is, along any fixed solution,
the homogeneous linear equation
\begin{equation}\label{b:transition:drift}
 A_\chi:\nabla^2 f+b\cdot\nabla f-c_f f=0,
 \qquad
 b=-\frac{3+p}{3+pv}\frac l{\sqrt D}K(\nabla f,\cdot),
 \quad c_f=\frac{\eta_N\sqrt D}{l}\ge\frac{\eta_N}{e}>0.
\end{equation}
Here $l\le e$ follows directly from its definition.  Thus retaining $K$
introduces a drift, but no independent forcing of the graph.  An arbitrary
nonmaximal tensor would leave the term $\tr_gK$ and would not have this
property.

\begin{lemma}[The maximal end equation]\label{b:transition:end}
Suppose $1<q<2$, $0<\delta<q-1$,
$g-\delta_{\mathrm{Eucl}}=O_2(r^{-q})$ and
$K=O_1(r^{-1-q})$ on a four-dimensional end.  Assume
$\tr_gK=0$ and $R_g=O(r^{-4-\delta})$.
For fixed $N,\varepsilon$, consider smooth solutions on successively
larger spherical truncations, with $f=Z=0$ on the outer sphere, of
\eqref{b:transition:trace} and
\begin{equation}\label{b:transition:divergence}
 \operatorname{div}_gV
 =u\left[\delta_0\mathcal T+\rho+
       \delta_0\eta_N |w||\nabla f|
       -m_0(t)C_N(\rho_0+v\rho_1)\right].
\end{equation}
Here $0<\delta_0<1$, $\rho,\rho_0=O(r^{-4-\delta})$ with
their differentiated bounds, $\rho_1=O(r^{-2\gamma})$ for a fixed
$\gamma>1$, $m_0$ is a smooth bounded function, and $\mathcal T$ is
the quadratic term in the four-dimensional scalar identity.
Assume that the solutions have a common bounded range and that, for every
fixed derivative order, the background coefficients, weights and
solutions have uniform fixed-$N$ classical bounds on end unit patches,
including the outer boundary patches.  Assume also that $A_\chi$ is
uniformly elliptic there.  These bounds may depend on $N$ and the
derivative order; they are independent of the outer radius.  They are
separate assumptions from the displayed $O_2/O_1$ decay rates.
Then, uniformly on those truncations,
\[
 |\nabla^j f|\le C_j(N)e^{-c_j(N)r}\quad(j\ge0),
 \qquad |Z|+|t|\le C(N)r^{-2}.
\]
Every smooth local limit satisfies $Z,t=O_2(r^{-2})$.  Before taking
that limit, its scalar equation is
\begin{equation}\label{b:transition:zero-graph}
 \begin{split}
 3\Delta_g Z
 +3[p(Z)+2-\delta_0(p(Z)+1)]|\dd Z|^2
 ={}&\tfrac{\delta_0}{2}|K|^2+\rho
       -m_0(Z)C_N\rho_0+\mathcal E_N,
 \end{split}
\end{equation}
where $\mathcal E_N$ and each required fixed derivative have
exponential unit-patch bounds.
\end{lemma}

\begin{proof}
The coefficients in \eqref{b:transition:drift} are uniformly bounded
and elliptic at fixed $N$.  For sufficiently small $c>0$, substitution
of $e^{-cr}$ into its left side gives an upper bound
\[
 e^{-cr}\bigl(C(N)c^2+C(N)c-\eta_N/e\bigr)<0
\]
outside one fixed sphere.  A multiple of this positive supersolution
dominates $|f|$ on that sphere; the zero outer value is already dominated.
Comparison for $f$ and $-f$ gives exponential decay, uniformly in the
truncation radius.  The assumed unit-patch estimates, followed by the
smooth linear interior and zero-Dirichlet boundary estimates for
\eqref{b:transition:drift}, give the asserted differentiated decay.
These are the smooth-coefficient estimates of
\cite[Chs.~6--9]{GilbargTrudinger2001} and
\cite{AgmonDouglisNirenberg1959}.
The implicit formula for $t$ then shows that $t-Z$ and the graph metric
error have the same type of bounds.

It remains to identify what survives when those graph errors are removed.
Set the graph height and its first two derivatives equal to zero in
the coefficients, retaining the exponentially small difference as an
error.  Then $H^f=0$, $F=0$, $w=0$, $D=d=\chi=1$, $t=Z$ and
$u=L_0(Z)$.  The smooth scalar quadratic form at this zero graph is
\begin{equation}\label{b:transition:quadratic}
 \mathcal T=\tfrac12|K|^2+3(p(Z)+1)|\dd Z|^2.
\end{equation}
One can also verify this directly from its block formula.  Choose any
unit vector $e$, write $j=K(e,e)$, $M=K(e,\cdot)|_{e^\perp}$ and
$T=K|_{e^\perp}$, and use $\tr T=-j$.  The terms involving $K$ become
$\tfrac12|T^0|^2+|M|^2+\tfrac23j^2$.
Since $|T|^2=|T^0|^2+j^2/3$, this is exactly $|K|^2/2$.
The terms linear in $\dd Z$ contain $w$ and vanish, leaving the
second term in \eqref{b:transition:quadratic}.  In particular setting
the graph to zero does not set the tensor to zero.

Now $(\log L_0)'=p+2$.  Divide
\eqref{b:transition:divergence} by $L_0$, insert
\eqref{b:transition:quadratic}, and move the quadratic gradient term
to the left.  Smooth dependence on the bounded solution variables
absorbs the graph errors into $\mathcal E_N$ and gives
\eqref{b:transition:zero-graph}.

Define $\Phi(0)=0$ and
\[
 \Phi'(z)=\exp\!\left(\int_0^z
     [p(s)+2-\delta_0(p(s)+1)]\,\dd s\right).
\]
On the common bounded range of $Z$, both $\Phi'$ and its reciprocal
are bounded at fixed $N$.  The chain rule gives
\[
 \Delta_g\Phi(Z)=\frac{\Phi'(Z)}3
 \left[\tfrac{\delta_0}{2}|K|^2+\rho
       -m_0(Z)C_N\rho_0+\mathcal E_N\right].
\]
The right side is $O_N(r^{-4-\delta})$, because
$|K|^2=O(r^{-2-2q})$ and $2+2q>4+\delta$.
With $\alpha=\delta/2$, set $b_0=r^{-2}-r^{-2-\alpha}$.
For sufficiently large $r$, this is positive and
\[
 -\Delta_g b_0
 =\alpha(2+\alpha)r^{-4-\alpha}+O(r^{-4-q})
 \ge c r^{-4-\alpha}.
\]
A fixed multiple of $b_0$ dominates both the absolute source and
$|\Phi(Z)|$ on the fixed inner sphere.  Since $\Phi(Z)=0$ at the outer
sphere, comparison for both signs gives $|\Phi(Z)|\le C(N)r^{-2}$,
and hence $|Z|+|t|\le C(N)r^{-2}$.  These estimates hold before
exhaustion, so the local limit has the asserted normalization at
infinity.

Finally rescale annuli $R<r<2R$ to unit size.  The metric has uniformly
controlled scaled $C^{1,\beta}$ coefficients for each fixed $0<\beta<1$.
Interior $W^{2,p_1}$ estimates with $p_1>4$ first give
$|\dd Z|=O_N(r^{-3})$.  The first derivative bound on $K$, the
differentiated weights and the exponential error then give a controlled
scaled H\"older norm for the source.  Schauder estimates yield
$\nabla^2Z=O_N(r^{-4})$.  The same follows for $t$ from the graph error.
\end{proof}

The term $|K|^2/2$ has two consequences for the rest of the construction.
First, the end proof must use the original decay of $K$; it cannot invoke
a result whose hypothesis is $K=0$ outside a compact set.  Second, the
continuation must control the corresponding noncompact drift before
these fixed-parameter estimates are available.  The appropriate tail
weights follow from a short calculation.  Choose
\[
 1<\gamma<\min\{2,q-\delta/2\},\qquad
 \rho_0\asymp r^{-4-\delta},\quad \rho_1\asymp r^{-2\gamma}.
\]
This interval is nonempty because $\delta<q-1$.
For $a=|w|$ one has
\begin{equation}\label{b:transition:tail-budget}
 a r^{-2-q}
 \le\tfrac12a^2r^{-2\gamma}
       +\tfrac12r^{-4-2q+2\gamma}
 \le C(v\rho_1+\rho_0).
\end{equation}
The first inequality is Young's inequality; the second uses
$2\gamma<2q-\delta$.  Also $|K|^2\le C\rho_0$.
Thus the drift error controlled by
$\Pi_N(|K|^2+a|\nabla K|)$ fits the same kind of floor penalty as
in the boundary construction, with exponents adapted to the original
end.  Since $4\gamma>4$, the square of $\rho_1$ is integrable.
That last fact is exactly what permits its contribution to the flux
loss to be absorbed with an error tending to zero.

The local comparison is now explicit.  Signed boundary faces and their
inner flux formula are shared.  The maximal route supplies its own
strict approximation and an all-black threshold boundary; it proves
the noncompact drift estimate and its complete boundary continuation
with the weights above.  Once these yield the fixed-$N$ estimates,
Lemma~\ref{b:transition:end} gives the correct end normalization.
In applying that lemma, unit-scale convolution of the strict end
coefficients supplies their uniform bounds at every fixed derivative
order; a subsequent trace-free projection preserves maximality.  This
is an actual preparation step, not a consequence of smoothness and
$O_2/O_1$ decay alone.  The displayed decay rates are the ones retained
from the original data and used in the final annular estimates.
The final comparison still uses the actual metric $\widehat g$, its
minimal enclosure, and its ADM flux.  Maximality preserves a method
with the original tensor tail; it does not replace those existence,
boundary, or enclosure arguments.

\section{The maximal comparison and its scope}
\label{four:maximal:section}

Maximality allows the boundary deformation to retain a decaying second
fundamental form all the way to infinity.  The reason is the homogeneous
trace equation in Section~\ref{b:transition}: when $\tr_gK=0$, the tensor
contributes a drift proportional to the graph gradient.  Thus the graph
can decay exponentially even though $K$ does not vanish outside a compact
set.  The conformal equation still sees $|K|^2/2$.  We now construct the
solutions for which that end calculation applies.

There are two geometric tasks before solving the equations.  We must
create strict dominant-energy slack without destroying maximality or
changing the limiting energy and enclosing area.  We must then select an
inner boundary whose outward collars control the large graph slopes.
The remaining analysis has the same input and output as the boundary
method: nonnegative scalar curvature, negative inner mean curvature,
a lower metric comparison on every cut, and an upper error in energy.
The noncompact drift requires the tail weights and continuation estimates
proved in this part.

The following theorem records the maximal-vacuum application considered
here.  Corollary~\ref{b:four-weak} gives a broader numerical comparison
under its one-ended weak-decay and dominant-energy hypotheses.  That
corollary first replaces the distant end and tensor.  The construction
below retains the decaying tensor tail through the deformation.

Throughout, four refers to the spatial dimension. The associated spacetime
dimension is five, and a horizon has dimension three. Write
\(\omega_3=|S^3|=2\pi^2\). The notation \(O_j(r^{-q})\) includes coordinate
derivatives through order \(j\), with the corresponding additional decay.

\begin{theorem}\label{four:maximal:thm:main}
Let \((M,g,K)\) be a smooth connected orientable initial-data exterior,
diffeomorphic to \(\R^4\) minus an open ball, complete up to its connected
compact boundary \(S\cong S^3\), and with one asymptotically flat end.
Here \(K\) is a symmetric covariant two-tensor. Assume the vacuum constraint
equations with zero cosmological constant and maximality:
\begin{equation}\label{four:maximal:intro:constraints}
 R_g+(\tr_gK)^2-|K|_g^2=0,\qquad
 \divg_g\bigl(K-(\tr_gK)g\bigr)=0,\qquad \tr_gK=0.
\end{equation}
Assume a smooth effective \(T^2=U(1)\times U(1)\) action preserving \(g,K\)
which, in an asymptotically Euclidean chart \(\R^4=\mathbb C^2\), is
\[
 (e^{i\alpha},e^{i\beta})\cdot(z_1,z_2)
       =(e^{i\alpha}z_1,e^{i\beta}z_2)
\]
at all sufficiently large radii. No cohomogeneity-one assumption is made.
In that chart suppose
\begin{equation}\label{four:maximal:intro:decay}
 g_{ij}-\delta_{ij}=O_2(r^{-q}),\qquad
 K_{ij}=O_1(r^{-q-1}),\qquad q>1.
\end{equation}
Suppose the following ADM limits exist and are finite:
\begin{align}
 E&=\frac{1}{6\omega_3}\lim_{r\to\infty}
     \int_{|x|=r}(\partial_jg_{ij}-\partial_ig_{jj})
                  \nu_\delta^i\dd A_\delta,\label{four:maximal:intro:energy}\\
 P_i&=\frac{1}{3\omega_3}\lim_{r\to\infty}
     \int_{|x|=r}\bigl(K_{ij}-(\tr_gK)g_{ij}\bigr)
                  \nu_\delta^j\dd A_\delta.\label{four:maximal:intro:momentum}
\end{align}
Assume \(P=0\) and \(E>0\).

Orient \(S\) by its unit normal \(\nu\) toward the end and set
\(H_S=\divg_S\nu\). Suppose
\begin{equation}\label{four:maximal:intro:mots}
 \Theta_K(S):=H_S+\tr_SK=0.
\end{equation}
Suppose \(S\) is outermost toward the end: there is no compact smooth
embedded two-sided hypersurface in the interior enclosing \(S\), oriented
toward the end, with \(\Theta_K\le0\) everywhere. Suppose also that \(S\)
is outer area-minimizing: every compact smooth embedded enclosing
hypersurface has three-volume at least \(A=|S|_g\). Disconnected and
noninvariant enclosing competitors are included.
Then
\begin{equation}\label{four:maximal:intro:main-inequality}
 E\ge\frac12\left(\frac{A}{2\pi^2}\right)^{2/3}.
\end{equation}
\end{theorem}

An enclosing hypersurface separates \(S\) from the selected end and bounds
with \(S\) a compact region. The original \(S\) is itself allowed when
taking the enclosing-area infimum. All hypersurface areas in this part
are three-dimensional volumes. In particular, outer area-minimization in
Theorem~\ref{four:maximal:thm:main} is not restricted by the torus action.
The torus action and outermostness specify the stated application class;
the construction below does not use them.  Ball-exterior topology supplies
the compact filling, and vacuum supplies the initial strictification.
Outer area-minimization identifies the final enclosing-area infimum with
\(A\).  Neither \(H_S=0\) nor \(\tr_SK=0\) is imposed.  The conditions
\(P=0\) and \(E>0\) identify the energy in
\eqref{four:maximal:intro:energy} with the ADM rest mass.

The normalization is sharp. The time-symmetric Schwarzschild--Tangherlini
family \cite{Tangherlini1963} has, in isotropic coordinates of radius
\(\varrho\), the spatial metric
\[
 g_m=\left(1+\frac{m}{2\varrho^2}\right)^2\delta,
 \qquad K=0,
 \qquad \varrho\ge\sqrt{m/2};
\]
see also \cite[p.~83]{BrayLee2009}. Its leading metric term is
\(m\varrho^{-2}\delta\), so \eqref{four:maximal:intro:energy} gives energy \(m\).
The boundary sphere has volume
\(8\omega_3(m/2)^{3/2}=\omega_3(2m)^{3/2}\), making the right side of
\eqref{four:maximal:intro:main-inequality} equal to \(m\).


Unless a different metric is indicated, norms, contractions, covariant
derivatives and index raising use the background metric \(g\). We use
\(\Delta=\divg\grad\). For symmetric two-tensors define
\begin{align}
 P_D&=D-(\tr_gD)g, &
 \mathsf{B}(D_1,D_2)&=D_1:D_2-(\tr_gD_1)(\tr_gD_2),\label{four:maximal:intro:bilinear}\\
 \mu&=\tfrac12\bigl(R_g-\mathsf{B}(K,K)\bigr), &
 J&=\divg_gP_K.\label{four:maximal:intro:densities}
\end{align}
Thus the vacuum constraints give \(\mu=J=0\).
For an oriented hypersurface the expansion is
\(\Theta_K=H+\tr_{\mathrm{tan}}K\), with the normal toward the
specified exterior. On an inner boundary this normal points into the
exterior domain; on a distant coordinate sphere the outward normal points
toward infinity. We retain these conventions in every boundary flux.

For an inner cut \(S_0\), let
\[
 \mathcal A_g(S_0)=\inf\{|\Sigma|_g:
              \Sigma\text{ is a smooth enclosing cut of }S_0\},
\]
using the full class of cuts in Definition~\ref{four:def:cuts} and
including \(S_0\) when it is smooth. The proof uses this infimum until
the final application of the outer area-minimizing hypothesis.


\subsection{The geometric output and the route to it}

It is useful to state what the construction will produce for strict
data before introducing its equations.  This is the point at which
the analytic problem rejoins the mass--area comparison.

\begin{proposition}[Maximal deformation output]
\label{four:maximal:prop:geometric-output}
Let $(M,g,K)$ be a smooth connected orientable four-dimensional
ball exterior, complete including its compact boundary $S$, with one
Euclidean coordinate end.  Let $r\ge1$ be a fixed smooth positive
extension of its Euclidean radius.  Suppose $1<q<2$, $0<\delta<q-1$,
\[
 \tr_gK=0,\quad g-\delta_{\rm Eucl}=O_2(r^{-q}),\quad
 K=O_1(r^{-1-q}),\quad R_g=O(r^{-4-\delta}),
 \qquad \mu-|J|\ge c_0r^{-4-\delta}>0,
\]
and suppose $H_S+\tr_SK<0$.  Assume finite ADM energy and uniform
bounds of every fixed derivative order on end patches of unit size.
Write $A_* =\mathcal A_g(S)$ for the infimum over all smooth enclosing
cuts, including disconnected cuts.

There is a smooth cut $B$ enclosing $S$, with connected one-ended
exterior $\Omega$, such that for each $0<\eps<1$ and all sufficiently
large integers $N$ there is a smooth metric $\widehat g_N$ on
$\overline\Omega$ satisfying
\[
 \widehat g_N\ge e^{-2\eps}g,\qquad
 R_{\widehat g_N}\ge0,\qquad \widehat H_B<0,
 \qquad
 E_{\widehat g_N}\le E_g+
       \frac{\Pi_N}{3\omega_3}
            (e^{-\eps N}+e^{-\eps N/2}).
\]
Here $\Pi_N$ has polynomial growth at fixed strict data and $\eps$.
The new metric is complete including $B$, has
$\widehat g_N-\delta_{\rm Eucl}=O_2(r^{-q})$ and
$R_{\widehat g_N}=O(r^{-4-\delta})$, and every smooth cut enclosing
$B$ has $\widehat g_N$-area at least $e^{-3\eps}A_*$.
\end{proposition}

The proof occupies Sections~\ref{four:maximal:sec:preparation}--
\ref{four:maximal:sec:comparison}.  First we obtain $B$ and its outward
collars.  Next we prove a lower bound for $t$, followed by height,
boundary-flux and range estimates for every possible continuation
solution.  Local regularity and an unrestricted scalar solve then
define the compact continuation map.  Only after existence and
fixed-$N$ exhaustion do we use the end calculation of
Lemma~\ref{b:transition:end}.  The final flux estimate proves the
displayed energy bound.  This order separates the polynomial estimates
needed for the mass limit from arbitrary fixed-$N$ regularity constants.

The trace $\tr_gK$ vanishes throughout this construction.  Its
positive regularizing parameter is $\eta_N=\ell^{3/2}$; this parameter
is distinct from the trace.
The scalar weight is denoted by $L_0=e^{2t}l(t)$.
First fix the strict geometry and its decay exponents, then $\eps>0$ and
a sufficiently large $N$.  Exhaust the outer radius at fixed $N$, send
$N\to\infty$ using polynomial constants only, then send $\eps\downarrow0$
and remove the strict approximation.  Arbitrary fixed-$N$ regularity
constants do not enter the final mass limit.


\section{Strict data and the geometric inner boundary}
\label{four:maximal:sec:preparation}

This section produces the strict data and the inner boundary used in the
deformation.  The approximation is controlled in ADM energy and in the
infimum over all enclosing areas.  The geometric construction also supplies
a comparison principle for closed sets described by smooth exterior
supports; this will be used before any derivative estimates for the
deformation are available.

\subsection{Strict conformal approximation}

The spacetime Poisson construction of
Jaracz~\cite{Jaracz2025StrictDEC} is the model for this reduction.
Here the conformal exponents are four-dimensional, the inner Neumann
datum is nonzero, and we prove preservation of maximality and comparison
for the full enclosing-area infimum.

Decrease the decay exponent, if necessary, so that \(1<q<2\), and fix
\(0<\delta<q-1\).  Extend the end radius to a positive smooth function
\(r\ge1\) on the exterior.  We use the enclosing-area infimum
\(\mathcal A_g(S)\) defined above.

\begin{proposition}[Reduction to strict data]
\label{four:maximal:prop:strict-data}
The data of Theorem~\ref{four:maximal:thm:main} admit smooth approximations
\((g_j,K_j)\) on the same exterior such that
\begin{equation}
\begin{aligned}
 \tr_{g_j}K_j&=0,&
 \mu_j-|J_j|_{g_j}&\ge c_j r^{-4-\delta},\\
 \Theta_{K_j}(S)&<0,&
 R_{g_j}&=O(r^{-4-\delta}),
 \qquad c_j>0.
\end{aligned}
\label{four:maximal:prep:strict-properties}
\end{equation}
They have the prescribed \(O_2(r^{-q})\) metric decay and
\(O_1(r^{-1-q})\) tensor decay, and all their coordinate derivatives
are bounded uniformly over end patches of unit size, with constants
allowed to depend on \(j\) and the derivative order.  Moreover,
\begin{equation}
 E_{g_j}\longrightarrow E_g,
 \qquad \mathcal A_{g_j}(S)\longrightarrow\mathcal A_g(S).
\label{four:maximal:prep:approx-limits}
\end{equation}
Consequently, it suffices to prove
\(E_g\ge\tfrac12(\mathcal A_g(S)/\omega_3)^{2/3}\) for such strict maximal
data on this exterior.  Neither vacuum nor marginality is required of
the intermediate approximations.
\end{proposition}

\begin{proof}
Choose a positive smooth function \(b\ge|K|_g\) which equals a
constant multiple of \(r^{-1-q}\) sufficiently far out.  We first solve
\begin{equation}
 -\Delta_g\phi
   =b\sqrt{|\dd\phi|_g^2+r^{-6}}+r^{-4-\delta},
 \qquad \partial_\nu\phi=-1\text{ on }S,
 \qquad \phi\longrightarrow0\text{ at infinity},
\label{four:maximal:prep:phi-equation}
\end{equation}
where \(\nu\) points into the exterior.  In particular, the normal
derivative in this equation is the negative of the domain's outward
normal derivative at its inner boundary.

Here are the existence and decay details.  On the compact truncation
\(M_R\), impose \(\phi=0\) on the outer coordinate sphere and replace
\(b\) on the right by \(\lambda b\), \(0\le\lambda\le1\).
The linearization, with homogeneous boundary data, is
\[
 v\longmapsto-\Delta_gv
 -\lambda b\frac{\langle\dd\phi,\dd v\rangle_g}
                        {\sqrt{|\dd\phi|_g^2+r^{-6}}}.
\]
Its drift is bounded by \(b\).  Its kernel is zero by the strong
maximum principle and the boundary point lemma, since the outer
Dirichlet face is nonempty and the inner face has homogeneous normal
derivative.  The mixed operator has Fredholm index zero, by deformation
to the mixed Laplacian, and hence is invertible.  The two boundary
components are disjoint, smooth, and compact, so ordinary scalar
elliptic estimates apply without a junction between boundary types;
see \cite[Chapters~3, 6, and~9]{GilbargTrudinger2001} and
\cite[Theorems~7.3 and~15.2--15.3]{AgmonDouglisNirenberg1959}.

For fixed \(R\), the global mixed \(W^{2,p}\) estimate, \(p>4\),
and interpolation give
\begin{equation}
 \|\phi\|_{W^{2,p}(M_R)}
 \le C_{R,p}\bigl(1+\|\dd\phi\|_{L^p(M_R)}
                         +\|\phi\|_{L^p(M_R)}\bigr)
 \le C'_{R,p}(1+\|\phi\|_\infty).
\label{four:maximal:prep:phi-estimate}
\end{equation}
These constants are uniform in \(\lambda\).  If the suprema were
unbounded, divide by their sup norms and take a subsequence.  Compact
embedding of \(W^{2,p}\) in \(C^{1,\alpha}\), with a smaller
\(\alpha>0\), gives a nonzero limit \(v\) satisfying
\[
 -\Delta_gv=\lambda_*b|\dd v|_g,
 \qquad \partial_\nu v=0\text{ on }S,
 \qquad v=0\text{ on the outer face}.
\]
This is a homogeneous linear equation with a bounded measurable drift:
set the drift to \(\lambda_*b\grad v/|\dd v|\) where
\(\dd v\ne0\), and to zero otherwise.  The maximum principle and
boundary point lemma again exclude a nonzero solution.  Thus
\eqref{four:maximal:prep:phi-estimate} is bounded, and Schauder estimates and
bootstrapping give closedness in the continuity argument.  The
\(\lambda=0\) mixed Poisson problem is solvable, so the implicit
function theorem and compactness give a smooth solution for
\(\lambda=1\).  It is nonnegative: a negative minimum is excluded
in the interior by \(-\Delta\phi>0\), on the outer face by its
zero value, and on the inner face by \(\partial_\nu\phi=-1\).

We next bound these solutions independently of \(R\).  On the end
consider
\[
 T(r)=r^{-2}(1-Ar^{-\delta}).
\]
Choose a fixed \(A>0\), then a large fixed \(r_0\), so that \(T>0\)
and
\begin{equation}
 (-\Delta_g-b|\dd\,\cdot\,|)T
 \ge c\,r^{-4-\delta}\qquad(r\ge r_0).
\label{four:maximal:prep:radial-barrier}
\end{equation}
Indeed,
\(-\Delta_\R T=A\delta(2+\delta)r^{-4-\delta}\), while both
the metric error and \(b|\dd T|\) are \(O(r^{-4-q})\).
Also
\[
 -\Delta_g\phi-b|\dd\phi|
 \le b r^{-3}+r^{-4-\delta}\le C r^{-4-\delta}.
\]
Comparison on \(\{r_0\le r\le R\}\), including the zero outer
value, therefore yields
\begin{equation}
 0\le\phi_R(x)\le C(1+m_R)r(x)^{-2},
 \qquad m_R=\max_{\{r\le r_0\}}\phi_R.
\label{four:maximal:prep:uniform-tail}
\end{equation}
If \(m_R\) were unbounded along expanding truncations, divide by
\(m_R\).  Local versions of \eqref{four:maximal:prep:phi-estimate}, including
the inner boundary estimates, give a nonnegative limit attaining
value one in the fixed core.  It solves
\(-\Delta v=b|\dd v|\), has homogeneous inner normal derivative,
and is \(O(r^{-2})\) by \eqref{four:maximal:prep:uniform-tail}.  Its positive
maximum is therefore attained at a finite point.  The strong maximum
principle or the boundary point lemma excludes it.  We have a uniform
core bound, and diagonal compactness gives a smooth solution of
\eqref{four:maximal:prep:phi-equation}.

Rescale annuli by their radii.  The bound \(\phi=O(r^{-2})\),
the scaled metric bounds furnished by \(g-\delta=O_2(r^{-q})\),
and the linear gradient growth in \eqref{four:maximal:prep:phi-equation} first
give scaled \(W^{2,p}\) estimates by interpolation.  They give
H\"older control of the gradient; the right side then has the
scaled H\"older control needed for Schauder estimates.  Hence
\begin{equation}
 \phi=O_2(r^{-2}).
\label{four:maximal:prep:phi-decay}
\end{equation}
The source in \eqref{four:maximal:prep:phi-equation} is
\(O(r^{-4-q})+O(r^{-4-\delta})\), and is integrable.  The
divergence theorem shows that
\[
 I_\phi=\lim_{R\to\infty}\int_{S_R}\partial_{\nu_g}\phi\dd A_g
\]
exists.  Replacing this by its Euclidean-coordinate flux changes it
by \(O(R^{-q})\), which tends to zero.

For small \(s>0\) set
\(g_s=e^{2s\phi}g\) and \(K_s=e^{s\phi}K\).
The trace remains zero.  The connection change is
\[
 \Gamma(g_s)^k_{ij}-\Gamma(g)^k_{ij}
 =s(\delta_i^k\phi_j+\delta_j^k\phi_i-g_{ij}\phi^k).
\]
Contracting it against the scaled trace-free tensor, and using the
four-dimensional scalar-curvature formula, gives
\begin{equation}
\begin{split}
 e^{2s\phi}\mu_s
   &=\mu-3s\Delta_g\phi-3s^2|\dd\phi|^2,\\
 e^{s\phi}J_s
   &=J+3sK(\grad\phi,\cdot),\\
 \Theta_{K_s}(S)
   &=e^{-s\phi}\bigl(\Theta_K(S)+3s\partial_\nu\phi\bigr).
\end{split}
\label{four:maximal:prep:conformal-constraints}
\end{equation}
The last coefficient is three because the boundary has dimension
three.  A covector norm contributes one more factor \(e^{-s\phi}\),
so the first two identities and \(\mu=J=0\) imply
\begin{align}
 e^{2s\phi}(\mu_s-|J_s|_{g_s})
 &\ge 3s\bigl[-\Delta\phi-|K||\dd\phi|-s|\dd\phi|^2\bigr]
 \notag\\
 &\ge3s\bigl[r^{-4-\delta}-s|\dd\phi|^2\bigr].
\label{four:maximal:prep:strict-gap}
\end{align}
The ratio \(|\dd\phi|^2/r^{-4-\delta}\) is bounded, by
\eqref{four:maximal:prep:phi-decay} and \(\delta<1\).  Thus small \(s>0\)
gives a positive gap of the asserted order.  The boundary expansion
is \(-3s e^{-s\phi}<0\).  Moreover,
\[
 R_{g_s}=e^{-2s\phi}
 \bigl(|K|^2-6s\Delta\phi-6s^2|\dd\phi|^2\bigr)
       =O(r^{-4-\delta}),
\]
because \(2+2q>4+\delta\).  The original metric and tensor decay
is preserved, since \(q<2\).

Only \(-6s\partial_i\phi\) survives from the conformal change
in the ADM integrand.  Products with a metric error or with \(\phi\)
have vanishing integrated flux by \eqref{four:maximal:prep:phi-decay}.  Thus
\begin{equation}
 E_{g_s}=E_g-\frac{s}{\omega_3}I_\phi\longrightarrow E_g.
\label{four:maximal:prep:conformal-mass}
\end{equation}
Since \(\phi\) is bounded, the comparison
\(e^{-2s\|\phi\|_\infty}g\le g_s\le
e^{2s\|\phi\|_\infty}g\) holds on the entire exterior.  Its
three-dimensional area comparison holds for every enclosing cut,
and therefore for their infima.  In particular,
\(\mathcal A_{g_s}(S)\to\mathcal A_g(S)\), with no compactness assumption on
minimizing cuts.

\paragraph{Smoothing the end.}
The conformal step has produced strict dominance and convergence of energy
and enclosing-area infima.  The differentiated end estimates will also
need bounds of every fixed derivative order on patches of unit size.
The original \(O_2/O_1\) hypotheses do not supply those bounds; we obtain
them by smoothing each fixed strict pair sufficiently far out.

Fix one of these strict pairs and temporarily denote it by \((g,K)\).
Let \(h=g-\delta\) in the end chart.  Convolve \(h\) and \(K\)
with a nonnegative smooth kernel of integral one supported in a unit
ball.  Interpolate to the original data with a cutoff which is zero
for \(r\le R_{\rm sm}\), one for \(r\ge2R_{\rm sm}\), and has
derivatives of order \(i\) bounded by \(C_iR_{\rm sm}^{-i}\).
Call the resulting metric and tensor \(\widetilde g,\widetilde K\),
and project the tensor to its \(\widetilde g\)-trace-free part:
\[
 K^{\rm sm}=\widetilde K-\tfrac14
                (\tr_{\widetilde g}\widetilde K)\widetilde g.
\]
For sufficiently large \(R_{\rm sm}\), the metric remains positive
definite.  The available original derivatives give
\[
 h*\kappa-h=O(r^{-q-1}),\qquad
 D(h*\kappa-h)=O(r^{-q-2}),\qquad
 K*\kappa-K=O(r^{-q-2}).
\]
Here \(D\) denotes coordinate differentiation only in this paragraph.
For the trace projection, use the original trace-free identity in the
form
\[
 \tr_\delta K=(\delta^{ij}-g^{ij})K_{ij}
                       =O_1(r^{-1-2q}).
\]
Convolution and interpolation preserve this first-derivative bound.
Moreover
\[
 \tr_{\widetilde g}\widetilde K
 =\tr_\delta\widetilde K
      +(\widetilde g^{ij}-\delta^{ij})\widetilde K_{ij}
 =O_1(r^{-1-2q}).
\]
Differentiating the last product uses only the available first derivatives
of the metric and tensor.  Thus the projection changes the tensor by
\(O_1(r^{-1-2q})\), without requiring a second derivative of the original
\(K\).

For trace-free data with these decay bounds, the constraint pair has
the flat linear principal expressions
\begin{equation}
 (\mu,J)
 =\left(\tfrac12(\partial_i\partial_jh_{ij}
                      -\Delta_\R h_{ii}),\,
                   \partial_jK_{ij}\right)
     +O(r^{-2-2q}).
\label{four:maximal:prep:linear-constraints}
\end{equation}
The scalar error consists of \(hD^2h\), \((Dh)^2\), and
\(|K|^2\); the momentum error consists of \(hDK\) and \((Dh)K\).
The displayed linear operators commute with convolution.  Cutoff
derivatives multiply the preceding convolution differences and
contribute \(O(r^{-q-3})\).  The trace projection changes the
momentum by \(O(r^{-2-2q})\).  Hence the new constraint pair is the
interpolated average of the old pair, up to
\begin{equation}
 O(r^{-2-2q}+r^{-q-3})=o(r^{-4-\delta}).
\label{four:maximal:prep:smoothing-errors}
\end{equation}
Replacing the old momentum norm by its Euclidean norm, and the averaged
Euclidean momentum norm by the new metric norm, each costs
\(O(r^{-2-2q})\): both metrics differ from the Euclidean metric by
\(O(r^{-q})\), while both momenta are \(O(r^{-2-q})\).
The triangle inequality and positivity
of the kernel therefore preserve a fixed positive fraction of the strict
dominance gap once \(R_{\rm sm}\) is sufficiently large.  This radius is
chosen separately for each fixed strict pair, so the little-o errors are
smaller than its particular positive gap coefficient.  The two needed
exponent inequalities are
\[
 2q-2-\delta>0,\qquad q-1-\delta>0.
\]
The same calculation shows that the scalar curvature is still
\(O(r^{-4-\delta})\).

Beyond \(2R_{\rm sm}\), derivatives of the convolution kernel
bound every higher derivative on unit patches, using the original
zeroth, first, and second derivative bounds.  The transition region
is smooth and compact.  At infinity the metric's first derivative
changes by \(O(r^{-q-2})\), whose integrated ADM contribution is
\(O(r^{1-q})\to0\).  Thus this smoothing leaves the energy exactly
unchanged.  As \(R_{\rm sm}\to\infty\), its uniform metric ratio
tends to one, so the enclosing infimum converges as well.  Choosing
first \(s\downarrow0\), and then a sufficiently large smoothing
radius for each \(s\), proves \eqref{four:maximal:prep:strict-properties} and
\eqref{four:maximal:prep:approx-limits}.
\end{proof}

For the remainder of the proof fix strict data from
Proposition~\ref{four:maximal:prop:strict-data}, and write
\(A_*=\mathcal A_g(S)\).  Constants may depend on this fixed approximation.
Write \(c_0>0\) for a fixed constant such that
\(\mu-|J|\ge c_0r^{-4-\delta}\) for these data.

\subsection{Threshold regions and outward collars}

We use the following part of the total trapped-region theorem:
on a complete smooth asymptotically flat initial data set of spatial
dimension \(2\) through \(7\), the nonempty union of bounded regions
with smooth boundary and outward expansion at most zero has a smooth
embedded outermost stable MOTS as its outer frontier.  We take its
closed, filled representative.  Filling a bounded complementary
component deletes boundary components and preserves admissibility.
This is the existence and regularity conclusion of
\cite[Section~4.3, Theorem~4.6]{AnderssonEichmairMetzger2011}; it requires no energy
condition.  We do not use the additional topology conclusions of that
theorem.

\begin{lemma}[Threshold selection and collars]
\label{four:maximal:lem:threshold-collars}
There is a number \(c\) with \(\max_S\Theta_K<c<0\), and a
closed filled compact region \(\mathcal D_c\), containing a smooth
compact filling of \(S\) and an exterior collar of \(S\), with the
following properties.
Its boundary \(B\) is a smooth embedded stable threshold surface,
\(\Theta_K(B)=c\).  It belongs to an increasing family
\(\mathcal D_b\), \(\max_S\Theta_K<b<0\), of total regions
at threshold \(b\), all lying in a fixed compact set, and
\begin{equation}
 \mathcal D_{b}\longrightarrow\mathcal D_c
       \quad\text{in Hausdorff distance as }b\downarrow c.
\label{four:maximal:prep:right-continuity}
\end{equation}
The exterior \(\Omega\) of \(\mathcal D_c\) is connected and
one-ended.  Every smooth cut enclosing \(B\) has \(g\)-area at
least \(A_*\).  Each component of \(B\) has a smooth outward
collar with leaf coordinate \(s\ge0\), \(s=0\) at \(B\),
\(|\dd s|>0\), whose leaves satisfy
\begin{equation}
 \Theta_K(\{s=\text{constant}\})\ge c,
 \qquad \nu_s=\frac{\grad s}{|\dd s|}.
\label{four:maximal:prep:outward-collars}
\end{equation}
\end{lemma}

\begin{proof}
The given diffeomorphism of the exterior permits a smooth compact
filling behind \(S\).  Extend \(g,K\) smoothly over it; the
extended metric can be chosen positive definite by interpolation
with any interior metric away from a collar.  This extension is used
only in the present geometric construction.

Fix a large radius \(R_0\) beyond which every coordinate sphere has
\(\Theta_K=3/r+O(r^{-1-q})>0\).  Outside a still larger fixed
radius, smoothly cut the metric to the Euclidean metric and the
unshifted tensor to zero.  Dilated cutoffs preserve the bounds
\(g-\delta=O_2(r^{-q})\), \(K=O_1(r^{-1-q})\), and consequently
the positive expansion of all spheres in the cutoff annuli.  Denote
these complete auxiliary data by \((g^{\rm aux},K^{\rm aux})\).
Choose a fixed nonnegative smooth cutoff \(\xi\), equal to one
on the filling and throughout a large neighborhood of \(\{r\le
R_0\}\), and zero far out.  For
\(\max_S\Theta_K<b<0\), use the tensor
\[
 K_b^{\rm aux}=K^{\rm aux}-\tfrac b3\xi g^{\rm aux}.
\]
Where \(\xi=1\) and the data are uncut, the exact identity is
\begin{equation}
 \Theta_{K_b^{\rm aux}}(\Sigma)=\Theta_K(\Sigma)-b.
\label{four:maximal:prep:threshold-shift}
\end{equation}
In the cutoff region the added tangential trace is \(-b\xi\ge0\),
so all the distant spheres still have positive expansion, uniformly
in \(b\).

A smooth weakly trapped bounded region cannot meet those distant
spheres: at a point where its boundary reaches its greatest radius,
its tangent plane and outward normal agree with those of the
coordinate sphere, and its mean curvature is at least the sphere's
mean curvature.  The tensor traces there agree.  This contradicts
nonpositive expansion.  The same argument applies to the outer
frontier furnished by the total-region theorem.  Thus all total
regions are contained in the uncut part of the construction.

For every indicated \(b\), strictness at \(S\) gives a short
outward collar whose leaves have expansion less than \(b\).
The filling together with this collar is a nonempty trapped region
for \(K_b^{\rm aux}\).  The theorem therefore gives a total closed
region \(\mathcal D_b\) with smooth stable frontier.  Its frontier
lies strictly outside \(S\), in the original data, and has
\(\Theta_K=b\) by \eqref{four:maximal:prep:threshold-shift}.
We used one filling and one set of cutoffs for the entire interval
of \(b\).  Since \(\xi\ge0\), increasing \(b\) decreases
the shifted expansion on every candidate hypersurface.  It follows
that
\begin{equation}
 b_1<b_2\quad\Longrightarrow\quad
                 \mathcal D_{b_1}\subset\mathcal D_{b_2}.
\label{four:maximal:prep:threshold-monotonicity}
\end{equation}

Choose a compact enlargement \(Q\) containing all these regions.
Lemma~\ref{n3:domains:right-continuity} applies to the increasing family
of closed subsets \(E_b=\mathcal D_b\) of this fixed compact metric
space, with \(\max_S\Theta_K<b<0\).  The preceding construction proves
both the common compact containment and monotonicity; every member is
nonempty.  Choose \(c\) outside the lemma's countable exceptional set.
Its conclusion is exactly \eqref{four:maximal:prep:right-continuity}.
Filling bounded complementary
pockets ensures that the remaining exterior is the connected
component containing the sole end.  Any cut enclosing its entire
boundary \(B\) also encloses \(S\), proving the area assertion.

It remains to construct the collars.  Apply
Lemma~\ref{n3:domains:stable-collar} with ambient dimension four,
\(Q=K\), threshold \(c\), and each connected component of \(B\).
The total-region theorem gives a compact smooth stable frontier;
its principal eigenvalue is nonnegative
\cite[Proposition~5.1]{AnderssonMarsSimon2008}.  On an open neighborhood
of this frontier the auxiliary cutoff \(\xi\) equals one and the data
are uncut, so the relevant shifted tensor is exactly \(K-(c/3)g\).
Its expansion differs from \(\Theta_K\) by the constant \(c\), with
zero derivative under graph variations.  Replacing one component by
a sufficiently short outward graph of expansion at most \(c\)
would strictly enlarge the total region while preserving every other
component.  Thus the local maximality hypothesis of the collar lemma
holds in this application.

For clarity, its zero-mode construction uses the expansion
\(\mathcal E(v)\) of the outward graph on this one component and
\(L=D\mathcal E(0)\).  The common augmented inverse specializes to
\begin{equation}
 (v,a)\longmapsto\left(Lv-a,\ \int_Bv\dd A_g\right),
 \qquad C^{2,\alpha}(B)\times\R
             \longrightarrow C^{0,\alpha}(B)\times\R.
\label{four:maximal:prep:augmented-inverse}
\end{equation}
Here the integrals are over the chosen component.  The positive
principal and adjoint eigenfunctions, their simple kernels, and the
positive pairing used to invert this map are those in
\cite[Section~4, Lemma~4.1]{AnderssonMarsSimon2008}, as in the common
proof.  That proof gives strictly larger expansion on every positive
outward leaf, including the zero-eigenvalue case; stability alone
would not give this conclusion.  Finitely many components permit
disjoint shortened collars, proving
\eqref{four:maximal:prep:outward-collars}.
\end{proof}

\subsection{Exterior supports of nonsmooth sets}

The threshold regions control smooth surfaces.  The graph-height limit
will give only a closed sublevel set, so we need the same comparison for
smooth surfaces that touch such a set from outside.  The next argument
transports the supporting tangent plane exactly; this prevents an
uncontrolled curvature of the support from entering the error.

For a closed set \(A\), a smooth exterior support at
\(x\in\partial A\) is a regular level \(\{\psi=0\}\) in a
neighborhood of \(x\), where \(\psi(x)=0\), \(\dd\psi(x)\ne0\),
and \(A\subset\{\psi\le0\}\) locally.  Its normal is
\(\nu_\psi=\grad\psi/|\dd\psi|\), and its expansion is
\begin{equation}
 \Theta_K[\psi]
 =\frac{\tr_g\Hess\psi-\Hess\psi(\nu_\psi,\nu_\psi)}
              {|\dd\psi|}
      +\tr_gK-K(\nu_\psi,\nu_\psi).
\label{four:maximal:prep:support-expansion}
\end{equation}
Only supports that exist are tested; no regularity is assumed of
\(A\) itself.

\begin{lemma}[Enclosure from exterior supports]
\label{four:maximal:lem:support-enclosure}
Let \(A\) be a compact closed set containing \(\mathcal D_c\).
Suppose that every smooth exterior support to its frontier has
expansion at most \(k'\), where \(c\le k'<0\).  Then
\[
 A\subset\mathcal D_{c'}\qquad\text{for every }k'<c'<0.
\]
All expansions at the frontier of \(A\) are computed in the
original strict exterior data.
\end{lemma}

\begin{proof}
First, every set under consideration lies inside one fixed large
sphere.  Otherwise the coordinate sphere at its maximum radius is
an exterior support with positive expansion, contrary to \(k'<0\).
Choose a compact enlargement containing this sphere and a small
metric neighborhood of it.  All constants below depend only on
the smooth geometry and tensor on this enlargement.  Distances
may be computed in the complete smooth filling, since the contact
points and sufficiently short segments used below lie in the
original exterior: \(A\) contains the filling and a collar of
\(S\).

For \(z>0\), put \(A_z=\{y:\dist(y,A)\le z\}\).
Apply Lemma~\ref{n3:domains:parallel-support} in ambient dimension
four, with its set \(D=A\), tensor \(Q=K\), and threshold
\(k'\).  The compact neighborhood and short segments required by
that local lemma are precisely those just fixed.  The definition
\eqref{four:maximal:prep:support-expansion} is its normalized
expansion.  It follows that every exterior support to \(A_z\)
has expansion at most
\begin{equation}
 k'+Cz,
\label{four:maximal:prep:parallel-support-bound}
\end{equation}
where \(C\) is independent of that support's curvature.

The exact geometric feature of this application is worth recording.
For a nearest pair \(x\in A\), \(x'\in\partial A_z\), let
\(P:T_xM\to T_{x'}M\) be parallel transport along the short
minimizing segment.  The Jacobi construction in the common proof gives
\begin{equation}
 F_z(x)=x',\qquad \dd F_z|_x=P,
 \qquad |\nabla\dd F_z|_x\le Cz.
\label{four:maximal:prep:exact-transport}
\end{equation}
If \(\psi\) defines the support at \(x'\), then
\(\psi\circ F_z\) supports \(A\) at \(x\), and
\begin{equation}
\begin{split}
 \Hess(\psi\circ F_z)_x(X,Y)
  ={}&\Hess\psi_{x'}(PX,PY)\\
    &+\dd\psi_{x'}\bigl((\nabla\dd F_z)_x(X,Y)\bigr).
\end{split}
\label{four:maximal:prep:hessian-pullback}
\end{equation}
The exact tangent isometry preserves the normalized Hessian trace;
only the \(O(z)\) second-map derivative and the \(C^1\) change
of \(K\) enter the expansion error.  Thus no term multiplying
\(\Hess\psi\) has been discarded.  This local transport step
uses compact ambient bounds, while the following enclosure argument
still uses the auxiliary total regions specific to the maximal end.

Fix \(k'<c'<0\), and take \(\beta>0\) below the preceding
injectivity and collar scales, so small that
\(k'+C\beta<c'\).  Smooth approximation of the distance function,
followed by a regular-value choice, gives a smooth compact region
\(U\) with
\[
 A_{\beta/4}\subset\operatorname{int}U,
 \qquad U\subset A_{3\beta/4}.
\]
For example, approximate the distance uniformly within
\(\beta/32\), and choose a regular value between
\(3\beta/8\) and \(5\beta/8\).  Choose a smooth nonnegative
function \(\zeta\) equal to one near \(\partial U\) with
\(\supp\zeta\subset\{\dist(\cdot,A)<\beta\}\).
For sufficiently large finite \(\Lambda\), replacing the
auxiliary threshold tensor by
\[
 K_{c'}^{\rm aux}-\Lambda\zeta g^{\rm aux}
\]
makes \(\partial U\) strictly trapped: the added term decreases
its expansion by \(3\Lambda\).  Apply the total-region theorem
to these smoothly modified complete data.  Let \(\widetilde{\mathcal
D}\) be the resulting closed filled total region and
\(\widetilde B\) its smooth frontier.  It contains \(U\).
All cutoffs were placed outside the fixed enlargement above;
outside the support of \(\zeta\), large spheres retain positive
expansion.  The maximum-radius argument confines
\(\widetilde B\) to the uncut exterior.  There it satisfies
\begin{equation}
 \Theta_K(\widetilde B)=c'+3\Lambda\zeta\ge c'.
\label{four:maximal:prep:modified-frontier}
\end{equation}

Let \(z=\dist(A,\widetilde B)\), the global minimum distance.
Because \(\widetilde{\mathcal D}\) contains
\(A_{\beta/4}\), one has \(z\ge\beta/4>0\).
If \(z<\beta\), every path of length less than \(z\) starting
in \(A\) remains on the inner side of \(\widetilde B\): its
first crossing would contradict the definition of \(z\).
Taking closures proves \(A_z\subset\widetilde{\mathcal D}\).
At a pair attaining the distance, \(\widetilde B\) is therefore
an exterior support to \(A_z\).  Equations
\eqref{four:maximal:prep:parallel-support-bound} and
\eqref{four:maximal:prep:modified-frontier} imply
\[
 c'\le\Theta_K(\widetilde B)\le k'+Cz
                  \le k'+C\beta<c',
\]
a contradiction.  Thus \(z\ge\beta\).  Since this is the
global minimum, \(\zeta\) vanishes at every point of
\(\widetilde B\), including when the minimum equals \(\beta\).
It follows that \(\Theta_K(\widetilde B)=c'\) everywhere.
Its filled inner region is an admissible threshold-\(c'\) region
for the original auxiliary data, so
\(A\subset\widetilde{\mathcal D}\subset\mathcal D_{c'}\),
as required.
\end{proof}

\subsection{Fixed collar data}

Set \(b_*=c\).  On the disjoint collars supplied by
Lemma~\ref{four:maximal:lem:threshold-collars}, choose a nonpositive smooth
compactly supported function \(C\), equal to \(-\kappa s\)
near each component of \(B\), with \(\kappa>0\) sufficiently
small.  Indeed, multiply \(-\kappa s\) by nonnegative collar
cutoffs.  The strict gap has a positive minimum on their compact
support, so \(\kappa\) can be chosen small enough that, for a
fixed smooth positive function \(\rho\) equal to a positive
multiple of \(r^{-4-\delta}\) on the end,
\begin{equation}
 |\dd C|+\rho\le\mu-|J|.
\label{four:maximal:prep:offset-slack}
\end{equation}
For instance first choose \(\rho\le(c_0/4)r^{-4-\delta}\),
then bound \(|\dd C|\) by half the remaining strict gap on its
support.  These data, the collars, the threshold \(b_*\), and
the region \(\Omega\) are fixed before choosing any of the
deformation parameters below.

\section{The scalar-curvature deformation}\label{four:maximal:sec:deformation}

We now fix the strict data and the exterior $\Omega$ constructed in
Section~\ref{four:maximal:sec:preparation}.  Thus $B=\partial\Omega$ has expansion
$H_B+\tr_BK=b_*<0$, the tensor is maximal, and the fixed functions $C\le0$
and $\rho>0$ satisfy
\begin{equation}\label{four:maximal:def:slack}
 |dC|+\rho\le\mu-|J|.
\end{equation}
The normal $\nu$ at $B$ points into $\Omega$.  The compactly supported
function $C$ vanishes on $B$ and equals $-\kappa s$ in its fixed collars.
Throughout this section the base dimension is four and a transverse block
has dimension three.

\subsection{Variables and the exact scalar identity}

The pointwise identities used here are those of
Section~\ref{four:sec:system}.  We retain their variables and
four-dimensional density conventions below.  Their proofs are local
and impose no support condition on $K$; maximality will enter through
$\tr_gK=0$ in the equations and the estimates for the tensor tail.

Fix $0<\eps<1$.  The integer $N\ge4$ is an auxiliary parameter.  Fix once
and for all a smooth nonincreasing function $\vartheta:\R\to[0,1]$ equal
to one on $(-\infty,0]$ and zero on $[1,\infty)$, and set
\begin{equation}\label{four:maximal:def:warp}
 \begin{gathered}
 p(t)=N\vartheta(Nt),\qquad
 l(t)=\exp\left(\int_0^t p(s)\dd s\right),\qquad
 L_0(t)=e^{2t}l(t),\\
 \ell=e^{-\eps N},\qquad \eta_N=\ell^{3/2}.
 \end{gathered}
\end{equation}
In particular $0\le p\le N$, $|p'|\le\|\vartheta'\|_\infty N^2$,
$l(t)=e^{Nt}$ for $t\le0$,
and $l\le e$ everywhere.  A symbol $\Pi_N$ denotes a positive polynomial
bound in $N$; its coefficients may depend on the fixed strict data,
collars, cutoffs, $\eps$, and subsequently fixed small constants, but not
on a solution, the outer radius, or a homotopy parameter.  Its value may
increase from one estimate to the next.  Arbitrary constants at fixed
$N$ will instead be denoted by $C(N)$.

On $\Omega_R=\Omega\cap\{r<R\}$ the unknowns are $f,Z$.  Define $t$
implicitly and the remaining variables explicitly by
\begin{align}
 \sigma&=|\grad f|,& D&=1+l(t)^2\sigma^2,&d&=D^{-1},
 & Z&=t+\tfrac16\log D,\label{four:maximal:def:implicit}\\
 w&=\frac{l\grad f}{\sqrt D},&a&=|w|,&v&=a^2=1-d,
 &\chi&=\frac{3d}{3+pv},\label{four:maximal:def:axial-variables}\\
 h&=\eta_N f,&U&=l\sqrt D,&u&=L_0\sqrt D.
 \label{four:maximal:def:weights}
\end{align}
At fixed $\sigma$, the derivative of the right side of
\eqref{four:maximal:def:implicit} with respect to $t$ is $1+pv/3>0$.  As $t\to-\infty$
it tends to $-\infty$, and as $t\to\infty$ it tends to $\infty$.
Consequently $t=t(Z,df)$ exists uniquely and smoothly for all finite
arguments.  Notice that $t\ge-\eps$ implies $\ell\le l\le e$.

For $\sigma>0$ put $e_f=\grad f/\sigma$ and
$A_j=\Id+(j-1)e_f\otimes e_f$.  The two matrices used below have the
globally smooth formulas
\begin{equation}\label{four:maximal:def:matrices}
 \Ad=\Id-w\otimes w,\qquad
 \Achi=\Id-\frac{3+p}{3+pv}w\otimes w,
 \qquad \frac{3d}{3+N}\le\chi\le d\le1.
\end{equation}
They are positive definite; the displayed formulas remove the apparent
ambiguity of the axis when $df=0$.  We identify vectors and covectors
using $g$, write $|\xi|_A^2=A(\xi,\xi)$, and put
\begin{equation}\label{four:maximal:def:graph-fields}
 \begin{gathered}
 \gbar=g+l^2df\otimes df,\quad
 \ghat=e^{2t}\gbar,\quad
 H^f=\frac{l}{\sqrt D}\Hess f,\quad S_f=K+H^f,\\
 V=u\Achi\bigl(3\grad Z+K(w,\cdot)\bigr).
 \end{gathered}
\end{equation}
The trace equation will always be
\begin{equation}\label{four:maximal:def:trace}
 \tr_{\Achi}S_f=F.
\end{equation}
Here $F$ is the prescribed right side, including its dependence on $f$,
$Z$, and $df$ when a system is specified.

For symmetric tensors recall
$P_A=A-(\tr A)g$ and
$\mathsf{B}(A_1,A_2)=A_1:A_2-(\tr A_1)(\tr A_2)$.
The constraint quantities are
$\mu=(R_g-\mathsf{B}(K,K))/2$ and $J=\divg P_K$.
At a point with $\sigma>0$, decompose $S_f$ relative to
$e_f^\perp\oplus\R e_f$ into its transverse symmetric block $T$, mixed
covector $M$, and axial entry $j$.  The letter $M$ in the formulas below
denotes only this block.  Put
\begin{equation}\label{four:maximal:def:blocks}
 x=\partial_{e_f}t,\qquad y=(\grad t)_\perp,
 \qquad T^0=T-\tfrac13(F-\chi j)\Id_\perp.
\end{equation}
The trace equation gives $\tr T=F-\chi j$.

\begin{proposition}[Scalar identity]\label{four:maximal:prop:scalar-identity}
For smooth $f,t$ satisfying \eqref{four:maximal:def:trace},
\begin{equation}\label{four:maximal:def:scalar-identity}
 \frac12 e^{2t}R_{\ghat}
 =\mu+J(w)+\cT+w(F)-F\tr K-u^{-1}\divg V,
\end{equation}
where
\begin{align}\label{four:maximal:def:quadratic-form}
 \cT={}&\tfrac12|T^0|^2+|M+(p+1)ay|^2-v|y|^2
       +3(p+1)(|y|^2+dx^2)\notag\\
 &-\tfrac13(F-\chi j)^2+\chi j^2-\chi Fj+F^2
       +2(p+1)\chi jax.
\end{align}
This expression extends smoothly across $df=0$.
\end{proposition}

\begin{proof}
This is Proposition~\ref{four:sys:scalar-proposition}, with its axis
$e=e_f$ and the same graph variables and density normalization.
Its arbitrary trace $\tau$ is zero here.

We record the differential relations needed in the boundary and drift
calculations.  Differentiating the implicit definition of $Z$ gives
\begin{equation}\label{four:maximal:def:Z-differential}
 3\,dZ=(3+pv)dt+H^f(w,\cdot).
\end{equation}
If $k=-H^f-p(dt\otimes w^\flat+w^\flat\otimes dt)$ and $Q=K-k$, then
\begin{align}
 V/u&=P_Qw+3\Ad\grad t+Fw,
                    \label{four:maximal:def:flux-rewrite}\\
 u^{-1}\divg(uw)
 &=F+(1-\chi)j-\tr_gK+2(p+1)w(t).
                    \label{four:maximal:def:w-divergence}
\end{align}
These are the identities in the proof of the cited proposition.
They follow by inserting $\tr T=F-\chi j$ and
$H^f(w,\cdot)=v\,d\log\sigma$ into its invariant formula; their
smooth tensor expressions also hold at $df=0$.  The quadratic expression
there is precisely the displayed $\cT$.
\end{proof}

\begin{lemma}[Polynomial coercivity]\label{four:maximal:lem:coercivity}
The quadratic form in \eqref{four:maximal:def:quadratic-form} is nonnegative, and
\begin{equation}\label{four:maximal:def:coercivity}
 |T|^2+|M|^2+dj^2+F^2+|dt|_{\Ad}^2\le\Pi_N\cT.
\end{equation}
At $dt=0$ it has the stronger equivalence
\begin{equation}\label{four:maximal:def:coercivity-stationary}
 \cT\asymp |T^0|^2+|M|^2+\chi j^2+F^2,
\end{equation}
with both constants absolute, independent of $N,d,\chi$.
\end{lemma}

\begin{proof}
Lemma~\ref{four:sys:coercivity-lemma} applies with the same dictionary.
For clarity, its exact square completion is
\begin{align}\label{four:maximal:def:square-completion}
 \cT={}&\tfrac12|T^0|^2+|M+(p+1)ay|^2
       +\bigl(3(p+1)-v\bigr)|y|^2\notag\\
 &+3(p+1)d\left(x+\frac{\chi aj}{3d}\right)^2
       +\tfrac23\left(F-\frac{\chi j}{4}\right)^2
       +\frac{d(48-3d)}{8(3+pv)^2}j^2.
\end{align}
The last coefficient is at least $45d/[8(3+N)^2]$ and the coefficient
of $|y|^2$ is at least two.  Undoing the shifted squares therefore
costs only a polynomial in $N$.  When $dt=0$, the normal terms instead
reduce to
\[
 \tfrac23(F-\chi j/4)^2+\chi(1-3\chi/8)j^2.
\]
The bounds $5/8\le1-3\chi/8\le1$ and $\chi^2\le\chi$ give the
absolute equivalence claimed in \eqref{four:maximal:def:coercivity-stationary}.
\end{proof}

\begin{lemma}[Boundary identity]\label{four:maximal:lem:boundary-identity}
On a hypersurface where $f$ is constant, let $H$ denote its $g$-mean
curvature for a chosen normal $\nu$, and put $P_B=\tr_BK$.  Then
\begin{align}
 V_\nu/u&=d(H+3\partial_\nu t)-H+w_\nu(F-P_B),
 \label{four:maximal:def:boundary-flux}\\
 \widehat H&=e^{-t}\sqrt d\,(H+3\partial_\nu t).
 \label{four:maximal:def:boundary-curvature}
\end{align}
The identities allow either sign of $\partial_\nu f$, including zero.
\end{lemma}

\begin{proof}
Apply Lemma~\ref{four:sys:boundary-lemma} with the chosen normal $\nu$
and trace relation \eqref{four:maximal:def:trace}.  Its formulas use
the signed component $w_\nu$ and extend smoothly to $df=0$, so they
also apply when the normal graph slope changes sign.
\end{proof}

\subsection{The equations and their homotopy}\label{four:maximal:def:homotopy}

Choose
\begin{equation}\label{four:maximal:def:end-weights}
 1<\gamma<\min(2,q-\delta/2),\qquad
 \rho_0=r^{-4-\delta},\qquad \rho_1=r^{-2\gamma}.
\end{equation}
Here $r$ is the fixed positive smooth extension of the end radius.
Fix smooth cutoffs $m_0^*,j_0:\R\to[0,1]$ with $m_0^*=1$ on
$(-\infty,0]$, $m_0^*=0$ on $[1,\infty)$, $j_0=0$ on
$(-\infty,0]$, and $j_0=1$ on $[1,\infty)$.  Define
\begin{equation}\label{four:maximal:def:penalty}
 \mathcal P=C_N(\rho_0+v\rho_1),\qquad
 m_0(t)=m_0^*(N(t+\eps)),\qquad
 \Xi=\delta_0\cT+\rho+\delta_0\eta_N a\sigma-m_0(t)\mathcal P.
\end{equation}
Proposition~\ref{four:maximal:prop:floor} below fixes a single $\delta_0>0$ independent
of $N$ and then a sufficiently large polynomial $C_N$.
Choose a smooth $N_B>1+\sup_B|H|$ on $B$ and write
\begin{equation}\label{four:maximal:def:boundary-datum}
 \mathcal B=-H+w_\nu(F-P_B)-N_Bd,\qquad P_B=\tr_BK.
\end{equation}
The system whose solution we seek is
\begin{equation}\label{four:maximal:def:system}
 \begin{cases}
  \tr_{\Achi}(K+H^f)=F=h+C,&\text{in }\Omega_R,\\
  \divg V=u\Xi,&\text{in }\Omega_R,\\
  h=b_*,\quad V_\nu/u=\mathcal B,&\text{on }B,\\
  f=Z=0,&\text{on }S_R.
 \end{cases}
\end{equation}
In the boundary datum, $F$ always means the prescribed trace right side,
so the datum contains no unprescribed second derivatives.

The two geometric signs follow directly for any smooth solution of
\eqref{four:maximal:def:system}.  Maximality and Proposition~\ref{four:maximal:prop:scalar-identity}
give the exact decomposition
\begin{equation}\label{four:maximal:def:scalar-sign}
 \frac12e^{2t}R_{\ghat}
 =\bigl(\mu+J(w)+w(C)-\rho\bigr)
  +(1-\delta_0)\cT+(1-\delta_0)\eta_N a\sigma+m_0\mathcal P.
\end{equation}
The bracket is nonnegative by \eqref{four:maximal:def:slack} and $|w|\le1$, and all
other terms are nonnegative when $\delta_0<1$.  On $B$, subtracting
\eqref{four:maximal:def:boundary-datum} from \eqref{four:maximal:def:boundary-flux} gives
$H+3\partial_\nu t=-N_B$, and therefore
\begin{equation}\label{four:maximal:def:geometric-signs}
 R_{\ghat}\ge0,\qquad \widehat H=-e^{-t}\sqrt d\,N_B<0.
\end{equation}
These conclusions use the actual boundary expansion $H+P_B=b_*$.
They require no assumption that $H$ vanishes.

Here is a homotopy to a scalar system with a trivial endpoint.
Let
\[
 V_\lambda=u\Achi\bigl(3\grad Z+\lambda K(w,\cdot)\bigr),
 \qquad 0\le\lambda\le1.
\]
In the first three stages impose $\divg V_\lambda=u\Xi$ and
\begin{equation}\label{four:maximal:def:homotopy-boundary}
 \frac{(V_\lambda)_\nu}{u}
 =\bigl[1-(1-\lambda)j_0(t)\bigr]
   \bigl[\mathcal B-(1-\lambda)
                (\Achi K(w,\cdot))_\nu\bigr].
\end{equation}
All outer values remain $f=Z=0$.  Let $\nu_s=\grad s/|ds|$ in the fixed
collars.  Choose smooth bounded functions $D_0(x,w)$ and $D_1(x,w)$,
supported in those collars and with uniformly bounded derivatives for
$|w|\le1$, as follows:
\begin{itemize}
 \item $D_0\le0$, $D_0=0$ when $w\cdot \nu_s\ge0$, and
 $D_0(x,-\nu)<-2\sup_B|H|$ on $B$;
 \item $D_1=A_1\zeta_B(x)w\cdot \nu_s$, where $\zeta_B=1$ near $B$ is a
 fixed collar cutoff and $A_1>1+\sup_B|H|$.
\end{itemize}
Such $D_0$ is obtained by a fixed smooth function of $w\cdot \nu_s$, zero
on the nonnegative half-line and sufficiently negative at $-1$, times
a collar cutoff.  In particular $D_0(x,0)=D_1(x,0)=0$.
The stages, with their endpoints included, are:
\begin{enumerate}[label=\arabic*.]
 \item Decrease $\lambda$ from one to zero, keeping $F=h+C$ and
 $h|_B=b_*$.
 \item Keep $\lambda=0$ and increase $\alpha$ from zero to one, with
 $F=h+C+\alpha D_0$ and $h|_B=b_*$.
 \item Keep $\lambda=0$ and increase $\zeta$ from zero to one, with
 \begin{equation}\label{four:maximal:def:stage-three}
  F=h+(1-\zeta)(C+D_0)
       +\zeta\bigl(\tr_{\Achi}K+D_1\bigr),
  \qquad h|_B=(1-\zeta)b_*.
 \end{equation}
 \item Keep the terminal trace problem
 \begin{equation}\label{four:maximal:def:terminal-trace}
  \tr_{\Achi}H^f=h+D_1,\qquad f|_{B\cup S_R}=0,
 \end{equation}
 and multiply both $\Xi$ and the right side of
 \eqref{four:maximal:def:homotopy-boundary} by a common factor $\eta$, decreasing
 from one to zero.
\end{enumerate}
The prescriptions agree at every junction.  In every stage the
derivative of $F$ with respect to $h$ is one.  The form $\cT$ is always
computed from the actual trace equation and the actual tensor $K+H^f$,
including when $\lambda=0$.

Two features of this choice will be used in the estimates.
First, in stage three the limiting values at $w=\nu,-\nu$, with
$\chi=0$ and the assigned face height, obey
\begin{equation}\label{four:maximal:def:face-direction-signs}
 F(x,\nu)\ge P_B+H,\qquad F(x,-\nu)\le P_B-H.
\end{equation}
At $\zeta=0$ the first equality follows from $b_*=P_B+H$ and
$D_0(x,\nu)=C|_B=0$; the second follows from the choice of $D_0$.
At $\zeta=1$ the two values are $P_B\pm A_1$.  Convex combination proves
\eqref{four:maximal:def:face-direction-signs} for intermediate $\zeta$.
Second, the terminal trace equation forces $f=0$ for every
$C^{1,\alpha}$ trial $Z$ and its classical trace solution: at a positive interior maximum the left side is nonpositive and
the right side is $\eta_N f>0$, while at a negative interior minimum the
signs reverse.  Boundary values are zero.  Thus in stage four
\begin{equation}\label{four:maximal:def:terminal-reduction}
 f=0,\quad t=Z,\quad d=\chi=1,\quad u=L_0(t),\quad
 \cT=\tfrac12|K|^2+3(p+1)|dt|^2.
\end{equation}

\begin{lemma}[Height and initial end bounds]\label{four:maximal:lem:height-end}
Every smooth homotopy solution satisfies $|h|+|F|\le C_0$ for a fixed
constant independent of $N,R$ and the homotopy parameter.  There are
fixed $r_0,C_2$ such that, for all sufficiently large $R$ depending on
$N,\eps$, the first three stages satisfy
\begin{equation}\label{four:maximal:def:height-end}
 |h|\le C_2(r^{-\gamma}-R^{-\gamma})\quad(r_0\le r\le R),
 \qquad |\grad f|\big|_{S_R}\le C_3\eta_N^{-1}R^{-1-\gamma}.
\end{equation}
Moreover $t>-\eps$ on $S_R$.  None of these assertions assumes an
interior lower bound for $t$.
\end{lemma}

\begin{proof}
At an interior extremum of $h$, we have $w=0$, $\Achi=\Id$,
$D_0=D_1=0$, and $\tr K=0$.  The trace equation and its positive
Hessian coefficient show that a minimum has $h\ge0$ and that a maximum
has $h\le\|C\|_\infty$.  The boundary heights lie in $[b_*,0]$.
Consequently
\begin{equation}\label{four:maximal:def:height-range}
 b_*\le h\le\|C\|_\infty
\end{equation}
in stages one through three; stage four has $f=0$.
All remaining terms in the trace prescriptions are uniformly bounded
for $|w|\le1$, $0\le\chi\le1$, so $F$ has a fixed bound as well.

Outside the compact supports the trace equation takes the form
\[
 \beta\tr_{\Achi}\Hess h+s_K\tr_{\Achi}K=h,
 \qquad \beta=\frac{l}{\eta_N\sqrt D},\qquad 0\le s_K\le1.
\]
Here $s_K=1$ in the first two stages and $s_K=1-\zeta$ in the third.
At a nonzero test gradient $\beta=a/|\grad h|$.
Test with $\psi=C_2(r^{-\gamma}-R^{-\gamma})$.  Since the test axis is
radial, the end metric estimates give, uniformly for $0<\chi\le1$,
\[
 \frac{\tr_{\Achi}\Hess\psi}{|\grad\psi|}
 =\frac{-3+(\gamma+1)\chi}{r}+O(r^{-1-q})
 \le-\frac{c_1}{r},\qquad c_1=(2-\gamma)/2>0,
\]
after fixing $r_0$ sufficiently large.  In addition, maximality gives
\begin{equation}\label{four:maximal:def:K-trace-estimate}
 |\tr_{\Achi}K|
 \le K_0r^{-1-q}\min(1,(1+N/3)a^2).
\end{equation}
Indeed $\tr_{\Achi}K=(\chi-1)K(e_f,e_f)$ and
$1-\chi=(3+p)v/(3+pv)\le\min(1,(1+N/3)v)$.
These estimates hold for every value of $t$.

Put $c_\gamma=1-2^{-\gamma}$ and
$H_0=\max(|b_*|,\|C\|_\infty)$.  Choose the fixed coefficient $C_2$ so that
\[
 C_2\ge\max\left\{1,
 \frac{2H_0r_0^\gamma}{c_\gamma},
 \frac{2K_0r_0^{\gamma-1-q}}{c_\gamma}\right\}.
\]
For $R\ge2r_0$, this dominates the heights at $r=r_0$.
On $r\le R/2$, the tensor term in \eqref{four:maximal:def:K-trace-estimate} is at
most $\psi/2$, since $\gamma<1+q$.  On $r\ge R/2$, use
$\min(1,(1+N/3)a^2)\le\sqrt{1+N/3}\,a$ and require
\[
 R\ge2\left(\frac{2K_0\sqrt{1+N/3}}{c_1}\right)^{1/q}.
\]
The tensor error is then at most $c_1a/(2r)$.

If $h-\psi$ had a positive maximum, it would occur in the interior of
this annulus.  The gradients there agree and the Hessian of $h$ is no
larger than the Hessian of $\psi$.  The coefficients use the actual $Z$
and this common gradient, so they agree at the contact regardless of
the two heights.  On the first half of the annulus the equation gives
$h\le-c_1a/r+\psi/2<\psi$; on the second half it gives
$h\le-c_1a/(2r)<0$.  Both are contradictions.
At a negative minimum of $h+\psi$ the Hessian inequality reverses,
while its axis and $\chi$ are unchanged on reversing the gradient;
the same estimates give the lower comparison.  This proves the height
bound in \eqref{four:maximal:def:height-end}.

Because $h=0$ at $S_R$, the two comparisons bound its normal derivative
there by $C_3R^{-1-\gamma}$, with $C_3=2C_2\gamma$ after increasing $r_0$
so that $|dr|\le2$.  Tangential derivatives vanish, giving the asserted
bound for $|\grad f|$.  Finally, at $S_R$ the strictly increasing function
$G(t)=t+\frac16\log(1+l(t)^2\sigma^2)$ has $G(t)=Z=0$.
Since $l(-\eps)=\ell$, its root exceeds $-\eps$ whenever
$\ell^2\sigma^2<e^{6\eps}-1$.  It suffices to require
\begin{equation}\label{four:maximal:def:outer-radius-floor}
 R^{2\gamma+2}>
 \frac{C_3^2e^{\eps N}}{e^{6\eps}-1}.
\end{equation}
In stage four the conclusion follows instead from $f=0$ and $Z=0$ on
$S_R$.  Every restriction on $R$ made here depends only on the fixed
data, $N,\eps$, and is compatible with $R\to\infty$ at fixed $N$.
\end{proof}

\subsection{Exclusion of the lower boundary of the admissible range}

The next argument concerns any smooth homotopy solution whose minimum
value of $t$ is $-\eps$.  It supplies the strict exclusion needed for
continuation before any global gradient or Hessian estimate is available.

\begin{lemma}[Gauss comparison at a minimum]\label{four:maximal:def:gauss-minimum}
At an interior minimum of $t$,
\begin{equation}\label{four:maximal:def:graph-gauss}
 \tfrac12e^{2t}R_{\ghat}
 \le\tfrac12R_g-v\Ric_g(e_f,e_f)+\mathcal S(H^f),
\end{equation}
where, for a symmetric tensor $A$,
\begin{equation}\label{four:maximal:def:gauss-form}
 \mathcal S(A)=\tfrac13(\tr_\perp A)^2
 -\tfrac12|A_\perp^0|^2
 +dA_{ee}\tr_\perp A-d|A_{e\perp}|^2.
\end{equation}
At every point where $dt=0$, there is an absolute $c_2>0$ such that
\begin{equation}\label{four:maximal:def:floor-comparison}
 \cT-\mathcal S(H^f)
 \ge c_2\cT-\Pi_N\bigl(vF^2+|K|^2\bigr).
\end{equation}
At $df=0$ the terms multiplied by $v$ vanish and any axis may be used.
\end{lemma}

\begin{proof}
Lemma~\ref{four:sys:gauss-lemma}, with $e=e_f$, uses exactly
\eqref{four:maximal:def:gauss-form}.  At a minimum $\Hess t\ge0$,
so its stationary Gauss identity gives
\eqref{four:maximal:def:graph-gauss}.

The two algebraic identities used by its floor estimate specialize to
\begin{align}\label{four:maximal:def:gauss-difference}
 \cT-\mathcal S(S_f)
 ={}&|T^0|^2+(1+d)|M|^2\notag\\
 &+\chi(1+d-2\chi/3)j^2+(\chi/3-d)Fj+F^2/3
\end{align}
and
\begin{equation}\label{four:maximal:def:d-chi-ratio}
 \frac{d^2}{\chi}=\frac{d(3+pv)}3\le\frac{N+3}3.
\end{equation}
They are \eqref{four:sys:stationary-difference} and the first ratio
in \eqref{four:sys:polynomial-ratios}.  The common proof splits at
$(1+p)v$ small, uses stationary coercivity, and replaces $S_f$ by
$H^f=S_f-K$.  It yields an absolute positive $c_2$ and a polynomial
$\Pi_N$ in precisely \eqref{four:maximal:def:floor-comparison},
including the smooth extension to $df=0$.
\end{proof}

\begin{lemma}[Coefficient derivatives]\label{four:maximal:def:coefficient-lemma}
The following estimates hold along every smooth trace solution:
\begin{equation}\label{four:maximal:def:derivative-bounds}
 |\nabla w|+|\nabla\Achi|+|d\log u|_{\Achi}
 \le\Pi_N(\sqrt{\cT}+|K|).
\end{equation}
Consequently, for each fixed $\eta_1>0$,
\begin{equation}\label{four:maximal:def:drift-divergence}
 \frac{|\divg(u\Achi K(w,\cdot))|}{u}
 \le\eta_1\cT+\Pi_N\bigl(|K|^2+a|\nabla K|\bigr).
\end{equation}
In each of the first three trace stages,
\begin{equation}\label{four:maximal:def:F-derivative}
 w(F)\ge\eta_N a\sigma-\eta_1\cT
   -\Pi_N\bigl(\mathbf1_{\mathcal K}+|K|^2+a|\nabla K|\bigr),
\end{equation}
where $\mathcal K$ is a fixed compact enlargement of the supports of
$C,D_0,D_1$.  For fixed $\eta_1$ the bounds remain polynomial in $N$.
\end{lemma}

\begin{proof}
Differentiating $w=l\grad f/\sqrt D$ covariantly gives
\begin{equation}\label{four:maximal:def:w-derivative-exact}
 \nabla_iw=\Ad(H^f_{i\cdot})+pd\,w\,\partial_i t,
 \qquad
 d\log u=(p+pv+2)dt+H^f(w,\cdot).
\end{equation}
In the first formula, $d$ is the scalar $D^{-1}$.  Relative to the axis,
the $\Ad$ factor multiplies the axial output of $H^f$ by $d$.
The transverse and mixed entries are controlled by
\eqref{four:maximal:def:coercivity}; the possibly large axial entry occurs as
$d(j-K_{ee})$ and is also controlled because $d\le\sqrt d$.
Furthermore $d|dt|\le|dt|_{\Ad}$.  This proves the estimate for
$\nabla w$ without an inverse power of $l$ or $d$.
For the second formula use $\chi\le d$: the weighted norm of its first
term is at most $(2N+2)|dt|_{\Ad}$, while its last term is bounded by
$a(|M-K_{e\perp}|+\sqrt\chi\,|j-K_{ee}|)$.

Write $\Achi=\Id-c(p,v)w\otimes w$, where
$c(p,v)=(3+p)/(3+pv)$.  The relevant derivatives are
\[
 \partial_pc=\frac{3d}{(3+pv)^2},\qquad
 \partial_vc=-\frac{p(3+p)}{(3+pv)^2},\qquad
 |\nabla v|\le2a|\nabla w|.
\]
The $p'$ contribution retains its factor $d$, and is bounded by a
polynomial times $d|dt|$.  All other terms are bounded by a polynomial
times $|\nabla w|$.  This proves \eqref{four:maximal:def:derivative-bounds}.

In \eqref{four:maximal:def:drift-divergence}, differentiating $u$ contributes at most
\[
 |d\log u|_{\Achi}|K(w,\cdot)|_{\Achi}.
\]
The other derivatives contribute at most
\[
 |K|(|\nabla\Achi|+|\nabla w|)+a|\nabla K|.
\]
Apply \eqref{four:maximal:def:derivative-bounds} and then Young's inequality to
$\Pi_N|K|\sqrt\cT$.  Squaring a polynomial remains polynomial, proving
\eqref{four:maximal:def:drift-divergence} with the asserted dependence.

For \eqref{four:maximal:def:F-derivative}, the common height term gives exactly
$w(h)=\eta_N a\sigma$.  Derivatives of $C,D_0,D_1$ cost, on $\mathcal K$,
a fixed multiple of $1+|\nabla w|$, hence an arbitrarily small fixed
multiple of $\cT$ plus a polynomial compact error and $\Pi_N|K|^2$.
The only other term is $w(\tr_{\Achi}K)$ in stage three; it is bounded
in absolute value by a fixed multiple of
$a|\nabla K|+a|K||\nabla\Achi|$, handled in the same way.
All stage parameters lie in $[0,1]$, so the estimates are uniform in
them.  This proves the last assertion.
\end{proof}

\begin{proposition}[Floor exclusion]\label{four:maximal:prop:floor}
There is a choice $0<\delta_0<1/4$ independent of $N$ and a positive
polynomial $C_N$ such that, for every sufficiently large allowed $R$,
no smooth solution of any of the four homotopy stages has
\[
 \min_{\overline{\Omega_R}}t=-\eps.
\]
These choices use the fixed data and the preceding pointwise and height
estimates only.  In particular they do not use existence, an upper bound
for $Z$, or a uniform Hessian estimate.
\end{proposition}

\begin{proof}
We first record all end losses that will occur.  There is a fixed
constant $C_*$ with
\begin{equation}\label{four:maximal:def:floor-weight-domination}
 \mathbf1_{\mathcal K}+|K|^2+vF^2+v|\Ric_g|+a|\nabla K|
 \le C_*(\rho_0+v\rho_1).
\end{equation}
On a fixed compact set this follows from positivity of $\rho_0$ and
Lemma~\ref{four:maximal:lem:height-end}.  Outside that set,
$|F|\le |h|+|K|$ in all three trace stages, and the same lemma controls
$h^2$ by a fixed multiple of $r^{-2\gamma}$.
The metric and tensor decay give
$|\Ric_g|=O(r^{-2-q})$, $|K|^2=O(r^{-2-2q})$ and
$|\nabla K|=O(r^{-2-q})$.  Since $2\gamma<2+q$ and $2+2q>4+\delta$,
the curvature and squared tensor terms have the stated bounds.
For the derivative term the precise estimate is
\[
 ar^{-2-q}
 \le\tfrac12a^2r^{-2\gamma}
       +\tfrac12r^{-4-2q+2\gamma}
 \le\tfrac12v\rho_1+\tfrac12\rho_0,
\]
where $2(q-\gamma)>\delta$ was used.  This proves
\eqref{four:maximal:def:floor-weight-domination}.  Also $\rho\le C_\rho\rho_0$ for a
fixed $C_\rho$, by its prescribed tail and compact positivity of $\rho_0$.

Suppose first that the minimum occurs in the interior in one of the
first three stages.  Then $dt=0$ and $\Hess t\ge0$ there.
Proposition~\ref{four:maximal:prop:scalar-identity} and
\eqref{four:maximal:def:graph-gauss} imply, using $\tr K=0$,
\[
 u^{-1}\divg V
 \ge\mu-\tfrac12R_g+J(w)+v\Ric_g(e_f,e_f)
       +\cT-\mathcal S(H^f)+w(F).
\]
Here $\mu-R_g/2=-|K|^2/2$ and
$|J(w)|\le C_*a|\nabla K|$.
Use \eqref{four:maximal:def:floor-comparison}, \eqref{four:maximal:def:F-derivative}, and
\eqref{four:maximal:def:floor-weight-domination}; then pass from $V$ to
$V_\lambda=V-(1-\lambda)u\Achi K(w,\cdot)$ by
\eqref{four:maximal:def:drift-divergence}.  Choose the two fixed Young errors smaller
than $c_2/4$.  This yields an absolute $c_3>0$ and a polynomial bound
$\Pi_N^{\mathrm f}$ such that at the proposed minimum
\begin{equation}\label{four:maximal:def:floor-divergence}
 u^{-1}\divg V_\lambda
 \ge c_3\cT+\eta_N a\sigma
       -\Pi_N^{\mathrm f}(\rho_0+v\rho_1).
\end{equation}
In particular $c_3$ is independent of $N$.

Fix $\delta_0<\min(c_3/2,1/4)$.  At the floor $m_0=1$, and the prescribed
right side in the first three stages is
$\delta_0\cT+\rho+\delta_0\eta_N a\sigma-C_N(\rho_0+v\rho_1)$.
The lower bound in \eqref{four:maximal:def:floor-divergence} exceeds it by at least
\[
 (c_3-\delta_0)\cT+(1-\delta_0)\eta_N a\sigma
 +(C_N-\Pi_N^{\mathrm f})(\rho_0+v\rho_1)-\rho.
\]
Choose $C_N\ge\Pi_N^{\mathrm f}+C_\rho+1$.  The last display is strictly
positive because $\rho_0>0$, contradicting the equation.  This choice is
polynomial and independent of $R$ and of the solution.

In stage four use \eqref{four:maximal:def:terminal-reduction}.  At an interior floor
minimum the left side is
$u^{-1}\divg(3u\grad t)=3\Delta t\ge0$,
while the right side is
\[
 \eta\bigl(\delta_0|K|^2/2+\rho-C_N\rho_0\bigr).
\]
Increasing $C_N$ by a fixed constant makes the bracket strictly negative
everywhere, since $|K|^2\le C_*\rho_0$.  Thus an interior floor minimum
is impossible whenever $\eta>0$.

An inner boundary minimum in the first three stages is also impossible.
At $t=-\eps$ we have $j_0=0$.  Since
$V_\lambda=V-(1-\lambda)u\Achi K(w,\cdot)$, the drift correction in
\eqref{four:maximal:def:homotopy-boundary} cancels, leaving $V_\nu/u=\mathcal B$.
Lemma~\ref{four:maximal:lem:boundary-identity} therefore gives
\[
 3\partial_\nu t=-H-N_B<0.
\]
This contradicts the nonnegative derivative into $\Omega_R$ required
at a boundary minimum.  In stage four, where $f=0$, the corresponding
formula is $3\partial_\nu t=\eta(-H-N_B)<0$ for $\eta>0$.
The outer boundary is excluded by Lemma~\ref{four:maximal:lem:height-end}.

It remains to consider $\eta=0$.  The last-stage equation and boundary
conditions then reduce to
\[
 \divg(3L_0(Z)\grad Z)=0,\qquad
 \partial_\nu Z|_B=0,\qquad Z|_{S_R}=0.
\]
Multiplying by $Z$ and integrating gives
$\int_{\Omega_R}3L_0(Z)|\grad Z|^2\dd V_g=0$.
Thus $Z=t=0$, which is strictly above the floor.  This treats every
stage and every possible location of a floor minimum.
\end{proof}

The order of the choices is now fixed: strict background data and
collars, $\eps$ and the cutoffs, the absolute small $\delta_0$, the
polynomial $C_N$, and finally sufficiently large $N$ and allowed $R$.
The estimates in the next Section concern arbitrary smooth admissible
solutions with $t\ge-\eps$.  Proposition~\ref{four:maximal:prop:floor} separately
ensures that a homotopy solution cannot meet the boundary of this range.

\section{Global a priori bounds}\label{four:maximal:sec:bounds}

Throughout this section the strict background data, collars, and
$\eps>0$ are fixed. We consider arbitrary smooth solutions of the
homotopy on $\Omega_R$ satisfying $t\ge-\eps$. Existence is not assumed:
each assertion is an estimate for every solution that might occur.
The constants denoted by $\Pi_N$ are bounded by a polynomial in $N$,
independently of the solution, homotopy parameter, and allowed radius.
Constants denoted by $C(N)$ may depend arbitrarily on the fixed integer
$N$. All integrals without an indicated graph measure use $g$.

\subsection{Separation from the threshold region}

\begin{lemma}[Compact separation]\label{four:maximal:lem:separation}
Let $Q_0$ be a compact subset of $\overline\Omega$ disjoint from
$\mathcal D_c$. There are $\eta_2>0$, $N_0$, and allowed radii
$R_{\min}(N)\to\infty$ such that, in stages~1 and~2,
\begin{equation}\label{four:maximal:bounds:separation}
 h\ge c+\eta_2\quad\hbox{on }Q_0
 \qquad(N\ge N_0,\ R\ge R_{\min}(N)).
\end{equation}
The constants are uniform over both stages.
\end{lemma}

\begin{proof}
We first identify the geometric inequality satisfied by a relaxed
limit. Take any sequence of admissible solutions with $N_i\to\infty$
and allowed $R_i\to\infty$, and write
\[
 \underline h(x)=\lim_{j\to\infty}
 \inf\{h_i(y):i\ge j,\ \dist(x,y)<j^{-1}\}.
\]
On compact sets this is a finite lower semicontinuous function, by
Lemma~\ref{four:maximal:lem:height-end}. Let a smooth function $\phi$ touch
$\underline h$ from below at an interior point $x$, with
$\grad\phi(x)\ne0$. Subtracting a small multiple of the fourth power
of the distance makes the contact strict without changing its two-jet.
Minimization on a small closed ball then gives a subsequence of local
lower contacts of $\phi$ plus constants with $h_i$, at points $x_i\to x$,
whose contact values converge to $\underline h(x)$.

At these contacts, $\eta_{N_i}\grad f_i=\grad\phi$, so
\[
 l_i\sigma_i\ge\frac{\ell_i}{\eta_{N_i}}|\grad\phi(x_i)|\longrightarrow\infty,
 \qquad d_i,\chi_i\longrightarrow0,\quad a_i\longrightarrow1.
\]
In particular this conclusion uses the floor and the fixed test
gradient, rather than an estimate for the solution gradients.
Put $\nu_\phi=\grad\phi/|\grad\phi|$. Positivity of $\Achi$ and the
Hessian inequality at a lower contact give
\[
 \tr_{A_{\chi_i}}K+
 \frac{a_i}{|\grad\phi|}
       \tr_{A_{\chi_i}}\Hess\phi
 \le F_i\le h_i.
\]
The last inequality holds throughout stages~1 and~2 because
$C\le0$ and $D_0\le0$. Passing to the limit proves
\begin{equation}\label{four:maximal:bounds:lower-test}
 \frac{\tr_{\nu_\phi^\perp}\Hess\phi}{|\grad\phi|}
       +\tr_{\nu_\phi^\perp}K\le\underline h(x).
\end{equation}
The left side is the outward expansion of the regular level of $\phi$.
Neither convergence of the homotopy parameters nor continuity bounds
for the solution coefficients were needed.

Fix $c<k<k'<0$.  Apply Lemma~\ref{n3:apriori:sublevel-support}
on $U=\operatorname{int}\Omega$ with $u=\underline h$, $Q=K$ and
$q_0=k'$.  The relaxed limit is finite and lower semicontinuous.
At a nonzero-gradient lower test whose contact value is less than $k'$,
\eqref{four:maximal:bounds:lower-test} bounds its expansion by that
value, hence by $k'$.  The lemma therefore shows that every smooth
exterior support to the closed sublevel $\{\underline h\le k\}$
at an interior frontier point has expansion at most $k'$.  In particular,
this conclusion also applies when the sublevel has a plateau at height $k$.

Lemma~\ref{four:maximal:lem:height-end} makes this sublevel bounded: for a large
fixed $r$, its lower height bound exceeds the fixed negative number
$k$. Adjoin $\mathcal D_c$ to obtain a compact closed set. At a frontier
point on $B$, each exterior support also supports $\mathcal D_c$;
tangency comparison with its smooth boundary gives expansion at most
$c$. Lemma~\ref{four:maximal:lem:support-enclosure} therefore places the union in
$\mathcal D_{c'}$ for every $k'<c'<0$.

Hausdorff right continuity at $c$ and the positive distance of $Q_0$
from $\mathcal D_c$ allow us to choose $c<k<k'<c'<0$ so that
$\mathcal D_{c'}\cap Q_0=\varnothing$. If
\eqref{four:maximal:bounds:separation} failed, choose $N_i\to\infty$, allowed
$R_i\ge i$, and $x_i\in Q_0$ with $h_i(x_i)\le c+i^{-1}$.
A subsequence has $x_i\to x\in Q_0$, and then
$\underline h(x)\le c<k$, contradicting that containment.
This also proves uniformity over the two stages. Enlarging the allowed
radius function never invalidates the assertion.
\end{proof}

\subsection{Collar barriers and polynomial trace estimates}

\begin{lemma}[Inner collar slopes]\label{four:maximal:lem:collar-slopes}
For all sufficiently large $N$ and allowed $R$, there are fixed
constants $0<c_5<C_6$ such that
\begin{equation}\label{four:maximal:bounds:collar-slopes}
 \frac{c_5}{\eta_N}\le\partial_\nu f\le\frac{C_6}{\eta_N}
 \quad\hbox{in stages 1 and 2},\qquad
 |\partial_\nu f|\le\frac{C_6}{\eta_N}
 \quad\hbox{in stage 3}.
\end{equation}
Here $\nu$ points into $\Omega$. In stage~4, $f=0$.
\end{lemma}

\begin{proof}
Use a fixed short collar $0\le s\le s_0$ from
Lemma~\ref{four:maximal:lem:threshold-collars}, with $C=-\kappa s$,
$\nu_s=\grad s/|\dd s|$, and $0<c_s\le|\dd s|\le C_s$.
All its geometric derivatives are bounded, and
$P_s+H_s\ge c$, where $P_s=\tr_{\{s=\mathrm{constant}\}}K$.
The finitely many outer collar faces satisfy
Lemma~\ref{four:maximal:lem:separation}.

For a positive-slope test $h_0=c+\psi(s)$, the trace operator, with
$t$ computed from the solution's $Z$ and the common contact gradient,
is exactly
\begin{equation}\label{four:maximal:bounds:collar-operator}
 P_s+aH_s+\chi K(\nu_s,\nu_s)
 +a\chi\left(
    \frac{\Hess s(\nu_s,\nu_s)}{|\dd s|}
       +|\dd s|\frac{\psi''}{\psi'}\right).
\end{equation}
For a test with negative slope, replace $a$ by $-a$ in this formula.
If $|\psi'|$ lies between fixed positive constants, the floor gives
\begin{equation}\label{four:maximal:bounds:steep-coefficients}
 d\le C(\eta_N/\ell)^2=Ce^{-\eps N},\qquad
 1-a\le d,\qquad
 \frac{N\chi}{d}=\frac{3N}{3+pv}\ge\frac{3N}{3+N}.
\end{equation}
Thus $a\ge1/2$ and $N\chi\ge c_4d$ for large $N$.

Choose a fixed smooth $\vartheta_1:[0,\infty)\to[0,1]$ equal to one
on $[0,1]$ and zero on $[2,\infty)$, and put
$b_N(s)=\Lambda_* N\vartheta_1(Ns)$, where $\Lambda_*>0$ is a fixed collar
constant, independent of $N$. The constant will be
chosen large; its crucial property is
\[
 0\le\int_0^{s_0}b_N(s)\dd s\le2\Lambda_*.
\]
For the lower barrier in stages~1 and~2, prescribe
$\psi'(s)=q_-\exp(\int_0^s b_N)$ and $\psi(0)=0$.
For the chosen value of $\Lambda_*$, take $q_->0$ so small that $\psi'\le\kappa/2$ and
$\psi(s_0)<\eta_2$. Positive normal direction makes $D_0=0$.
Subtracting the prescribed trace right side from
\eqref{four:maximal:bounds:collar-operator} gives a residual at least
\begin{equation}\label{four:maximal:bounds:lower-residual}
 \kappa s/2-Cd+a\chi|\dd s|b_N(s).
\end{equation}
The constant $C$ in this residual depends only on collar geometry.
We choose $\Lambda_*$ so large that the last term dominates $2Cd$ on
$[0,1/N]$. On $[1/N,s_0]$, the first term dominates $Cd$ for large
$N$, because $Ne^{-\eps N}\to0$, and the last term has a favorable
sign. This is a strict lower barrier. Its outer value is below $h$
by separation and its inner value agrees with $h$.

For the upper barrier take
$\psi'(s)=q_+\exp(-\int_0^s b_N)$.
Choose its fixed minimum slope larger than twice the bounded linear
variation of $P_s+H_s-c+\kappa s$, and choose $q_+$ still larger if
needed to dominate the height bound at $s_0$. The residual is now at
most $-c_7s+Cd-a\chi|\dd s|b_N$, for fixed $c_7>0$.
The same two intervals prove strict negativity. The bounded integral
of $b_N$ leaves all lower and upper test slopes between fixed positive
constants, independently of $N$.

For completeness, at an interior positive maximum of
$h_0-h$, the test Hessian is at most the solution Hessian, whereas its
height is larger. The positive principal matrix and the strictly
increasing dependence of $F$ on $h$ contradict a positive test residual.
This proves lower comparison; reversing the difference proves upper
comparison. In both comparisons $Z$ and the gradient are the same
at contact, so all coefficients other than the height agree.
Taking inward derivatives at $s=0$ gives the first assertion of
\eqref{four:maximal:bounds:collar-slopes}.

In stage~3 the face height is $b_\zeta=(1-\zeta)c$.
Use $b_\zeta\pm\psi(s)$ with fixed large positive slopes and
$\psi''/\psi'=-b_N$. At $s=0$ the two unit-direction inequalities
in \eqref{four:maximal:def:face-direction-signs} say that the positive-slope residual is nonpositive
and the negative-slope residual is nonnegative. Replacing $\pm \nu_s$ by
$\pm a\nu_s$ costs $O(d)$, including the term $\chi K(\nu_s,\nu_s)$; moving to
distance $s$ costs $O(s)$, uniformly in $\zeta$. The change of height
then supplies a strict linear margin, with the desired sign, when the
minimum test slope exceeds the fixed spatial error constant.
The term involving $b_N$ has the correct sign for both tests and
dominates the $O(d)$ error on $[0,1/N]$. On the remainder the linear
margin dominates it. The height bound handles $s=s_0$.
These two comparisons bound the slope in absolute value in stage~3.
The terminal trace equation has $f=0$, as established in
Section~\ref{four:maximal:sec:deformation}.
\end{proof}

\begin{lemma}[Weighted trace and boundary flux]\label{four:maximal:lem:weighted-trace}
Let $\mathcal C$ be a fixed union of collars. In stages~1 and~2, every
nonnegative $C^1$ function $\varphi$ satisfies
\begin{align}
 \int_B L_0\varphi
 &\le\Pi_N\int_{\mathcal C}
  u\bigl((1+\sqrt{\cT})\varphi+|\dd\varphi|_{\Achi}\bigr),
       \label{four:maximal:bounds:trace-weighted}\\
 \int_B\varphi
 &\le\Pi_N\int_{\mathcal C}
  \sqrt D\bigl((1+\sqrt{\cT})\varphi+|\dd\varphi|_{\Achi}\bigr).
       \label{four:maximal:bounds:trace-unweighted}
\end{align}
The integrands vanish outside the support of $\varphi$ and its
derivative; thus both statements localize to any base patch.
Moreover
\begin{equation}\label{four:maximal:bounds:boundary-flux}
 |(V_\lambda)_\nu|\le C_7ud=C_7L_0\sqrt d
                     \le C_8(\eta_N/\ell)L_0.
\end{equation}
Consequently,
\begin{equation}\label{four:maximal:bounds:boundary-absorption}
 \int_B|(V_\lambda)_\nu|
 \le \Pi_N\frac{\eta_N}{\ell}
             \int_{\mathcal C}u(\rho+\delta_0\cT).
\end{equation}
All constants in these assertions have the stated polynomial control.
\end{lemma}

\begin{proof}
Set $W=l\grad_{\gbar}f=\sqrt d\,w$. For each covector $\xi$,
\[
 |\xi(W)|\le|\xi|_{\Ad}
       \le\sqrt{(3+N)/3}\,|\xi|_{\Achi},\qquad uW=L_0w.
\]
The second inequality follows from $d/\chi=(3+pv)/3$.
Lemma~\ref{four:maximal:lem:collar-slopes} gives $w\cdot\nu=a\ge1/2$ on $B$.
Direct differentiation yields
\begin{align*}
 \divg w&=\tr_{\Ad}H^f+dp\,w(t),\\
 u^{-1}\divg(uW)
 &=\sqrt d\bigl(\tr_{\Ad}H^f+(dp+p+2)w(t)\bigr),\\
 D^{-1/2}\divg(\sqrt D W)
 &=\sqrt d\bigl(\tr_{\Ad}H^f+dp\,w(t)\bigr).
\end{align*}
These last two expressions have absolute value at most
$\Pi_N(1+\sqrt{\cT})$ on the collars. Indeed the axial tensor term is
$d^{3/2}H^f_{ee}$, controlled by the $d j^2$ term of
Lemma~\ref{four:maximal:lem:coercivity}, and
$\sqrt d\,|w(t)|\le|\dd t|_{\Ad}$; all background $K$ terms are
bounded there. No inverse power of $l$ has occurred.

Take a fixed collar cutoff $\eta$, equal to one near $B$ and zero
near the outer collar faces. Apply the divergence theorem to
$\eta\varphi uW$. The domain's outward normal on $B$ is $-\nu$,
so the boundary integral is $-\int_B L_0a\varphi$.
Taking absolute values of the three interior terms, and using
$a\ge1/2$, proves \eqref{four:maximal:bounds:trace-weighted}.
The derivative of $\eta$ costs only a fixed multiple of $\varphi$.
Using $\sqrt D W=w$ instead proves \eqref{four:maximal:bounds:trace-unweighted}.
The same computation proves the asserted localization.

On $B$, positive slope gives $w_\nu=a$, $D_0=0$, and
$F=c=P_B+H$. Thus the expression inside the boundary multiplier of
the homotopy is
\[
 (a-1)H-N_Bd-(1-\lambda)\chi aK(\nu,\nu).
\]
The multiplier lies in $[0,1]$, and $1-a\le d$, $\chi\le d$.
This proves the first inequality in \eqref{four:maximal:bounds:boundary-flux}.
The lower collar slope gives
$\sqrt d\le(l|\partial_\nu f|)^{-1}\le C\eta_N/\ell$,
proving the rest. Finally apply \eqref{four:maximal:bounds:trace-weighted} with
$\varphi=1$. On the fixed collars, $\rho$ has a positive minimum, so
$1+\sqrt{\cT}\le C(\rho+\delta_0\cT)$ with fixed $C$.
This proves \eqref{four:maximal:bounds:boundary-absorption}.
\end{proof}

\subsection{The global high-level energy bound}

The collar estimates have already made the inner flux small with
polynomial constants.  We now obtain a bounded range for the conformal
variable at each fixed parameter.  That is a separate task: its constants
may be arbitrary at fixed parameter because the mass estimate will
return to the earlier collar bounds.

We now allow arbitrary constants $C(N)$. At sufficiently high levels
$Z>Z_0(N)>0$, the unscaled source satisfies
\begin{equation}\label{four:maximal:bounds:high-source}
 \Xi\ge\delta_0\cT+\rho+\tfrac12\delta_0\eta_N a\sigma.
\end{equation}
To see this, $\mathcal P$ is bounded at fixed $N$, and on the admissible
support of $m_0$ we have
$-\eps\le t\le-\eps+1/N$. There
$D=e^{6(Z-t)}$ and
$\eta_N a\sigma=(\eta_N/l)(\sqrt D-D^{-1/2})$ tends uniformly to infinity
as $Z\to\infty$. Outside this band the penalty vanishes.
On the full range we also have $\Xi\ge\delta_0\cT-C(N)$.

Only stage~1 requires an energy argument to handle its nonzero drift.
Choose a smooth nondecreasing $k_0$ which is zero below $Z_0$, one
above $Z_0+1$, and has bounded derivative. Testing the divergence
equation with $k_0(Z)$, whose outer trace is zero, gives
\[
 \int_{\Omega_R}uk_0\Xi
       +3\int_{\Omega_R}uk_0'|\dd Z|_{\Achi}^2
 =-\int_B k_0(V_\lambda)_\nu
       -\lambda\int_{\Omega_R}uk_0'
                   \langle K(w,\cdot),\dd Z\rangle_{\Achi}.
\]
Young's inequality absorbs half the gradient term and leaves
$C\int uk_0'|K|^2$. This is at most $C(N)$: on the derivative strip,
\[
 u=l e^{3Z-t}\le e^{1+3(Z_0+1)+\eps},
\]
and $|K|^2$ is integrable since $q>1$.
By \eqref{four:maximal:bounds:boundary-flux} and \eqref{four:maximal:bounds:trace-weighted},
the boundary cost is at most
\[
 \Pi_N\frac{\eta_N}{\ell}\int_{\mathcal C}
   u\bigl((1+\sqrt{\cT})k_0+k_0'|\dd Z|_{\Achi}\bigr).
\]
The first term is absorbed by a fixed fraction of
$\int uk_0(\delta_0\cT+\rho)$ for large $N$, because
$\Pi_N\eta_N/\ell\to0$. Young's inequality absorbs the second term
into $\int uk_0'|\dd Z|_{\Achi}^2$ and a bounded derivative-strip
integral over $\mathcal C$. We obtain
\begin{equation}\label{four:maximal:bounds:global-energy}
 \int_{\Omega_R}u\bigl[
       k_0(\cT+\rho+\eta_N a\sigma)
                    +k_0'|\dd Z|_{\Achi}^2\bigr]\le C(N).
\end{equation}
This is uniform in $R$ and in the stage~1 parameter.

Let $\Gamma$ denote the ordinary graph of $f$ in $(\Omega,g)\times\R$
with product metric $g+\dd b^2$. This is a four-dimensional graph.
Its measure and gradient satisfy, at fixed $N$,
\begin{equation}\label{four:maximal:bounds:graph-comparison}
 \dd\Gamma\asymp_N\sqrt D\dd V_g,\qquad
 |\grad_\Gamma\phi|\asymp_N|\dd\phi|_{\Achi}.
\end{equation}
These comparisons use only $\ell\le l\le e$ and
$3d/(3+N)\le\chi\le d$. Also
\[
 D^{1/3}\le C(N)(1+\eta_N a\sigma),\qquad
 e^{2Z}\sqrt D=l^{-1}uD^{1/3}.
\]
For the first inequality, put $y=l\sigma$. It is immediate for
$y\le1$; for $y\ge1$ use
$\eta_N a\sigma\ge(\eta_N/(\sqrt2 l))y$ and $y^{2/3}\le y$.
On any fixed compact base set $Q$, the positive minimum of $\rho$
and \eqref{four:maximal:bounds:global-energy} control the integral of
$u(1+\eta_N a\sigma)$ above $Z_0+1$. Below that level the floor and
$D=e^{6(Z-t)}$ bound every integrand on $Q$. Consequently
\begin{equation}\label{four:maximal:bounds:graph-l2}
 \int_{\Gamma\cap\pi^{-1}(Q)} e^{2Z}\dd\Gamma\le C(N,Q).
\end{equation}
Here $\pi$ is the graph projection. The estimate is local in the base;
the global estimate from which it follows is
\eqref{four:maximal:bounds:global-energy}.

\subsection{Graph Sobolev inequality and the exponential test}

We continue with stage~1 for the next two estimates.
For a fixed compact base patch and a function $\omega$ supported over
it, away from $S_R$ but possibly meeting $B$, we claim
\begin{equation}\label{four:maximal:bounds:graph-sobolev}
 \|\omega\|_{L^4(\dd\Gamma)}^2
 \le C(N)\int_\Gamma
       \bigl(|\grad_\Gamma\omega|^2+(1+\cT)\omega^2\bigr)\dd\Gamma.
\end{equation}
Extend a compact enlargement of the base metric smoothly to a closed
manifold and use a smooth isometric Euclidean embedding
\cite[Theorem~2, p.~59]{Nash1956}. Its product with $\R$ gives a Euclidean
immersion of the graph. The contribution to its mean curvature from
the base embedding is bounded by a fixed constant: it is a trace of
the fixed second fundamental form on a four-dimensional tangent
plane. The remaining product mean curvature has magnitude
\[
 \frac{\sqrt D}{l\sqrt{1+\sigma^2}}
 \left|\tr_\perp H^f+(1+\sigma^2)^{-1}H^f_{ee}\right|
 \le C(N)(1+\sqrt{\cT}).
\]
Indeed the prefactor is bounded at fixed $N$, while
$(1+\sigma^2)^{-1}\le C(N)d$ and the axial term is controlled by
$d j^2$ in Lemma~\ref{four:maximal:lem:coercivity}.

The Euclidean submanifold Sobolev inequality in dimension four gives
\[
 \left(\int_\Gamma |\omega|^4\dd\Gamma\right)^{3/4}
 \le C(N)\int_\Gamma
    \bigl(|\omega|^2|\grad_\Gamma\omega|
               +(1+\sqrt{\cT})|\omega|^3\bigr)\dd\Gamma
           +C\int_B|\omega|^3.
\]
We have applied the $L^1$ Sobolev inequality to $|\omega|^3$;
see Michael--Simon~\cite{MichaelSimon1973} and the precise
rectifiable-submanifold statement in
\cite[Chapter~4, \S5, Theorem~5.7, p.~98]{Simon2014}.
Its hypotheses hold for this smooth multiplicity-one immersion with
locally integrable mean curvature and compactly supported test.
To justify the displayed boundary term from the interior statement,
extend each individual smooth graph across $B$ and cut off in an
exterior strip whose width tends to zero. The integral of the cutoff
gradient tends to the boundary integral, while the strip's other
integrals tend to zero. This requires smoothness of each graph, not a
uniform curvature estimate for extensions. The induced boundary
metric is $g|_B$, since $f$ is constant there.

Apply \eqref{four:maximal:bounds:trace-unweighted} to $|\omega|^3$ and then
\eqref{four:maximal:bounds:graph-comparison}. It bounds the last term by the same
two interior terms. Each such term is at most
\[
 C(N)\left(\int_\Gamma|\omega|^4\dd\Gamma\right)^{1/2}
 \left(\int_\Gamma
       (|\grad_\Gamma\omega|^2+(1+\cT)\omega^2)\dd\Gamma\right)^{1/2}
\]
by Cauchy--Schwarz. Dividing, and squaring, proves
\eqref{four:maximal:bounds:graph-sobolev}; the case $\omega=0$ is immediate.

Next let $\psi$ be a smooth compact base cutoff of the same kind.
There is $k_*(N)\ge1$ such that for every real $k\ge k_*(N)$,
\begin{equation}\label{four:maximal:bounds:exponential-energy}
 \int_\Gamma e^{2kZ}\psi^2
        (\cT+k|\grad_\Gamma Z|^2)\dd\Gamma
 \le C(N)k\int_\Gamma e^{2kZ}
               (\psi^2+|\dd\psi|_g^2)\dd\Gamma.
\end{equation}
The constant $C(N)$ here is independent of $k$. To prove this, test
the stage~1 divergence equation with
$\varphi=\psi^2e^{2kZ}/L_0$. Write
$Q=\psi^2e^{2kZ}$ and $b=\lambda K(w,\cdot)$.
After integration, cancellation of $L_0$ gives the exact identity
\begin{align*}
 &\int_{\Omega_R}\sqrt D\,Q
       \bigl(\Xi+6k|\dd Z|_{\Achi}^2\bigr)\\
 &=-\int_B(V_\lambda)_\nu Q/L_0
   -\int_{\Omega_R}\sqrt D\,e^{2kZ}
       \bigl[2\psi\langle3\dd Z+b,\dd\psi\rangle_{\Achi}
          +2k\psi^2\langle b,\dd Z\rangle_{\Achi}\bigr]\\
 &\hspace{1cm}
   +\int_{\Omega_R}\sqrt D\,Q(p+2)
                    \langle3\dd Z+b,\dd t\rangle_{\Achi}.
\end{align*}
Use $\Xi\ge\delta_0\cT-C(N)$, $|b|_{\Achi}\le|K|$, and
Lemma~\ref{four:maximal:lem:coercivity}. For example,
\[
 (p+2)|\dd t|_{\Achi}(|\dd Z|_{\Achi}+|K|)
 \le\frac{\delta_0}{8}\cT
               +C(N)(|\dd Z|_{\Achi}^2+|K|^2).
\]
Young's inequality handles the drift and cutoff products with a
fixed fraction of $k|\dd Z|_{\Achi}^2$ and a remainder bounded by
$C(N)k e^{2kZ}(\psi^2+|\dd\psi|^2)$. Choose $k_*(N)$ large enough
to absorb the preceding $C(N)|\dd Z|_{\Achi}^2$ term.

By \eqref{four:maximal:bounds:boundary-flux} the boundary contribution is at most
$C\int_B Q$ (enlarging a fixed constant since $\eta_N/\ell\le1$).
Applying \eqref{four:maximal:bounds:trace-unweighted} to $Q$ gives, only on
$\mathcal C\cap\supp\psi$,
\[
 C(N)\int\sqrt D\,e^{2kZ}
   \bigl[(1+\sqrt{\cT})\psi^2
       +2|\psi||\dd\psi|+2k\psi^2|\dd Z|_{\Achi}\bigr].
\]
The $\sqrt{\cT}$ and $k|\dd Z|_{\Achi}$ terms are absorbed by
small fixed fractions of $\delta_0\psi^2\cT$ and
$k\psi^2|\dd Z|_{\Achi}^2$, leaving at most
$C(N)k e^{2kZ}(\psi^2+|\dd\psi|^2)$.
Finally \eqref{four:maximal:bounds:graph-comparison} proves
\eqref{four:maximal:bounds:exponential-energy}. In particular, neither the trace
cutoff nor differentiation of $1/L_0$ introduces a constant exponential
in $k$.

\subsection{Moser iteration and the full range}

\begin{proposition}[Bounded variables]\label{four:maximal:prop:bounded-range}
For every sufficiently large fixed $N$ and all sufficiently large
allowed $R$, every admissible smooth homotopy solution satisfies
\begin{equation}\label{four:maximal:bounds:bounded-range}
 |f|\le C_0/\eta_N,\qquad
 -\eps\le t\le Z\le C(N),\qquad |\grad f|\le C(N).
\end{equation}
The constants are independent of $R$, the solution, and the homotopy
parameter. No uniform Hessian or coefficient continuity estimate is
used in obtaining these bounds.
\end{proposition}

\begin{proof}
We first finish stage~1. Apply \eqref{four:maximal:bounds:graph-sobolev} to
$\omega=\psi e^{kZ}$ and then use
\eqref{four:maximal:bounds:exponential-energy}. If $\psi=1$ on a smaller patch
and $|\dd\psi|\le C/\mathrm{gap}$, the resulting norm step is
\[
 \|e^Z\|_{L^{4k}(\text{smaller graph})}
 \le\bigl[C(N)k^2(1+\mathrm{gap}^{-2})\bigr]^{1/(2k)}
       \|e^Z\|_{L^{2k}(\text{larger graph})}.
\]
Choose $k_j=2^j k_*(N)$ and geometrically decreasing gaps between
fixed nested base patches. The logarithms of the factors sum because
$\sum_j (1+j)/2^j<\infty$. Writing $S(r)=\sup_{Q_r}e^Z$ for a
nested family of compact patches gives
\[
 S(r')\le C(N)(r''-r')^{-A(N)}
          \|e^Z\|_{L^{2k_*}(\Gamma\cap\pi^{-1}(Q_{r''}))}.
\]
Interpolation with \eqref{four:maximal:bounds:graph-l2} yields
\[
 \|e^Z\|_{2k_*}
     \le S(r'')^{1-1/k_*}\|e^Z\|_2^{1/k_*}.
\]
If $k_*=1$ this already proves the bound. Otherwise Young's inequality
gives, for every $\alpha>0$,
\[
 S(r')\le\alpha S(r'')
       +C(N,\alpha)(r''-r')^{-A'(N)}.
\]
Iterate over radii increasing geometrically to a fixed larger radius
and take $\alpha<2^{-A'(N)-1}$. The sum of the remainder terms is
finite. The final term $\alpha^j S(r_j)$ tends to zero for each
individual smooth solution, since all $Q_{r_j}$ lie in a fixed compact
patch. The resulting bound is independent of that individual
supremum. Finitely many such patches prove an upper bound for $Z$
on any fixed compact base set, including $B$.

To extend the bound along the end, use \eqref{four:maximal:def:drift-divergence},
with a fixed small coefficient:
\[
 u^{-1}|\divg(u\Achi K(w,\cdot))|
 \le\frac{\delta_0}{4}\cT
       +\Pi_N(|K|^2+a|\grad K|).
\]
Since $a\sigma= l\sigma^2/\sqrt D\ge e^{-1}a^2$,
Young's inequality bounds the second term by
\[
 \frac{\delta_0}{4}\eta_N a\sigma
       +C(N)(|K|^2+|\grad K|^2).
\]
The last expression, apart from its favorable term, is $o(\rho)$
on the end. Indeed $2+2q>4+\delta$ and
$|\grad K|^2=O(r^{-4-2q})$.
Beyond a sufficiently large fixed-$N$ sphere, the high-level source
\eqref{four:maximal:bounds:high-source} therefore implies
\[
 \divg(3u\Achi\grad Z)>0
 \quad\hbox{where }Z>Z_0(N).
\]
At an interior maximum its left side is
$3u\Achi:\Hess Z\le0$, a contradiction.
The compact bound controls the inner sphere and $Z=0$ on $S_R$.
This proves the global upper bound in stage~1.

In stages~2 and~3, $\lambda=0$. The same high-level source is
strictly positive, so it excludes interior high maxima directly.
The boundary slope bound implies
\[
 Z=t+\tfrac16\log(1+l^2|\partial_\nu f|^2)
       \le t+\tfrac16\log(1+e^2 C_6^2/\eta_N^2)
       \quad\hbox{on }B.
\]
Thus sufficiently high boundary $Z$ forces $t\ge1$. At such points
$j_0(t)=1$, and the boundary prescription gives zero conormal flux.
The boundary point principle excludes a nonconstant high maximum
there: in a neighborhood of the maximum the individual smooth
operator is elliptic and $\divg(3u\Achi\grad Z)>0$, so its outward
conormal derivative at a boundary maximum is strictly positive,
contrary to the zero datum. This uses ellipticity only for each
individual solution. The outer Dirichlet value is zero, so the upper
bound follows uniformly through these stages.

In stage~4, $f=0$ and $t=Z$. For positive scaling $\eta$, sufficiently
high levels lie outside the penalty band and have source
$\eta(\delta_0\cT+\rho)>0$; the same maximum and boundary point
argument applies. At zero scaling, multiplying the homogeneous
equation by $Z$ gives $\int 3L_0|\dd Z|^2=0$, and the outer value yields
$Z=0$. These bounds are uniform even as $\eta\downarrow0$.

Finally the height bound is Lemma~\ref{four:maximal:lem:height-end}, and
$t\le Z$ follows from $D\ge1$. The implicit equation gives
\[
 D=e^{6(Z-t)}\le e^{6(C(N)+\eps)},\qquad
 |\grad f|^2=(D-1)/l^2\le\ell^{-2}e^{6(C(N)+\eps)}.
\]
This proves \eqref{four:maximal:bounds:bounded-range}.
\end{proof}

\begin{remark}\label{four:maximal:bounds:parameter-order}
Separation, the collar slopes, and
\eqref{four:maximal:bounds:trace-weighted}--\eqref{four:maximal:bounds:boundary-absorption}
were proved before introducing any arbitrary $C(N)$.
The graph comparisons, Moser iteration, and subsequent regularity
may have arbitrary fixed-$N$ losses. The final mass comparison uses
the earlier weighted trace and boundary-flux inequalities directly;
none of these later constants enters its limit as $N\to\infty$.
\end{remark}

\section{Local regularity of the coupled system}\label{four:maximal:sec:regularity}

The bounded-range estimate makes the trace equation uniformly elliptic
at fixed $N$.  The second equation still contains $|\grad Z|^2$ and
$|\Hess f|^2$, so classical estimates cannot yet be applied to it.
We first use Theorem~\ref{four:thm:axial-regularity} to control $df$
and the local integral of $|\Hess f|^2$.  The scalar measure-source
estimate of Lemma~\ref{b:lem:natural-growth} then gives H\"older
continuity of $Z$, including at both boundary types.  These two gains
start the classical bootstrap.

All constants in this section may depend arbitrarily on the fixed data
and $N$.  The earlier polynomial bounds will control the final energy
error.  In coordinate calculations $D$ denotes ordinary differentiation;
the graph scalar $D=1+l^2|\grad f|^2$ retains its earlier meaning in
geometric formulas.

\begin{proposition}\label{four:maximal:prop:classical-estimates}
Fix the strict data, $\eps$, a sufficiently large $N$, and any of the
homotopy stages of Section~\ref{four:maximal:sec:deformation}.  Every smooth
solution on an allowed finite truncation with $t\ge-\eps$ has uniform
local bounds for every finite number of derivatives of $f$ and $Z$,
including up to the inner and outer boundary faces.  On the uniform
unit patches on the end these constants are independent of the
outer truncation radius and of the stage parameter.  They may depend
arbitrarily on $N$.  In particular, on a fixed finite truncation they
bound the $C^{1,\alpha}$ norm of $Z$ for each fixed $0<\alpha<1$.
\end{proposition}

\begin{proof}
Work on the uniform interior and boundary-normal patches of the
prepared geometry.

\paragraph{The axial estimate and the Hessian measure.}
Proposition~\ref{four:maximal:prop:bounded-range} bounds $f,Z,t,\grad f$ and
keeps $l/\sqrt D$ bounded above and away from zero at fixed $N$.
Dividing the trace equation by this factor gives
\begin{equation}\label{four:maximal:reg:f-normalized}
 \Achi:\Hess f=G_f,\qquad
 G_f=\frac{\sqrt D}{l}
             \bigl(F-\tr_{\Achi}K\bigr),\qquad
 |G_f|\le C(N),\quad \chi\ge\chi _0(N)>0.
\end{equation}
The geometry of each patch has the bounds required by
Theorem~\ref{four:thm:axial-regularity}, and each face value of $f$
is constant.  Each individual solution is smooth; the theorem's
constants depend only on the bounded gradient, source, geometry and
ellipticity.  It gives a uniform H\"older bound for $\grad f$ and
\begin{equation}\label{four:maximal:reg:f-hessian-measure}
 \int_{B_r(x)\cap\Omega_R}|\Hess f|^2\dd V_g
                              \le C(N)r^{2+2\alpha}.
\end{equation}
Coordinate Hessians satisfy the same estimate, since their difference
from covariant Hessians is bounded.  At a face the theorem uses the
$C^1$ odd graph double and controls its interface flux explicitly.
It has no support hypothesis on $K$; the prepared tensor tail and
geometry supply the uniform patch bounds used here.

\paragraph{The scalar source and its boundary reflection.}
In coordinates, absorb the volume density in the second equation and
write it as
\begin{equation}\label{four:maximal:reg:natural-equation}
 \operatorname{div}(\mathcal A DZ+B_1)=E.
\end{equation}
Here $\mathcal A$ and $B_1$ are the coordinate versions of $3u\Achi$
and $\lambda u\Achi K(w,\cdot)$, respectively, with the appropriate
homotopy factors in the source.  On the established range
$\mathcal A$ is bounded, symmetric and uniformly elliptic, and $B_1$
is bounded.  The implicit relation for $t$ has derivative
$1+pv/3>0$, so $t=t(x,Z,Df)$ is smooth on that range and
\begin{equation}\label{four:maximal:reg:t-chain}
 |Dt|\le C(N)(1+|DZ|+|D^2f|).
\end{equation}
The quadratic expression for $\cT$, together with the bounded
lower-order terms in each stage, therefore gives
\begin{equation}\label{four:maximal:reg:coupled-growth}
 |E|\le C(N)(1+|DZ|^2+|D^2f|^2).
\end{equation}
This estimate does not presume a bound for $DZ$.  The prescribed
inner conormal value concerns the total flux and is bounded on this
range.  The outer value of $Z$ is zero.

To apply Lemma~\ref{b:lem:natural-growth} at a boundary, flatten the
face to $x_4=0$, with domain $x_4>0$, and write
$Q=\mathcal A DZ+B_1$.  Put $R=\operatorname{diag}(1,1,1,-1)$.
For $x_4<0$ set
\[
 \widetilde Z(x)=sZ(Rx),\quad
 \widetilde{\mathcal A}(x)=R\mathcal A(Rx)R,\quad
 \widetilde B_1(x)=sRB_1(Rx),\qquad s\in\{1,-1\}.
\]
Then $\widetilde Q(x)=sRQ(Rx)$ below the plane.  Its distributional
divergence has reflected bulk source $sE(Rx)$ and plane source
\begin{equation}\label{four:maximal:reg:reflection-source}
 (1+s)Q_4^+\dd\mathcal H^3\big|_{\{x_4=0\}}.
\end{equation}
At the inner conormal face use $s=1$.  The added signed density is
bounded because $Q_4^+$ is the prescribed total conormal flux in
these coordinates.  At the outer zero Dirichlet face use $s=-1$.
The reflected $Z$ is continuous and the normal total flux matches,
so there is no plane term.  The coefficient matrix remains measurable,
symmetric and uniformly elliptic.  Each reflected solution is continuous
and smooth on its two closed sides.  Integration on those sides gives
the weak exponential product rule required by the shared lemma,
including the full measure in \eqref{four:maximal:reg:reflection-source}.

The absolute source on the doubled patch is consequently dominated by
\[
 C(N)(1+|DZ|^2)\dd x+\mu _0,\qquad
 \mu _0=C(N)|D^2f|^2\dd x
          +2|Q_4^+|\dd\mathcal H^3\big|_{\{x_4=0\}},
\]
with the bulk density reflected and the plane term omitted where absent.
Equation~\eqref{four:maximal:reg:f-hessian-measure} and the area of a
three-dimensional plane section give, at every local center,
\[
 \mu _0(B_r(x))\le C(N)r^{2+\beta},\qquad
                         \beta=\min(2\alpha,1)>0.
\]
On an interior patch only the bulk term is needed.  Thus all the
hypotheses of Lemma~\ref{b:lem:natural-growth}, with dimension four,
hold.  Its H\"older representative agrees with the existing continuous
$Z$, giving a uniform H\"older bound up to either face.

\paragraph{Schauder for the trace equation and the normal datum.}
The actual coefficients in \eqref{four:maximal:reg:f-normalized} are smooth
functions of $x,Z,Df$, including at $Df=0$, since
\[
 \Achi=\Id-\frac{3+p}{3+pv}\,w\otimes w.
\]
Its source is smooth in the same variables and in $f$.  The
H\"older bounds just proved therefore make them uniformly
$C^{0,\theta}$ for some $\theta>0$.  Interior and constant
Dirichlet Schauder estimates give
\begin{equation}\label{four:maximal:reg:f-schauder}
 \|f\|_{C^{2,\theta}(U')}\le C(N)
\end{equation}
on smaller patches; the hypotheses are precisely uniform
ellipticity, $C^{0,\theta}$ leading coefficients and source, smooth
faces and constant Dirichlet data
\cite[Sections~6.1--6.2]{GilbargTrudinger2001}.

Expand the divergence equation for $Z$ using
\eqref{four:maximal:reg:f-schauder}.  A derivative of a coefficient depending on
$x,Z,Df$ is a bounded term plus bounded multiples of $DZ,D^2f$.
Thus it takes the nondivergence form
\begin{equation}\label{four:maximal:reg:Z-nondivergence}
 a^{ij}(x)D_{ij}Z=R_Z,\qquad
 |R_Z|\le C(N)(1+|DZ|^2),
\end{equation}
with uniformly elliptic $C^{0,\theta}$ leading coefficients.  No
derivative of $Z$ of order two has been placed on its right side.

On an inner face, $f$ is constant and $Df$ is normal, so
$\Achi\nu=\chi\nu$.  Dividing the flux condition by $3u\chi>0$
gives
\begin{equation}\label{four:maximal:reg:normal-datum}
 \partial_\nu Z=b(x,Z,f,Df),
\end{equation}
where $b$ is a smooth bounded-range function, separately on each
smooth stage of the homotopy.  The formulas agree at stage junctions
and have uniform bounds.  At zero gradient the smooth formula for
$\Achi$ above gives the same conclusion.  A smooth collar
extension of the right hand side has derivative bounded by
\[
 C(N)(1+|DZ|+|Df|+|D^2f|)
                         \le C(N)(1+|DZ|).
\]
The trace theorem therefore bounds its
$W^{1-1/p,p}$ norm on the face by
$C(N)(1+\|Z\|_{W^{1,p}})$ on a larger collar patch.  In particular,
differentiating this extension costs $DZ$ and $D^2f$, and does not
cost $D^2Z$.

\paragraph{Absorbing the quadratic gradient term.}
Fix $p_1>4$.  The linear local $W^{2,p_1}$ estimates for
\eqref{four:maximal:reg:Z-nondivergence}, with zero Dirichlet data or
\eqref{four:maximal:reg:normal-datum}, give on nested patches
\begin{equation}\label{four:maximal:reg:Z-linear-estimate}
 \|Z\|_{W^{2,p_1}(U')}
 \le C(N)\left(1+\|DZ\|_{L^{2p_1}(U)}^2
                          +\|Z\|_{W^{1,p_1}(U)}\right).
\end{equation}
We specify why the coefficient regularity at this point suffices.
Freeze the continuous leading matrix on small patches.  The constant
coefficient interior Hessian estimate follows from the Laplacian
estimate by a linear coordinate change.  At a boundary, use normal
coordinates.  On that face the matrix has no mixed tangential-normal
entry, by the normal-axis property above.  After removing the
boundary datum by its Sobolev trace extension, odd Dirichlet or even
homogeneous-normal reflection gives the constant coefficient
estimate.  The term
$(a(x)-a(x_0)):D^2Z$ is absorbed when its oscillation is sufficiently
small.  Cutoff terms are first order.  This proves the local
estimate with the $W^{1-1/p_1,p_1}$ normal trace norm just estimated;
it is the scalar Dirichlet/oblique specialization of
\cite[Theorems~15.1--15.3, pp.~702--706, especially
Theorem~15.2, estimate~(15.5)]{AgmonDouglisNirenberg1959}, together with the interior $L^p$ estimate
\cite[Section~9.5]{GilbargTrudinger2001}.  Uniform H\"older control fixes the patch
size uniformly.  The complementing condition can also be checked
directly: a nonzero decaying Fourier mode for the frozen scalar
operator in a half-space has nonzero value and nonzero normal
derivative at its boundary.  Thus neither the Dirichlet nor the
normal operator annihilates that mode.  The normal field is smooth
and has unit transverse component; the boundary is smooth and the
leading coefficients are continuous, as required for this $W^{2,p_1}$
specialization.  The inner and outer faces are disjoint, so no
change of boundary type occurs within such a patch.

For a ball or smooth half-patch $Q_h$ of radius $h$, subtract a
constant from $Z$.  Nirenberg's interpolation inequality, with
derivative orders $1,2$, endpoint exponent $\infty$, Hessian
exponent $p_1$ and interpolation parameter $1/2$, gives
\begin{equation}\label{four:maximal:reg:nirenberg-scaled}
 \|DZ\|_{L^{2p_1}(Q_h)}^2
 \le C\operatorname{osc}_{Q_h}Z\,
                       \|D^2Z\|_{L^{p_1}(Q_h)}
       +Ch^{4/p_1-2}(\operatorname{osc}_{Q_h}Z)^2.
\end{equation}
This is the bounded-domain form of
\cite[Lecture~II, (2.2) and Remark~5,
pp.~125--126]{Nirenberg1959}; extension on half-patches preserves
the estimate.  The power of $h$ follows by rescaling from dimension
four.  Our choice $p_1>4$, $j=1$, $m=2$ avoids the exceptional
endpoint cases in that inequality.

The H\"older estimate for $Z$ makes its oscillation at most $Ch^\theta$.
Cover each larger patch in \eqref{four:maximal:reg:Z-linear-estimate} by such
$Q_h$ with bounded overlap.  Taking the $\ell^{p_1}$ norm of
\eqref{four:maximal:reg:nirenberg-scaled} over this cover gives
\[
 \|DZ\|_{L^{2p_1}(U)}^2
 \le C\sup_{Q_h}\operatorname{osc}_{Q_h}Z\,
                 \|D^2Z\|_{L^{p_1}(U^+)}+C(N,h),
\]
where $U^+$ is a fixed enlargement.  The number of small patches
may affect the harmless constant $C(N,h)$; their overlap, hence the
coefficient of the Hessian norm, is bounded independently of the
truncation.  Choose $h$ so that this coefficient, including the
constant in \eqref{four:maximal:reg:Z-linear-estimate} and the fixed enlargement
factor, is less than $1/4$.  Ordinary interpolation handles the
$W^{1,p_1}$ term with another coefficient less than $1/4$.

To make the absorption on enlarged patches precise, use a uniform
radius cover of the whole finite truncation and let $S$ be the
supremum of the local $W^{2,p_1}$ norms on its smaller patches.
Each fixed enlargement is covered by a bounded number of these
patches, independently of the total number of patches.  The last
two estimates give $S\le C(N)+\tfrac12S$.  For each individual
smooth solution on a finite truncation $S$ is finite before this
inequality is used.  Thus $S\le2C(N)$ uniformly.  The same argument
on nested compact patches gives the local estimate when no global
cover is needed.

\paragraph{Bootstrap and the continuation interface.}
Since $p_1>4$, the bound just proved gives
$Z\in C^{1,\theta'}$ for some $\theta'>0$.  The trace equation now
has $C^{1,\theta''}$ coefficients and source after
\eqref{four:maximal:reg:f-schauder}, so one differentiation and Dirichlet
Schauder give $f\in C^{3,\theta''}$, with
$0<\theta''\le\min(\theta,\theta')$.  Equation
\eqref{four:maximal:reg:Z-nondivergence} consequently has a
$C^{0,\theta''}$ source, and \eqref{four:maximal:reg:normal-datum} has a
$C^{1,\theta''}$ datum.  The scalar oblique/Dirichlet Schauder
estimate gives $Z\in C^{2,\theta''}$; the boundary estimate is
\cite[Theorem~7.3, estimate~(7.8), p.~668]{AgmonDouglisNirenberg1959}, with the
complementing condition verified above.
Repeating these two steps in this order raises the differentiability
by one each time, because the trace equation contains $Z$ without
its derivatives and the second equation contains at most $D^2f$
and $DZ$ in its lower order terms.  The normal datum has the same
derivative count.  Smooth prepared geometry and smooth data complete
the bootstrap to every finite order.

For the scalar solve with trial $Z\in C^{1,\alpha}$ used in
Section~\ref{four:maximal:sec:existence}, this argument makes no assertion that
$D^2Z$ exists.  Its preliminary individual third derivatives for
$f$ require only one differentiation of the trace equation, hence
$DZ$.  After the scalar gradient and axial estimates, bounded
$D^2f$ makes $Df$ Lipschitz; smooth compositions with the trial
$C^{1,\alpha}$ input then give $C^{0,\alpha}$ scalar coefficients,
Dirichlet $C^{2,\alpha}$ estimates, and one more differentiation
gives $C^{3,\alpha}$.  The independent scalar gradient bound and
solvability are proved there.  Finally, for actual coupled solutions
the uniform $C^2$ bound already controls every fixed
$C^{1,\alpha}$ norm, $\alpha<1$, regardless of the first H\"older
exponent obtained above.

The prepared end has uniform smooth geometry on unit patches.
Large outer coordinate spheres admit boundary normal patches with
the same bounds; their curvature and higher local chart derivatives
are bounded uniformly.  All the local constants and covering
overlap numbers in this proof can therefore be chosen independently
of the outer radius.  This proves the proposition.
\end{proof}

\section{Scalar solvability and compact continuation}\label{four:maximal:sec:existence}

Fix the strict data and $0<\eps<1$.  In this section $N\ge4$ and an
allowed finite truncation $\Omega_R$ are fixed.  The boundary components
$B$ and $S_R$ are smooth and disjoint.  Fix also $0<\alpha<1$, with no
requirement that $\alpha$ equal the exponent in
Theorem~\ref{four:thm:axial-regularity}.  We parametrize the four successive
stages of Section~\ref{four:maximal:def:homotopy} by $s\in[0,4]$: $s=0$ is the
desired system and $s=4$ is the final zero-scale system.  This parameter
is unrelated to a collar coordinate.  Constants in the scalar
construction below may depend on $N,\eps,R,\alpha$ and on a specified
bound for the trial input.  They need not be polynomial in $N$.

\subsection{The trace equation on unrestricted trial inputs}

\begin{proposition}[Scalar solution operator]\label{four:maximal:prop:scalar-solve}
For every $Z\in C^{1,\alpha}(\overline\Omega_R)$ and every homotopy
parameter $s$, the trace equation with its assigned constant Dirichlet
values has a unique solution $f=f_s[Z]\in
C^{3,\alpha}(\overline\Omega_R)$.  No lower bound on
$t=t(Z,\grad f)$ is assumed.  On bounded sets of trial inputs, these
solutions have bounded $C^{3,\alpha}$ norms, uniformly in $s$.  The
map $(s,Z)\mapsto f_s[Z]$ is continuous into $C^{3,\alpha}$ and is
locally obtained by inversion of the scalar Dirichlet operator.
\end{proposition}

\begin{proof}
Write the trace prescription at parameter $s$ as
\[
 F_s=\eta_N f+F_s^0(x,Z,\grad f),\qquad
 Q=\frac{\sqrt D}{l},\qquad
 G_s=Q\bigl(F_s-\tr_{\Achi}K\bigr).
\]
Here $F_s^0$ is smooth in its displayed finite-dimensional arguments.
For example, in stage three it is
$(1-\zeta)(C+D_0)+\zeta(\tr_{\Achi}K+D_1)$.
The functions $D_0,D_1$ are evaluated on $w$, whose length is less than
one.  Thus $F_s^0-\tr_{\Achi}K$ is bounded on the finite truncation,
uniformly in the gradient and in $s$.  The normalized scalar equation is
\begin{equation}\label{four:maximal:exist:normalized-trace}
 \Achi:\Hess f=G_s(x,Z,f,\grad f).
\end{equation}
Its matrix is independent of the value of $f$, and
$\partial_fG_s=\eta_N Q>0$.

For the scalar method of continuity introduce a further parameter
$\theta\in[0,1]$, retaining the assigned Dirichlet values and using
\begin{align}
 A^{(\theta)}&=(1-\theta)\Id+\theta\Achi
       =\Id-(1-\chi_\theta)e_f\otimes e_f,
       &\chi_\theta&=1-\theta+\theta\chi,\label{four:maximal:exist:scalar-interpolation}\\
 G_s^{(\theta)}&=(1-\theta)\eta_N f+\theta G_s,
       &A^{(\theta)}:\Hess f&=G_s^{(\theta)}.
       \label{four:maximal:exist:interpolated-equation}
\end{align}
At zero gradient these expressions are interpreted by their smooth
tensor formulas.  In particular,
\begin{equation}\label{four:maximal:exist:strict-height}
 \partial_fG_s^{(\theta)}
       =\eta_N\bigl(1-\theta+\theta Q\bigr)>0.
\end{equation}
Every trace stage has coefficient one on $h=\eta_N f$, so this strict
sign holds throughout both parameters.

\paragraph{Height and large-gradient behavior.}
At a critical point of $f$, $w=0$, $\Achi=\Id$, $t=Z$,
$D_0(x,0)=D_1(x,0)=0$, and $\tr K=0$.
The signs at extrema used in Lemma~\ref{four:maximal:lem:height-end} therefore
persist in \eqref{four:maximal:exist:interpolated-equation}: a fixed interval
containing the assigned boundary values and the zeros of its scalar
right side contains every solution.  Denote its resulting bound by
\begin{equation}\label{four:maximal:exist:scalar-height}
 |f|\le M_f.
\end{equation}
This argument uses neither the divergence equation nor the floor.

Suppose $\|Z\|_{C^0}\le M_Z$ and put $\sigma=|\grad f|$.
The implicit relation $Z=t+\frac16\log(1+l(t)^2\sigma^2)$ implies
$t\to-\infty$ uniformly as $\sigma\to\infty$: if $t$ stayed bounded
below, its right side would tend to infinity.  Eventually $p=N$ and
$l=e^{Nt}$, whence exactly
\begin{equation}\label{four:maximal:exist:large-gradient-identity}
 e^{6Z}=e^{6t}+e^{(6+2N)t}\sigma^2.
\end{equation}
Since the first term tends uniformly to zero and $e^{6Z}$ stays between
two positive constants, this gives
\begin{equation}\label{four:maximal:exist:large-gradient-asymptotics}
 \begin{split}
 t&=-\frac{\log\sigma}{N+3}+O(1),\\
 d=e^{6(t-Z)}&\asymp\sigma^{-6/(N+3)},\qquad
 \chi=\frac{3d}{3+N(1-d)}\asymp d,\\
 Q=\sqrt{\sigma^2+l^{-2}}&\asymp\sigma.
 \end{split}
\end{equation}
The constants here are uniform for bounded $Z$ at the fixed parameters.
On bounded gradient sets the implicit equation has bounded $t$ and
positive $l,d,\chi$.  Since $N\ge4$ makes $6/(N+3)<1$, combining the
two ranges yields constants $c_*>0$ and $C_*>0$ such that
\begin{equation}\label{four:maximal:exist:linear-growth}
 \chi_\theta\ge\chi\ge\frac{c_*}{1+\sigma},\qquad
 |G_s^{(\theta)}|\le C_*(1+\sigma)
       \quad\text{when } |f|\le M_f.
\end{equation}
No uniform ellipticity has yet been asserted.

\paragraph{Logarithmic collar barriers.}
Extend the metric smoothly across all faces to a compact enlargement.
Choose a short distance scale below its injectivity radius, the tubular
radii of the faces, and their positive mutual separation.  The geometry
of all these collars is bounded.  For the inward distance $r$ from a
face with value $b$, consider
\begin{equation}\label{four:maximal:exist:log-barrier}
\begin{gathered}
 b\pm\psi(r),\qquad
 \psi(r)=\frac1k\log(1+B_0r),\\
 q_1(r)=\psi'(r)=\frac{B_0}{k(1+B_0r)},\qquad
 \psi''=-kq_1^2.
\end{gathered}
\end{equation}
Here $k$ and $B_0$ are constants, not any of the geometric tensors.
At a comparison contact the gradient is $\pm q_1\grad r$.
The distance Hessian vanishes in that axial direction, and the three
transverse eigenvalues of $A^{(\theta)}$ are one.  Its contraction on
the positive test is consequently
\[
 \chi_\theta\psi''+
      q_1\sum_{a=1}^3\Hess r(E_a,E_a),
\]
where $E_1,E_2,E_3$ are orthonormal tangent directions to the distance
leaf.  If $q_1\ge1$, \eqref{four:maximal:exist:linear-growth} bounds this above by
$-c_*kq_1/2+Cq_1$; the negative test has the reverse bound.
Choose $k$ so large that this axial term dominates both the distance
Hessian bound and $C_*(1+q_1)$, and also the pair-comparison constants
below.  Then choose a width $\delta_2<1/(4k)$ below all the preceding
geometric scales.  Finally choose $B_0$ large enough that
\begin{equation}\label{four:maximal:exist:barrier-choices}
 q_1(\delta_2)\ge2,\qquad \psi(\delta_2)>2M_f+1.
\end{equation}
These choices are compatible: as $B_0\to\infty$ the two quantities
tend respectively to $1/(k\delta_2)>4$ and to infinity.

At a positive interior maximum of $f-b-\psi(r)$ the solution and test
have the same gradient and the solution Hessian is no larger.  Evaluating
the source at the \emph{solution} height, which obeys
\eqref{four:maximal:exist:scalar-height}, the preceding strict upper bound
contradicts \eqref{four:maximal:exist:interpolated-equation}.
At the outer end of the collar the height inequality follows from
\eqref{four:maximal:exist:barrier-choices}, and at its inner end there is equality.
The negative test is handled with the reversed inequalities.  Thus
\begin{equation}\label{four:maximal:exist:collar-comparison}
 |f(x)-b|\le\psi\bigl(\dist(x,\text{face})\bigr)
       \qquad (\dist(x,\text{face})\le\delta_2).
\end{equation}
The tests themselves need not stay in the bounded height interval;
the source estimate was used only at solution heights.

\paragraph{The full two-point comparison.}
The face barriers control graph slopes only at the boundary.  An interior
two-point comparison supplies the missing global gradient bound.  Its
key feature is that the three transverse eigenvalues agree at both
points; the two axial eigenvalues need only have the same lower bound.

Let $r(x,y)$ be distance in the smooth extension.  On the compact set
$\{(x,y)\in\overline\Omega_R^2:r(x,y)\le\delta_2\}$ consider
\[
 f(x)-f(y)-\psi(r(x,y)).
\]
It is zero on the diagonal and negative when $r=\delta_2$.
It cannot have a positive maximum with a boundary endpoint: if, for
example, $x$ belongs to a face of height $b$, then $y$ lies in that
same face's collar, its distance to the face is at most $r(x,y)$,
and \eqref{four:maximal:exist:collar-comparison} and monotonicity of $\psi$ give
$b-f(y)\le\psi(r(x,y))$.  The case with $y$ on a face is identical.
The separation choice excludes different faces at such distances.

Suppose then that a positive maximum occurs at an interior pair with
$0<r<\delta_2$.  The short geodesic $\gamma$ from $y$ to $x$ is unique
in the extension; the distance is smooth at this pair.  Write
$e_y=\dot\gamma(0)$ and $e_x=\dot\gamma(r)$, with unit speed.  First
variation gives
\begin{equation}\label{four:maximal:exist:pair-gradients}
 \grad f(x)=q_1e_x,\qquad \grad f(y)=q_1e_y,
       \qquad q_1=\psi'(r)\ge1.
\end{equation}
Parallel translate an orthonormal transverse frame $E_1,E_2,E_3$
along $\gamma$.  Vary both endpoints simultaneously in the corresponding
$E_a$ directions, using geodesic endpoint curves.  The first variation
of the distance is zero.  The second variation is bounded above by
$Cr$: the parallel comparison field has index form
\[
 -\int_0^r
 \bigl\langle\operatorname{Rm}(E_a,\dot\gamma)\dot\gamma,E_a
 \bigr\rangle\dd v,
\]
whose absolute value is at most $Cr$, and the Jacobi field with the
same endpoints minimizes this index form on the short segment.
The second derivative test at the pair maximum therefore gives
\begin{equation}\label{four:maximal:exist:pair-transverse}
 \sum_{a=1}^3\bigl[
  \Hess f_x(E_a(r),E_a(r))-
  \Hess f_y(E_a(0),E_a(0))\bigr]\le Cq_1r.
\end{equation}
Varying only $x$ in direction $e_x$, and only $y$ in direction $e_y$,
the distance changes linearly and its second derivative vanishes.
Hence
\[
 \Hess f_x(e_x,e_x)\le\psi'',\qquad
 \Hess f_y(e_y,e_y)\ge-\psi''.
\]
The two axial coefficients need not be equal.  Positivity and
\eqref{four:maximal:exist:linear-growth} nevertheless imply
\begin{equation}\label{four:maximal:exist:pair-axial}
 \begin{split}
 &\chi_\theta(x)\Hess f_x(e_x,e_x)
        -\chi_\theta(y)\Hess f_y(e_y,e_y)\\
 &\hspace{12mm}\le
    \bigl(\chi_\theta(x)+\chi_\theta(y)\bigr)\psi''
    \le-\frac{2c_*kq_1^2}{1+q_1}\le-c_*kq_1.
 \end{split}
\end{equation}
There are exactly three transverse eigenvalues at \emph{each} endpoint,
all equal to one.  Thus \eqref{four:maximal:exist:pair-transverse} and
\eqref{four:maximal:exist:pair-axial}, followed by the two scalar equations, give
\[
 G_s^{(\theta)}(x)-G_s^{(\theta)}(y)
       \le Cq_1r-c_*kq_1.
\]
On the other hand, the source growth bound and
\eqref{four:maximal:exist:pair-gradients} give a lower bound
$-2C_*(1+q_1)\ge-4C_*q_1$.  Taking $k$ above so that
$c_*k>C+4C_*$ and $\delta_2<1$ is a contradiction.
This uses no continuity estimate on the two axial coefficients or on
their difference.  The connecting segment is permitted to leave
$\Omega_R$: only endpoint variations and the extension's bounded
geometry enter the argument, and the endpoints are interior.

It follows that $|f(x)-f(y)|\le\psi(r(x,y))$ for short pairs.  Letting
$y\to x$, and also using the collar comparison at a face, proves
\begin{equation}\label{four:maximal:exist:scalar-gradient}
 \|\grad f\|_{C^0(\overline\Omega_R)}\le\psi'(0)=B_0/k.
\end{equation}
This estimate is uniform in $s,\theta$ and on bounded trial-input sets,
and was obtained without assuming the floor or a positive uniform
ellipticity constant in advance.

\paragraph{Regularity and the scalar method of continuity.}
The gradient bound and the implicit relation now confine all scalar
coefficients to a compact range.  In particular
$0<\chi_0\le\chi_\theta\le1$, and the normalized source is bounded.
There is a small regularity point before applying
Theorem~\ref{four:thm:axial-regularity}.  An individual
$C^{2,\alpha}$ solution of \eqref{four:maximal:exist:interpolated-equation} is
$C^{3,\alpha}$ for a trial $Z\in C^{1,\alpha}$.
Indeed one differentiation of the equation gives a linear equation
for a first derivative of $f$.  Its coefficients and source involve
$Z,\grad Z,\grad f,\Hess f$ and smooth geometric coefficients;
they involve no second derivative of $Z$.
Difference quotients justify this differentiation initially.
In a flattened boundary chart tangential differentiation preserves
the zero Dirichlet value after subtracting the constant face height.
The linear boundary Schauder estimate controls these tangential
derivatives through order two.  Solving the original equation for
the pure normal second derivative, whose coefficient is positive,
and differentiating that identity in the normal direction recovers
the remaining pure normal third derivative.  Every other third
derivative already has a tangential index.  This gives individual
third-order regularity, without presuming a uniform estimate.

Theorem~\ref{four:thm:axial-regularity} now applies to
\eqref{four:maximal:exist:interpolated-equation}, including its constant-Dirichlet
faces.  It bounds $\grad f$ in $C^{0,\beta}$ for some $\beta>0$.
Since the trial inputs are bounded in $C^{1,\alpha}$, the coefficients
and source of the scalar equation are consequently bounded in
$C^{0,\beta'}$ for a positive $\beta'$.  Dirichlet Schauder estimates
give a uniform $C^{2,\beta'}$ bound for $f$.
In particular $\grad f$ is now uniformly Lipschitz.  Smooth composition
with $Z$, which also has bounded first derivatives, shows that those
same coefficients and sources have bounded $C^{0,\alpha}$ norms for
the originally chosen, arbitrary $\alpha<1$.
Schauder estimates upgrade $f$ to $C^{2,\alpha}$.
Their coefficients and sources then have bounded $C^{1,\alpha}$
norms, and the preceding differentiation argument gives
\begin{equation}\label{four:maximal:exist:scalar-schauder}
 \|f\|_{C^{3,\alpha}(\overline\Omega_R)}\le C.
\end{equation}
The linear estimates used here are the interior and Dirichlet boundary
Schauder estimates and their differentiated versions
\cite[Sections~6.1--6.4]{GilbargTrudinger2001}.  Their hypotheses hold on the finite
smooth domain: positive uniform ellipticity on the established range,
the indicated H\"older coefficient bounds, and smooth constant
Dirichlet data on each disjoint component.

For completeness, the local inverse has the required sign.
Linearizing $A^{(\theta)}:\Hess f-G_s^{(\theta)}$ in the $f$
variable gives
\begin{equation}\label{four:maximal:exist:scalar-linearization}
 \mathcal L\varphi
    = A^{(\theta)ij}\nabla_i\nabla_j\varphi
           +b^i\nabla_i\varphi
           -\eta_N(1-\theta+\theta Q)\varphi,
      \qquad \varphi|_{\partial\Omega_R}=0.
\end{equation}
Its first-order coefficients are bounded H\"older functions along each
solution.  Differentiating the gradient dependence of $A^{(\theta)}$
contributes terms involving $\Hess f$ only to this first-order part.
The strictly negative zeroth-order coefficient and the maximum
principle give a zero kernel.  The linear Dirichlet solvability theorem
and Schauder estimate give an isomorphism
$C^{2,\alpha}_0\to C^{0,\alpha}$
\cite[Sections~6.2--6.3]{GilbargTrudinger2001}; here the subscript means zero
trace on \emph{all} boundary components.  Equivalently, one can
continue the linear operator to $\Delta-\eta_N$, using its Schauder
estimate and the maximum principle at each point of that linear
continuation.  With the additional coefficient regularity it is also
an isomorphism $C^{3,\alpha}_0\to C^{1,\alpha}$.

The endpoint $\theta=0$ is the solvable linear Dirichlet problem
$\Delta f=\eta_N f$.  Subtract a smooth extension of the assigned face
values when working on spaces with homogeneous boundary data.
The implicit function theorem and \eqref{four:maximal:exist:scalar-linearization}
give openness in $\theta$.  The bound
\eqref{four:maximal:exist:scalar-schauder} gives compactness in $C^{2,\alpha}$,
so passage to a limiting equation gives closedness.  Thus the solution
exists at $\theta=1$.

If $f_1-f_2$ had a positive maximum in the interior, their gradients
and therefore their principal matrices and $Q$ would agree there.
Their Hessian difference would have nonpositive contraction, whereas
the difference of the sources would be
$\eta_N Q(f_1-f_2)>0$.  This contradiction, and its reverse, prove
uniqueness.  Finally use \eqref{four:maximal:exist:scalar-linearization} in
$C^{3,\alpha}_0\to C^{1,\alpha}$ with $Z$ as a
$C^{1,\alpha}$ parameter.  Smooth composition and the implicit
function theorem show that the unique scalar solution depends
continuously on that input in $C^{3,\alpha}$.  The same statement
holds along each homotopy stage; the explicit prescriptions and
boundary heights agree at their junctions, so it holds across those
junctions as well.
\end{proof}

\subsection{The frozen linear return problem}

The trace equation is now solvable for every trial $Z$, including trials
below the intended floor.  We can therefore define a compact return map
on a fixed function space.  The floor becomes a restriction on the
degree domain, rather than an unproved prerequisite for defining the map.

Set
\[
 X=\{Z\in C^{1,\alpha}(\overline\Omega_R):Z|_{S_R}=0\}.
\]
Given $(s,Z)\in[0,4]\times X$, first compute $f=f_s[Z]$ by
Proposition~\ref{four:maximal:prop:scalar-solve} and then all of $t,l,u,w,\chi$,
$F$ and $\cT$ from this pair.  Let $\lambda_s$ denote the drift
parameter, let $q_s=1$ in the first three stages and $q_s=\eta$ in
the last, and let $b_s$ denote the prescribed right side of
$(V_{\lambda_s})_\nu/u$ in the full homotopy, including its last-stage
scale.  Define $T_sZ=\widetilde Z$ by
\begin{equation}\label{four:maximal:exist:return-problem}
 \begin{cases}
 \divg\bigl(3u\Achi\grad\widetilde Z
                  +u\Achi\lambda_sK(w,\cdot)\bigr)=q_su\Xi
                         &\text{in }\Omega_R,\\
 \bigl(3u\Achi\grad\widetilde Z
                  +u\Achi\lambda_sK(w,\cdot)\bigr)_\nu=ub_s
                         &\text{on }B,\\
 \widetilde Z=0          &\text{on }S_R.
 \end{cases}
\end{equation}
Every quantity in this display except $\grad\widetilde Z$ is frozen
at the input pair, including every occurrence of $\grad Z$ in $\Xi$.

The coefficient $\mathcal A=3u\Achi$ is symmetric and positive
definite.  On a bounded subset of $X$ its ellipticity constants are
uniformly positive, by Proposition~\ref{four:maximal:prop:scalar-solve} and the
smooth implicit equation for $t$.  Let
$\mathcal D=u\Achi\lambda_sK(w,\cdot)$.
On the Hilbert space of $H^1$ functions vanishing on $S_R$, the weak
form of \eqref{four:maximal:exist:return-problem} is
\begin{equation}\label{four:maximal:exist:return-weak}
 \int_{\Omega_R}\mathcal A\grad\widetilde Z\cdot\grad\varphi
 =-\int_{\Omega_R}q_su\Xi\varphi
  -\int_{\Omega_R}\mathcal D\cdot\grad\varphi
  -\int_Bub_s\varphi.
\end{equation}
The last sign uses the inward normal $\nu$ on $B$.
Poincar\'e's inequality holds because $S_R$ is nonempty and the
domain is connected.  Consequently the left side is a coercive
bilinear form and the right side is a bounded linear functional,
by the trace inequality.  To solve it directly, minimize one half
of this quadratic form minus the right-side functional on that
Hilbert space.  Coercivity bounds a minimizing sequence in $H^1$;
weak compactness and lower semicontinuity give a minimizer.  Its first
variation is \eqref{four:maximal:exist:return-weak}, and strict convexity gives
uniqueness.  There is no Neumann compatibility constraint.

We check the precise regularity of this linear problem.  The formula
\[
 \Achi=\Id-\frac{3+p}{3+pv}\,w\otimes w
\]
is smooth also at zero gradient, and the derivative of the implicit
relation in $t$ is $1+pv/3>0$.
Thus $t,u,\Achi,\mathcal D$ have bounded $C^{1,\alpha}$ norms on
bounded input sets.  Differentiating $t=t(x,Z,\grad f)$ requires only
$\grad Z$ and $\Hess f$.  The expression for $\cT$ is quadratic in
$\Hess f$ and $\grad t$ with smooth bounded-range coefficients,
so $q_su\Xi$ has a bounded $C^{0,\alpha}$ norm.  The prescribed
boundary expression $ub_s$ has a bounded $C^{1,\alpha}$ norm.
At $B$ the scalar Dirichlet data make $\grad f$ normal, hence
$\Achi\nu=\chi\nu$; the boundary equation is equivalently
\begin{equation}\label{four:maximal:exist:return-normal}
 \partial_\nu\widetilde Z
  =\frac{b_s-\lambda_s(\Achi K(w,\cdot))_\nu}{3\chi}.
\end{equation}
This is a $C^{1,\alpha}$ normal derivative datum, with a denominator
bounded away from zero on the input set.
After expanding the divergence, interior and boundary Schauder
estimates therefore give a bounded $C^{2,\alpha}$ norm for
$\widetilde Z$.  The estimates apply locally with Dirichlet data on
$S_R$ and a uniformly oblique, here normal, derivative on $B$;
these faces do not meet.  We use the Dirichlet estimates just cited
and the oblique derivative estimates of
\cite[Theorem~6.30, p.~127]{GilbargTrudinger2001}; the weak existence argument above fixes
the additive constant that is present in a pure Neumann problem.

Subtracting two return equations and applying the same linear
estimates proves continuity in $(s,Z)$ into $C^{2,\alpha}$.
The parameter continuity is uniform on bounded input sets.  Indeed
the negative zeroth-order bound in \eqref{four:maximal:exist:scalar-linearization}
and the Schauder estimate bound the scalar inverse uniformly on
such a set.  Differentiation in a stage parameter, after extending
its Dirichlet data, then bounds the parameter derivative of $f_s[Z]$;
the return equation gives the same conclusion for $T_sZ$.
The finitely many stage junctions preserve continuity.
The inclusion $C^{2,\alpha}\hookrightarrow C^{1,\alpha}$ on the
fixed smooth compact domain is compact.  We have thus constructed
a continuous homotopy $T_s:X\to X$ that is compact on bounded sets,
jointly in $s$, before making any floor restriction.
A fixed point satisfies both original equations and their exact
boundary conditions.  It is smooth: the first gain is
$Z\in C^{2,\alpha}$, and successive differentiation of the scalar
equation, the divergence equation and its normal boundary condition
then alternately improves $f$ and $Z$ to every finite order.

\subsection{Degree on the open floor sections}

\begin{proposition}[Existence and exhaustion]\label{four:maximal:prop:existence}
For the fixed strict data and every $0<\eps<1$ there are $N_0$ and
$R_{\min}(N)$, with $R_{\min}(N)\to\infty$, such that for every
integer $N\ge N_0$ and every $R\ge R_{\min}(N)$ the desired coupled
system has a smooth solution on $\overline\Omega_R$ with
$t>-\eps$.  At each such fixed $N$, a subsequence as $R\to\infty$
converges smoothly on compact subsets of $\overline\Omega$ to a
solution of the equations and inner boundary conditions with
$t\ge-\eps$.  The previously proved height, range and collar
estimates hold for this limit.
\end{proposition}

\begin{proof}
All estimates from Sections~\ref{four:maximal:sec:deformation}--\ref{four:maximal:sec:regularity}
were conditional estimates for arbitrary smooth homotopy solutions.
Specifically, Lemma~\ref{four:maximal:lem:height-end} and Proposition~\ref{four:maximal:prop:floor}
precede the bounded-range and derivative estimates;
Lemma~\ref{four:maximal:lem:separation} and Lemma~\ref{four:maximal:lem:collar-slopes} apply to
solutions satisfying the floor;
Proposition~\ref{four:maximal:prop:bounded-range} then bounds their variables;
and Proposition~\ref{four:maximal:prop:classical-estimates} bounds their
derivatives.  None presupposes a fixed point of the map just defined.
On this finite truncation these results give, uniformly in $s$,
\[
 \|Z\|_{C^{1,\alpha}}\le M_*
 \quad\text{for every smooth fixed point with }\min t\ge-\eps.
\]
The arbitrary exponent $\alpha<1$ is permitted because the classical
estimates give a $C^2$ bound, and higher bounds as needed.  Choose
$M_0>M_*+1$.

Let
\begin{equation}\label{four:maximal:exist:floor-sections}
 \begin{split}
 m(s,Z)&=\min_{\overline\Omega_R}t\bigl(x,Z,\grad f_s[Z]\bigr),\\
 \mathcal W&=\{(s,Z)\in[0,4]\times X:
            m(s,Z)>-\eps, \|Z\|_{C^{1,\alpha}}<M_0\},\\
 \mathcal W_s&=\{Z:(s,Z)\in\mathcal W\}.
 \end{split}
\end{equation}
The minimum is continuous by Proposition~\ref{four:maximal:prop:scalar-solve},
so $\mathcal W$ is relatively open and every section is a bounded
open subset of $X$.  Consider
\[
 \mathcal K=
 \{(s,Z)\in\overline{\mathcal W}:Z=T_sZ\}.
\]
This is compact: any sequence has a convergent parameter subsequence,
and joint compactness of $T_s$ on the bounded input ball makes its
fixed points converge in $X$; continuity preserves the fixed-point
equation and membership in the closed set.
Every point of $\mathcal K$ is a smooth solution with
$m(s,Z)\ge-\eps$.  Its norm is at most $M_*<M_0$, and
Proposition~\ref{four:maximal:prop:floor} excludes $m(s,Z)=-\eps$.
Therefore
\begin{equation}\label{four:maximal:exist:no-boundary-fixed-point}
 \mathcal K\subset\mathcal W.
\end{equation}

Apply Lemma~\ref{b:lem:moving-degree} with $I=[0,4]$, the Banach
space $X$ above, the jointly compact homotopy $(s,Z)\mapsto T_sZ$, and
$\mathcal U=\mathcal W$.  The uniform norm bound makes $\mathcal W$
bounded; continuity of $m$ makes it relatively open.  The preceding
compactness argument and \eqref{four:maximal:exist:no-boundary-fixed-point}
verify that its closed fixed-point set is compact and contained in
$\mathcal W$, so no fixed point lies on the relative boundary.
The lemma gives constancy of
\[
 \deg(\Id-T_s,\mathcal W_s,0)\qquad (0\le s\le4).
\]
It applies to these varying sections without a convexity assumption.

At $s=4$ the scalar boundary values are zero and the final trace
equation is
\[
 \tr_{\Achi}H^f=\eta_N f+D_1(x,w).
\]
At a positive or negative extremum $w=0$ and $D_1(x,0)=0$, so
the strict height sign forces $f=0$.  This holds for every trial
$Z$, not just at fixed points.  Thus $t=Z$, $\chi=1$ and $u=L_0(Z)$.
Both the drift and the prescribed source and boundary flux in
\eqref{four:maximal:exist:return-problem} vanish.  Testing its homogeneous
equation by $\widetilde Z$ gives
\[
 \int_{\Omega_R}3L_0(Z)|\grad\widetilde Z|^2=0,
\]
so $T_4Z=0$ for every input.  At its sole fixed point $Z=0$ one
has $t=0$, so $0\in\mathcal W_4$ and
$\deg(\Id-T_4,\mathcal W_4,0)=1$
\cite[Chapter~II, \S8, p.~55]{LeraySchauder1934}.
The same degree at $s=0$ yields a desired fixed point in
$\mathcal W_0$, proving the assertion on the finite truncation.

We finally make the parameter order explicit.  Fix the strict
background, $\eps$, the collars and the end weights first.
The pointwise floor argument chooses the small $\delta_0$ and a
polynomially large penalty coefficient $C_N$ without derivative
estimates.  The separation statement is uniform along arbitrary
expanding sequences of the first two homotopy stages; applied at
the finitely many fixed collar ends, it supplies the eventual
separation needed for the slopes.  Increase $N_0$ to meet this
requirement and the finitely many conditions of the form
$\Pi_N\eta_N/\ell<\text{constant}$.  Increase $R_{\min}(N)$ to
meet all the fixed-$N$ outer-radius requirements, requiring also
$R_{\min}(N)\ge N$.  These choices give the estimates for every
$N\ge N_0$ and every $R\ge R_{\min}(N)$.
All subsequent scalar, degree-radius and classical regularity
constants may depend arbitrarily on the now fixed $N$.

For that fixed $N$, take $R_j\to\infty$ and the desired fixed
points just obtained.  Proposition~\ref{four:maximal:prop:classical-estimates}
gives uniform bounds of every finite order on each fixed compact
subset, including patches meeting $B$.
Successive compact extractions on a countable exhaustion produce
a subsequence converging in $C^k$ on each such compact set for
every finite $k$.  All equations, inner data and collar inequalities
pass to the limit.  The strict floor may become non-strict, which
is the asserted $t\ge-\eps$.
The normalization at infinity is supplied by the uniform tail
estimates in Proposition~\ref{four:maximal:prop:end-control}; it is not inferred
from compact convergence alone.  Only after this fixed-$N$
exhaustion will $N$ tend to infinity in the flux comparison.
The latter uses the earlier polynomial estimates directly, and
does not use the arbitrary constants of the present section.
\end{proof}

\section{End normalization and the mass--area comparison}
\label{four:maximal:sec:comparison}

Fix the strict approximation data, the threshold exterior $\Omega$ with
boundary $B$, and $0<\eps<1$.  Take $N$ sufficiently large as in
Proposition~\ref{four:maximal:prop:existence}.  In the first part of this section, $N$
remains fixed while the truncation radius tends to infinity.  Constants
in end estimates may depend arbitrarily on these fixed parameters.  In
the flux deficit estimate, by contrast, we will return directly to the
polynomial bounds of Lemma~\ref{four:maximal:lem:weighted-trace}.

\subsection{Normalization inherited from the truncations}

\begin{proposition}[Control of the asymptotic end]\label{four:maximal:prop:end-control}
The exhaustion solutions of Proposition~\ref{four:maximal:prop:existence} admit a
smooth local limit $(f,Z)$ on $\overline\Omega$ satisfying the desired
system and $t\ge-\eps$.  Every such limit has
\begin{equation}\label{four:maximal:end:decay}
 |\nabla^j f|\le C_j(N)e^{-c_j(N)r}\quad(j\ge0),
 \qquad Z,t=O_2(r^{-2}).
\end{equation}
The zeroth-order estimates producing this normalization are uniform on
the finite truncations.  The resulting metric
$\ghat=e^{2t}(g+l^2df\otimes df)$ is smooth, complete with its boundary
included, and asymptotically flat with
\begin{equation}\label{four:maximal:end:geometry}
 \ghat_{ij}-\delta_{ij}=O_2(r^{-q}),\qquad
 0\le R_{\ghat}=O(r^{-4-\delta}),\qquad \widehat H_B<0.
\end{equation}
Here the normal defining $\widehat H_B$ points toward infinity.
\end{proposition}

\begin{proof}
The local compactness and the inner boundary conditions follow from
Propositions~\ref{four:maximal:prop:classical-estimates} and
\ref{four:maximal:prop:existence}.  To determine the value at infinity,
apply Lemma~\ref{b:transition:end} to the finite truncations before
taking their local limit.  We verify its hypotheses explicitly.

The strict background has $1<q<2$, $0<\delta<q-1$,
$g-\delta_{\mathrm{Eucl}}=O_2(r^{-q})$,
$K=O_1(r^{-1-q})$, $\tr_gK=0$, and
$R_g=O(r^{-4-\delta})$ by
Proposition~\ref{four:maximal:prop:strict-data}.  Outside the compact
support of $C$, the trace right side is $\eta_N f$.
The divergence equation is exactly
\eqref{b:transition:divergence}, with
$\rho_0=r^{-4-\delta}$, $\rho_1=r^{-2\gamma}$ and the smooth penalty
switch defined in \eqref{four:maximal:def:penalty}.  Their differentiated
end bounds follow from the fixed smooth radial choices.
Propositions~\ref{four:maximal:prop:bounded-range} and
\ref{four:maximal:prop:classical-estimates} supply the common bounded
range and all fixed-$N$ classical estimates, on end unit patches and on
the zero-Dirichlet outer faces, uniformly in the truncation radius.
The leading matrix is uniformly elliptic on that range.

The lemma therefore gives exponential decay of $f$ and its fixed
derivatives, a truncation-uniform bound $|Z|+|t|\le C(N)r^{-2}$, and
$Z,t=O_2(r^{-2})$ for every local limit.  In particular normalization
at infinity follows from a comparison on the actual truncations, not
from local convergence.  The equation used for that comparison is
\begin{align}\label{four:maximal:end:zero-graph}
 3\Delta_g Z
 &+3\bigl[p(Z)+2-\delta_0(p(Z)+1)\bigr]|dZ|^2\notag\\
 &=\tfrac{\delta_0}{2}|K|^2+\rho
       -m_0(Z)C_N\rho_0+\mathcal E_N,
\end{align}
where the graph error $\mathcal E_N$ has exponential unit-patch bounds.
The $|K|^2/2$ term survives because the tensor is retained; its decay is
integrable since $2+2q>4+\delta$.  All constants in this application
are allowed to depend on the fixed $N$.  They will not be used to
estimate the later limit in $N$.

The signs can now be recorded directly from the desired system.
Since $F=\eta_N f+C$ and $w(f)=a\sigma$, the scalar identity becomes
\begin{align}\label{four:maximal:end:scalar-sign}
 \tfrac12e^{2t}R_{\ghat}
 ={}&[\mu+J(w)+w(C)-\rho]
       +(1-\delta_0)\cT\notag\\
    &+(1-\delta_0)\eta_N a\sigma+m_0(t)\mathcal P.
\end{align}
The bracket is nonnegative by $|w|\le1$ and
$\mu-|J|\ge|dC|+\rho$; all remaining terms are nonnegative.
The boundary identity and the prescribed conormal datum give
\begin{equation}\label{four:maximal:end:boundary-sign}
 H_B+3\partial_\nu t=-N_B,
 \qquad \widehat H_B=-e^{-t}\sqrt d\,N_B<0.
\end{equation}
This uses the normal into $\Omega$ on every component of $B$.

As quadratic forms $\ghat\ge e^{-2\eps}g$.  A $\ghat$-Cauchy sequence
is therefore $g$-Cauchy; completeness of the background exterior and
local equivalence of the smooth metrics give completeness with the
boundary included.  Since $q<2$, the conformal $O_2(r^{-2})$ term and
the exponential graph error retain the background decay
$\ghat-\delta=O_2(r^{-q})$.  Equation~\eqref{four:maximal:end:zero-graph} gives
$\Delta_g t=O_N(r^{-4-\delta})$, since $|dt|^2=O_N(r^{-6})$ and
$0<\delta<1$.  Finally,
\[
 R_{\ghat}
  =e^{-2t}\bigl(R_{\gbar}-6\Delta_{\gbar}t
                                  -6|dt|_{\gbar}^2\bigr)
\]
and the exponential graph errors, together with
$R_g=O(r^{-4-\delta})$, give the scalar decay in
\eqref{four:maximal:end:geometry}.
\end{proof}

\subsection{The ADM change and a polynomial flux deficit}

All integrals in the next Proposition use the background metric $g$.
At the inner boundary $\nu$ points into $\Omega$, and at a coordinate
sphere $\nu_g$ points toward increasing radius.

\begin{proposition}[Flux and mass comparison]\label{four:maximal:prop:flux}
For the limiting solution at fixed $N$, the flux
\[
 \mathfrak F=\lim_{r\to\infty}\int_{S_r}V\cdot\nu_g\dd A_g
\]
exists and satisfies
\begin{equation}\label{four:maximal:end:mass-flux}
 E_{\ghat}=E_g-\frac{\mathfrak F}{3\omega_3},
 \qquad \omega_3=2\pi^2.
\end{equation}
There is a polynomial bound $\Pi_N$, independent of the exhaustion
radius and of the fixed-$N$ regularity constants, such that
\begin{equation}\label{four:maximal:end:flux-deficit}
 \mathfrak F\ge-\Pi_N\left(\ell+\frac{\ell^2}{\eta_N}\right).
\end{equation}
Consequently, for fixed strict data and $\eps>0$,
$E_{\ghat}\le E_g+o(1)$ as $N\to\infty$.
\end{proposition}

\begin{proof}
The end estimates imply
\begin{align*}
 \cT&=\tfrac12|K|^2+3(p(t)+1)|dt|^2+O_N(e^{-cr}),\\
 u&=1+O_N(r^{-2}),\qquad
 V=3\grad_g t+O_N(r^{-5})+O_N(e^{-cr}).
\end{align*}
In particular $u\Xi$ is absolutely integrable.  Its nonexponential
terms have decay $r^{-2-2q}$, $r^{-6}$, or $r^{-4-\delta}$.
The possibly slower weight $\rho_1$ is multiplied by the exponentially
small $v$ on this fixed-$N$ end.  The divergence theorem gives
\[
 \int_{S_r}V\cdot\nu_g\dd A_g
    -\int_B V\cdot\nu\dd A_g
       =\int_{\Omega_r}u\Xi\dd V_g.
\]
Here $\Omega_r$ denotes the part of $\Omega$ inside $S_r$.
Absolute integrability proves existence of the limit and yields the
orientation-sensitive formula
\begin{equation}\label{four:maximal:end:divergence-flux}
 \mathfrak F=\int_\Omega u\Xi\dd V_g
                        +\int_B V\cdot\nu\dd A_g.
\end{equation}

We next compute the mass change in the same asymptotic chart.  The
linear conformal perturbation is $2t\delta_{ij}$, and in dimension
four its contribution to the ADM numerator is exactly
\[
 \partial_j(2t\delta_{ij})-\partial_i(2t\delta_{jj})
       =-6\partial_i t.
\]
The differentiated remainders are
$O_N(r^{-3-q})+O_N(r^{-5})+O_N(e^{-cr})$ and have vanishing sphere
integrals.  They include products of $t$ with $g-\delta$, quadratic
conformal terms, and the graph perturbation.  Replacing Euclidean
surface flux for $dt$ by $g$-surface flux changes it by
$O_N(r^{-q})$ after integration.  Since the ADM denominator is
$6\omega_3$, we obtain
\[
 E_{\ghat}-E_g
   =-\frac1{\omega_3}\lim_{r\to\infty}
                  \int_{S_r}\partial_{\nu_g}t\dd A_g
   =-\frac{\mathfrak F}{3\omega_3}.
\]
This also proves existence of the ADM energy of $\ghat$ from the
assumed existence of $E_g$; no interchange of the $r$ and $N$ limits
is involved.

To estimate \eqref{four:maximal:end:divergence-flux}, use the boundary flux bound
\eqref{four:maximal:bounds:boundary-flux} and the weighted trace estimate
\eqref{four:maximal:bounds:trace-weighted} with test function one.  They give
\begin{align}\label{four:maximal:end:boundary-cost}
 \int_B|V\cdot\nu|\dd A_g
 &\le C\frac{\eta_N}{\ell}\int_B L_0\dd A_g\notag\\
 &\le\Pi_N\frac{\eta_N}{\ell}
      \int_{\mathrm{collars}}u(1+\sqrt{\cT})\dd V_g\notag\\
 &\le\Pi_N\frac{\eta_N}{\ell}
      \int_{\mathrm{collars}}u(\rho+\delta_0\cT)\dd V_g.
\end{align}
For the last step, $\rho$ has a positive minimum on these fixed compact
collars, and $\sqrt{\cT}\le c\cT+C_c$ with fixed $c>0$.
Only the polynomial constants in those two earlier estimates enter
\eqref{four:maximal:end:boundary-cost}.  Since $\eta_N/\ell=e^{-\eps N/2}$, increasing
$N$ absorbs this entire cost into, for example, one half of the
$u(\rho+\delta_0\cT)$ contribution to the bulk integral.

The negative bulk term is $um_0(t)\mathcal P$.  On its support,
\[
 -\eps\le t\le-\eps+1/N,
 \qquad \ell\le l\le e\ell,\qquad c\ell\le L_0\le C\ell.
\]
These constants are uniform for $0<\eps<1$ and sufficiently large
$N$.  From the definitions, rather than from a derivative estimate,
one obtains
\begin{equation}\label{four:maximal:end:band-identities}
 u\le C(\ell+\ell^2\sigma),\qquad
 uv\le C\ell^2\sigma,\qquad
 ua\sigma=L_0l\sigma^2\ge c\ell^2\sigma^2.
\end{equation}
For the middle inequality use
$v\sqrt D=l^2\sigma^2/\sqrt{1+l^2\sigma^2}\le l\sigma$.
Thus Young's inequality, with $C_N$ polynomial and $0\le m_0\le1$,
gives
\begin{align}\label{four:maximal:end:penalty-loss}
 um_0\mathcal P
 &\le \Pi_N\bigl[\ell\rho_0+
                         \ell^2\sigma(\rho_0+\rho_1)\bigr]\notag\\
 &\le\tfrac12\delta_0\eta_N ua\sigma
       +\Pi_N\left[\ell\rho_0+
                  \frac{\ell^2}{\eta_N}(\rho_0^2+\rho_1^2)\right].
\end{align}
Off the band the left side vanishes, so the final inequality holds
everywhere.  All weights in the loss are integrable:
\[
 \rho_0=r^{-4-\delta}\in L^1,\qquad
 \rho_0^2\in L^1,\qquad
 \rho_1^2=r^{-4\gamma}\in L^1
\]
on the four-dimensional end, using $\delta>0$ and $\gamma>1$.
Their smooth extensions are bounded on the compact interior.  In
particular no integrability of $\rho_1$ itself is being used.

Insert \eqref{four:maximal:end:boundary-cost} and \eqref{four:maximal:end:penalty-loss} into
\eqref{four:maximal:end:divergence-flux}, absorb the indicated nonnegative terms,
and discard the remaining nonnegative bulk terms.  This proves
\eqref{four:maximal:end:flux-deficit}.  Finally,
\[
 \ell+\frac{\ell^2}{\eta_N}
       =e^{-\eps N}+e^{-\eps N/2},
\]
which tends to zero against every polynomial in $N$.  The constants
used to prove \eqref{four:maximal:end:decay} have not entered this estimate.
\end{proof}

\begin{proof}[Completion of Proposition~\ref{four:maximal:prop:geometric-output}]
The threshold exterior and its all-cut comparison are supplied by
Lemma~\ref{four:maximal:lem:threshold-collars}.  Propositions~
\ref{four:maximal:prop:existence} and \ref{four:maximal:prop:end-control}
construct the complete metric with the stated signs and decay.
Its lower metric bound gives the asserted three-dimensional area
comparison on each cut separately.  Proposition~\ref{four:maximal:prop:flux}
gives the energy estimate, since $\ell=e^{-\eps N}$ and
$\eta_N=\ell^{3/2}$.  Every stated output is therefore established
before the Riemannian Penrose inequality is applied.
\end{proof}

\subsection{A minimal enclosing hypersurface}

The boundary $B$ has strictly negative mean curvature in $\ghat$.
We obtain the minimal boundary required for the Riemannian theorem
from the full-perimeter enclosure proposition.  Perimeter here means the full
area of the measure-theoretic boundary of a set, including any portion
coinciding with an obstacle.

\begin{lemma}[Minimal enclosure]\label{four:maximal:lem:minimal-enclosure}
The metric $\ghat$ of Proposition~\ref{four:maximal:prop:end-control} admits a
nonempty compact smooth embedded minimal hypersurface $\Sigma$
enclosing $B$.  Its exterior is connected, complete with boundary,
and has the same single asymptotically flat end.  The hypersurface is
outer area-minimizing, with disconnected competitors permitted, and
its area $\widehat A=|\Sigma|_{\ghat}$ satisfies
\begin{equation}\label{four:maximal:end:area-comparison}
 \widehat A\ge e^{-3\eps}A_*.
\end{equation}
\end{lemma}

\begin{proof}
For the classical codimension-one regularity background, see
\cite[Sections~5.3--5.4]{Federer1969} and \cite{Federer1970}; the
smooth-Riemannian formulation is
\cite[Chapter~7, \S5, Theorem~5.8, p.~200]{Simon2014}.
Apply Proposition~\ref{four:input:enclosure} to
$(X,g_X)=(\Omega,\ghat)$.  The threshold construction gives a smooth
connected orientable four-dimensional exterior with nonempty compact
smooth boundary $B$ and exactly one end.  Proposition~\ref{four:maximal:prop:end-control}
gives completeness with $B$ included and
$\ghat-\delta=O_2(r^{-q})$ with $q>1$.  These are the full geometric
hypotheses of the cited proposition; it imposes no scalar-curvature
condition and does not require decay exponent two.  Its perimeter counts
the entire frontier, including contact with the obstacle.

The strict inequality $\widehat H_B<0$ places the minimizing frontier
$\Sigma$ strictly outside $B$ and makes it a nonempty compact smooth
embedded minimal hypersurface.  The same proposition gives its connected
complete one-ended exterior and outer area-minimization with every
component counted, including disconnected enclosing competitors.

Every such cut encloses the original obstacle.  The cut comparison in
Lemma~\ref{four:maximal:lem:threshold-collars} gives $|\Sigma|_g\ge A_*$.
Restriction of $\ghat\ge e^{-2\eps}g$ to each three-dimensional tangent
plane gives
\[
 |\Sigma|_{\ghat}\ge e^{-3\eps}|\Sigma|_g
                    \ge e^{-3\eps}A_*.
\]
This proves the asserted area comparison.  The argument uses the
general full-perimeter proposition directly, with its full $q>1$ decay range.  An enclosure result requiring
$O_2(r^{-2})$ decay would not suffice for this application.
\end{proof}

\subsection{The Riemannian inequality and the ordered limits}

We use the following dimension-four instance of the boundary
Riemannian Penrose inequality
\cite[Theorem~1.4, p.~84]{BrayLee2009}.  Let $(\mathscr E,g_{\mathscr E})$ be a smooth
connected four-dimensional Riemannian manifold, complete with its
nonempty compact boundary included, having one asymptotically flat
end.  Suppose in that end
\[
 (g_{\mathscr E})_{ij}-\delta_{ij}=O_2(r^{-p_0}),\quad p_0>1,
 \qquad R_{g_{\mathscr E}}=O(r^{-s_0}),\quad s_0>4,
\]
and $R_{g_{\mathscr E}}\ge0$ everywhere.  If its boundary is minimal and outer
area-minimizing, then
\begin{equation}\label{four:maximal:end:bray-lee}
 E_{g_{\mathscr E}}\ge\frac12
          \left(\frac{|\partial\mathscr E|_{g_{\mathscr E}}}{\omega_3}\right)^{2/3}.
\end{equation}
The boundary may have one or more components, and the ADM
normalization is $1/(6\omega_3)$.  The numerical inequality does not
require a spin assumption; the additional spin hypothesis in the
published theorem concerns its equality conclusion.  We use only
\eqref{four:maximal:end:bray-lee}.

\begin{proof}[Proof of Theorem~\ref{four:maximal:thm:main}]
First keep one set of strict approximation data fixed.  For fixed
$\eps>0$ and each sufficiently large $N$, take the exhaustion limit
of Proposition~\ref{four:maximal:prop:end-control}.  Apply
\eqref{four:maximal:end:bray-lee} to the exterior of $\Sigma$ from
Lemma~\ref{four:maximal:lem:minimal-enclosure}, with metric $\ghat$.
That exterior is smooth and complete with boundary, has one end,
and has nonnegative scalar curvature.  Its decay exponents are
$p_0=q>1$ and $s_0=4+\delta>4$ by \eqref{four:maximal:end:geometry}.  Its boundary
is compact, minimal, and outer area-minimizing by the lemma.  Its
ADM energy exists by Proposition~\ref{four:maximal:prop:flux} and is unchanged
by restriction past the compact region.  Thus every hypothesis of
the Riemannian theorem has been checked for this new exterior.

Equations~\eqref{four:maximal:end:mass-flux}, \eqref{four:maximal:end:flux-deficit}, and
\eqref{four:maximal:end:area-comparison} give
\begin{align*}
 E_g+\frac{\Pi_N}{3\omega_3}
                  \left(\ell+\frac{\ell^2}{\eta_N}\right)
 &\ge E_{\ghat}
 \ge\frac12\left(\frac{\widehat A}{\omega_3}\right)^{2/3}\\
 &\ge\frac12e^{-2\eps}
                       \left(\frac{A_*}{\omega_3}\right)^{2/3}.
\end{align*}
Let $N\to\infty$ with these strict data and $\eps$ fixed.  The error
vanishes by its polynomial bound.  Letting $\eps\downarrow0$ next
yields
\[
 E_g\ge\frac12\left(\frac{A_*}{\omega_3}\right)^{2/3}
\]
for the strict data.  No convergence of the metrics $\ghat$ as
$N\to\infty$ has been assumed: only their individual inequalities
and the uniform deficit are used.

Finally remove the strict approximation using
Proposition~\ref{four:maximal:prop:strict-data}.  The approximation energies
converge to the original $E$, and their enclosing infima $A_*$
converge to the original enclosing infimum.  The given outer
area-minimizing hypothesis, including the admissibility of $S$
itself and all disconnected enclosing competitors, identifies this
infimum with $A=|S|_g$.  Therefore
\[
 \boxed{\displaystyle
 E\ge\frac12\left(\frac{A}{2\pi^2}\right)^{2/3}.}
\]
The prescribed $P=0$ and $E>0$ identify this energy with the ADM rest
mass $m=\sqrt{E^2-|P|^2}=E$.  The original boundary was used with its
future MOTS condition; no minimality or vanishing tangential trace
of $K$ on that boundary was imposed.
\end{proof}

\bibliographystyle{plainnat}
\bibliography{references}
\end{document}
